\documentclass[11pt,a4paper]{amsart}

\usepackage{amsmath,amssymb,amsthm,mathtools}
\usepackage[T1]{fontenc}
\usepackage[utf8]{inputenc}
\usepackage{lmodern}
\usepackage{microtype}
\usepackage{hyperref}
\usepackage{orcidlink}
\usepackage{enumitem}
\usepackage[margin=1in]{geometry}

\theoremstyle{plain}
\newtheorem{theorem}{Theorem}[section]
\newtheorem{proposition}[theorem]{Proposition}
\newtheorem{lemma}[theorem]{Lemma}
\newtheorem{corollary}[theorem]{Corollary}

\theoremstyle{definition}
\newtheorem{definition}[theorem]{Definition}
\newtheorem{assumption}[theorem]{Assumption}

\theoremstyle{remark}
\newtheorem{remark}[theorem]{Remark}

\DeclareMathOperator{\KL}{KL}
\DeclareMathOperator{\Var}{Var}
\DeclareMathOperator{\TV}{TV}
\newcommand{\R}{\mathbb{R}}
\newcommand{\N}{\mathbb{N}}
\newcommand{\PP}{\mathbb{P}}
\newcommand{\EE}{\mathbb{E}}
\newcommand{\Acal}{\mathcal A}
\newcommand{\as}{\mathrm{a.s.}}
\newcommand{\dd}{\,d}
\newcommand{\one}{\mathbf 1}

\begin{document}

\title[Dynamic Observability of Latent Contagion]{%
Dynamic Observability of Latent Contagion\\[0.3em]
{\small Sequential Revelation and Observable Screening}}
\author[J.~Vidal Llaurad\'o]{J.~Vidal Llaurad\'o\,\orcidlink{0009-0002-1918-5458}}
\thanks{DOI: \href{https://doi.org/10.2139/ssrn.6539199}{10.2139/ssrn.6539199}.}
\email{vidal.llaurado@gmail.com}
\subjclass[2020]{60G22, 62M15, 91G80, 47B06}
\keywords{rough volatility, latent volatility contagion, Gaussian Volterra models, dynamic observability, sequential revelation, short-horizon prediction, observable screening, fractional transfer operators}
\date{April 2026}

\begin{abstract}
A latent Volterra contagion channel has no single visibility level.  It can accumulate
path-space evidence, improve oracle short-horizon prediction, and still be screened after
conditioning on the target asset's own observable past.  Dynamic observability is formulated as the
residual predictive component left after latent drivers are replaced by observable Gaussian
channels.  In a bivariate Gaussian Volterra model this separates three experiments:
finite-resolution path-space comparison, oracle prediction with latent drivers, and observable
prediction in the target filtration.

In the smoothing regime \(H_{XY}>H_Y\), the relative covariance operator generates a canonical
finite-resolution Gaussian experiment.  Exact likelihood, Kullback--Leibler, Hellinger, and
Bayes-error formulas recover the path-detectability boundary
\(H_{XY}=H_Y+\tfrac14\) as the critical rate at which evidence accumulates with resolution.  At the
oracle latent-driver level, short-horizon prediction has a different boundary, namely
\(H_{XY}=H_Y\).  Rougher contagion improves prediction at leading order, smoother contagion is
dynamically latent at leading order, and the boundary case places both channels on the same scale.

At the observable level, the oracle channel is compressed again.  Passing from latent drivers to
observed Gaussian channels sends dynamic information through a hidden transfer operator.  In the
rougher regime \(0<H_{XY}<H_Y<\tfrac12\), a screening theorem is first proved under an
\(h^\alpha\)-stable boundary-cutoff hypothesis.  The exact-correction closure verifies that
hypothesis for the model classes considered here under the stated \(C^2\) near-diagonal regularity
assumptions.  The endpoint-flat closure proves that the target's own past screens the leading
oracle rough gain:
\[
\|\mathcal F_{t,h}\|_{\mathfrak G_t^*}=o(h^{H_{XY}}),
\qquad
\widetilde G_{X\to Y}(t,h)=o(h^{2H_{XY}}).
\]
Pricing visibility, path-space detectability, and dynamic observability are therefore distinct
visibility levels.  The screened observable component is the inferential object passed to
projected covariance-score inference.
\end{abstract}

\maketitle

\section{Introduction}\label{sec:intro}

The threshold analysis of \cite{vidal2026latent} separates pricing visibility from path-space
detectability for smoothing cross-asset Volterra channels.  That static separation leaves open
which part of a latent channel survives once prediction is conditioned on observable histories.
Dynamic observability is this residual predictive object.

Observable filtrations replace latent drivers by conditional laws generated by asset histories.
Evidence accumulation, oracle short-horizon prediction, and target-past screened prediction are
therefore different experiments.  A channel may be detectable at finite resolution, may improve
oracle prediction when latent drivers are resolved, and may still vanish from the leading gain after
conditioning on the target asset's own past.

The bivariate Gaussian Volterra model keeps the covariance structure explicit enough to make these
reductions exact.  Asset histories enter through Gaussian channels; latent contagion enters through
a cross Volterra kernel.  The dynamic problem asks which component of the oracle signal remains in
the observable filtration, rather than restating the smoothing threshold of \cite{vidal2026latent} in
time-indexed form.

The probabilistic inputs are classical, but each one controls a different level.  The
finite-resolution comparison uses Gaussian experiment theory in the tradition of
Le~Cam~\cite{LeCam86}, Le~Cam and Yang~\cite{LeCamYang00}, van der Vaart~\cite{vdV98}, and
Ibragimov and Has'minskii~\cite{IH81}, together with fractional LAN benchmarks such as
Dahlhaus~\cite{Dahl89} and Brouste--Fukasawa~\cite{BF18}.  The prediction problem uses fractional
filtering results of Gripenberg and Norros~\cite{GN96}, Kleptsyna and
Le~Breton~\cite{KLB02}, Nuzman and Poor~\cite{NP00}, Cai, Chigansky, and
Kleptsyna~\cite{CCK16}, and Afterman, Chigansky, Kleptsyna, and
Marushkevych~\cite{ACKM22}.  The observable reduction is a Volterra transfer problem, using the
operator theory of Gripenberg, Londen, and Staffans~\cite{GLS90}, trace-ideal estimates in the
sense of Simon~\cite{Sim05}, and the fractional covariance representation of Dzhaparidze and van
Zanten~\cite{DvZ05}.

The argument passes through three non-equivalent reductions.  Theorem~\ref{thm:ratesN} converts
the static path-space threshold into a finite-resolution Gaussian experiment.  In the smoothing
regime \(H_{XY}>H_Y\), the relative covariance operator controls likelihood,
Kullback--Leibler, Hellinger, and Bayes-error asymptotics.  The formulas recover
\[
H_{XY}=H_Y+\tfrac14
\]
as the critical evidence-accumulation boundary.

Theorem~\ref{thm:prediction-threshold} separates static evidence from short-horizon prediction.
At the oracle latent-driver level \(H_Y+\tfrac14\) no longer governs.  The prediction boundary is
\[
H_{XY}=H_Y.
\]
Rougher contagion improves prediction at leading order, smoother contagion is dynamically latent at
leading order, and the boundary case puts both channels on the same scale.

Observable screening is the final reduction.  The target past is already contaminated by the
contagion channel, so the source cannot be appended as an independent predictor.  The screening
theorem, Theorem~\ref{thm:observable-screening-smooth}, isolates the relation once
\(h^\alpha\)-stable boundary cutoffs are available.  The exact-correction closure theorem,
Theorem~\ref{thm:observable-screening-closed}, verifies the cutoff hypothesis for the model
classes considered here under the stated \(C^2\) near-diagonal regularity assumptions.  In the
rougher regime \(0<H_{XY}<H_Y<\tfrac12\), the endpoint-flat closure gives
\[
\|\mathcal F_{t,h}\|_{\mathfrak G_t^*}=o(h^{H_{XY}})
\qquad\text{and}\qquad
\widetilde G_{X\to Y}(t,h)=o(h^{2H_{XY}}).
\]
Pricing visibility, path-space detectability, and dynamic observability therefore need not
coincide.

The screened component is an estimand before it is an estimator.  Recent work on rough-volatility
inference has sharpened what can be learned from high-frequency Gaussian or rough observations
\cite{BF18,CHLRS24,CD24}.  The object defined here precedes the choice of econometric procedure.  A
rough channel may exist while its observable content is smaller than both its pricing signature and
its oracle predictive content.

The assumptions strengthen with the object.  The static threshold model uses near-diagonal Volterra
modulation at the \(C^1\) level.  Closed observable screening uses a \(C^2\) near-diagonal expansion
together with endpoint-flat cutoffs.  Projected covariance-score inference requires that
second-order screening structure plus the sampling and proxy-dominance conditions of the projected
experiment.  The hierarchy records assumption import across objects rather than an inconsistency between
the threshold and screening statements.

\subsection*{Results by reduction}
The finite-resolution Gaussian experiment comes first.  It gives quantitative content to the
path-detectability threshold of \cite{vidal2026latent} and determines the rate at which statistical
evidence accumulates with observation resolution.

The oracle prediction theorem separates prediction from path-space comparison.  It identifies the
dynamic threshold \(H_{XY}=H_Y\) and proves that the short-horizon predictive problem is not the
static detectability problem rewritten in time.

The observable screening block completes the reduction.  It defines the inferential target used by
feasible source-screened inference: the part of the oracle signal that remains after conditioning on
the target channel and restricting information to observables.  The finite-resolution experiment
fixes the evidence scale.  The oracle theorem gives the latent predictive benchmark.  The screening
theorem under \(h^\alpha\)-stable boundary cutoffs isolates the screening relation, and the
exact-correction theorem, Theorem~\ref{thm:observable-screening-closed}, verifies the cutoff
structure for the model classes treated here.  The endpoint-flat closure fixes the observable object
to be estimated.

\subsection*{Organization}
Section~\ref{sec:finres} develops the finite-resolution Gaussian experiment and its information
functionals in the smoothing regime.  Sections~\ref{sec:prediction} and~\ref{sec:observed}
formulate the oracle and observable short-horizon prediction problems.  Sections~\ref{sec:geometry}
through~\ref{sec:micro-solver} build the exact Volterra reduction and the frozen fractional model.
Sections~\ref{sec:micro-regularity} through~\ref{sec:screening} prove the screening theorem under
\(h^\alpha\)-stable boundary cutoffs.  The final sections verify those cutoffs by exact correction,
establish the closed observable theorem in the model classes treated here, and record the
consequences in Section~\ref{sec:conclusion}.

\section{Finite-resolution information geometry}\label{sec:finres}

The first reduction turns the static smoothing threshold of \cite{vidal2026latent} into a finite
Gaussian experiment.  The governing object is the positive compact relative covariance operator
\[
\mathcal R_{Y,XY}:=C_0^{-1/2}(C_1-C_0)C_0^{-1/2}
\]
and, in the smoothing regime, it satisfies
\[
\tau_n(\mathcal R_{Y,XY})\asymp n^{-2\delta},
\qquad
\delta:=H_{XY}-H_Y>0,
\]
where \(\tau_n(\mathcal R_{Y,XY})\) denotes the \(n\)th eigenvalue in decreasing order.  This
eigenvalue order fixes the information scale of the finite-resolution experiment.

The same spectral profile generates the canonical finite-resolution Gaussian experiment.  It also
fixes the rate at which statistical evidence accumulates with the resolved dimension $N$ or,
equivalently, with an observation scale $\Delta$ such that $N_\Delta\asymp \Delta^{-1}$.

\begin{definition}[Canonical $N$-mode experiment]\label{def:EN}
Let $\{\tau_n\}_{n\ge1}$ be the eigenvalues of \(\mathcal R_{Y,XY}\), repeated with multiplicity and ordered
decreasingly. For each $N\in\N$, define the centred Gaussian laws
\[
\nu_{0,N}:=\bigotimes_{n=1}^N \mathcal N(0,1),
\qquad
\nu_{1,N}:=\bigotimes_{n=1}^N \mathcal N(0,1+\tau_n)
\]
on $\R^N$.  The experiment
\[
\mathcal E_N := (\R^N,\nu_{0,N},\nu_{1,N})
\]
is the \emph{canonical $N$-mode experiment} associated with \(\mathcal R_{Y,XY}\).
\end{definition}

\begin{remark}[Interpretation of the normal form]
Definition~\ref{def:EN} is the spectral normal form associated with the relative covariance
operator \(\mathcal R_{Y,XY}\).  In the Gaussian comparison problem, whitening and diagonalisation of
the pair $(C_0,C_1)$ reduce the finite-resolution comparison to the same eigenvalue profile
$\{\tau_n\}$.  All finite-resolution information quantities used below depend only on that profile.
The experiment $\mathcal E_N$ is therefore the canonical object to analyse; see, for example,
Le~Cam~\cite{LeCam86}, Le~Cam and Yang~\cite[Ch.~2]{LeCamYang00}, van~der~Vaart~\cite{vdV98}, and
Ibragimov and Has'minskii~\cite{IH81}.
\end{remark}

For later reference set
\[
        E_N:=\sum_{n=1}^N\tau_n,
        \qquad
        I_N:=\sum_{n=1}^N\tau_n^2.
\]
The first quantity is a finite-resolution latent-energy diagnostic.  The second is the quadratic
information scale of the product Gaussian experiment.  The evidence boundary below is governed by
\(I_N\), not by \(E_N\), because likelihood, Kullback--Leibler, Hellinger, and Bayes-error
asymptotics all reduce to the quadratic profile.

\begin{proposition}[Likelihood ratio in normal form]\label{prop:lr}
For every $N\in\N$, the law $\nu_{1,N}$ is absolutely continuous with respect to
$\nu_{0,N}$, with Radon--Nikodym derivative
\begin{equation}\label{eq:lr}
\frac{d\nu_{1,N}}{d\nu_{0,N}}(x)
=
\prod_{n=1}^N (1+\tau_n)^{-1/2}
\exp\!\left(
\frac12 \sum_{n=1}^N \frac{\tau_n}{1+\tau_n}x_n^2
\right),
\qquad x=(x_1,\dots,x_N)\in\R^N.
\end{equation}
Equivalently,
\begin{equation}\label{eq:loglr}
\log \frac{d\nu_{1,N}}{d\nu_{0,N}}(x)
=
-\frac12\sum_{n=1}^N \log(1+\tau_n)
+ \frac12\sum_{n=1}^N \frac{\tau_n}{1+\tau_n}x_n^2.
\end{equation}
\end{proposition}

\begin{proof}
For one coordinate,
\[
\frac{d\mathcal N(0,1+\tau_n)}{d\mathcal N(0,1)}(x_n)
=
(1+\tau_n)^{-1/2}
\exp\!\left(
-\frac{x_n^2}{2(1+\tau_n)}+\frac{x_n^2}{2}
\right)
\]
for $\tau_n\ge0$. Since
\[
-\frac{1}{2(1+\tau_n)}+\frac12
=
\frac{\tau_n}{2(1+\tau_n)},
\]
the one-dimensional density ratio is
\[
(1+\tau_n)^{-1/2}
\exp\!\left(\frac{\tau_n}{2(1+\tau_n)}x_n^2\right).
\]
Independence of the coordinates under both product laws gives~\eqref{eq:lr}, and taking
logarithms yields~\eqref{eq:loglr}.
\end{proof}

\begin{proposition}[Exact divergence functionals]\label{prop:div}
For every $N\in\N$:
\begin{enumerate}[label=(\roman*)]
\item the two Kullback--Leibler divergences are
\begin{align}
\KL(\nu_{1,N}\,\|\,\nu_{0,N})
&=
\frac12\sum_{n=1}^N \Bigl(\tau_n-\log(1+\tau_n)\Bigr), \label{eq:kl10}\\
\KL(\nu_{0,N}\,\|\,\nu_{1,N})
&=
\frac12\sum_{n=1}^N \left(\log(1+\tau_n)-\frac{\tau_n}{1+\tau_n}\right); \label{eq:kl01}
\end{align}
\item the Hellinger affinity is
\begin{equation}\label{eq:affinity}
\Acal_N
:=
\int_{\R^N}\sqrt{\frac{d\nu_{1,N}}{dx}\frac{d\nu_{0,N}}{dx}}\,dx
=
\prod_{n=1}^N
\frac{\sqrt{2}\,(1+\tau_n)^{1/4}}{(2+\tau_n)^{1/2}};
\end{equation}
\item if $\pi_N^\ast$ denotes the optimal Bayes misclassification probability for equal
priors, then
\begin{equation}\label{eq:bayes-bounds}
\frac12\Bigl(1-\sqrt{1-\Acal_N^2}\Bigr)
\le
\pi_N^\ast
\le
\frac12 \Acal_N.
\end{equation}
\end{enumerate}
\end{proposition}

\begin{proof}
For item~(i), a centred Gaussian law $\mathcal N(0,\sigma^2)$ satisfies
\[
\KL(\mathcal N(0,\sigma^2)\,\|\,\mathcal N(0,1))
=
\frac12(\sigma^2-1-\log \sigma^2).
\]
Applying this coordinatewise with $\sigma_n^2=1+\tau_n$ and using independence yields
~\eqref{eq:kl10}. Likewise,
\[
\KL(\mathcal N(0,1)\,\|\,\mathcal N(0,\sigma^2))
=
\frac12\left(\log \sigma^2+\frac{1}{\sigma^2}-1\right),
\]
which gives~\eqref{eq:kl01} after substituting $\sigma_n^2=1+\tau_n$ and simplifying
\[
\frac{1}{1+\tau_n}-1=-\frac{\tau_n}{1+\tau_n}.
\]

For item~(ii), in one dimension the Hellinger affinity between
$\mathcal N(0,\sigma_0^2)$ and $\mathcal N(0,\sigma_1^2)$ is
\[
\sqrt{\frac{2\sigma_0\sigma_1}{\sigma_0^2+\sigma_1^2}}.
\]
With $\sigma_0=1$ and $\sigma_1=\sqrt{1+\tau_n}$ this becomes
\[
\frac{\sqrt{2}\,(1+\tau_n)^{1/4}}{(2+\tau_n)^{1/2}}.
\]
Multiplying over independent coordinates gives~\eqref{eq:affinity}.

For item~(iii), the optimal Bayes error with equal priors is
\[
\pi_N^\ast = \frac12\bigl(1-\TV(\nu_{1,N},\nu_{0,N})\bigr).
\]
Also,
\[
\TV(\nu_{1,N},\nu_{0,N})
=
\frac12\int_{\R^N} \left|
\sqrt{\frac{d\nu_{1,N}}{dx}}-\sqrt{\frac{d\nu_{0,N}}{dx}}
\right|
\left(
\sqrt{\frac{d\nu_{1,N}}{dx}}+\sqrt{\frac{d\nu_{0,N}}{dx}}
\right)dx.
\]
By Cauchy--Schwarz,
\[
\TV(\nu_{1,N},\nu_{0,N})
\le
\sqrt{1-\Acal_N^2},
\]
which gives the lower bound in~\eqref{eq:bayes-bounds}. For the upper bound, using
$\min(a,b)\le \sqrt{ab}$ for $a,b\ge0$,
\[
2\pi_N^\ast
=
\int_{\R^N} \min\!\left(\frac{d\nu_{1,N}}{dx},\frac{d\nu_{0,N}}{dx}\right)dx
\le
\int_{\R^N}\sqrt{\frac{d\nu_{1,N}}{dx}\frac{d\nu_{0,N}}{dx}}\,dx
=
\Acal_N.
\]
This proves~\eqref{eq:bayes-bounds}.
\end{proof}

\begin{lemma}[Quadratic control of the divergence profiles]\label{lem:quadratic}
There exist constants $0<c_1\le c_2<\infty$ such that for every $\tau\in[0,1]$,
\begin{align}
c_1\tau^2
\le
\tau-\log(1+\tau)
\le
c_2\tau^2, \label{eq:quad-kl10}\\
c_1\tau^2
\le
\log(1+\tau)-\frac{\tau}{1+\tau}
\le
c_2\tau^2, \label{eq:quad-kl01}\\
c_1\tau^2
\le
-\log\!\left(
\frac{\sqrt{2}\,(1+\tau)^{1/4}}{(2+\tau)^{1/2}}
\right)
\le
c_2\tau^2. \label{eq:quad-aff}
\end{align}
\end{lemma}

\begin{proof}
Define
\[
\phi_1(\tau):=\tau-\log(1+\tau),
\qquad
\phi_0(\tau):=\log(1+\tau)-\frac{\tau}{1+\tau},
\]
and
\[
\psi(\tau):=
-\log\!\left(
\frac{\sqrt{2}\,(1+\tau)^{1/4}}{(2+\tau)^{1/2}}
\right).
\]
All three functions vanish at $\tau=0$ and satisfy
\[
\phi_1(\tau)=\frac12\tau^2+O(\tau^3),
\qquad
\phi_0(\tau)=\frac12\tau^2+O(\tau^3),
\qquad
\psi(\tau)=\frac{1}{16}\tau^2+O(\tau^3)
\]
as $\tau\downarrow0$. Therefore the quotients
\[
\frac{\phi_1(\tau)}{\tau^2},
\qquad
\frac{\phi_0(\tau)}{\tau^2},
\qquad
\frac{\psi(\tau)}{\tau^2}
\]
extend continuously to $\tau=0$ with positive finite limits. By continuity on the compact
interval $[0,1]$, each quotient is bounded above and below by positive constants, which
proves~\eqref{eq:quad-kl10}--\eqref{eq:quad-aff}.
\end{proof}

\begin{theorem}[Sharp information accumulation at resolution $N$]
\label{thm:ratesN}
Assume there exist $n_0\in\N$, $\delta>0$, and constants
$0<c_-\le c_+<\infty$ such that
\begin{equation}\label{eq:tau-asymp}
c_- n^{-2\delta}\le \tau_n \le c_+ n^{-2\delta},
\qquad n\ge n_0.
\end{equation}
Then there exist constants $0<C_1\le C_2<\infty$ such that for all $N\ge n_0$:
\begin{align}
C_1\sum_{n=1}^N n^{-4\delta}
\le
\KL(\nu_{1,N}\,\|\,\nu_{0,N})
\le
C_2\sum_{n=1}^N n^{-4\delta}, \label{eq:rate-kl10}\\
C_1\sum_{n=1}^N n^{-4\delta}
\le
\KL(\nu_{0,N}\,\|\,\nu_{1,N})
\le
C_2\sum_{n=1}^N n^{-4\delta}, \label{eq:rate-kl01}\\
C_1\sum_{n=1}^N n^{-4\delta}
\le
-\log \Acal_N
\le
C_2\sum_{n=1}^N n^{-4\delta}. \label{eq:rate-aff}
\end{align}
Consequently, as $N\to\infty$,
\begin{equation}\label{eq:phase-rates}
\KL(\nu_{1,N}\,\|\,\nu_{0,N}),
\;
\KL(\nu_{0,N}\,\|\,\nu_{1,N}),
\;
-\log \Acal_N
\asymp
\begin{cases}
N^{1-4\delta}, & 0<\delta<\tfrac14,\\[4pt]
\log N, & \delta=\tfrac14,\\[4pt]
1, & \delta>\tfrac14.
\end{cases}
\end{equation}
\end{theorem}

\begin{proof}
Because \(\tau_n\to0\), choose \(n_1\ge n_0\) such that \(\tau_n\le1\) for every
\(n\ge n_1\).  Applying Lemma~\ref{lem:quadratic} termwise in the exact formulas
\eqref{eq:kl10}--\eqref{eq:affinity}, there exist constants
\(0<C_1'\le C_2'<\infty\) such that, for every \(N\ge n_1\),
\begin{align*}
C_1' \sum_{n=n_1}^N \tau_n^2
\le
\KL(\nu_{1,N}\,\|\,\nu_{0,N})
\le
C_2' \sum_{n=n_1}^N \tau_n^2,\\
C_1' \sum_{n=n_1}^N \tau_n^2
\le
\KL(\nu_{0,N}\,\|\,\nu_{1,N})
\le
C_2' \sum_{n=n_1}^N \tau_n^2,\\
C_1' \sum_{n=n_1}^N \tau_n^2
\le
-\log \Acal_N
\le
C_2' \sum_{n=n_1}^N \tau_n^2.
\end{align*}
The eigenvalue bounds give
\[
c_-^2 \sum_{n=n_1}^N n^{-4\delta}
\le
\sum_{n=n_1}^N \tau_n^2
\le
c_+^2 \sum_{n=n_1}^N n^{-4\delta}.
\]
The finite prefix \(\{1,\dots,n_1-1\}\) changes both the information functionals and the comparison
series only by finite constants.  Adjusting the constants on the finite set
\(n_0\le N<n_1\), and then on the finite prefix for \(N\ge n_1\), gives
\eqref{eq:rate-kl10}--\eqref{eq:rate-aff} for every \(N\ge n_0\).

The remaining step is the integral-test asymptotic
\[
\sum_{n=1}^N n^{-4\delta}
\asymp
\begin{cases}
N^{1-4\delta}, & 0<\delta<\tfrac14,\\[4pt]
\log N, & \delta=\tfrac14,\\[4pt]
1, & \delta>\tfrac14.
\end{cases}
\]
Substitution into \eqref{eq:rate-kl10}--\eqref{eq:rate-aff} gives \eqref{eq:phase-rates}.
\end{proof}

\begin{corollary}[Observation-scale formulation]\label{cor:delta}
Let $\Delta\in(0,1)$ be an observation scale and set
\[
N_\Delta := \lfloor \Delta^{-1}\rfloor.
\]
Under the assumptions of Theorem~\ref{thm:ratesN},
\begin{equation}\label{eq:delta-rates}
\KL(\nu_{1,N_\Delta}\,\|\,\nu_{0,N_\Delta}),
\;
\KL(\nu_{0,N_\Delta}\,\|\,\nu_{1,N_\Delta}),
\;
-\log \Acal_{N_\Delta}
\asymp
\begin{cases}
\Delta^{-(1-4\delta)}, & 0<\delta<\tfrac14,\\[4pt]
\log(\Delta^{-1}), & \delta=\tfrac14,\\[4pt]
1, & \delta>\tfrac14,
\end{cases}
\end{equation}
as $\Delta\downarrow0$.
\end{corollary}

\begin{proof}
The relation \(N_\Delta\asymp \Delta^{-1}\) transfers the three cases in
\eqref{eq:phase-rates} to the observation scale.  If \(0<\delta<\tfrac14\), then
\(N_\Delta^{1-4\delta}\asymp \Delta^{-(1-4\delta)}\).  If \(\delta=\tfrac14\), then
\(\log N_\Delta\asymp \log(\Delta^{-1})\).  If \(\delta>\tfrac14\), the rate remains of order
one.  This gives \eqref{eq:delta-rates}.
\end{proof}

\begin{corollary}[Finite-resolution discrimination threshold]\label{cor:bayes}
Under the assumptions of Theorem~\ref{thm:ratesN}, the optimal Bayes error
$\pi_N^\ast$ satisfies:
\begin{enumerate}[label=(\roman*)]
\item if $0<\delta\le\tfrac14$, then $\Acal_N\to0$ and therefore
\[
\pi_N^\ast \to 0
\qquad\text{as }N\to\infty;
\]
\item if $\delta>\tfrac14$, then $\Acal_N\to \Acal_\infty$ for some
$\Acal_\infty\in(0,1]$, and
\[
\liminf_{N\to\infty}\pi_N^\ast
\ge
\frac12\Bigl(1-\sqrt{1-\Acal_\infty^2}\Bigr)
>0.
\]
\end{enumerate}
\end{corollary}

\begin{proof}
If $0<\delta\le\tfrac14$, then~\eqref{eq:rate-aff} implies
$-\log\Acal_N\to\infty$, hence $\Acal_N\to0$. The upper bound in~\eqref{eq:bayes-bounds}
then yields $\pi_N^\ast\to0$.

If $\delta>\tfrac14$, then~\eqref{eq:rate-aff} shows that $-\log \Acal_N$ remains bounded
as $N\to\infty$. Each Hellinger factor belongs to \(0,1]\), so \(\Acal_N\) is decreasing.
The boundedness of \(-\log \Acal_N\) keeps the limit strictly positive, and therefore
\(\Acal_N\to\Acal_\infty\in(0,1]\). The lower bound in~\eqref{eq:bayes-bounds} therefore gives
\[
\liminf_{N\to\infty}\pi_N^\ast
\ge
\frac12\Bigl(1-\sqrt{1-\Acal_\infty^2}\Bigr).
\]
Because \eqref{eq:tau-asymp} gives $\tau_n>0$ for every $n\ge n_0$, the Hellinger factor at
$n_0$ is strictly smaller than $1$. Hence $\Acal_N\le \Acal_{n_0}<1$ for every
$N\ge n_0$, and therefore $\Acal_\infty\le \Acal_{n_0}<1$. Since the boundedness of
$-\log\Acal_N$ gives $\Acal_\infty>0$, the right-hand side is strictly positive.
\end{proof}

\section{Short-horizon sequential prediction}\label{sec:prediction}

The second reduction is predictive.  The bivariate Volterra state already contains a short-horizon
boundary before pricing or observable screening enters.  The competition between the idiosyncratic
exponent $H_Y$ and the contagion exponent $H_{XY}$ appears in prediction risk itself.  For the
prediction and filtering background used here, see in particular Gripenberg and Norros~\cite{GN96},
Kleptsyna and Le~Breton~\cite{KLB02}, Nuzman and Poor~\cite{NP00}, and Afterman, Chigansky,
Kleptsyna, and Marushkevych~\cite{ACKM22}.

\begin{assumption}[Bivariate Volterra volatility factor]\label{ass:model}
Fix $T>0$. Let $W^\perp$ and $W^X$ be independent standard Brownian motions. For
$i\in\{Y,XY\}$ let $H_i\in(0,\tfrac12)$ and let $L_i$ be defined on
\[
\Delta_T:=\{(t,s)\in[0,T]^2:0\le s\le t\le T\},
\]
with the following properties:
\begin{enumerate}[label=(\roman*)]
\item $L_i$ extends to a $C^1$ function on a neighbourhood of $\Delta_T$;
\item $L_i$ and its first partial derivatives are bounded on $\Delta_T$;
\item $L_i(t,t)>0$ for every $t\in[0,T]$.
\end{enumerate}
Define the kernels
\[
K_i(t,s):=(t-s)^{H_i-\frac12}L_i(t,s)\one_{\{0\le s<t\le T\}},
\qquad i\in\{Y,XY\}.
\]
Let $\nu_Y>0$ and $\eta\ge0$, and define the Gaussian volatility factor
\begin{equation}\label{eq:Yfactor}
\mathcal Y_t
:=
\nu_Y\int_0^t K_Y(t,s)\,\dd W_s^\perp
\;+\;
\eta\int_0^t K_{XY}(t,s)\,\dd W_s^X,
\qquad 0\le t\le T.
\end{equation}
\end{assumption}

\begin{definition}[Oracle intrinsic information filtrations]\label{def:intrinsic-filtrations}
For $0\le t\le T$, define
\[
\mathcal G_t^Y:=\sigma(W_s^\perp:0\le s\le t),
\qquad
\mathcal G_t^{X,Y}:=\sigma(W_s^\perp,W_s^X:0\le s\le t).
\]
For $0\le t<t+h\le T$, set
\[
\Delta_h \mathcal Y_t := \mathcal Y_{t+h}-\mathcal Y_t.
\]
The associated mean-square prediction risks are
\begin{align}
R_Y(t,h)
&:=
\EE\!\left[
\bigl(\Delta_h\mathcal Y_t-\EE[\Delta_h\mathcal Y_t\mid \mathcal G_t^Y]\bigr)^2
\right], \label{eq:RY}\\
R_{XY}(t,h)
&:=
\EE\!\left[
\bigl(\Delta_h\mathcal Y_t-\EE[\Delta_h\mathcal Y_t\mid \mathcal G_t^{X,Y}]\bigr)^2
\right]. \label{eq:RXY}
\end{align}
We also write
\begin{equation}\label{eq:gain}
G_{X\to Y}(t,h):=R_Y(t,h)-R_{XY}(t,h)
\end{equation}
for the prediction gain generated by the contagion channel.
\end{definition}

\begin{remark}[Oracle character of the intrinsic filtration block]
The filtrations of Definition~\ref{def:intrinsic-filtrations} separate the idiosyncratic latent
driver from the enlarged latent filtration that also contains the contagion driver up to time $t$.
They are \emph{not} yet the sigma-fields generated by observable asset prices or
realised-volatility proxies.  The section identifies the intrinsic short-horizon prediction gain
available when the cross-asset latent channel is resolved exactly.  This is the benchmark for the
observable-filtration problem.
\end{remark}

\begin{definition}[Universal short-horizon constants]\label{def:constants}
For $H\in(0,\tfrac12)$, define
\begin{equation}\label{eq:aH}
a_H:=\frac{1}{2H},
\end{equation}
and
\begin{equation}\label{eq:bH}
b_H
:=
\int_0^\infty \Bigl((1+u)^{H-\frac12}-u^{H-\frac12}\Bigr)^2\,\dd u.
\end{equation}
Also set
\begin{equation}\label{eq:qH}
q_H:=a_H+b_H.
\end{equation}
\end{definition}

\begin{lemma}[Finiteness of the past-resolving constant]\label{lem:bH}
For every $H\in(0,\tfrac12)$ one has $b_H\in(0,\infty)$.
\end{lemma}

\begin{proof}
Near $u=0$, since $H-\tfrac12<0$,
\[
\Bigl((1+u)^{H-\frac12}-u^{H-\frac12}\Bigr)^2
\asymp
u^{2H-1},
\]
and $u^{2H-1}$ is integrable at $0$ because $2H-1>-1$.

As $u\to\infty$, the mean value theorem gives
\[
(1+u)^{H-\frac12}-u^{H-\frac12}
=
\left(H-\frac12\right)\xi_u^{H-\frac32}
\]
for some $\xi_u\in(u,u+1)$, hence
\[
\Bigl((1+u)^{H-\frac12}-u^{H-\frac12}\Bigr)^2
\lesssim
u^{2H-3},
\]
and $u^{2H-3}$ is integrable at $\infty$ because $2H-3<-2$.
The integrand is not identically zero, so $b_H>0$.
\end{proof}

\begin{proposition}[Exact oracle projection formulas for the prediction risks]
\label{prop:exact-risks}
Fix $0\le t<t+h\le T$. Then
\begin{align}
R_{XY}(t,h)
&=
\nu_Y^2 \int_t^{t+h} K_Y(t+h,s)^2\,\dd s
+ \eta^2 \int_t^{t+h} K_{XY}(t+h,s)^2\,\dd s, \label{eq:RXY-exact}\\
R_Y(t,h)
&=
\nu_Y^2 \int_t^{t+h} K_Y(t+h,s)^2\,\dd s
+ \eta^2 \Var(\Delta_h Z_t), \label{eq:RY-exact}\\
G_{X\to Y}(t,h)
&=
\eta^2 \int_0^t \bigl(K_{XY}(t+h,s)-K_{XY}(t,s)\bigr)^2\,\dd s, \label{eq:G-exact}
\end{align}
where
\[
Z_t:=\int_0^t K_{XY}(t,s)\,\dd W_s^X.
\]
Equivalently,
\begin{equation}\label{eq:Var-increment}
\Var(\Delta_h Z_t)
=
\int_0^t \bigl(K_{XY}(t+h,s)-K_{XY}(t,s)\bigr)^2\,\dd s
+ \int_t^{t+h} K_{XY}(t+h,s)^2\,\dd s.
\end{equation}
\end{proposition}

\begin{proof}
Define
\[
V_t:=\int_0^t K_Y(t,s)\,\dd W_s^\perp,
\qquad
Z_t:=\int_0^t K_{XY}(t,s)\,\dd W_s^X.
\]
By~\eqref{eq:Yfactor},
\[
\Delta_h \mathcal Y_t
=
\nu_Y \Delta_h V_t + \eta \Delta_h Z_t,
\]
where
\[
\Delta_h V_t
=
\int_0^t \bigl(K_Y(t+h,s)-K_Y(t,s)\bigr)\,\dd W_s^\perp
+ \int_t^{t+h} K_Y(t+h,s)\,\dd W_s^\perp,
\]
and
\[
\Delta_h Z_t
=
\int_0^t \bigl(K_{XY}(t+h,s)-K_{XY}(t,s)\bigr)\,\dd W_s^X
+ \int_t^{t+h} K_{XY}(t+h,s)\,\dd W_s^X.
\]

Conditioning on $\mathcal G_t^{X,Y}$, the integrals over $[0,t]$ are measurable, while
the integrals over $(t,t+h]$ are independent centred Gaussians. Therefore
\[
\Delta_h \mathcal Y_t-\EE[\Delta_h\mathcal Y_t\mid \mathcal G_t^{X,Y}]
=
\nu_Y\int_t^{t+h} K_Y(t+h,s)\,\dd W_s^\perp
+ \eta\int_t^{t+h} K_{XY}(t+h,s)\,\dd W_s^X.
\]
Applying It\^o isometry and independence yields~\eqref{eq:RXY-exact}.

Conditioning instead on $\mathcal G_t^Y$, the term $\Delta_h Z_t$ is a centred functional of
\(W^X\) and is independent of the sigma-field generated by \(W^\perp\) up to time \(t\), so
\[
\Delta_h \mathcal Y_t-\EE[\Delta_h\mathcal Y_t\mid \mathcal G_t^Y]
=
\nu_Y\int_t^{t+h} K_Y(t+h,s)\,\dd W_s^\perp
+ \eta\,\Delta_h Z_t.
\]
The two displayed summands are independent, so It\^o isometry gives
\[
R_Y(t,h)
=
\nu_Y^2 \int_t^{t+h} K_Y(t+h,s)^2\,\dd s
+ \eta^2 \Var(\Delta_h Z_t),
\]
which is~\eqref{eq:RY-exact}.

Finally, subtracting~\eqref{eq:RXY-exact} from~\eqref{eq:RY-exact} and using the
decomposition of $\Delta_h Z_t$ into the past-resolved and future-innovation pieces gives
\eqref{eq:G-exact} and~\eqref{eq:Var-increment}.
\end{proof}

\begin{proposition}[Innovation variance asymptotics]\label{prop:new}
Let $H\in(0,\tfrac12)$ and let $L$ satisfy the regularity assumptions of
Assumption~\ref{ass:model}. Define
\[
K_H(t,s):=(t-s)^{H-\frac12}L(t,s)\one_{\{0\le s<t\le T\}}.
\]
Then for every fixed $t\in[0,T)$,
\begin{equation}\label{eq:new-asymp}
\int_t^{t+h} K_H(t+h,s)^2\,\dd s
=
L(t,t)^2 a_H h^{2H}+o(h^{2H}),
\qquad h\downarrow0.
\end{equation}
\end{proposition}

\begin{proof}
By the change of variables $s=t+h-hu$,
\[
\int_t^{t+h} K_H(t+h,s)^2\,\dd s
=
h^{2H}\int_0^1 u^{2H-1}L(t+h,t+h-hu)^2\,\dd u.
\]
Since $L$ is continuous on a neighbourhood of the diagonal,
\[
L(t+h,t+h-hu)\to L(t,t)
\qquad\text{for every }u\in[0,1]
\]
as $h\downarrow0$, and the integrand is dominated by
$C u^{2H-1}$ for some constant $C<\infty$. Because $u^{2H-1}$ is integrable on $(0,1)$,
dominated convergence yields
\[
h^{-2H}\int_t^{t+h} K_H(t+h,s)^2\,\dd s
\longrightarrow
L(t,t)^2\int_0^1 u^{2H-1}\,\dd u
=
L(t,t)^2 a_H.
\]
This proves~\eqref{eq:new-asymp}.
\end{proof}

\begin{proposition}[Past-resolved variance asymptotics]\label{prop:old}
Let $H\in(0,\tfrac12)$ and let $L$ satisfy the regularity assumptions of
Assumption~\ref{ass:model}. Define $K_H$ as in Proposition~\ref{prop:new}. Then for every
fixed $t\in(0,T)$,
\begin{equation}\label{eq:old-asymp}
\int_0^t \bigl(K_H(t+h,s)-K_H(t,s)\bigr)^2\,\dd s
=
L(t,t)^2 b_H h^{2H}+o(h^{2H}),
\qquad h\downarrow0.
\end{equation}
\end{proposition}

\begin{proof}
Write $\alpha:=H-\tfrac12\in(-\tfrac12,0)$ and change variables $u=t-s$. Then
\[
I_h
:=
\int_0^t \bigl(K_H(t+h,s)-K_H(t,s)\bigr)^2\,\dd s
=
\int_0^t \Bigl((u+h)^\alpha L(t+h,t-u)-u^\alpha L(t,t-u)\Bigr)^2\,\dd u.
\]

Split $I_h=I_h^{\mathrm{near}}+I_h^{\mathrm{far}}$, where
\[
I_h^{\mathrm{near}}:=\int_0^{h^{1/2}} (\cdots)^2\,\dd u,
\qquad
I_h^{\mathrm{far}}:=\int_{h^{1/2}}^t (\cdots)^2\,\dd u.
\]

\emph{Near-diagonal part.}
Since $L$ is $C^1$, there exists $C<\infty$ such that for all small $h$ and all
$u\in[0,h^{1/2}]$,
\[
|L(t+h,t-u)-L(t,t)|+|L(t,t-u)-L(t,t)|\le C(h+u)\le 2C h^{1/2}.
\]
Hence
\begin{align*}
&(u+h)^\alpha L(t+h,t-u)-u^\alpha L(t,t-u) \\
&\qquad=
L(t,t)\bigl((u+h)^\alpha-u^\alpha\bigr)+r_h(u),
\end{align*}
with
\[
|r_h(u)|\le C h^{1/2}\bigl((u+h)^\alpha+u^\alpha\bigr).
\]
Therefore
\[
I_h^{\mathrm{near}}
=
L(t,t)^2 \int_0^{h^{1/2}} \bigl((u+h)^\alpha-u^\alpha\bigr)^2\,\dd u
+ o(h^{2H}),
\]
because
\[
\int_0^{h^{1/2}} r_h(u)^2\,\dd u
\lesssim
h\int_0^{h^{1/2}} \bigl((u+h)^{2H-1}+u^{2H-1}\bigr)\,\dd u
=
O(h^{1+H})
=
o(h^{2H}).
\]
The main model term is $O(h^{2H})$ by Lemma~\ref{lem:bH}, so the corresponding cross term is
$o(h^{2H})$ by Cauchy--Schwarz.

Now substitute $u=hv$ in the main term:
\[
\int_0^{h^{1/2}} \bigl((u+h)^\alpha-u^\alpha\bigr)^2\,\dd u
=
h^{2H}\int_0^{h^{-1/2}} \bigl((1+v)^\alpha-v^\alpha\bigr)^2\,\dd v.
\]
By Lemma~\ref{lem:bH}, the integrand is integrable on $(0,\infty)$, so
\[
h^{-2H}I_h^{\mathrm{near}}
\longrightarrow
L(t,t)^2 b_H.
\]

\emph{Far-from-diagonal part.}
For $u\ge h^{1/2}$, define
\[
F_h(u):=(u+h)^\alpha L(t+h,t-u).
\]
The mean value theorem in $h$ gives
\[
|F_h(u)-F_0(u)|
\le
h\sup_{0\le r\le h}
\left|
\partial_r\bigl((u+r)^\alpha L(t+r,t-u)\bigr)
\right|.
\]
Since $u\ge h^{1/2}$ and $L,\partial_1L$ are bounded,
\[
|F_h(u)-F_0(u)|
\le
Ch\bigl(u^{\alpha-1}+u^\alpha\bigr)
\]
for some constant $C$. Consequently,
\begin{align*}
I_h^{\mathrm{far}}
&\lesssim
h^2\int_{h^{1/2}}^t \bigl(u^{2\alpha-2}+u^{2\alpha}\bigr)\,\dd u \\
&=
O\!\left(h^2\int_{h^{1/2}}^t u^{2H-3}\,\dd u\right)+O(h^2) \\
&=
O(h^{1+H})+O(h^2)
=
o(h^{2H}),
\end{align*}
because $H<\tfrac12$.

Combining the near and far parts proves~\eqref{eq:old-asymp}.
\end{proof}

\begin{corollary}[Short-horizon increment asymptotics]\label{cor:increment}
Under the assumptions of Proposition~\ref{prop:old},
\begin{equation}\label{eq:increment-asymp}
\Var(X_{t+h}-X_t)
=
L(t,t)^2 q_H h^{2H}+o(h^{2H}),
\qquad h\downarrow0,
\end{equation}
for the Gaussian Volterra process
\[
X_t:=\int_0^t K_H(t,s)\,\dd W_s.
\]
\end{corollary}

\begin{proof}
The increment decomposes as
\[
X_{t+h}-X_t
=
\int_0^t\bigl(K_H(t+h,s)-K_H(t,s)\bigr)\,\dd W_s
+
\int_t^{t+h}K_H(t+h,s)\,\dd W_s.
\]
The two integrals are independent centred Gaussian variables because they are supported on disjoint
Brownian time intervals.  It\^o isometry therefore gives the sum of the two variances, and
Propositions~\ref{prop:new} and~\ref{prop:old} give the constants \(a_H\) and \(b_H\).  Since
\(q_H=a_H+b_H\), \eqref{eq:increment-asymp} follows.
\end{proof}

\begin{lemma}[Off-diagonal pointwise increment bound]\label{lem:offdiag-pointwise}
Let $H\in(0,\tfrac12)$ and let $L$ satisfy the regularity assumptions of
Assumption~\ref{ass:model}. Define $K_H$ as in Proposition~\ref{prop:new}. Then for every fixed
$t\in(0,T)$ there exist constants $C<\infty$ and $h_0\in(0,t]$ such that for every
$0<h\le h_0$ and every $s\in[0,t-h]$,
\begin{equation}\label{eq:offdiag-pointwise}
|K_H(t+h,s)-K_H(t,s)|
\le
Ch\bigl((t-s)^{H-\frac32}+(t-s)^{H-\frac12}\bigr).
\end{equation}
\end{lemma}

\begin{proof}
Write $\alpha:=H-\tfrac12\in(-\tfrac12,0)$ and define
\[
F(r,s):=(r-s)^\alpha L(r,s),
\qquad 0\le s<r\le T.
\]
Fix $t\in(0,T)$ and choose
\[
h_0\in\left(0,\min\{t,T-t,1\}\right].
\]
By the regularity assumptions on $L$, there exists $C<\infty$ such that for all
$0\le s\le t$ and all $r\in[t,t+h_0]$ with $r>s$,
\[
|L(r,s)|+|\partial_1L(r,s)|\le C.
\]
For $s\in[0,t-h]$, the mean value theorem gives
\[
K_H(t+h,s)-K_H(t,s)
=
F(t+h,s)-F(t,s)
=
h\,\partial_1F(\xi,s)
\]
for some $\xi\in(t,t+h)$. Because $s\le t-h<\xi$, one has
\[
\partial_1F(\xi,s)
=
\alpha(\xi-s)^{\alpha-1}L(\xi,s)
+
(\xi-s)^\alpha \partial_1L(\xi,s),
\]
and therefore
\[
|\partial_1F(\xi,s)|
\le
C\bigl((\xi-s)^{H-\frac32}+(\xi-s)^{H-\frac12}\bigr).
\]
Since $H-\tfrac32<0$ and $H-\tfrac12<0$, while $\xi-s\ge t-s$, we may absorb the resulting
comparison into the constant and obtain
\[
|\partial_1F(\xi,s)|
\le
C\bigl((t-s)^{H-\frac32}+(t-s)^{H-\frac12}\bigr).
\]
Multiplying by $h$ proves \eqref{eq:offdiag-pointwise}.
\end{proof}

\begin{lemma}[$L^2$-continuity of Volterra channels]\label{lem:L2cont}
Under the assumptions of Proposition~\ref{prop:old}, the Gaussian Volterra process
\[
X_t:=\int_0^t K_H(t,s)\,\dd W_s,
\qquad 0\le t\le T,
\]
is continuous from $[0,T]$ into $L^2(\Omega)$.
\end{lemma}

\begin{proof}
For $t\in[0,T]$, define
\[
k_t(s):=K_H(t,s)\one_{\{0\le s\le t\}},
\qquad 0\le s\le T.
\]
Then $X_t=\int_0^T k_t(s)\,\dd W_s$, so by It\^o isometry
\[
\|X_u-X_t\|_{L^2(\Omega)}^2=\|k_u-k_t\|_{L^2([0,T])}^2.
\]
It therefore suffices to show that $t\mapsto k_t$ is continuous in $L^2([0,T])$.

If $t=0$, Proposition~\ref{prop:new} gives
\[
\|k_u\|_{L^2([0,T])}^2
=
\int_0^u K_H(u,s)^2\,\dd s
=
L(0,0)^2a_Hu^{2H}+o(u^{2H})
\longrightarrow 0
\]
as $u\downarrow0$.

Fix now $t\in(0,T]$ and let $u\to t$. Choose $\delta\in(0,t)$ so small that
$u\in[t-\delta,t+\delta]\cap[0,T]$ for all $u$ sufficiently close to $t$. Then
\begin{align*}
\|k_u-k_t\|_{L^2([0,T])}^2
&=
\int_0^{t-\delta}\!\!\bigl(K_H(u,s)-K_H(t,s)\bigr)^2\,\dd s \\
&\quad+
\int_{t-\delta}^{\max\{u,t\}}
\bigl(k_u(s)-k_t(s)\bigr)^2\,\dd s.
\end{align*}
On $[0,t-\delta]$, the map $(r,s)\mapsto K_H(r,s)$ is continuous and bounded on the compact
set $[t-\delta,t+\delta]\times[0,t-\delta]$, so dominated convergence yields
\[
\int_0^{t-\delta}\bigl(K_H(u,s)-K_H(t,s)\bigr)^2\,\dd s\longrightarrow0.
\]

For the near-diagonal part, there exists $C<\infty$ such that for all
$r\in[t-\delta,t+\delta]\cap[0,T]$ and all $s<r$,
\[
|K_H(r,s)|\le C(r-s)^{H-\frac12}.
\]
Hence
\begin{align*}
\int_{t-\delta}^{\max\{u,t\}}\bigl(k_u(s)-k_t(s)\bigr)^2\,\dd s
&\le
2\int_{t-\delta}^{\max\{u,t\}}k_u(s)^2\,\dd s
+2\int_{t-\delta}^{t}k_t(s)^2\,\dd s \\
&\le
C\int_{t-\delta}^{\max\{u,t\}}
\bigl((u-s)_+^{2H-1}+(t-s)_+^{2H-1}\bigr)\,\dd s \\
&\le C\delta^{2H}.
\end{align*}
Taking $\limsup_{u\to t}$ and then letting $\delta\downarrow0$ proves
$\|k_u-k_t\|_{L^2([0,T])}\to0$.
\end{proof}

\begin{theorem}[Oracle dynamic prediction threshold]\label{thm:prediction-threshold}
Assume the bivariate Volterra model of Assumption~\ref{ass:model}. Fix $t\in(0,T)$. Then, as
$h\downarrow0$,
\begin{align}
R_{XY}(t,h)
&=
\nu_Y^2 L_Y(t,t)^2 a_{H_Y} h^{2H_Y}
+ \eta^2 L_{XY}(t,t)^2 a_{H_{XY}} h^{2H_{XY}}
+ o\!\bigl(h^{2H_Y}+h^{2H_{XY}}\bigr), \label{eq:RXY-asymp}\\
G_{X\to Y}(t,h)
&=
\eta^2 L_{XY}(t,t)^2 b_{H_{XY}} h^{2H_{XY}}
+ o(h^{2H_{XY}}), \label{eq:G-asymp}\\
R_Y(t,h)
&=
\nu_Y^2 L_Y(t,t)^2 a_{H_Y} h^{2H_Y}
+ \eta^2 L_{XY}(t,t)^2 q_{H_{XY}} h^{2H_{XY}}
+ o\!\bigl(h^{2H_Y}+h^{2H_{XY}}\bigr). \label{eq:RY-asymp}
\end{align}
If \(\eta>0\), then
\begin{equation}\label{eq:gain-ratio}
\frac{G_{X\to Y}(t,h)}{R_Y(t,h)}
\longrightarrow
\begin{cases}
\dfrac{b_{H_{XY}}}{q_{H_{XY}}}, & H_{XY}<H_Y,\\[12pt]
\dfrac{\eta^2L_{XY}(t,t)^2 b_{H_Y}}
{\nu_Y^2L_Y(t,t)^2 a_{H_Y}+\eta^2L_{XY}(t,t)^2 q_{H_Y}},
& H_{XY}=H_Y,\\[12pt]
0, & H_{XY}>H_Y.
\end{cases}
\end{equation}
If \(\eta=0\), then \(G_{X\to Y}(t,h)\equiv0\). Thus, for a nonzero contagion channel, the
dynamic prediction threshold is
\[
H_{XY}=H_Y.
\]
Below it the contagion channel improves prediction at leading order, at it the gain is of leading
order but mixed with the idiosyncratic term, and above it the gain is asymptotically negligible at
leading order.
\end{theorem}

\begin{proof}
The exact formulas~\eqref{eq:RXY-exact}--\eqref{eq:G-exact}, together with
Propositions~\ref{prop:new} and~\ref{prop:old}, give the asymptotic expansions
\eqref{eq:RXY-asymp}--\eqref{eq:RY-asymp}. If \(\eta=0\), equation~\eqref{eq:G-exact}
gives \(G_{X\to Y}(t,h)\equiv0\), so the null-channel statement follows. Hence assume
\(\eta>0\) for the gain-ratio limits.

If $H_{XY}<H_Y$, then $h^{2H_{XY}}$ decays more slowly than $h^{2H_Y}$, so the
contagion contribution is the leading term in $R_Y(t,h)$ and
\[
R_Y(t,h)
=
\eta^2 L_{XY}(t,t)^2 q_{H_{XY}} h^{2H_{XY}}(1+o(1)).
\]
Combining this with~\eqref{eq:G-asymp} yields the first limit in~\eqref{eq:gain-ratio}.

If $H_{XY}=H_Y$, both terms are of the same order and the second limit in
\eqref{eq:gain-ratio} follows by dividing numerator and denominator by
$h^{2H_Y}$ and retaining the two coefficients.

If $H_{XY}>H_Y$, then
\[
R_Y(t,h)
=
\nu_Y^2 L_Y(t,t)^2 a_{H_Y} h^{2H_Y}(1+o(1)),
\]
while $G_{X\to Y}(t,h)=O(h^{2H_{XY}})$. Therefore
\[
\frac{G_{X\to Y}(t,h)}{R_Y(t,h)}
=
O(h^{2(H_{XY}-H_Y)})
\longrightarrow 0.
\]
This proves~\eqref{eq:gain-ratio}.
\end{proof}

\begin{remark}[Interpretation of the oracle prediction block]
Theorem~\ref{thm:prediction-threshold} is the time-dynamic analogue of the pricing
threshold of \cite{vidal2026latent} at the \emph{oracle latent-driver level}.  It proves that,
at short horizons, the predictive value of the contagion channel is governed by the same exponent
comparison $H_{XY}$ versus $H_Y$.  Rougher contagion is dynamically visible at leading order,
smoother contagion is dynamically latent at leading order, and the boundary case places both
channels on the same scale.
\end{remark}

\begin{remark}[Size of the leading-order gain]
Since $a_H>0$ and $b_H>0$, one has
\[
0<\frac{b_H}{q_H}<1,
\qquad q_H=a_H+b_H.
\]
Thus in the rougher-contagion regime $H_{XY}<H_Y$, the enlarged oracle filtration removes
a nontrivial but not total fraction of the leading-order prediction risk: the future
innovation on $(t,t+h]$ remains unpredictable even after the past contagion driver is
resolved.
\end{remark}

\section{Observed-channel Gaussian reduction}\label{sec:observed}

Section~\ref{sec:prediction} resolves the latent contagion driver directly.  The observed-channel
reduction changes the information set.  The latent source driver is replaced by an observed
Gaussian source channel, and the problem becomes the posterior-variance reduction induced by
observing the target past alone.

\begin{assumption}[Observed-channel identifiability]
\label{ass:obs-source}
Let $H_X\in(0,\tfrac12)$ and let $L_X$ satisfy the same regularity conditions as in
Assumption~\ref{ass:model}. Define
\[
K_X^{\mathrm{obs}}(t,s):=(t-s)^{H_X-\frac12}L_X(t,s)\one_{\{0\le s<t\le T\}}
\]
and the observed source channel
\begin{equation}\label{eq:Xobs}
X_t^{\mathrm{obs}}
:=
\int_0^t K_X^{\mathrm{obs}}(t,s)\,\dd W_s^X,
\qquad 0\le t\le T.
\end{equation}
For each $t\in(0,T]$, let
\[
\mathcal K_{X,t}^{\mathrm{obs}}:L^2([0,t])\to L^2([0,t])
\]
be the associated Volterra operator
\[
(\mathcal K_{X,t}^{\mathrm{obs}}f)(r)
:=
\int_0^r K_X^{\mathrm{obs}}(r,s)f(s)\,\dd s,
\qquad 0\le r\le t.
\]
Also, for each $t\in(0,T]$, let
\[
\mathcal K_{Y,t}:L^2([0,t])\to L^2([0,t]),
\qquad
(\mathcal K_{Y,t}f)(r):=\int_0^r K_Y(r,s)f(s)\,\dd s.
\]
Assume that both $\mathcal K_{X,t}^{\mathrm{obs}}$ and $\mathcal K_{Y,t}$ are injective
for every $t\in(0,T]$.
\end{assumption}

\begin{remark}[Interpretation of the identifiability hypotheses]
Injectivity of $\mathcal K_{X,t}^{\mathrm{obs}}$ means that the observed source channel
does not lose first-chaos information about the source driver on $[0,t]$, while
injectivity of $\mathcal K_{Y,t}$ means that once the contagion component is removed from
the target, the remaining idiosyncratic channel still determines the past Brownian driver
$W^\perp$ on $[0,t]$. These are exactly the hypotheses needed to identify the enlarged
observed filtration with the enlarged oracle filtration.
\end{remark}

\begin{lemma}[Closedness of measurable $L^2$ subspaces]\label{lem:L2meas}
Let $\mathcal F$ be a sigma-field on $\Omega$. Then
\[
L^2(\mathcal F)
:=
\left\{
X\in L^2(\Omega): X \text{ admits an }\mathcal F\text{-measurable version}
\right\}
\]
is a closed linear subspace of $L^2(\Omega)$.
\end{lemma}

\begin{proof}
Linearity follows by taking measurable versions and forming finite linear combinations.  To prove
closedness, let \(X_n\in L^2(\mathcal F)\) and suppose that
\[
X_n\to X
\qquad\text{in }L^2(\Omega).
\]
Choose \(\mathcal F\)-measurable versions \(\widetilde X_n\) of \(X_n\).  The same convergence holds
for \(\widetilde X_n\), so a subsequence \(\widetilde X_{n_k}\) converges to \(X\) almost surely.
The almost-sure limit of \(\mathcal F\)-measurable random variables is \(\mathcal F\)-measurable
after redefining on a null set.  Thus \(X\) admits an \(\mathcal F\)-measurable version, and
\(X\in L^2(\mathcal F)\).
\end{proof}

\begin{lemma}[Gaussian first-chaos criterion]\label{lem:first-chaos}
Let $\{G_\alpha\}_{\alpha\in A}$ and $\{H_\beta\}_{\beta\in B}$ be centred Gaussian
families on the same probability space, and let
\[
\mathfrak H_G:=\overline{\mathrm{span}}\{G_\alpha:\alpha\in A\},
\qquad
\mathfrak H_H:=\overline{\mathrm{span}}\{H_\beta:\beta\in B\}
\]
be the corresponding closed subspaces of $L^2(\Omega)$. If $\mathfrak H_G=\mathfrak H_H$,
then
\[
\sigma(G_\alpha:\alpha\in A)=\sigma(H_\beta:\beta\in B)
\qquad\text{mod }\PP.
\]
\end{lemma}

\begin{proof}
Every finite linear combination of the \(G_\alpha\) is measurable with respect to
\(\sigma(G_\alpha:\alpha\in A)\).  Lemma~\ref{lem:L2meas} therefore implies that the closed span
\(\mathfrak H_G\) is contained in
\[
L^2\bigl(\sigma(G_\alpha:\alpha\in A)\bigr),
\]
and hence
\[
\sigma(\mathfrak H_G)\subseteq \sigma(G_\alpha:\alpha\in A).
\]
Conversely, each \(G_\alpha\) belongs to \(\mathfrak H_G\), so
\[
\sigma(G_\alpha:\alpha\in A)\subseteq \sigma(\mathfrak H_G).
\]
Thus \(\sigma(G_\alpha:\alpha\in A)=\sigma(\mathfrak H_G)\) modulo null sets.  The same
argument applies to the \(H_\beta\)-family.  If \(\mathfrak H_G=\mathfrak H_H\), the generated
sigma-fields coincide modulo \(\PP\).
\end{proof}

\begin{lemma}[$L^2$-continuous families and integral spans]\label{lem:integral-span}
Let $t>0$ and let $G:[0,t]\to L^2(\Omega)$ be continuous. Then
\[
\overline{\mathrm{span}}^{L^2(\Omega)}\{G(r):0\le r\le t\}
=
\overline{\left\{
\int_0^t f(r)G(r)\,\dd r : f\in L^2([0,t])
\right\}}^{L^2(\Omega)}.
\]
\end{lemma}

\begin{proof}
Every Bochner integral $\int_0^t f(r)G(r)\,\dd r$ belongs to the closed linear span of
$\{G(r):0\le r\le t\}$, so only the reverse inclusion needs proof.

Fix $r_0\in[0,t]$. Choose a sequence of intervals $I_n\subset[0,t]$ such that
$r_0\in I_n$, $|I_n|>0$, and $|I_n|\downarrow0$. Let
\[
f_n(r):=\frac{\one_{I_n}(r)}{|I_n|}.
\]
Then $f_n\in L^2([0,t])$ and
\[
\int_0^t f_n(r)G(r)\,\dd r
=
\frac{1}{|I_n|}\int_{I_n} G(r)\,\dd r.
\]
By Jensen's inequality and $L^2$-continuity of $G$ at $r_0$,
\[
\left\|
\frac{1}{|I_n|}\int_{I_n} G(r)\,\dd r-G(r_0)
\right\|_{L^2(\Omega)}^2
\le
\frac{1}{|I_n|}\int_{I_n}\|G(r)-G(r_0)\|_{L^2(\Omega)}^2\,\dd r
\longrightarrow 0.
\]
Hence each $G(r_0)$ lies in the closed linear span of the integral family, proving the
claim.
\end{proof}

\begin{proposition}[Single-driver Volterra filtration transfer]
\label{prop:filtration-transfer}
Fix $t\in(0,T]$, let $B$ be a standard Brownian motion, and let
\[
\mathcal K_t:L^2([0,t])\to L^2([0,t])
\]
be a Volterra operator of the form
\[
(\mathcal K_tf)(r):=\int_0^r K(r,s)f(s)\,\dd s,
\qquad 0\le r\le t.
\]
Assume that $\mathcal K_t$ is injective, and define
\[
U_r:=\int_0^r K(r,s)\,\dd B_s,
\qquad 0\le r\le t.
\]
Assume in addition that $r\mapsto U_r$ is continuous from $[0,t]$ into $L^2(\Omega)$.
Then
\begin{equation}\label{eq:filtration-transfer}
\sigma(U_r:0\le r\le t)
=
\sigma(B_r:0\le r\le t)
\qquad\text{mod }\PP.
\end{equation}
\end{proposition}

\begin{proof}
Let
\[
I_t^B(f):=\int_0^t f(s)\,\dd B_s,
\qquad f\in L^2([0,t]).
\]
By It\^o isometry, \(I_t^B\) is an isometric isomorphism from \(L^2([0,t])\) onto the first chaos
generated by \(\{B_r:0\le r\le t\}\).

For \(f\in L^2([0,t])\), boundedness of \(\mathcal K_t\) gives
\(\mathcal K_t^*f\in L^2([0,t])\).  The square-integrable stochastic Fubini theorem therefore
applies and gives
\begin{align*}
\int_0^t f(r)U_r\,\dd r
&=
\int_0^t f(r)\int_0^r K(r,s)\,\dd B_s\,\dd r \\
&=
\int_0^t \left(\int_s^t f(r)K(r,s)\,\dd r\right)\dd B_s \\
&=
I_t^B\bigl(\mathcal K_t^*f\bigr).
\end{align*}
Lemma~\ref{lem:integral-span} therefore gives
\[
\overline{\mathrm{span}}^{L^2(\Omega)}\{U_r:0\le r\le t\}
=
I_t^B\bigl(\overline{\mathrm{Ran}(\mathcal K_t^*)}\bigr).
\]
Because \(\mathcal K_t\) is injective,
\[
\overline{\mathrm{Ran}(\mathcal K_t^*)}
=
(\ker \mathcal K_t)^\perp
=
L^2([0,t]).
\]
Thus the first chaos generated by \(\{U_r:0\le r\le t\}\) equals the first chaos generated by
\(\{B_r:0\le r\le t\}\).  Lemma~\ref{lem:first-chaos} yields \eqref{eq:filtration-transfer}.
\end{proof}

\begin{definition}[Observed-channel filtrations]\label{def:observed-filtrations}
For $0\le t\le T$, define
\[
\mathcal F_t^Y:=\sigma(\mathcal Y_r:0\le r\le t),
\qquad
\mathcal F_t^{X,Y}:=\sigma(X_r^{\mathrm{obs}},\mathcal Y_r:0\le r\le t).
\]
For $0\le t<t+h\le T$, define the observed-channel prediction risks
\begin{align}
\widetilde R_Y(t,h)
&:=
\EE\!\left[
\Var(\Delta_h\mathcal Y_t\mid \mathcal F_t^Y)
\right], \label{eq:RtildeY}\\
\widetilde R_{XY}(t,h)
&:=
\EE\!\left[
\Var(\Delta_h\mathcal Y_t\mid \mathcal F_t^{X,Y})
\right], \label{eq:RtildeXY}
\end{align}
and the associated gain
\begin{equation}\label{eq:Gtilde}
\widetilde G_{X\to Y}(t,h):=
\widetilde R_Y(t,h)-\widetilde R_{XY}(t,h).
\end{equation}
\end{definition}

\begin{definition}[Past-state operators]\label{def:operators}
Fix $t\in(0,T]$ and write $H_t:=L^2([0,t])$ and $H_t^\oplus:=H_t\oplus H_t$.
Define the truncated Volterra operators
\[
(\mathcal K_{Y,t}f)(r):=\int_0^r K_Y(r,s)f(s)\,\dd s,
\qquad
(\mathcal K_{XY,t}g)(r):=\int_0^r K_{XY}(r,s)g(s)\,\dd s
\]
for $0\le r\le t$, and the observation operator
\begin{equation}\label{eq:Ot}
\mathcal O_t:H_t^\oplus\to H_t,
\qquad
\mathcal O_t(f,g):=\nu_Y\mathcal K_{Y,t}f+\eta\mathcal K_{XY,t}g.
\end{equation}
Let
\[
I_t^\oplus(f,g):=\int_0^t f(s)\,\dd W_s^\perp+\int_0^t g(s)\,\dd W_s^X
\]
be the isonormal map on $H_t^\oplus$.

For $0\le t<t+h\le T$, set
\begin{align}
g_{Y,t,h}(s)&:=K_Y(t+h,s)-K_Y(t,s), \label{eq:gY}\\
g_{XY,t,h}(s)&:=K_{XY}(t+h,s)-K_{XY}(t,s), \label{eq:gXY}
\end{align}
and define the increment test vector
\begin{equation}\label{eq:gamma}
\gamma_{t,h}:=\bigl(\nu_Y g_{Y,t,h},\, \eta g_{XY,t,h}\bigr)\in H_t^\oplus.
\end{equation}
Let $\Pi_t$ denote the orthogonal projection in $H_t^\oplus$ onto
\[
\overline{\mathrm{Ran}(\mathcal O_t^*)},
\]
and let
\begin{equation}\label{eq:Sigma}
\Sigma_t:=I_{H_t^\oplus}-\Pi_t.
\end{equation}
\end{definition}

\begin{remark}[Interpretation of $\Sigma_t$]
The operator $\Sigma_t$ is the posterior covariance operator of the hidden two-component past state
\((W^\perp_{[0,t]},W^X_{[0,t]})\) after observing the target past filtration $\mathcal F_t^Y$.
The exact observed-channel gain is the quadratic form of $\Sigma_t$ against the past-resolved
increment vector $\gamma_{t,h}$.
\end{remark}

\begin{proposition}[Exact observed-channel reduction]
\label{prop:observed-reduction}
Assume Assumptions~\ref{ass:model} and~\ref{ass:obs-source}. Fix $0<t<t+h\le T$ and
define the past-resolved term
\begin{equation}\label{eq:past-term}
P_{t,h}:=I_t^\oplus(\gamma_{t,h})
\end{equation}
and the future innovation
\begin{equation}\label{eq:future-term}
M_{t,h}
:=
\nu_Y\int_t^{t+h}K_Y(t+h,s)\,\dd W_s^\perp
+ \eta\int_t^{t+h}K_{XY}(t+h,s)\,\dd W_s^X.
\end{equation}
Then:
\begin{enumerate}[label=(\roman*)]
\item
\begin{equation}\label{eq:decomp}
\Delta_h\mathcal Y_t=P_{t,h}+M_{t,h},
\end{equation}
with $M_{t,h}$ independent of $\mathcal F_t^Y$ and of $\mathcal F_t^{X,Y}$;
\item
\begin{equation}\label{eq:RtildeXY-exact}
\widetilde R_{XY}(t,h)
=
\nu_Y^2 \int_t^{t+h} K_Y(t+h,s)^2\,\dd s
+ \eta^2 \int_t^{t+h} K_{XY}(t+h,s)^2\,\dd s;
\end{equation}
\item the target-past first chaos is
\begin{equation}\label{eq:first-chaos-Y}
\overline{\mathrm{span}}^{L^2(\Omega)}\{\mathcal Y_r:0\le r\le t\}
=
I_t^\oplus\bigl(\overline{\mathrm{Ran}(\mathcal O_t^*)}\bigr);
\end{equation}
\item the observed-channel prediction gain is exactly
\begin{equation}\label{eq:Gtilde-exact}
\widetilde G_{X\to Y}(t,h)
=
\EE\!\left[\Var(P_{t,h}\mid \mathcal F_t^Y)\right]
=
\|\Sigma_t \gamma_{t,h}\|_{H_t^\oplus}^2.
\end{equation}
\end{enumerate}
\end{proposition}

\begin{proof}
By the definitions of $g_{Y,t,h}$ and $g_{XY,t,h}$,
\[
P_{t,h}
=
\nu_Y\int_0^t \bigl(K_Y(t+h,s)-K_Y(t,s)\bigr)\,\dd W_s^\perp
+ \eta\int_0^t \bigl(K_{XY}(t+h,s)-K_{XY}(t,s)\bigr)\,\dd W_s^X,
\]
so \eqref{eq:decomp} follows by splitting the increment $\Delta_h\mathcal Y_t$ into its
past-resolved and future-innovation parts. The future innovation $M_{t,h}$ depends only on
the Brownian increments on $(t,t+h]$, hence is independent of the sigma-field generated by
\[
\{W_s^\perp,W_s^X:0\le s\le t\},
\]
and therefore of both $\mathcal F_t^Y$ and $\mathcal F_t^{X,Y}$.

To prove \eqref{eq:RtildeXY-exact}, first note from Lemma~\ref{lem:L2cont} applied to the
observed source channel that $r\mapsto X_r^{\mathrm{obs}}$ is continuous from $[0,t]$ into
$L^2(\Omega)$. Proposition~\ref{prop:filtration-transfer} applied to that
channel therefore yields
\[
\mathcal F_t^{X,Y}
=
\sigma(W_r^X,\mathcal Y_r:0\le r\le t)
\qquad\text{mod }\PP.
\]
Hence the past contagion process
\[
Z_r:=\int_0^r K_{XY}(r,s)\,\dd W_s^X,
\qquad 0\le r\le t,
\]
is $\mathcal F_t^{X,Y}$-measurable. Since
\[
\mathcal Y_r=\nu_Y V_r+\eta Z_r,
\qquad
V_r:=\int_0^r K_Y(r,s)\,\dd W_s^\perp,
\]
it follows that the full idiosyncratic past process $\{V_r:0\le r\le t\}$ is also
$\mathcal F_t^{X,Y}$-measurable. By Lemma~\ref{lem:L2cont}, the map $r\mapsto V_r$ is
continuous from $[0,t]$ into $L^2(\Omega)$. Applying
Proposition~\ref{prop:filtration-transfer} to the channel $\{V_r:0\le r\le t\}$ and the
Brownian motion $W^\perp$, and using the injectivity of $\mathcal K_{Y,t}$ from
Assumption~\ref{ass:obs-source}, we obtain
\[
\sigma(V_r:0\le r\le t)
=
\sigma(W_r^\perp:0\le r\le t)
\qquad\text{mod }\PP.
\]
Hence
\[
\mathcal F_t^{X,Y}
=
\sigma(W_r^\perp,W_r^X:0\le r\le t)
=
\mathcal G_t^{X,Y}
\qquad\text{mod }\PP.
\]
Therefore $P_{t,h}$ is $\mathcal F_t^{X,Y}$-measurable. Using \eqref{eq:decomp} and
independence of $M_{t,h}$,
\[
\Var(\Delta_h\mathcal Y_t\mid \mathcal F_t^{X,Y})=\Var(M_{t,h}),
\]
almost surely. It\^o isometry gives \eqref{eq:RtildeXY-exact}.

For \eqref{eq:first-chaos-Y}, first note from Lemma~\ref{lem:L2cont} applied to the
two Volterra components of $\mathcal Y$ that $r\mapsto \mathcal Y_r$ is continuous from
$[0,t]$ into $L^2(\Omega)$. For $f\in H_t$,
\begin{align*}
\int_0^t f(r)\mathcal Y_r\,\dd r
&=
\nu_Y\int_0^t f(r)\int_0^r K_Y(r,s)\,\dd W_s^\perp\,\dd r
+ \eta\int_0^t f(r)\int_0^r K_{XY}(r,s)\,\dd W_s^X\,\dd r \\
&=
I_t^\oplus\bigl(\mathcal O_t^* f\bigr),
\end{align*}
where
\[
\mathcal O_t^*f=\bigl(\nu_Y\mathcal K_{Y,t}^*f,\eta\mathcal K_{XY,t}^*f\bigr).
\]
Lemma~\ref{lem:integral-span} therefore implies that the closed first chaos generated by
the target past equals
$I_t^\oplus(\overline{\mathrm{Ran}(\mathcal O_t^*)})$.

Finally, because $M_{t,h}$ is independent of $\mathcal F_t^Y$,
\[
\Var(\Delta_h\mathcal Y_t\mid \mathcal F_t^Y)
=
\Var(P_{t,h}\mid \mathcal F_t^Y)+\Var(M_{t,h}).
\]
Taking expectations and subtracting \eqref{eq:RtildeXY-exact} gives
\[
\widetilde G_{X\to Y}(t,h)
=
\EE\!\left[\Var(P_{t,h}\mid \mathcal F_t^Y)\right].
\]
Now $P_{t,h}=I_t^\oplus(\gamma_{t,h})$ lies in the hidden first chaos $I_t^\oplus(H_t^\oplus)$,
while the target-past first chaos is $I_t^\oplus(\overline{\mathrm{Ran}(\mathcal O_t^*)})$.
By Lemma~\ref{lem:first-chaos},
\[
\mathcal F_t^Y
=
\sigma\bigl(I_t^\oplus(h):h\in \overline{\mathrm{Ran}(\mathcal O_t^*)}\bigr)
\qquad\text{mod }\PP.
\]
All variables are centred and jointly Gaussian.  Conditioning on this sigma-field is therefore the
orthogonal projection onto its closed first chaos, and
\[
\EE[P_{t,h}\mid \mathcal F_t^Y]
=
I_t^\oplus(\Pi_t\gamma_{t,h}),
\]
and hence
\[
\EE\!\left[\Var(P_{t,h}\mid \mathcal F_t^Y)\right]
=
\| (I_{H_t^\oplus}-\Pi_t)\gamma_{t,h}\|_{H_t^\oplus}^2
=
\|\Sigma_t\gamma_{t,h}\|_{H_t^\oplus}^2.
\]
This proves \eqref{eq:Gtilde-exact}.
\end{proof}

\section{Deterministic geometry of the observed gain}\label{sec:geometry}

Section~\ref{sec:observed} reduces observable prediction to a deterministic distance in the
hidden past-state space.  Variational and nullspace formulas describe the same distance.  The
obstruction is the failure of the target past to reproduce the enlarged oracle gain.

\begin{definition}[Normal operator and score vector]\label{def:normal}
Fix $t\in(0,T]$. Define the positive self-adjoint operator
\begin{equation}\label{eq:Ct}
\mathcal C_t:=\mathcal O_t\mathcal O_t^*
=
\nu_Y^2\mathcal K_{Y,t}\mathcal K_{Y,t}^*
+ \eta^2\mathcal K_{XY,t}\mathcal K_{XY,t}^*
\end{equation}
on $H_t=L^2([0,t])$. For $0\le t<t+h\le T$, define the score vector
\begin{equation}\label{eq:rt}
r_{t,h}:=\mathcal O_t\gamma_{t,h}
=
\nu_Y^2\mathcal K_{Y,t}g_{Y,t,h}
+ \eta^2\mathcal K_{XY,t}g_{XY,t,h}.
\end{equation}
\end{definition}

\begin{proposition}[Variational representation of the observed gain]
\label{prop:variational}
Assume Assumptions~\ref{ass:model} and~\ref{ass:obs-source}. Fix $0<t<t+h\le T$. Then
\begin{equation}\label{eq:variational}
\widetilde G_{X\to Y}(t,h)
=
\inf_{f\in H_t}
\left\{
\nu_Y^2\|g_{Y,t,h}-\mathcal K_{Y,t}^*f\|_{H_t}^2
+
\eta^2\|g_{XY,t,h}-\mathcal K_{XY,t}^*f\|_{H_t}^2
\right\}.
\end{equation}
Equivalently,
\begin{equation}\label{eq:dual-variational}
\widetilde G_{X\to Y}(t,h)
=
\|\gamma_{t,h}\|_{H_t^\oplus}^2
- \sup_{f\in H_t}
\left\{
2\langle r_{t,h},f\rangle_{H_t}
- \langle \mathcal C_t f,f\rangle_{H_t}
\right\}.
\end{equation}
\end{proposition}

\begin{proof}
By Proposition~\ref{prop:observed-reduction},
\[
\widetilde G_{X\to Y}(t,h)=\|\Sigma_t\gamma_{t,h}\|_{H_t^\oplus}^2.
\]
Since $\Sigma_t$ is the orthogonal projection onto
$\overline{\mathrm{Ran}(\mathcal O_t^*)}^\perp$, this equals the squared distance from
$\gamma_{t,h}$ to $\overline{\mathrm{Ran}(\mathcal O_t^*)}$. Therefore
\[
\widetilde G_{X\to Y}(t,h)
=
\inf_{f\in H_t}\|\gamma_{t,h}-\mathcal O_t^*f\|_{H_t^\oplus}^2.
\]
Using \eqref{eq:Ot} and \eqref{eq:gamma}, we obtain
\[
\gamma_{t,h}-\mathcal O_t^*f
=
\bigl(\nu_Y(g_{Y,t,h}-\mathcal K_{Y,t}^*f),\,
\eta(g_{XY,t,h}-\mathcal K_{XY,t}^*f)\bigr),
\]
which gives \eqref{eq:variational}.

Expanding the square,
\begin{align*}
\|\gamma_{t,h}-\mathcal O_t^*f\|_{H_t^\oplus}^2
&=
\|\gamma_{t,h}\|_{H_t^\oplus}^2
-2\langle \gamma_{t,h},\mathcal O_t^*f\rangle_{H_t^\oplus}
+\|\mathcal O_t^*f\|_{H_t^\oplus}^2 \\
&=
\|\gamma_{t,h}\|_{H_t^\oplus}^2
-2\langle \mathcal O_t\gamma_{t,h},f\rangle_{H_t}
+\langle \mathcal O_t\mathcal O_t^*f,f\rangle_{H_t} \\
&=
\|\gamma_{t,h}\|_{H_t^\oplus}^2
-2\langle r_{t,h},f\rangle_{H_t}
+\langle \mathcal C_t f,f\rangle_{H_t}.
\end{align*}
Taking the infimum over $f$ yields \eqref{eq:dual-variational}.
\end{proof}

\begin{proposition}[Nullspace representation of the observed gain]
\label{prop:nullspace}
Assume Assumptions~\ref{ass:model} and~\ref{ass:obs-source}. Fix $0<t<t+h\le T$. Then
\begin{equation}\label{eq:kernelOt}
\ker \mathcal O_t
=
\left\{
(u,v)\in H_t^\oplus:
\nu_Y\mathcal K_{Y,t}u+\eta\mathcal K_{XY,t}v=0
\right\}.
\end{equation}
Moreover, with the convention that the supremum is zero when $\ker\mathcal O_t=\{0\}$,
\begin{equation}\label{eq:nullspace-rep}
\widetilde G_{X\to Y}(t,h)
=
\sup\left\{
\bigl|\langle \gamma_{t,h},w\rangle_{H_t^\oplus}\bigr|^2:
w\in\ker \mathcal O_t,\ \|w\|_{H_t^\oplus}=1
\right\}.
\end{equation}
\end{proposition}

\begin{proof}
Identity \eqref{eq:kernelOt} is the kernel condition for $\mathcal O_t$.
Because
\[
(\overline{\mathrm{Ran}(\mathcal O_t^*)})^\perp=\ker \mathcal O_t,
\]
the orthogonal projection $\Sigma_t$ is the projection onto $\ker \mathcal O_t$. Hence
\[
\widetilde G_{X\to Y}(t,h)=\|\Sigma_t\gamma_{t,h}\|_{H_t^\oplus}^2.
\]
If $\ker\mathcal O_t=\{0\}$, then $\Sigma_t=0$ and the convention gives
\eqref{eq:nullspace-rep}. Otherwise, for any orthogonal projection $P$ on a Hilbert space and any
vector $x$,
\[
\|Px\|^2
=
\sup\left\{
|\langle x,w\rangle|^2 : w\in \mathrm{Ran}(P),\ \|w\|=1
\right\}.
\]
Applying this with $P=\Sigma_t$ and $\mathrm{Ran}(\Sigma_t)=\ker \mathcal O_t$ gives
\eqref{eq:nullspace-rep}.
\end{proof}

\begin{corollary}[Zero-gain criterion]\label{cor:zero-gain}
Assume Assumptions~\ref{ass:model} and~\ref{ass:obs-source}. Fix $0<t<t+h\le T$. Then the
following are equivalent:
\begin{enumerate}[label=(\roman*)]
\item $\widetilde G_{X\to Y}(t,h)=0$;
\item $\gamma_{t,h}\in\overline{\mathrm{Ran}(\mathcal O_t^*)}$;
\item there exists a sequence $(f_n)$ in $H_t$ such that
\[
\mathcal K_{Y,t}^*f_n\to g_{Y,t,h},
\qquad
\eta\,\mathcal K_{XY,t}^*f_n\to \eta\,g_{XY,t,h}
\]
in $H_t$.
\end{enumerate}
\end{corollary}

\begin{proof}
By Proposition~\ref{prop:observed-reduction},
\[
\widetilde G_{X\to Y}(t,h)=\|\Sigma_t\gamma_{t,h}\|_{H_t^\oplus}^2.
\]
Hence $\widetilde G_{X\to Y}(t,h)=0$ if and only if $\Sigma_t\gamma_{t,h}=0$, which is
equivalent to $\gamma_{t,h}\in\overline{\mathrm{Ran}(\mathcal O_t^*)}$. This proves the
equivalence of (i) and (ii).

Condition (ii) is equivalent to the existence of $f_n\in H_t$ such that
\[
\mathcal O_t^*f_n\to \gamma_{t,h}
\qquad\text{in }H_t^\oplus.
\]
Using the explicit form of $\mathcal O_t^*$ and $\gamma_{t,h}$ gives condition~(iii).
\end{proof}

\begin{remark}[Interpretation of the geometric reduction]
Section~\ref{sec:observed} gives the probabilistic reduction.  Section~\ref{sec:geometry} gives
the deterministic geometry.  The observable gain is the distance of the increment vector
$\gamma_{t,h}$ from the subspace $\overline{\mathrm{Ran}(\mathcal O_t^*)}$, or equivalently the
energy of its component along the hidden null directions of the target-past observation operator.
Asymptotic screening therefore depends on whether the near-diagonal increment vector approaches
that subspace or retains a nontrivial projection on $\ker \mathcal O_t$ as $h\downarrow0$.
\end{remark}

\section{Transfer-operator reduction}\label{sec:transfer}

The observation nullspace first lives on the two-component hidden state space $H_t^\oplus$.
Injectivity of the idiosyncratic Volterra operator removes one component and leaves the contagion
component alone.  The hidden cancellation geometry is therefore encoded by the graph of a single
closed operator.

\begin{definition}[Hidden transfer operator]\label{def:transfer}
Assume $\eta>0$ and fix $t\in(0,T]$. Since $\mathcal K_{Y,t}$ is injective, its inverse
\[
\mathcal K_{Y,t}^{-1}:\mathrm{Ran}(\mathcal K_{Y,t})\to H_t
\]
is well defined. Define the domain
\begin{equation}\label{eq:domainS}
\mathcal D(\mathcal S_t)
:=
\left\{
v\in H_t:\mathcal K_{XY,t}v\in \mathrm{Ran}(\mathcal K_{Y,t})
\right\},
\end{equation}
and the transfer operator
\begin{equation}\label{eq:St}
\mathcal S_t:\mathcal D(\mathcal S_t)\subset H_t\to H_t,
\qquad
\mathcal S_t v:=\mathcal K_{Y,t}^{-1}\mathcal K_{XY,t}v.
\end{equation}
\noindent
This is the operator-theoretic form of the observable screening problem.  The Volterra-equation
background comes from \cite{GLS90}; the compactness and graph arguments below are carried in
trace-ideal terms \cite{Sim05}.
\end{definition}

\begin{proposition}[Closedness and graph form of the hidden nullspace]
\label{prop:transfer-graph}
Assume Assumptions~\ref{ass:model} and~\ref{ass:obs-source}, and let $\eta>0$. Fix
$t\in(0,T]$. Then $\mathcal S_t$ is a linear closed operator on $H_t$, and
\begin{equation}\label{eq:kernel-graph}
\ker \mathcal O_t
=
\left\{
\left(-\frac{\eta}{\nu_Y}\mathcal S_t v,\; v\right): v\in\mathcal D(\mathcal S_t)
\right\}.
\end{equation}
\end{proposition}

\begin{proof}
The domain \(\mathcal D(\mathcal S_t)\) is linear because \(\mathrm{Ran}(\mathcal K_{Y,t})\) is
linear and both Volterra operators are linear.  On this domain, \(\mathcal S_t\) is linear by the
linearity of \(\mathcal K_{XY,t}\) and of \(\mathcal K_{Y,t}^{-1}\) on
\(\mathrm{Ran}(\mathcal K_{Y,t})\).

To prove closedness, let \(v_n\in\mathcal D(\mathcal S_t)\) satisfy
\[
v_n\to v,
\qquad
\mathcal S_t v_n\to u
\]
in \(H_t\).  Since \(\mathcal K_{XY,t}\) and \(\mathcal K_{Y,t}\) are bounded,
\[
\mathcal K_{XY,t}v_n\to \mathcal K_{XY,t}v,
\qquad
\mathcal K_{Y,t}\mathcal S_t v_n\to \mathcal K_{Y,t}u
\]
in \(H_t\).  The identity \(\mathcal K_{Y,t}\mathcal S_t v_n=\mathcal K_{XY,t}v_n\) passes to the
limit, so
\[
\mathcal K_{XY,t}v=\mathcal K_{Y,t}u.
\]
Hence \(\mathcal K_{XY,t}v\in\mathrm{Ran}(\mathcal K_{Y,t})\), so
\(v\in\mathcal D(\mathcal S_t)\), and injectivity of \(\mathcal K_{Y,t}\) gives
\(\mathcal S_t v=u\).  Thus \(\mathcal S_t\) is closed.

Now let \((u,v)\in H_t^\oplus\).  By definition of \(\mathcal O_t\),
\[
(u,v)\in\ker \mathcal O_t
\iff
\nu_Y\mathcal K_{Y,t}u+\eta\mathcal K_{XY,t}v=0.
\]
Since \(\eta>0\), this is equivalent to
\[
\mathcal K_{XY,t}v=-\frac{\nu_Y}{\eta}\mathcal K_{Y,t}u.
\]
Therefore \(v\in\mathcal D(\mathcal S_t)\) and
\[
\mathcal S_t v
=
\mathcal K_{Y,t}^{-1}\mathcal K_{XY,t}v
=
-\frac{\nu_Y}{\eta}u,
\]
so \(u=-(\eta/\nu_Y)\mathcal S_t v\).  This gives one inclusion in
\eqref{eq:kernel-graph}.  Conversely, if \(v\in\mathcal D(\mathcal S_t)\) and
\(u:=-(\eta/\nu_Y)\mathcal S_t v\), then
\[
\nu_Y\mathcal K_{Y,t}u+\eta\mathcal K_{XY,t}v
=
-\eta\mathcal K_{Y,t}\mathcal S_t v+\eta\mathcal K_{XY,t}v
=
0,
\]
so \((u,v)\in\ker\mathcal O_t\).  This proves \eqref{eq:kernel-graph}.
\end{proof}

\begin{proposition}[Transfer-operator formula for the observed gain]
\label{prop:transfer-gain}
Assume Assumptions~\ref{ass:model} and~\ref{ass:obs-source}, let $\eta>0$, and fix
$0<t<t+h\le T$. Then
\begin{equation}\label{eq:transfer-gain}
\widetilde G_{X\to Y}(t,h)
=
\eta^2
\sup_{v\in\mathcal D(\mathcal S_t)\setminus\{0\}}
\frac{
\bigl|
\langle g_{XY,t,h},v\rangle_{H_t}
-\langle g_{Y,t,h},\mathcal S_t v\rangle_{H_t}
\bigr|^2
}{
\|v\|_{H_t}^2+\frac{\eta^2}{\nu_Y^2}\|\mathcal S_t v\|_{H_t}^2
},
\end{equation}
with the convention that the supremum is $0$ if $\mathcal D(\mathcal S_t)=\{0\}$.
\end{proposition}

\begin{proof}
If $\mathcal D(\mathcal S_t)=\{0\}$, then Proposition~\ref{prop:transfer-graph} gives
$\ker \mathcal O_t=\{0\}$. Proposition~\ref{prop:nullspace} therefore yields
$\widetilde G_{X\to Y}(t,h)=0$, which agrees with the convention in
\eqref{eq:transfer-gain}.

Assume now that $\mathcal D(\mathcal S_t)\neq\{0\}$. By
Proposition~\ref{prop:transfer-graph}, every nonzero vector in $\ker \mathcal O_t$ can be
written uniquely as
\[
w_v:=\left(-\frac{\eta}{\nu_Y}\mathcal S_t v,\; v\right),
\qquad
v\in\mathcal D(\mathcal S_t)\setminus\{0\}.
\]
Its norm is
\begin{equation}\label{eq:wvnorm}
\|w_v\|_{H_t^\oplus}^2
=
\frac{\eta^2}{\nu_Y^2}\|\mathcal S_t v\|_{H_t}^2+\|v\|_{H_t}^2.
\end{equation}
Using \eqref{eq:gamma},
\begin{align}
\langle \gamma_{t,h},w_v\rangle_{H_t^\oplus}
&=
\left\langle
(\nu_Y g_{Y,t,h},\eta g_{XY,t,h}),
\left(-\frac{\eta}{\nu_Y}\mathcal S_t v,\; v\right)
\right\rangle_{H_t^\oplus} \notag\\
&=
\eta\Bigl(
\langle g_{XY,t,h},v\rangle_{H_t}
-\langle g_{Y,t,h},\mathcal S_t v\rangle_{H_t}
\Bigr). \label{eq:transfer-pairing}
\end{align}
Applying Proposition~\ref{prop:nullspace} and parametrising unit vectors in
$\ker \mathcal O_t$ by $w_v/\|w_v\|_{H_t^\oplus}$ gives
\[
\widetilde G_{X\to Y}(t,h)
=
\sup_{v\in\mathcal D(\mathcal S_t)\setminus\{0\}}
\frac{
|\langle \gamma_{t,h},w_v\rangle_{H_t^\oplus}|^2
}{
\|w_v\|_{H_t^\oplus}^2
}.
\]
Substituting \eqref{eq:wvnorm} and \eqref{eq:transfer-pairing} yields
\eqref{eq:transfer-gain}.
\end{proof}

\begin{corollary}[Transfer-operator zero-gain criterion]
\label{cor:transfer-zero}
Assume Assumptions~\ref{ass:model} and~\ref{ass:obs-source}, let $\eta>0$, and fix
$0<t<t+h\le T$. Then the following are equivalent:
\begin{enumerate}[label=(\roman*)]
\item $\widetilde G_{X\to Y}(t,h)=0$;
\item for every $v\in\mathcal D(\mathcal S_t)$,
\begin{equation}\label{eq:transfer-zero}
\langle g_{XY,t,h},v\rangle_{H_t}
=
\langle g_{Y,t,h},\mathcal S_t v\rangle_{H_t};
\end{equation}
\item for every $(u,v)\in\ker \mathcal O_t$,
\begin{equation}\label{eq:kernel-orth}
\nu_Y\langle g_{Y,t,h},u\rangle_{H_t}
+ \eta\langle g_{XY,t,h},v\rangle_{H_t}=0.
\end{equation}
\end{enumerate}
\end{corollary}

\begin{proof}
If $\mathcal D(\mathcal S_t)=\{0\}$, then Proposition~\ref{prop:transfer-gain} gives
$\widetilde G_{X\to Y}(t,h)=0$, and both (ii) and (iii) are vacuous. So suppose
$\mathcal D(\mathcal S_t)\neq\{0\}$.

By Proposition~\ref{prop:transfer-gain}, $\widetilde G_{X\to Y}(t,h)=0$ if and only if
\[
\langle g_{XY,t,h},v\rangle_{H_t}
-\langle g_{Y,t,h},\mathcal S_t v\rangle_{H_t}
=
0
\]
for every $v\in\mathcal D(\mathcal S_t)$, which proves the equivalence of (i) and (ii).

For (ii) versus (iii), Proposition~\ref{prop:transfer-graph} says that every
$(u,v)\in\ker \mathcal O_t$ is of the form
\[
(u,v)=\left(-\frac{\eta}{\nu_Y}\mathcal S_t v,\; v\right)
\]
with $v\in\mathcal D(\mathcal S_t)$. Substituting this identity into
\eqref{eq:kernel-orth} gives exactly \eqref{eq:transfer-zero}, so (ii) and (iii) are
equivalent.
\end{proof}

\begin{remark}[Interpretation of the transfer-operator reduction]
The operator $\mathcal S_t$ maps a contagion-side hidden direction into the unique
idiosyncratic hidden direction that cancels it in the target past. Thus the observable gain
is controlled by a single mismatch functional,
\[
v\mapsto
\langle g_{XY,t,h},v\rangle_{H_t}
-\langle g_{Y,t,h},\mathcal S_t v\rangle_{H_t},
\]
on the graph domain of $\mathcal S_t$. The asymptotic analysis therefore reduces to the
near-diagonal action of $\mathcal S_t$ on the singular increment profiles
$g_{Y,t,h}$ and $g_{XY,t,h}$.
\end{remark}

\section{Graph-space blow-up for the unbounded regime}\label{sec:unbounded}

The bounded-transfer case covers the regime in which the cancellation operator acts on all of
$H_t$.  Rougher contagion changes the space.  The transfer operator is expected to be unbounded,
the ambient object becomes the graph space of $\mathcal S_t$, and the observable gain becomes the
dual norm of a graph-space functional.  The near-diagonal increment profiles are the local objects
that survive the blow-up.

\begin{definition}[Graph space and graph functional]\label{def:graphspace}
Assume $\eta>0$ and fix $t\in(0,T]$. Let
\[
\mathfrak G_t:=\mathcal D(\mathcal S_t)
\]
equipped with the graph inner product
\begin{equation}\label{eq:graph-inner}
\langle v,w\rangle_{\mathfrak G_t}
:=
\langle v,w\rangle_{H_t}
+
\frac{\eta^2}{\nu_Y^2}\langle \mathcal S_t v,\mathcal S_t w\rangle_{H_t},
\qquad
v,w\in\mathfrak G_t.
\end{equation}
For $0\le t<t+h\le T$, define the graph functional
\begin{equation}\label{eq:graph-functional}
\mathcal F_{t,h}(v)
:=
\langle g_{XY,t,h},v\rangle_{H_t}
-\langle g_{Y,t,h},\mathcal S_t v\rangle_{H_t},
\qquad
v\in\mathfrak G_t.
\end{equation}
\end{definition}

\begin{proposition}[Exact graph-space dual formula]
\label{prop:graph-dual}
Assume Assumptions~\ref{ass:model} and~\ref{ass:obs-source}, let $\eta>0$, and fix
$0<t<t+h\le T$. Then:
\begin{enumerate}[label=(\roman*)]
\item $(\mathfrak G_t,\langle\cdot,\cdot\rangle_{\mathfrak G_t})$ is a Hilbert space;
\item $\mathcal F_{t,h}$ is a continuous linear functional on $\mathfrak G_t$;
\item the observed gain satisfies
\begin{equation}\label{eq:graph-dual}
\widetilde G_{X\to Y}(t,h)
=
\eta^2\|\mathcal F_{t,h}\|_{\mathfrak G_t^*}^2.
\end{equation}
\end{enumerate}
\end{proposition}

\begin{proof}
Because $\mathcal S_t$ is closed by Proposition~\ref{prop:transfer-graph}, its graph is
closed in $H_t^\oplus$. Therefore the graph norm
\[
\|v\|_{\mathfrak G_t}^2
=
\|v\|_{H_t}^2+\frac{\eta^2}{\nu_Y^2}\|\mathcal S_t v\|_{H_t}^2
\]
makes $\mathfrak G_t$ into a Hilbert space, proving (i).

For $v\in\mathfrak G_t$, Cauchy--Schwarz gives
\begin{align*}
|\mathcal F_{t,h}(v)|
&\le
\|g_{XY,t,h}\|_{H_t}\|v\|_{H_t}
+\|g_{Y,t,h}\|_{H_t}\|\mathcal S_t v\|_{H_t} \\
&\le
\left(
\|g_{XY,t,h}\|_{H_t}
+\frac{\nu_Y}{\eta}\|g_{Y,t,h}\|_{H_t}
\right)\|v\|_{\mathfrak G_t},
\end{align*}
so $\mathcal F_{t,h}\in\mathfrak G_t^*$, proving (ii).

Finally, if $\mathfrak G_t=\{0\}$, then Proposition~\ref{prop:transfer-gain} gives
$\widetilde G_{X\to Y}(t,h)=0$ and $\|\mathcal F_{t,h}\|_{\mathfrak G_t^*}=0$. Otherwise,
Proposition~\ref{prop:transfer-gain} states that
\[
\widetilde G_{X\to Y}(t,h)
=
\eta^2
\sup_{v\in\mathfrak G_t\setminus\{0\}}
\frac{|\mathcal F_{t,h}(v)|^2}{\|v\|_{\mathfrak G_t}^2},
\]
which is \eqref{eq:graph-dual}. This proves (iii).
\end{proof}

\begin{definition}[Dilation and universal profile]\label{def:dilation}
For $t\in(0,T]$ and $h\in(0,t]$, define the dilation map
\begin{equation}\label{eq:dilation}
(\mathcal D_{t,h}f)(u):=h^{1/2}f(t-hu)\one_{(0,t/h)}(u),
\qquad
u>0,
\end{equation}
from $L^2([0,t])$ into $L^2(\R_+)$.

For $H\in(0,\tfrac12)$, define the universal profile
\begin{equation}\label{eq:psiH}
\psi_H(u):=(1+u)^{H-\frac12}-u^{H-\frac12},
\qquad u>0,
\end{equation}
viewed as an element of $L^2(\R_+)$. This is legitimate by Lemma~\ref{lem:bH}.
\end{definition}

\begin{lemma}[Dilation isometry]\label{lem:dilation-isometry}
Fix $t\in(0,T]$ and $h\in(0,t]$. Then $\mathcal D_{t,h}:H_t\to L^2(\R_+)$ is a linear
isometric embedding. For every $f,g\in H_t$,
\begin{equation}\label{eq:dilation-isometry}
\langle \mathcal D_{t,h}f,\mathcal D_{t,h}g\rangle_{L^2(\R_+)}
=
\langle f,g\rangle_{H_t}.
\end{equation}
\end{lemma}

\begin{proof}
Linearity follows from the definition of \(\mathcal D_{t,h}\).  For \(f,g\in H_t\), the change of
variables \(s=t-hu\) yields
\begin{align*}
\langle \mathcal D_{t,h}f,\mathcal D_{t,h}g\rangle_{L^2(\R_+)}
&=
\int_0^{t/h} h\,f(t-hu)g(t-hu)\,du \\
&=
\int_0^t f(s)g(s)\,ds \\
&=
\langle f,g\rangle_{H_t}.
\end{align*}
This proves \eqref{eq:dilation-isometry}.
\end{proof}

\begin{proposition}[Universal profile limit for the increment kernels]
\label{prop:profile}
Assume Assumption~\ref{ass:model}. Fix $t\in(0,T)$ and let $i\in\{Y,XY\}$. Define
\begin{equation}\label{eq:profile-gh}
\widehat g_{i,t,h}
:=
h^{-H_i}\mathcal D_{t,h}g_{i,t,h}
\in L^2(\R_+),
\qquad 0<h\le t.
\end{equation}
Then
\begin{equation}\label{eq:profile-limit}
\widehat g_{i,t,h}\longrightarrow L_i(t,t)\psi_{H_i}
\qquad\text{in }L^2(\R_+)
\end{equation}
as $h\downarrow0$.
\end{proposition}

\begin{proof}
Fix $i\in\{Y,XY\}$, write $H:=H_i$, $L:=L_i$, and set $\alpha:=H-\tfrac12\in(-\tfrac12,0)$.
For $u\in(0,t/h)$ one has
\begin{align}
\widehat g_{i,t,h}(u)
&=
h^{\frac12-H}\Bigl(K_i(t+h,t-hu)-K_i(t,t-hu)\Bigr) \notag\\
&=
(1+u)^\alpha L(t+h,t-hu)-u^\alpha L(t,t-hu). \label{eq:profile-pointwise}
\end{align}
Also, $\widehat g_{i,t,h}(u)=0$ for $u\ge t/h$.

Let
\[
u_h:=h^{-1/2}.
\]
We split
\[
\|\widehat g_{i,t,h}-L(t,t)\psi_H\|_{L^2(\R_+)}^2
\le
I_h^{\mathrm{near}}+2I_h^{\mathrm{far}}+2I_h^{\mathrm{tail}},
\]
where
\begin{align*}
I_h^{\mathrm{near}}
&:=
\int_0^{u_h}
\bigl|\widehat g_{i,t,h}(u)-L(t,t)\psi_H(u)\bigr|^2\,du,\\
I_h^{\mathrm{far}}
&:=
\int_{u_h}^{t/h}\bigl|\widehat g_{i,t,h}(u)\bigr|^2\,du,\\
I_h^{\mathrm{tail}}
&:=
L(t,t)^2\int_{u_h}^\infty |\psi_H(u)|^2\,du.
\end{align*}

\emph{Near-diagonal region.}
For $u\in[0,u_h]$, one has $0\le hu\le h^{1/2}$. Since $L$ is $C^1$ near the diagonal,
there exists $C<\infty$ such that
\[
|L(t+h,t-hu)-L(t,t)|+|L(t,t-hu)-L(t,t)|\le C(h+hu)\le 2Ch^{1/2}.
\]
Hence by \eqref{eq:profile-pointwise},
\[
\widehat g_{i,t,h}(u)-L(t,t)\psi_H(u)
=
(1+u)^\alpha \Delta_{1,h}(u)-u^\alpha \Delta_{2,h}(u),
\]
with $|\Delta_{1,h}(u)|+|\Delta_{2,h}(u)|\le 2Ch^{1/2}$. Therefore
\[
\bigl|\widehat g_{i,t,h}(u)-L(t,t)\psi_H(u)\bigr|
\le
Ch^{1/2}\bigl((1+u)^\alpha+u^\alpha\bigr),
\]
and so
\[
I_h^{\mathrm{near}}
\lesssim
h\int_0^{u_h}\bigl((1+u)^{2H-1}+u^{2H-1}\bigr)\,du
=
O(h^{1-H})
\longrightarrow 0.
\]

\emph{Far region.}
Since $h^{1/2}\ge h$ for $0<h\le1$, Lemma~\ref{lem:offdiag-pointwise} yields a constant
$C<\infty$ such that for all sufficiently small $h$ and all $s\in[0,t-h^{1/2}]$,
\[
|K_i(t+h,s)-K_i(t,s)|
\le
Ch\bigl((t-s)^{H-\frac32}+(t-s)^{H-\frac12}\bigr).
\]
Since $u\mapsto t-hu$ maps $[u_h,t/h]$ onto $[0,t-h^{1/2}]$, it follows that
\begin{align*}
I_h^{\mathrm{far}}
&=
\int_{u_h}^{t/h}|\widehat g_{i,t,h}(u)|^2\,du \\
&=
h^{-2H}\int_0^{t-h^{1/2}} \bigl(K_i(t+h,s)-K_i(t,s)\bigr)^2\,ds \\
&\lesssim
h^{2-2H}\int_0^{t-h^{1/2}}
\bigl((t-s)^{2H-3}+(t-s)^{2H-1}\bigr)\,ds \\
&=
O(h^{1-H})+O(h^{2-2H})
\longrightarrow 0.
\end{align*}

\emph{Tail of the limit profile.}
Since $\psi_H\in L^2(\R_+)$ by Lemma~\ref{lem:bH}, one has
\[
I_h^{\mathrm{tail}}\longrightarrow 0
\qquad\text{as }h\downarrow0.
\]

Combining the three terms proves \eqref{eq:profile-limit}.
\end{proof}

\begin{corollary}[Profile pairing limit]\label{cor:profile-pairing}
Under the assumptions of Proposition~\ref{prop:profile}, for every
$\phi\in L^2(\R_+)$ one has
\begin{equation}\label{eq:profile-pairing-limit}
\left\langle \widehat g_{i,t,h},\phi\right\rangle_{L^2(\R_+)}
\longrightarrow
L_i(t,t)\left\langle \psi_{H_i},\phi\right\rangle_{L^2(\R_+)}
\end{equation}
as $h\downarrow0$.
\end{corollary}

\begin{proof}
By Proposition~\ref{prop:profile},
\[
\widehat g_{i,t,h}-L_i(t,t)\psi_{H_i}\to0
\qquad\text{in }L^2(\R_+).
\]
For fixed \(\phi\in L^2(\R_+)\), Cauchy--Schwarz gives
\[
\left|
\left\langle
\widehat g_{i,t,h}-L_i(t,t)\psi_{H_i},\phi
\right\rangle_{L^2(\R_+)}
\right|
\le
\left\|
\widehat g_{i,t,h}-L_i(t,t)\psi_{H_i}
\right\|_{L^2(\R_+)}
\|\phi\|_{L^2(\R_+)},
\]
and the right-hand side tends to zero.
\end{proof}

\section{Fractional model cancellation}\label{sec:model}

The profile block identifies the singular shapes carried by the increment kernels.  Cancellation
requires an operator with the same tangent scale.  On the half-line the natural candidate is a
fractional derivative of order $H_Y-H_{XY}$.  The exact model identity below is the cancellation law
that any graph blow-up must satisfy; the transfer-operator proof of that blow-up remains separate.

\begin{definition}[Liouville--Weyl fractional derivatives]\label{def:LW}
Let $\beta\in(0,1)$. For $\phi\in C_c^\infty(\R)$, define
\begin{align}
(\mathfrak D_+^\beta \phi)(x)
&:=
\frac{1}{\Gamma(1-\beta)}
\frac{d}{dx}
\int_{-\infty}^x (x-y)^{-\beta}\phi(y)\,dy,
\label{eq:Dplus}\\
(\mathfrak D_-^\beta \phi)(x)
&:=
-\frac{1}{\Gamma(1-\beta)}
\frac{d}{dx}
\int_x^\infty (y-x)^{-\beta}\phi(y)\,dy.
\label{eq:Dminus}
\end{align}
We extend $\mathfrak D_+^\beta$ to distributions by duality:
\begin{equation}\label{eq:Dplus-dual}
\langle \mathfrak D_+^\beta U,\phi\rangle
:=
\langle U,\mathfrak D_-^\beta\phi\rangle,
\qquad
U\in\mathcal D'(\R),\quad \phi\in C_c^\infty(\R).
\end{equation}
Finally, for $a\in\R$, let $\tau_a$ denote translation by
\[
(\tau_a f)(x):=f(x+a).
\]
\end{definition}

\begin{lemma}[Adjointness and translation invariance]\label{lem:LW-basic}
Let $\beta\in(0,1)$.
\begin{enumerate}[label=(\roman*)]
\item For every $f,g\in C_c^\infty(\R)$,
\begin{equation}\label{eq:LW-adjoint}
\int_\R (\mathfrak D_+^\beta f)(x)g(x)\,dx
=
\int_\R f(x)(\mathfrak D_-^\beta g)(x)\,dx.
\end{equation}
\item For every $a\in\R$ and every $f\in C_c^\infty(\R)$,
\begin{equation}\label{eq:LW-translation}
\mathfrak D_\pm^\beta(\tau_a f)
=
\tau_a(\mathfrak D_\pm^\beta f).
\end{equation}
\end{enumerate}
\end{lemma}

\begin{proof}
For $f\in C_c^\infty(\R)$, set
\[
(\mathfrak I_+^{1-\beta}f)(x)
:=
\frac{1}{\Gamma(1-\beta)}
\int_{-\infty}^x (x-y)^{-\beta}f(y)\,dy,
\]
so that $\mathfrak D_+^\beta f=(\mathfrak I_+^{1-\beta}f)'$. Likewise, for
$g\in C_c^\infty(\R)$,
\[
(\mathfrak I_-^{1-\beta}g)(x)
:=
\frac{1}{\Gamma(1-\beta)}
\int_x^\infty (y-x)^{-\beta}g(y)\,dy,
\]
so that $\mathfrak D_-^\beta g=-(\mathfrak I_-^{1-\beta}g)'$.

Because $f$ and $g$ are compactly supported, $\mathfrak I_+^{1-\beta}f$ and
$\mathfrak I_-^{1-\beta}g$ are absolutely continuous and vanish at $\pm\infty$. Hence
integration by parts gives
\begin{align*}
\int_\R (\mathfrak D_+^\beta f)(x)g(x)\,dx
&=
-\int_\R (\mathfrak I_+^{1-\beta}f)(x)g'(x)\,dx \\
&=
-\frac{1}{\Gamma(1-\beta)}
\int_\R\int_{-\infty}^x (x-y)^{-\beta}f(y)g'(x)\,dy\,dx.
\end{align*}
Fubini's theorem applies because the integrand is absolutely integrable. Therefore
\begin{align*}
\int_\R (\mathfrak D_+^\beta f)(x)g(x)\,dx
&=
-\frac{1}{\Gamma(1-\beta)}
\int_\R f(y)\left(\int_y^\infty (x-y)^{-\beta}g'(x)\,dx\right)dy.
\end{align*}
After the change of variables $x=y+r$,
\[
\frac{1}{\Gamma(1-\beta)}\int_y^\infty (x-y)^{-\beta}g'(x)\,dx
=
\frac{d}{dy}(\mathfrak I_-^{1-\beta}g)(y),
\]
so the last display equals
\[
\int_\R f(y)\bigl(-(\mathfrak I_-^{1-\beta}g)'(y)\bigr)\,dy
=
\int_\R f(y)(\mathfrak D_-^\beta g)(y)\,dy,
\]
which proves \eqref{eq:LW-adjoint}.

For translation invariance, let $a\in\R$. A change of variables gives
\begin{align*}
(\mathfrak D_+^\beta(\tau_a f))(x)
&=
\frac{1}{\Gamma(1-\beta)}
\frac{d}{dx}
\int_{-\infty}^{x+a} (x+a-z)^{-\beta}f(z)\,dz \\
&=
(\tau_a \mathfrak D_+^\beta f)(x).
\end{align*}
The proof for $\mathfrak D_-^\beta$ is identical.
\par
By duality, \eqref{eq:LW-translation} extends from test functions to distributions.
\end{proof}

\begin{proposition}[Fractional derivative of the one-sided power kernel]
\label{prop:power-kernel}
Let $H\in(0,\tfrac12)$ and $\beta\in(0,H+\tfrac12)$. Define
\begin{equation}\label{eq:kH}
k_H(x):=x_+^{H-\frac12},
\qquad x\in\R.
\end{equation}
Then, in $\mathcal D'(\R)$,
\begin{equation}\label{eq:power-kernel}
\mathfrak D_+^\beta k_H
=
\frac{\Gamma(H+\frac12)}{\Gamma(H-\beta+\frac12)}\,k_{H-\beta}.
\end{equation}
\end{proposition}

\begin{proof}
Write $\lambda:=H+\tfrac12\in(0,1)$. For $x<0$, both sides of \eqref{eq:power-kernel} are
zero. For $x>0$, one has
\begin{align*}
(\mathfrak I_+^{1-\beta}k_H)(x)
&=
\frac{1}{\Gamma(1-\beta)}
\int_0^x (x-y)^{-\beta}y^{\lambda-1}\,dy \\
&=
\frac{x^{\lambda-\beta}}{\Gamma(1-\beta)}
\int_0^1 (1-r)^{-\beta}r^{\lambda-1}\,dr \\
&=
\frac{\Gamma(\lambda)}{\Gamma(\lambda+1-\beta)}x^{\lambda-\beta},
\end{align*}
by the Beta-function identity. Differentiating yields
\[
(\mathfrak D_+^\beta k_H)(x)
=
\frac{\Gamma(\lambda)}{\Gamma(\lambda-\beta)}x^{\lambda-\beta-1}
=
\frac{\Gamma(H+\frac12)}{\Gamma(H-\beta+\frac12)}x^{H-\beta-\frac12},
\qquad x>0.
\]
Both sides are locally integrable on $\R$, so the almost-everywhere identity extends to the
distributional identity \eqref{eq:power-kernel}.
\end{proof}

\begin{proposition}[Exact fractional cancellation of the universal profiles]
\label{prop:model-cancel}
Assume
\[
0<H_{XY}<H_Y<\tfrac12,
\qquad
\beta:=H_Y-H_{XY}\in(0,H_Y).
\]
Fix $t\in(0,T)$ and define the model transfer coefficient
\begin{equation}\label{eq:model-coeff}
\mathbf c_t
:=
\frac{L_{XY}(t,t)}{L_Y(t,t)}
\frac{\Gamma(H_{XY}+\frac12)}{\Gamma(H_Y+\frac12)}.
\end{equation}
For $\phi\in C_c^\infty((0,\infty))$, set
\begin{equation}\label{eq:model-transfer}
(\mathcal T_{t,\beta}^{\mathrm{mod}}\phi)(u)
:=
\mathbf c_t\,(\mathfrak D_-^\beta \phi)(u),
\qquad u>0.
\end{equation}
Then for every $\phi\in C_c^\infty((0,\infty))$,
\begin{equation}\label{eq:model-cancel}
L_{XY}(t,t)\langle \psi_{H_{XY}},\phi\rangle_{L^2(\R_+)}
=
L_Y(t,t)
\left\langle
\psi_{H_Y},
\mathcal T_{t,\beta}^{\mathrm{mod}}\phi
\right\rangle_{L^2(\R_+)}.
\end{equation}
\end{proposition}

\begin{proof}
For $H\in(0,\tfrac12)$ and $u>0$, \eqref{eq:psiH} gives
\[
\psi_H(u)=k_H(u+1)-k_H(u)=((\tau_1-I)k_H)(u),
\]
with $k_H$ as in \eqref{eq:kH}. By Proposition~\ref{prop:power-kernel} and the translation
invariance from Lemma~\ref{lem:LW-basic},
\begin{align*}
\mathfrak D_+^\beta \psi_{H_Y}
&=
(\tau_1-I)\mathfrak D_+^\beta k_{H_Y} \\
&=
\frac{\Gamma(H_Y+\frac12)}{\Gamma(H_Y-\beta+\frac12)}
(\tau_1-I)k_{H_Y-\beta} \\
&=
\frac{\Gamma(H_Y+\frac12)}{\Gamma(H_{XY}+\frac12)}\psi_{H_{XY}}
\end{align*}
in $\mathcal D'((0,\infty))$.

Now let $\phi\in C_c^\infty((0,\infty))$. Since $\phi$ is compactly supported away from the
boundary, $\mathfrak D_-^\beta\phi\in L^2(\R_+)$. Hence the first and last pairings below
are ordinary $L^2$ integrals, while the middle one is distributional. Using the duality
\eqref{eq:Dplus-dual},
\begin{align*}
\left\langle
\psi_{H_Y},
\mathcal T_{t,\beta}^{\mathrm{mod}}\phi
\right\rangle_{L^2(\R_+)}
&=
\mathbf c_t\langle \psi_{H_Y},\mathfrak D_-^\beta\phi\rangle \\
&=
\mathbf c_t\langle \mathfrak D_+^\beta \psi_{H_Y},\phi\rangle \\
&=
\mathbf c_t
\frac{\Gamma(H_Y+\frac12)}{\Gamma(H_{XY}+\frac12)}
\langle \psi_{H_{XY}},\phi\rangle_{L^2(\R_+)} \\
&=
\frac{L_{XY}(t,t)}{L_Y(t,t)}
\langle \psi_{H_{XY}},\phi\rangle_{L^2(\R_+)},
\end{align*}
which is exactly \eqref{eq:model-cancel}.
\end{proof}

\begin{corollary}[Vanishing of the model leading functional]
\label{cor:model-functional}
Under the assumptions of Proposition~\ref{prop:model-cancel}, define
\begin{equation}\label{eq:model-functional}
\mathcal F_{t,\beta}^{\mathrm{mod}}(\phi)
:=
L_{XY}(t,t)\langle \psi_{H_{XY}},\phi\rangle_{L^2(\R_+)}
-L_Y(t,t)
\left\langle
\psi_{H_Y},
\mathcal T_{t,\beta}^{\mathrm{mod}}\phi
\right\rangle_{L^2(\R_+)},
\end{equation}
for $\phi\in C_c^\infty((0,\infty))$. Then
\[
\mathcal F_{t,\beta}^{\mathrm{mod}}(\phi)=0
\qquad\text{for every }\phi\in C_c^\infty((0,\infty)).
\]
\end{corollary}

\begin{proof}
The functional \(\mathcal F_{t,\beta}^{\mathrm{mod}}\) is the difference between the two sides of
\eqref{eq:model-cancel}.  Proposition~\ref{prop:model-cancel} identifies those sides for every
\(\phi\in C_c^\infty((0,\infty))\), so the difference is zero.
\end{proof}

\section{Directional cancellation under fractional blow-up}\label{sec:directional}

The model identity determines the only plausible leading-order cancellation law.  Localized smooth
profiles transfer that identity back to the original graph functional.  The conclusion is
directional; uniform control on the graph unit ball remains a separate problem.

\begin{definition}[Localized profile lift]\label{def:lift}
Fix $t\in(0,T)$ and $h\in(0,t]$. For $\phi\in L^2(\R_+)$, define
\begin{equation}\label{eq:lift}
(\mathcal L_{t,h}\phi)(s)
:=
h^{-1/2}\phi\!\left(\frac{t-s}{h}\right)\one_{[0,t]}(s),
\qquad 0\le s\le t.
\end{equation}
\end{definition}

\begin{lemma}[Dilation-lift duality]\label{lem:lift-duality}
Fix $t\in(0,T)$ and let $\phi\in C_c^\infty((0,\infty))$. Then for all sufficiently small
$h>0$ such that $\operatorname{supp}(\phi)\subset(0,t/h)$, one has:
\begin{enumerate}[label=(\roman*)]
\item $\mathcal L_{t,h}\phi\in H_t$ and
\[
\|\mathcal L_{t,h}\phi\|_{H_t}=\|\phi\|_{L^2(\R_+)};
\]
\item for every $f\in H_t$,
\begin{equation}\label{eq:lift-dual}
\langle f,\mathcal L_{t,h}\phi\rangle_{H_t}
=
\langle \mathcal D_{t,h}f,\phi\rangle_{L^2(\R_+)};
\end{equation}
\item for each $i\in\{Y,XY\}$,
\begin{equation}\label{eq:lift-gh}
h^{-H_i}\langle g_{i,t,h},\mathcal L_{t,h}\phi\rangle_{H_t}
=
\langle \widehat g_{i,t,h},\phi\rangle_{L^2(\R_+)}.
\end{equation}
\end{enumerate}
\end{lemma}

\begin{proof}
Choose $h>0$ so small that $\operatorname{supp}(\phi)\subset(0,t/h)$. Then
\[
\|\mathcal L_{t,h}\phi\|_{H_t}^2
=
\int_0^t h^{-1}\left|\phi\!\left(\frac{t-s}{h}\right)\right|^2ds
=
\int_0^{t/h} |\phi(u)|^2\,du
=
\|\phi\|_{L^2(\R_+)}^2,
\]
because $\phi$ vanishes outside $(0,t/h)$. This proves (i).

For $f\in H_t$, the change of variables $u=(t-s)/h$ gives
\begin{align*}
\langle f,\mathcal L_{t,h}\phi\rangle_{H_t}
&=
\int_0^t f(s)h^{-1/2}\phi\!\left(\frac{t-s}{h}\right)\,ds \\
&=
\int_0^{t/h} h^{1/2}f(t-hu)\phi(u)\,du \\
&=
\langle \mathcal D_{t,h}f,\phi\rangle_{L^2(\R_+)},
\end{align*}
which proves (ii). Finally, substituting $f=g_{i,t,h}$ into \eqref{eq:lift-dual} and using
\eqref{eq:profile-gh} yields (iii).
\end{proof}

\begin{definition}[Localized fractional blow-up at time $t$]
\label{def:localized-blowup}
Assume
\[
0<H_{XY}<H_Y<\tfrac12,
\qquad
\beta:=H_Y-H_{XY}.
\]
We say that the hidden transfer operator has \emph{localized fractional blow-up at time $t$}
if, for every $\phi\in C_c^\infty((0,\infty))$, there exists $h_\phi\in(0,t]$ such that
\[
\mathcal L_{t,h}\phi\in\mathcal D(\mathcal S_t)
\qquad\text{for all }0<h\le h_\phi,
\]
and
\begin{equation}\label{eq:localized-blowup}
h^{\beta}\mathcal D_{t,h}\bigl(\mathcal S_t\mathcal L_{t,h}\phi\bigr)
\longrightarrow
\mathcal T_{t,\beta}^{\mathrm{mod}}\phi
\qquad\text{in }L^2(\R_+)
\end{equation}
as $h\downarrow0$.
\end{definition}

\begin{theorem}[Directional leading-order cancellation]\label{thm:directional-cancel}
Assume Assumptions~\ref{ass:model} and~\ref{ass:obs-source}, let $\eta>0$, and fix
$t\in(0,T)$. Suppose
\[
0<H_{XY}<H_Y<\tfrac12
\]
and that $\mathcal S_t$ has localized fractional blow-up at time $t$ in the sense of
Definition~\ref{def:localized-blowup}. Then for every $\phi\in C_c^\infty((0,\infty))$,
\begin{equation}\label{eq:directional-cancel}
h^{-H_{XY}}\mathcal F_{t,h}(\mathcal L_{t,h}\phi)
\longrightarrow 0
\qquad\text{as }h\downarrow0.
\end{equation}
\end{theorem}

\begin{proof}
Fix $\phi\in C_c^\infty((0,\infty))$ and let $h>0$ be small enough that
$\mathcal L_{t,h}\phi\in\mathcal D(\mathcal S_t)$. Set
\[
v_h:=\mathcal L_{t,h}\phi,
\qquad
z_h:=h^{\beta}\mathcal D_{t,h}(\mathcal S_t v_h)\in L^2(\R_+).
\]
By \eqref{eq:graph-functional}, Lemma~\ref{lem:lift-duality}, and the identity
$H_{XY}=H_Y-\beta$,
\begin{align*}
h^{-H_{XY}}\mathcal F_{t,h}(v_h)
&=
h^{-H_{XY}}\langle g_{XY,t,h},v_h\rangle_{H_t}
-h^{-H_{XY}}\langle g_{Y,t,h},\mathcal S_t v_h\rangle_{H_t} \\
&=
\langle \widehat g_{XY,t,h},\phi\rangle_{L^2(\R_+)}
-h^{-H_Y+\beta}\langle g_{Y,t,h},\mathcal S_t v_h\rangle_{H_t}.
\end{align*}
Applying Lemma~\ref{lem:dilation-isometry} to the second term and using the definitions of
$\widehat g_{Y,t,h}$ and $z_h$ gives
\[
h^{-H_Y+\beta}\langle g_{Y,t,h},\mathcal S_t v_h\rangle_{H_t}
=
\left\langle
h^{-H_Y}\mathcal D_{t,h}g_{Y,t,h},
h^\beta \mathcal D_{t,h}(\mathcal S_t v_h)
\right\rangle_{L^2(\R_+)}
=
\langle \widehat g_{Y,t,h},z_h\rangle_{L^2(\R_+)},
\]
so
\[
h^{-H_{XY}}\mathcal F_{t,h}(v_h)
=
\langle \widehat g_{XY,t,h},\phi\rangle_{L^2(\R_+)}
-\langle \widehat g_{Y,t,h},z_h\rangle_{L^2(\R_+)}.
\]

By Corollary~\ref{cor:profile-pairing},
\[
\langle \widehat g_{XY,t,h},\phi\rangle_{L^2(\R_+)}
\longrightarrow
L_{XY}(t,t)\langle \psi_{H_{XY}},\phi\rangle_{L^2(\R_+)}.
\]
Also, Proposition~\ref{prop:profile} gives
\[
\widehat g_{Y,t,h}\to L_Y(t,t)\psi_{H_Y}
\qquad\text{in }L^2(\R_+),
\]
while Definition~\ref{def:localized-blowup} gives
\[
z_h\to \mathcal T_{t,\beta}^{\mathrm{mod}}\phi
\qquad\text{in }L^2(\R_+).
\]
Since $(z_h)$ is therefore bounded in $L^2(\R_+)$,
\begin{align*}
&\bigl|
\langle \widehat g_{Y,t,h},z_h\rangle_{L^2(\R_+)}
-L_Y(t,t)
\langle \psi_{H_Y},\mathcal T_{t,\beta}^{\mathrm{mod}}\phi\rangle_{L^2(\R_+)}
\bigr| \\
&\qquad\le
\|\widehat g_{Y,t,h}-L_Y(t,t)\psi_{H_Y}\|_{L^2(\R_+)}\|z_h\|_{L^2(\R_+)} \\
&\qquad\quad+
L_Y(t,t)\|\psi_{H_Y}\|_{L^2(\R_+)}
\|z_h-\mathcal T_{t,\beta}^{\mathrm{mod}}\phi\|_{L^2(\R_+)},
\end{align*}
and the right-hand side tends to $0$. Hence
\[
\langle \widehat g_{Y,t,h},z_h\rangle_{L^2(\R_+)}
\longrightarrow
L_Y(t,t)
\langle \psi_{H_Y},\mathcal T_{t,\beta}^{\mathrm{mod}}\phi\rangle_{L^2(\R_+)}.
\]
Applying Proposition~\ref{prop:model-cancel} to the limiting terms gives
\eqref{eq:directional-cancel}.
\end{proof}

\section{Uniform suppression from near-diagonal coercivity}\label{sec:uniform}

Directional cancellation does not control the graph-space supremum in
Proposition~\ref{prop:graph-dual}.  A separate coercive input gives the first uniform theorem.  If
graph-unit vectors cannot concentrate too much \(L^2\) mass inside the \(h^{1/2}\) boundary layer
near \(t\), the rough contagion increment loses its oracle scale uniformly.  The far part of that
increment is already of higher order and needs no further structure.

\begin{definition}[Near/far localization on $H_t$]\label{def:near-far}
Fix $t\in(0,T)$ and $h\in(0,t]$. Define the orthogonal projections
\begin{align}
(\Pi_{t,h}^{\mathrm{near}}f)(s)
&:=
f(s)\one_{[t-h^{1/2},\,t]}(s), \label{eq:Pi-near}\\
(\Pi_{t,h}^{\mathrm{far}}f)(s)
&:=
f(s)\one_{[0,\,t-h^{1/2})}(s), \label{eq:Pi-far}
\end{align}
for $f\in H_t=L^2([0,t])$.
\end{definition}

\begin{lemma}[Far-zone estimate for the increment kernels]\label{lem:far-kernel}
Assume Assumption~\ref{ass:model}. Fix $t\in(0,T)$ and let $i\in\{Y,XY\}$. Then
\begin{equation}\label{eq:far-kernel}
\|\Pi_{t,h}^{\mathrm{far}}g_{i,t,h}\|_{H_t}
=
O(h^{\frac{1+H_i}{2}})
\qquad\text{as }h\downarrow0.
\end{equation}
\end{lemma}

\begin{proof}
Write $H:=H_i$. By Lemma~\ref{lem:offdiag-pointwise}, there exist
constants $C<\infty$ and $h_0\in(0,t]$ such that for all $0<h\le h_0$ and all
$s\in[0,t-h^{1/2}]$,
\[
|g_{i,t,h}(s)|
\le
Ch\bigl((t-s)^{H-\frac32}+(t-s)^{H-\frac12}\bigr).
\]
Therefore
\begin{align*}
\|\Pi_{t,h}^{\mathrm{far}}g_{i,t,h}\|_{H_t}^2
&=
\int_0^{t-h^{1/2}} |g_{i,t,h}(s)|^2\,ds \\
&\lesssim
h^2\int_0^{t-h^{1/2}}
\bigl((t-s)^{2H-3}+(t-s)^{2H-1}\bigr)\,ds.
\end{align*}
With the change of variables $u=t-s$, this becomes
\[
\|\Pi_{t,h}^{\mathrm{far}}g_{i,t,h}\|_{H_t}^2
\lesssim
h^2\int_{h^{1/2}}^t (u^{2H-3}+u^{2H-1})\,du.
\]
Since $H\in(0,\tfrac12)$,
\[
\int_{h^{1/2}}^t u^{2H-3}\,du = O(h^{H-1}),
\qquad
\int_{h^{1/2}}^t u^{2H-1}\,du = O(1),
\]
and hence
\[
\|\Pi_{t,h}^{\mathrm{far}}g_{i,t,h}\|_{H_t}^2
=
O(h^{1+H})+O(h^2)
=
O(h^{1+H}).
\]
Taking square roots proves \eqref{eq:far-kernel}.
\end{proof}

\section{Exact Volterra localization of the transfer operator}\label{sec:local-transfer}

The boundary-layer mechanism is still expressed through cutoff and blow-up hypotheses.  The
Volterra structure replaces part of that abstraction by exact local operators.  Causality makes
right-end boundary cutoffs stable in graph norm.  After local dilation to the unit interval, the
true transfer operator induces a rescaled family whose kernels converge to the frozen fractional
model.  The problem becomes local transfer convergence.

\begin{definition}[Boundary-layer subspace and unit-interval dilation]
\label{def:unit-dilation}
Fix $t\in(0,T)$, $\alpha\in(0,1)$, and $h\in(0,t]$. Define the boundary-layer subspace
\[
H_{t,h,\alpha}^{\mathrm{near}}
:=
\left\{
f\in H_t:\operatorname{supp}(f)\subset[t-h^\alpha,t]
\right\}.
\]
Let
\[
E:=L^2([0,1]).
\]
Define the unit-interval dilation and lift by
\begin{align}
(\mathcal D_{t,h}^{(\alpha)}f)(u)
&:=
h^{\alpha/2}f(t-h^\alpha u),
\qquad 0\le u\le 1, \label{eq:unit-dilation}\\
(\mathcal L_{t,h}^{(\alpha)}\phi)(s)
&:=
h^{-\alpha/2}\phi\!\left(\frac{t-s}{h^\alpha}\right)\one_{[t-h^\alpha,t]}(s),
\qquad 0\le s\le t, \label{eq:unit-lift}
\end{align}
for $f\in H_{t,h,\alpha}^{\mathrm{near}}$ and $\phi\in E$.
\end{definition}

\begin{lemma}[Unit-interval dilation isometry]\label{lem:unit-dilation}
Fix $t\in(0,T)$, $\alpha\in(0,1)$, and $h\in(0,t]$. Then
\[
\mathcal D_{t,h}^{(\alpha)}:H_{t,h,\alpha}^{\mathrm{near}}\to E
\]
is an isometric isomorphism with inverse $\mathcal L_{t,h}^{(\alpha)}$. In particular, for
every $f\in H_{t,h,\alpha}^{\mathrm{near}}$ and $\phi\in E$,
\begin{align}
\|\mathcal D_{t,h}^{(\alpha)}f\|_E &= \|f\|_{H_t}, \label{eq:unit-iso}\\
\langle f,\mathcal L_{t,h}^{(\alpha)}\phi\rangle_{H_t}
&=
\langle \mathcal D_{t,h}^{(\alpha)}f,\phi\rangle_E. \label{eq:unit-dual}
\end{align}
\end{lemma}

\begin{proof}
The change of variables $s=t-h^\alpha u$ gives
\[
\|\mathcal D_{t,h}^{(\alpha)}f\|_E^2
=
\int_0^1 h^\alpha |f(t-h^\alpha u)|^2\,du
=
\int_{t-h^\alpha}^t |f(s)|^2\,ds
=
\|f\|_{H_t}^2,
\]
because $f$ is supported in $[t-h^\alpha,t]$. The same change of variables gives
\eqref{eq:unit-dual}. The definitions give
$\mathcal D_{t,h}^{(\alpha)}\mathcal L_{t,h}^{(\alpha)}=I_E$ and
$\mathcal L_{t,h}^{(\alpha)}\mathcal D_{t,h}^{(\alpha)}=I$ on
$H_{t,h,\alpha}^{\mathrm{near}}$.
\end{proof}

\begin{proposition}[Support preservation on boundary-supported inputs]
\label{prop:support-preserve}
Assume Assumptions~\ref{ass:model} and~\ref{ass:obs-source}, let $\eta>0$, fix $t\in(0,T)$,
let $\alpha\in(0,1)$, and set $a:=t-h^\alpha$. If $w\in\mathfrak G_t$ satisfies
\[
\operatorname{supp}(w)\subset[a,t],
\]
then
\begin{equation}\label{eq:support-preserve}
\operatorname{supp}(\mathcal S_t w)\subset[a,t].
\end{equation}
\end{proposition}

\begin{proof}
Let $w\in\mathfrak G_t$ satisfy $\operatorname{supp}(w)\subset[a,t]$, and set
$u:=\mathcal S_t w$. Then
\[
\mathcal K_{Y,t}u=\mathcal K_{XY,t}w.
\]
Because $w$ is supported in $[a,t]$, the Volterra form of $\mathcal K_{XY,t}$ yields
\[
(\mathcal K_{XY,t}w)(r)=0
\qquad\text{for }0\le r<a.
\]
Hence
\[
(\mathcal K_{Y,t}u)(r)=0
\qquad\text{for }0\le r<a.
\]
Equivalently,
\[
\int_0^r K_Y(r,s)u(s)\,ds=0
\qquad\text{for }0\le r<a.
\]
This means exactly that the truncated operator $\mathcal K_{Y,a}$ annihilates $u|_{[0,a]}$.
By Assumption~\ref{ass:obs-source}, $\mathcal K_{Y,a}$ is injective, so $u|_{[0,a]}=0$ almost
everywhere. Therefore \eqref{eq:support-preserve} holds.
\end{proof}

\begin{definition}[Frozen right-sided Volterra model on the unit interval]
\label{def:unit-model}
For $H\in(0,\tfrac12)$, define
\begin{equation}\label{eq:JH}
(\mathfrak J_H\phi)(u)
:=
\int_u^1 (v-u)^{H-\frac12}\phi(v)\,dv,
\qquad 0\le u\le 1,
\end{equation}
for $\phi\in E=L^2([0,1])$.
\end{definition}

\begin{definition}[Localized exact Volterra operators]\label{def:local-volterra}
Fix $t\in(0,T)$, $\alpha\in(0,1)$, $h\in(0,t]$, and $i\in\{Y,XY\}$. Define the operator
\[
\mathfrak A_{i,t,h,\alpha}:E\to E
\]
by
\begin{equation}\label{eq:Aih}
(\mathfrak A_{i,t,h,\alpha}\phi)(u)
:=
\int_u^1 (v-u)^{H_i-\frac12}
L_i(t-h^\alpha u,t-h^\alpha v)\phi(v)\,dv,
\qquad 0\le u\le 1.
\end{equation}
\end{definition}

\begin{proposition}[Exact localized representation and frozen-kernel limit]
\label{prop:local-volterra}
Assume Assumption~\ref{ass:model}. Fix $t\in(0,T)$, $\alpha\in(0,1)$, and let
$i\in\{Y,XY\}$. Then:
\begin{enumerate}[label=(\roman*)]
\item for every $\phi\in E$,
\begin{equation}\label{eq:local-representation}
\mathcal D_{t,h}^{(\alpha)}\mathcal K_{i,t}\mathcal L_{t,h}^{(\alpha)}\phi
=
h^{\alpha H_i}\mathfrak A_{i,t,h,\alpha}\phi;
\end{equation}
\item there exists $C<\infty$, depending only on $t,\alpha,i$, such that
\begin{equation}\label{eq:local-freeze}
\left\|
\mathfrak A_{i,t,h,\alpha}-L_i(t,t)\mathfrak J_{H_i}
\right\|_{\mathcal L(E)}
\le
Ch^\alpha
\end{equation}
for all sufficiently small $h>0$.
\end{enumerate}
\end{proposition}

\begin{proof}
Fix $\phi\in E$ and write $\mathcal L=\mathcal L_{t,h}^{(\alpha)}$,
$\mathcal D=\mathcal D_{t,h}^{(\alpha)}$. For $0\le u\le 1$,
\begin{align*}
(\mathcal D\mathcal K_{i,t}\mathcal L\phi)(u)
&=
h^{\alpha/2}
\int_0^{t-h^\alpha u}
K_i(t-h^\alpha u,s)\,
h^{-\alpha/2}\phi\!\left(\frac{t-s}{h^\alpha}\right)\one_{[t-h^\alpha,t]}(s)\,ds \\
&=
h^{\alpha/2}
\int_{t-h^\alpha}^{t-h^\alpha u}
(t-h^\alpha u-s)^{H_i-\frac12}
L_i(t-h^\alpha u,s)\,
h^{-\alpha/2}\phi\!\left(\frac{t-s}{h^\alpha}\right)\,ds.
\end{align*}
With the change of variables $s=t-h^\alpha v$, this becomes
\begin{align*}
(\mathcal D\mathcal K_{i,t}\mathcal L\phi)(u)
&=
h^{\alpha H_i}
\int_u^1
(v-u)^{H_i-\frac12}
L_i(t-h^\alpha u,t-h^\alpha v)\phi(v)\,dv \\
&=
h^{\alpha H_i}(\mathfrak A_{i,t,h,\alpha}\phi)(u),
\end{align*}
which proves \eqref{eq:local-representation}.

To prove \eqref{eq:local-freeze}, note that the difference kernel is
\[
R_{i,t,h,\alpha}(u,v)
:=
(v-u)^{H_i-\frac12}
\bigl(L_i(t-h^\alpha u,t-h^\alpha v)-L_i(t,t)\bigr)\one_{\{0\le u\le v\le 1\}}.
\]
Because $L_i$ is $C^1$ near the diagonal, there exists $C<\infty$ such that
\[
|L_i(t-h^\alpha u,t-h^\alpha v)-L_i(t,t)|
\le
Ch^\alpha
\qquad\text{for }0\le u\le v\le 1.
\]
Hence
\[
|R_{i,t,h,\alpha}(u,v)|
\le
Ch^\alpha (v-u)^{H_i-\frac12}\one_{\{0\le u\le v\le 1\}}.
\]
Since
\[
\int_0^1\int_0^v (v-u)^{2H_i-1}\,du\,dv<\infty
\]
for $H_i>0$, the Hilbert--Schmidt norm of the difference is $O(h^\alpha)$. This implies
\eqref{eq:local-freeze}.
\end{proof}

\begin{definition}[Exact local transfer family]\label{def:local-transfer}
Assume Assumptions~\ref{ass:model} and~\ref{ass:obs-source}, let $\eta>0$, fix $t\in(0,T)$,
$\alpha\in(0,1)$, and $h\in(0,t]$. Define
\[
\mathfrak G_{t,h,\alpha}^{\mathrm{near}}
:=
\mathcal D_{t,h}^{(\alpha)}
\bigl(
\mathfrak G_t\cap H_{t,h,\alpha}^{\mathrm{near}}
\bigr)
\subset E.
\]
The \emph{exact local transfer operator} is
\begin{equation}\label{eq:local-transfer}
\mathfrak T_{t,h,\alpha}:\mathfrak G_{t,h,\alpha}^{\mathrm{near}}\subset E\to E,
\qquad
\mathfrak T_{t,h,\alpha}\phi
:=
h^{\alpha(H_Y-H_{XY})}
\mathcal D_{t,h}^{(\alpha)}
\mathcal S_t
\mathcal L_{t,h}^{(\alpha)}\phi.
\end{equation}
\end{definition}

\begin{proposition}[Exact localized transfer identity]
\label{prop:local-transfer}
Assume Assumptions~\ref{ass:model} and~\ref{ass:obs-source}, let $\eta>0$, fix $t\in(0,T)$,
$\alpha\in(0,1)$, and let $h>0$ be sufficiently small. Then:
\begin{enumerate}[label=(\roman*)]
\item $\mathfrak T_{t,h,\alpha}$ is a well-defined closed operator on $E$;
\item for every $\phi\in\mathfrak G_{t,h,\alpha}^{\mathrm{near}}$,
\begin{equation}\label{eq:local-transfer-eq}
\mathfrak A_{Y,t,h,\alpha}\bigl(\mathfrak T_{t,h,\alpha}\phi\bigr)
=
\mathfrak A_{XY,t,h,\alpha}\phi.
\end{equation}
\end{enumerate}
\end{proposition}

\begin{proof}
Well-definedness follows from Proposition~\ref{prop:support-preserve}.  If
$\phi\in\mathfrak G_{t,h,\alpha}^{\mathrm{near}}$, then
$w:=\mathcal L_{t,h}^{(\alpha)}\phi\in \mathfrak G_t\cap H_{t,h,\alpha}^{\mathrm{near}}$, so
$\mathcal S_t w\in H_{t,h,\alpha}^{\mathrm{near}}$ and
$\mathcal D_{t,h}^{(\alpha)}\mathcal S_t w\in E$.

To prove closedness, let $\phi_n\to\phi$ in $E$ and
$\mathfrak T_{t,h,\alpha}\phi_n\to\psi$ in $E$, with
$\phi_n\in\mathfrak G_{t,h,\alpha}^{\mathrm{near}}$. Applying the isometric lifts gives
\[
w_n:=\mathcal L_{t,h}^{(\alpha)}\phi_n\to w:=\mathcal L_{t,h}^{(\alpha)}\phi
\qquad\text{in }H_t
\]
and
\[
h^{\alpha(H_Y-H_{XY})}\mathcal S_t w_n
=
\mathcal L_{t,h}^{(\alpha)}\mathfrak T_{t,h,\alpha}\phi_n
\longrightarrow
\mathcal L_{t,h}^{(\alpha)}\psi
\qquad\text{in }H_t.
\]
Since $\mathcal S_t$ is closed, $w\in\mathfrak G_t$ and
$h^{\alpha(H_Y-H_{XY})}\mathcal S_t w=\mathcal L_{t,h}^{(\alpha)}\psi$. Since
$w=\mathcal L_{t,h}^{(\alpha)}\phi$, one already has $w\in H_{t,h,\alpha}^{\mathrm{near}}$,
so $\phi=\mathcal D_{t,h}^{(\alpha)}w\in\mathfrak G_{t,h,\alpha}^{\mathrm{near}}$ and
$\mathfrak T_{t,h,\alpha}\phi=\psi$. Thus $\mathfrak T_{t,h,\alpha}$ is closed.

Now let $\phi\in\mathfrak G_{t,h,\alpha}^{\mathrm{near}}$ and set
$w:=\mathcal L_{t,h}^{(\alpha)}\phi$. Set
\[
\beta:=H_Y-H_{XY}.
\]
By definition,
\[
\mathcal K_{Y,t}\mathcal S_t w=\mathcal K_{XY,t}w.
\]
Applying Proposition~\ref{prop:local-volterra} to both sides and using
$H_Y-H_{XY}=\beta$ gives
\begin{align*}
h^{\alpha H_Y}
\mathfrak A_{Y,t,h,\alpha}
\bigl(\mathcal D_{t,h}^{(\alpha)}\mathcal S_t w\bigr)
&=
h^{\alpha H_{XY}}
\mathfrak A_{XY,t,h,\alpha}
\bigl(\mathcal D_{t,h}^{(\alpha)}w\bigr) \\
&=
h^{\alpha H_{XY}}
\mathfrak A_{XY,t,h,\alpha}\phi.
\end{align*}
Multiplying by $h^{\alpha\beta-\alpha H_Y}=h^{-\alpha H_{XY}}$ yields
\[
\mathfrak A_{Y,t,h,\alpha}
\bigl(
h^{\alpha\beta}\mathcal D_{t,h}^{(\alpha)}\mathcal S_t w
\bigr)
=
\mathfrak A_{XY,t,h,\alpha}\phi,
\]
which is exactly \eqref{eq:local-transfer-eq}.
\end{proof}

\begin{definition}[Finite-interval right fractional integral and derivative]
\label{def:unit-fractional}
For $\lambda\in(0,1)$ and $\phi\in E=L^2([0,1])$, define
\begin{equation}\label{eq:Iminus1}
(\mathfrak I_{-,1}^\lambda \phi)(u)
:=
\frac{1}{\Gamma(\lambda)}
\int_u^1 (v-u)^{\lambda-1}\phi(v)\,dv,
\qquad 0\le u\le 1.
\end{equation}
For $\beta\in(0,1)$ and $\phi\in C_c^\infty((0,1))$, define
\begin{equation}\label{eq:Dminus1}
(\mathfrak D_{-,1}^\beta \phi)(u)
:=
-\frac{d}{du}
\left[
\frac{1}{\Gamma(1-\beta)}
\int_u^1 (v-u)^{-\beta}\phi(v)\,dv
\right],
\qquad 0<u<1.
\end{equation}
These are the right-sided Riemann--Liouville operators; see
Samko, Kilbas, and Marichev~\cite[Ch.~2]{SKM93}.
\end{definition}

\begin{lemma}[Unit-interval semigroup and injectivity]\label{lem:unit-semigroup}
Let $\lambda,\mu\in(0,1)$ with $\lambda+\mu\le 1$.
\begin{enumerate}[label=(\roman*)]
\item For every $\phi\in E$,
\begin{equation}\label{eq:unit-semigroup}
\mathfrak I_{-,1}^\lambda \mathfrak I_{-,1}^\mu \phi
=
\mathfrak I_{-,1}^{\lambda+\mu}\phi.
\end{equation}
\item For every $H\in(0,\tfrac12)$, the operator $\mathfrak J_H$ of
Definition~\ref{def:unit-model} satisfies
\begin{equation}\label{eq:JH-fractional}
\mathfrak J_H
=
\Gamma\!\left(H+\frac12\right)\mathfrak I_{-,1}^{H+\frac12},
\end{equation}
and is injective on $E$.
\end{enumerate}
\end{lemma}

\begin{proof}
Let $\phi\in E$. For almost every $u\in[0,1]$, Fubini's theorem gives
\begin{align*}
(\mathfrak I_{-,1}^\lambda \mathfrak I_{-,1}^\mu \phi)(u)
&=
\frac{1}{\Gamma(\lambda)\Gamma(\mu)}
\int_u^1\int_r^1
(r-u)^{\lambda-1}(v-r)^{\mu-1}\phi(v)\,dv\,dr \\
&=
\frac{1}{\Gamma(\lambda)\Gamma(\mu)}
\int_u^1
\left(
\int_u^v (r-u)^{\lambda-1}(v-r)^{\mu-1}\,dr
\right)\phi(v)\,dv.
\end{align*}
With $r=u+(v-u)\theta$, the inner integral becomes
\[
(v-u)^{\lambda+\mu-1}
\int_0^1 \theta^{\lambda-1}(1-\theta)^{\mu-1}\,d\theta
=
(v-u)^{\lambda+\mu-1}
\frac{\Gamma(\lambda)\Gamma(\mu)}{\Gamma(\lambda+\mu)}.
\]
Substituting this identity proves \eqref{eq:unit-semigroup}.

Identity \eqref{eq:JH-fractional} follows from the definitions with
$\lambda=H+\tfrac12$. To prove injectivity, let $\mathfrak J_H\phi=0$. Set
$\lambda:=H+\tfrac12\in(0,1)$. By \eqref{eq:JH-fractional},
\[
\mathfrak I_{-,1}^\lambda \phi=0.
\]
Applying \eqref{eq:unit-semigroup} with $\mu=1-\lambda$ yields
\[
\mathfrak I_{-,1}^{1-\lambda}\mathfrak I_{-,1}^\lambda \phi
=
\mathfrak I_{-,1}^1\phi
=
0.
\]
But
\[
(\mathfrak I_{-,1}^1\phi)(u)=\int_u^1 \phi(v)\,dv,
\]
so differentiating almost everywhere gives $\phi=0$ almost everywhere. Thus $\mathfrak J_H$
is injective.
\end{proof}

\begin{definition}[Frozen transfer operator on the unit interval]
\label{def:frozen-transfer}
Assume
\[
0<H_{XY}<H_Y<\tfrac12,
\qquad
\beta:=H_Y-H_{XY}.
\]
Fix $t\in(0,T)$. Define the domain
\begin{equation}\label{eq:frozen-domain}
\mathcal D(\mathfrak T_t^{\mathrm{fr}})
:=
\left\{
\phi\in E:\mathfrak J_{H_{XY}}\phi\in\mathrm{Ran}(\mathfrak J_{H_Y})
\right\},
\end{equation}
and the frozen transfer operator
\begin{equation}\label{eq:frozen-transfer}
\mathfrak T_t^{\mathrm{fr}}:\mathcal D(\mathfrak T_t^{\mathrm{fr}})\subset E\to E,
\qquad
\mathfrak T_t^{\mathrm{fr}}\phi
:=
\frac{L_{XY}(t,t)}{L_Y(t,t)}
\mathfrak J_{H_Y}^{-1}\mathfrak J_{H_{XY}}\phi.
\end{equation}
\end{definition}

\begin{proposition}[Finite-interval fractional representation of the frozen transfer]
\label{prop:frozen-fractional}
Assume
\[
0<H_{XY}<H_Y<\tfrac12,
\qquad
\beta:=H_Y-H_{XY},
\]
and fix $t\in(0,T)$. Then:
\begin{enumerate}[label=(\roman*)]
\item for every $\phi\in\mathcal D(\mathfrak T_t^{\mathrm{fr}})$,
\begin{equation}\label{eq:frozen-identity}
L_Y(t,t)\mathfrak J_{H_Y}\bigl(\mathfrak T_t^{\mathrm{fr}}\phi\bigr)
=
L_{XY}(t,t)\mathfrak J_{H_{XY}}\phi;
\end{equation}
\item for every $\phi\in C_c^\infty((0,1))$,
\begin{equation}\label{eq:frozen-derivative}
\mathfrak T_t^{\mathrm{fr}}\phi
=
\mathbf c_t\,\mathfrak D_{-,1}^{\beta}\phi,
\end{equation}
where $\mathbf c_t$ is the coefficient from \eqref{eq:model-coeff};
\item if $\widetilde\phi$ denotes the zero extension of $\phi\in C_c^\infty((0,1))$ to $\R$,
then
\begin{equation}\label{eq:frozen-wholeline}
\mathfrak D_{-,1}^{\beta}\phi
=
\bigl(\mathfrak D_-^\beta \widetilde\phi\bigr)\big|_{(0,1)}.
\end{equation}
\end{enumerate}
\end{proposition}

\begin{proof}
Part (i) is exactly the definition of $\mathfrak T_t^{\mathrm{fr}}$ and injectivity of
$\mathfrak J_{H_Y}$ from Lemma~\ref{lem:unit-semigroup}.

For part (ii), let $\phi\in C_c^\infty((0,1))$ and set
\[
\psi:=\mathbf c_t\,\mathfrak D_{-,1}^{\beta}\phi.
\]
Since $\phi$ is compactly supported away from the boundary, $\psi\in E$. Using
\eqref{eq:JH-fractional}, the semigroup identity \eqref{eq:unit-semigroup}, and the left-inverse
property
\[
\mathfrak I_{-,1}^{1-\beta}\mathfrak I_{-,1}^{\beta}= \mathfrak I_{-,1}^{1},
\qquad
-\frac{d}{du}\mathfrak I_{-,1}^{1}f=f
\quad\text{a.e.},
\]
we obtain
\[
\mathfrak I_{-,1}^{H_Y+\frac12}\mathfrak D_{-,1}^{\beta}\phi
=
\mathfrak I_{-,1}^{H_{XY}+\frac12}\phi.
\]
Multiplying by the Gamma factors gives
\begin{align*}
L_Y(t,t)\mathfrak J_{H_Y}\psi
&=
L_Y(t,t)\mathbf c_t\Gamma\!\left(H_Y+\frac12\right)
\mathfrak I_{-,1}^{H_Y+\frac12}\mathfrak D_{-,1}^{\beta}\phi \\
&=
L_{XY}(t,t)\Gamma\!\left(H_{XY}+\frac12\right)
\mathfrak I_{-,1}^{H_{XY}+\frac12}\phi \\
&=
L_{XY}(t,t)\mathfrak J_{H_{XY}}\phi.
\end{align*}
By part (i) and injectivity of $\mathfrak J_{H_Y}$, this proves \eqref{eq:frozen-derivative}.

For part (iii), if $\widetilde\phi$ is the zero extension of $\phi$, then for $u\in(0,1)$,
\[
\frac{1}{\Gamma(1-\beta)}\int_u^\infty (v-u)^{-\beta}\widetilde\phi(v)\,dv
=
\frac{1}{\Gamma(1-\beta)}\int_u^1 (v-u)^{-\beta}\phi(v)\,dv.
\]
Differentiating yields \eqref{eq:frozen-wholeline}.
\end{proof}

\begin{proposition}[Weak local convergence under a uniform local bound]
\label{prop:local-weak}
Assume Assumptions~\ref{ass:model} and~\ref{ass:obs-source}, let $\eta>0$, fix $t\in(0,T)$
and $\alpha\in(0,1)$, and suppose
\[
0<H_{XY}<H_Y<\tfrac12.
\]
Let $\phi\in C_c^\infty((0,1))$. Assume that:
\begin{enumerate}[label=(\roman*)]
\item \(\phi\in \mathfrak G_{t,h,\alpha}^{\mathrm{near}}\) for all sufficiently small \(h>0\);
\item there exists \(C_\phi<\infty\) such that
\begin{equation}\label{eq:local-L2-bound}
\|\mathfrak T_{t,h,\alpha}\phi\|_E\le C_\phi
\end{equation}
for all sufficiently small \(h>0\).
\end{enumerate}
Then
\begin{equation}\label{eq:local-weak}
\mathfrak T_{t,h,\alpha}\phi
\rightharpoonup
\mathfrak T_t^{\mathrm{fr}}\phi
\qquad\text{weakly in }E=L^2([0,1])
\end{equation}
as \(h\downarrow0\).
\end{proposition}

\begin{proof}
Let \(h_n\downarrow0\). By \eqref{eq:local-L2-bound}, the sequence
\((\mathfrak T_{t,h_n,\alpha}\phi)\) is bounded in \(E\), so after passing to a subsequence we
may assume
\[
\mathfrak T_{t,h_n,\alpha}\phi \rightharpoonup \psi
\qquad\text{weakly in }E
\]
for some \(\psi\in E\).

By Proposition~\ref{prop:local-transfer},
\[
\mathfrak A_{Y,t,h_n,\alpha}\bigl(\mathfrak T_{t,h_n,\alpha}\phi\bigr)
=
\mathfrak A_{XY,t,h_n,\alpha}\phi.
\]
The right-hand side converges strongly in \(E\) to
\[
L_{XY}(t,t)\mathfrak J_{H_{XY}}\phi
\]
by Proposition~\ref{prop:local-volterra}. On the left-hand side, decompose
\begin{align*}
\mathfrak A_{Y,t,h_n,\alpha}\bigl(\mathfrak T_{t,h_n,\alpha}\phi\bigr)
&=
\bigl(\mathfrak A_{Y,t,h_n,\alpha}-L_Y(t,t)\mathfrak J_{H_Y}\bigr)
\bigl(\mathfrak T_{t,h_n,\alpha}\phi\bigr) \\
&\quad+
L_Y(t,t)\mathfrak J_{H_Y}\bigl(\mathfrak T_{t,h_n,\alpha}\phi\bigr).
\end{align*}
The first term tends to \(0\) strongly in \(E\) because
\eqref{eq:local-freeze} gives operator-norm convergence and the sequence
\((\mathfrak T_{t,h_n,\alpha}\phi)\) is bounded in \(E\). Therefore
\[
L_Y(t,t)\mathfrak J_{H_Y}\bigl(\mathfrak T_{t,h_n,\alpha}\phi\bigr)
\longrightarrow
L_{XY}(t,t)\mathfrak J_{H_{XY}}\phi
\qquad\text{strongly in }E.
\]
Since \(\mathfrak J_{H_Y}\) is bounded on \(E\), weak convergence of
\(\mathfrak T_{t,h_n,\alpha}\phi\) implies
\[
\mathfrak J_{H_Y}\bigl(\mathfrak T_{t,h_n,\alpha}\phi\bigr)
\rightharpoonup
\mathfrak J_{H_Y}\psi
\qquad\text{weakly in }E.
\]
Because the same sequence converges strongly, its weak limit must equal its strong limit, so
\[
L_Y(t,t)\mathfrak J_{H_Y}\psi
=
L_{XY}(t,t)\mathfrak J_{H_{XY}}\phi.
\]
By Proposition~\ref{prop:frozen-fractional}(i) and injectivity of \(\mathfrak J_{H_Y}\), this
implies
\[
\psi=\mathfrak T_t^{\mathrm{fr}}\phi.
\]
Thus every weakly convergent subsequence has the same weak limit, and \eqref{eq:local-weak}
follows.
\end{proof}

\begin{definition}[Local and frozen graph norms]\label{def:local-graph-norm}
Assume Assumptions~\ref{ass:model} and~\ref{ass:obs-source}, let \(\eta>0\), fix \(t\in(0,T)\),
let \(\alpha\in(0,1)\), and suppose
\[
0<H_{XY}<H_Y<\tfrac12.
\]
For \(h>0\) sufficiently small and \(\phi\in\mathfrak G_{t,h,\alpha}^{\mathrm{near}}\), define
\begin{equation}\label{eq:local-graph-norm}
\|\phi\|_{\mathfrak G_{t,h,\alpha}^{\mathrm{loc}}}^2
:=
\|\phi\|_E^2
+\frac{\eta^2}{\nu_Y^2}\|\mathfrak T_{t,h,\alpha}\phi\|_E^2.
\end{equation}
For \(\phi\in\mathcal D(\mathfrak T_t^{\mathrm{fr}})\), define
\begin{equation}\label{eq:frozen-graph-norm}
\|\phi\|_{\mathfrak G_t^{\mathrm{fr}}}^2
:=
\|\phi\|_E^2
+\frac{\eta^2}{\nu_Y^2}\|\mathfrak T_t^{\mathrm{fr}}\phi\|_E^2.
\end{equation}
\end{definition}

\begin{proposition}[Weak graph compactness of the exact local transfer family]
\label{prop:local-graph-compact}
Assume Assumptions~\ref{ass:model} and~\ref{ass:obs-source}, let \(\eta>0\), fix \(t\in(0,T)\)
and \(\alpha\in(0,1)\), and suppose
\[
0<H_{XY}<H_Y<\tfrac12.
\]
Let \(h_n\downarrow0\) and let
\[
\phi_n\in\mathfrak G_{t,h_n,\alpha}^{\mathrm{near}}
\]
be such that
\begin{equation}\label{eq:local-graph-bounded}
\sup_{n\ge1}\|\phi_n\|_{\mathfrak G_{t,h_n,\alpha}^{\mathrm{loc}}}<\infty.
\end{equation}
Then there exist a subsequence, still denoted \((\phi_n)\), and some
\[
\phi\in\mathcal D(\mathfrak T_t^{\mathrm{fr}})
\]
such that
\begin{align}
\phi_n &\rightharpoonup \phi
\qquad\text{weakly in }E, \label{eq:local-graph-weak-input}\\
\mathfrak T_{t,h_n,\alpha}\phi_n &\rightharpoonup \mathfrak T_t^{\mathrm{fr}}\phi
\qquad\text{weakly in }E. \label{eq:local-graph-weak-output}
\end{align}
Moreover,
\begin{equation}\label{eq:local-graph-liminf}
\|\phi\|_{\mathfrak G_t^{\mathrm{fr}}}
\le
\liminf_{n\to\infty}\|\phi_n\|_{\mathfrak G_{t,h_n,\alpha}^{\mathrm{loc}}}.
\end{equation}
\end{proposition}

\begin{proof}
By \eqref{eq:local-graph-bounded} and Definition~\ref{def:local-graph-norm}, both sequences
\((\phi_n)\) and \((\mathfrak T_{t,h_n,\alpha}\phi_n)\) are bounded in \(E\). Since \(E\) is a
Hilbert space, after passing to a subsequence we may assume
\[
\phi_n\rightharpoonup \phi,
\qquad
\mathfrak T_{t,h_n,\alpha}\phi_n\rightharpoonup \psi
\qquad\text{weakly in }E
\]
for some \(\phi,\psi\in E\).

By Proposition~\ref{prop:local-transfer},
\[
\mathfrak A_{Y,t,h_n,\alpha}\bigl(\mathfrak T_{t,h_n,\alpha}\phi_n\bigr)
=
\mathfrak A_{XY,t,h_n,\alpha}\phi_n.
\]
We claim that
\begin{align}
\mathfrak A_{XY,t,h_n,\alpha}\phi_n
&\rightharpoonup
L_{XY}(t,t)\mathfrak J_{H_{XY}}\phi
\qquad\text{weakly in }E, \label{eq:graph-compact-xy}\\
\mathfrak A_{Y,t,h_n,\alpha}\bigl(\mathfrak T_{t,h_n,\alpha}\phi_n\bigr)
&\rightharpoonup
L_Y(t,t)\mathfrak J_{H_Y}\psi
\qquad\text{weakly in }E. \label{eq:graph-compact-y}
\end{align}
The first convergence is a direct consequence of Proposition~\ref{prop:local-volterra}.  The
decomposition
\[
\mathfrak A_{XY,t,h_n,\alpha}
=
\bigl(\mathfrak A_{XY,t,h_n,\alpha}-L_{XY}(t,t)\mathfrak J_{H_{XY}}\bigr)
+L_{XY}(t,t)\mathfrak J_{H_{XY}}
\]
has an operator-norm error term tending to \(0\) by \eqref{eq:local-freeze}. Since
\((\phi_n)\) is bounded in \(E\), that error contributes strongly, hence weakly, to
\(0\); the frozen term converges weakly to \(L_{XY}(t,t)\mathfrak J_{H_{XY}}\phi\)
because \(\mathfrak J_{H_{XY}}\) is bounded on \(E\). This proves
\eqref{eq:graph-compact-xy}. The proof of \eqref{eq:graph-compact-y} is identical.

Passing to the weak limit in the exact identity therefore gives
\[
L_Y(t,t)\mathfrak J_{H_Y}\psi
=
L_{XY}(t,t)\mathfrak J_{H_{XY}}\phi.
\]
Since \(L_Y(t,t)>0\) by Assumption~\ref{ass:model}, this shows that
\(\mathfrak J_{H_{XY}}\phi\in\mathrm{Ran}(\mathfrak J_{H_Y})\), hence
\(\phi\in\mathcal D(\mathfrak T_t^{\mathrm{fr}})\) by Definition~\ref{def:frozen-transfer}. By
Proposition~\ref{prop:frozen-fractional}(i) and injectivity of \(\mathfrak J_{H_Y}\),
\[
\psi=\mathfrak T_t^{\mathrm{fr}}\phi.
\]
This proves \eqref{eq:local-graph-weak-input} and \eqref{eq:local-graph-weak-output}.

Finally, weak lower semicontinuity of the norm in \(E\) yields
\[
\|\phi\|_E\le \liminf_{n\to\infty}\|\phi_n\|_E,
\qquad
\|\mathfrak T_t^{\mathrm{fr}}\phi\|_E
=
\|\psi\|_E
\le
\liminf_{n\to\infty}\|\mathfrak T_{t,h_n,\alpha}\phi_n\|_E.
\]
Squaring both inequalities, adding them with weight \(\eta^2/\nu_Y^2\), and using
\[
\liminf_{n\to\infty}(a_n+b_n)\ge
\liminf_{n\to\infty}a_n+\liminf_{n\to\infty}b_n
\]
for nonnegative sequences gives
\[
\|\phi\|_{\mathfrak G_t^{\mathrm{fr}}}^2
\le
\liminf_{n\to\infty}\|\phi_n\|_{\mathfrak G_{t,h_n,\alpha}^{\mathrm{loc}}}^2.
\]
Taking square roots proves \eqref{eq:local-graph-liminf}.
\end{proof}

\begin{corollary}[Unit local graph balls converge to the frozen graph]
\label{cor:local-unit-ball}
Under the assumptions of Proposition~\ref{prop:local-graph-compact}, let \(h_n\downarrow0\) and
let
\[
\phi_n\in\mathfrak G_{t,h_n,\alpha}^{\mathrm{near}}
\]
satisfy
\[
\|\phi_n\|_{\mathfrak G_{t,h_n,\alpha}^{\mathrm{loc}}}\le 1
\qquad\text{for every }n.
\]
Then some subsequence converges weakly in \(E\) to an element
\[
\phi\in\mathcal D(\mathfrak T_t^{\mathrm{fr}})
\]
with
\[
\|\phi\|_{\mathfrak G_t^{\mathrm{fr}}}\le 1.
\]
\end{corollary}

\begin{proof}
Apply Proposition~\ref{prop:local-graph-compact} to the sequence \((\phi_n)\). The unit bound gives
weak compactness of the graph pairs in \(E\times E\), and the liminf inequality in that proposition
places the weak limit in the frozen graph with frozen graph norm at most one.
\end{proof}

\begin{definition}[Local graph recovery core]\label{def:local-recovery-core}
Assume Assumptions~\ref{ass:model} and~\ref{ass:obs-source}, let \(\eta>0\), fix \(t\in(0,T)\)
and \(\alpha\in(0,1)\), and suppose
\[
0<H_{XY}<H_Y<\tfrac12.
\]
A linear subspace
\[
\mathcal C_{t,\alpha}\subset \mathcal D(\mathfrak T_t^{\mathrm{fr}})
\]
is called a \emph{local graph recovery core} if:
\begin{enumerate}[label=(\roman*)]
\item \(\mathcal C_{t,\alpha}\) is dense in \(\mathcal D(\mathfrak T_t^{\mathrm{fr}})\) with respect
to the frozen graph norm \(\|\cdot\|_{\mathfrak G_t^{\mathrm{fr}}}\);
\item for every \(\phi\in\mathcal C_{t,\alpha}\), there exists \(h_\phi\in(0,t]\) such that
\[
\phi\in \mathfrak G_{t,h,\alpha}^{\mathrm{near}}
\qquad\text{for all }0<h\le h_\phi;
\]
\item for every \(\phi\in\mathcal C_{t,\alpha}\),
\begin{equation}\label{eq:recovery-strong}
\|\mathfrak T_{t,h,\alpha}\phi-\mathfrak T_t^{\mathrm{fr}}\phi\|_E\longrightarrow 0
\qquad\text{as }h\downarrow0.
\end{equation}
\end{enumerate}
\end{definition}

\begin{theorem}[Recovery sequences from a local graph recovery core]
\label{thm:local-recovery}
Assume the hypotheses of Proposition~\ref{prop:local-graph-compact}, and suppose that there
exists a local graph recovery core \(\mathcal C_{t,\alpha}\) in the sense of
Definition~\ref{def:local-recovery-core}. Let \(h_n\downarrow0\), and let
\[
\phi\in \mathcal D(\mathfrak T_t^{\mathrm{fr}}).
\]
Then there exists a sequence
\[
\phi_n\in \mathfrak G_{t,h_n,\alpha}^{\mathrm{near}}
\]
such that
\begin{align}
\|\phi_n-\phi\|_E &\longrightarrow 0, \label{eq:recovery-input}\\
\|\mathfrak T_{t,h_n,\alpha}\phi_n-\mathfrak T_t^{\mathrm{fr}}\phi\|_E
&\longrightarrow 0, \label{eq:recovery-output}
\end{align}
and consequently
\begin{equation}\label{eq:recovery-norm}
\|\phi_n\|_{\mathfrak G_{t,h_n,\alpha}^{\mathrm{loc}}}
\longrightarrow
\|\phi\|_{\mathfrak G_t^{\mathrm{fr}}}.
\end{equation}
\end{theorem}

\begin{proof}
By Definition~\ref{def:local-recovery-core}(i), choose a sequence
\[
\psi_m\in\mathcal C_{t,\alpha},
\qquad m\ge1,
\]
such that
\begin{equation}\label{eq:recovery-core-approx}
\|\psi_m-\phi\|_E+\|\mathfrak T_t^{\mathrm{fr}}\psi_m-\mathfrak T_t^{\mathrm{fr}}\phi\|_E
\le 2^{-m}
\qquad\text{for every }m\ge1.
\end{equation}
For each \(m\), Definition~\ref{def:local-recovery-core}(ii)--(iii) yields a number
\(\eta_m\in(0,t]\) such that for every \(0<h\le \eta_m\),
\begin{equation}\label{eq:recovery-core-smallh}
\psi_m\in \mathfrak G_{t,h,\alpha}^{\mathrm{near}}
\qquad\text{and}\qquad
\|\mathfrak T_{t,h,\alpha}\psi_m-\mathfrak T_t^{\mathrm{fr}}\psi_m\|_E\le 2^{-m}.
\end{equation}
After shrinking if necessary, we may assume that \((\eta_m)\) is nonincreasing.

For each \(n\), define
\[
m(n):=\max\{m\in\{1,\dots,n\}: h_n\le \eta_m\}.
\]
This set is nonempty for all sufficiently large \(n\) because \(h_n\downarrow0\) and
\(\eta_1>0\). Moreover, \(m(n)\to\infty\): for every fixed \(M\), one has
\(h_n\le \eta_M\) for all sufficiently large \(n\), hence \(m(n)\ge M\) eventually. Set
\[
\phi_n:=\psi_{m(n)}
\]
for all sufficiently large \(n\), and define the remaining finitely many terms by \(\phi_n:=0\).
By construction,
\(\phi_n\in \mathfrak G_{t,h_n,\alpha}^{\mathrm{near}}\) for every \(n\).

Using \eqref{eq:recovery-core-approx},
\[
\|\phi_n-\phi\|_E
=
\|\psi_{m(n)}-\phi\|_E
\le 2^{-m(n)}
\longrightarrow 0,
\]
which proves \eqref{eq:recovery-input}. Likewise, by
\eqref{eq:recovery-core-smallh} and \eqref{eq:recovery-core-approx},
\begin{align*}
\|\mathfrak T_{t,h_n,\alpha}\phi_n-\mathfrak T_t^{\mathrm{fr}}\phi\|_E
&\le
\|\mathfrak T_{t,h_n,\alpha}\psi_{m(n)}-\mathfrak T_t^{\mathrm{fr}}\psi_{m(n)}\|_E \\
&\quad+
\|\mathfrak T_t^{\mathrm{fr}}\psi_{m(n)}-\mathfrak T_t^{\mathrm{fr}}\phi\|_E \\
&\le 2^{-m(n)}+2^{-m(n)}
\longrightarrow 0,
\end{align*}
which proves \eqref{eq:recovery-output}.

Finally,
\[
\|\phi_n\|_E\to\|\phi\|_E,
\qquad
\|\mathfrak T_{t,h_n,\alpha}\phi_n\|_E\to\|\mathfrak T_t^{\mathrm{fr}}\phi\|_E,
\]
by \eqref{eq:recovery-input} and \eqref{eq:recovery-output}. Substituting into
Definition~\ref{def:local-graph-norm} and \eqref{eq:frozen-graph-norm} proves
\eqref{eq:recovery-norm}.
\end{proof}

\begin{corollary}[Sequential Mosco convergence of the local transfer graphs]
\label{cor:local-mosco}
Assume the hypotheses of Theorem~\ref{thm:local-recovery}. For \(h>0\), define the graph
\[
\Gamma_{t,h,\alpha}^{\mathrm{loc}}
:=
\left\{
(\phi,\mathfrak T_{t,h,\alpha}\phi):
\phi\in \mathfrak G_{t,h,\alpha}^{\mathrm{near}}
\right\}
\subset E\times E,
\]
and define the frozen graph
\[
\Gamma_t^{\mathrm{fr}}
:=
\left\{
(\phi,\mathfrak T_t^{\mathrm{fr}}\phi):
\phi\in \mathcal D(\mathfrak T_t^{\mathrm{fr}})
\right\}.
\]
Then along every sequence \(h_n\downarrow0\), the graphs
\(\Gamma_{t,h_n,\alpha}^{\mathrm{loc}}\) converge to \(\Gamma_t^{\mathrm{fr}}\) in the sequential
Mosco sense~\cite{Mosco69}:
\begin{enumerate}[label=(\roman*)]
\item if \((\phi_n,\mathfrak T_{t,h_n,\alpha}\phi_n)\in \Gamma_{t,h_n,\alpha}^{\mathrm{loc}}\) is
bounded in \(E\times E\), then some subsequence converges weakly in \(E\times E\) to an element
of \(\Gamma_t^{\mathrm{fr}}\);
\item every \((\phi,\mathfrak T_t^{\mathrm{fr}}\phi)\in \Gamma_t^{\mathrm{fr}}\) admits a sequence
\((\phi_n,\mathfrak T_{t,h_n,\alpha}\phi_n)\in \Gamma_{t,h_n,\alpha}^{\mathrm{loc}}\) converging
strongly in \(E\times E\) to \((\phi,\mathfrak T_t^{\mathrm{fr}}\phi)\).
\end{enumerate}
\end{corollary}

\begin{proof}
For item~(i), boundedness in \(E\times E\) implies boundedness of
\(\|\phi_n\|_{\mathfrak G_{t,h_n,\alpha}^{\mathrm{loc}}}\) by
Definition~\ref{def:local-graph-norm}. Proposition~\ref{prop:local-graph-compact} therefore
gives a subsequence and a limit \(\phi\in\mathcal D(\mathfrak T_t^{\mathrm{fr}})\) such that
\[
\phi_n\rightharpoonup \phi,
\qquad
\mathfrak T_{t,h_n,\alpha}\phi_n\rightharpoonup \mathfrak T_t^{\mathrm{fr}}\phi
\]
weakly in \(E\). This is weak convergence in \(E\times E\) toward an element of
\(\Gamma_t^{\mathrm{fr}}\).

Item~(ii) is exactly Theorem~\ref{thm:local-recovery}.
\end{proof}

\section{Strong local convergence on the smooth core}\label{sec:smooth-core}

The local graph limit must be populated by recovery sequences.  For the exact Volterra operators,
strong local convergence can be verified on the smooth core \(C_c^\infty((0,1))\).  This does not
prove graph density of that core in the frozen graph norm.  It proves the analytic part of the
recovery-core hypothesis: exact local solvability and strong convergence for fixed smooth profiles.

\begin{proposition}[Conjugated target-side operator]\label{prop:BY}
Assume Assumption~\ref{ass:model}, fix \(t\in(0,T)\) and \(\alpha\in(0,1)\), and set
\[
\lambda:=H_Y+\frac12\in(0,1).
\]
For \(h>0\), define
\[
a_{Y,t,h,\alpha}(u,v):=
L_Y(t-h^\alpha u,t-h^\alpha v),
\qquad
0\le u\le v\le 1.
\]
Then there exists a bounded operator
\[
\mathfrak B_{Y,t,h,\alpha}\in\mathcal L(E)
\]
such that:
\begin{enumerate}[label=(\roman*)]
\item for every \(\psi\in E\), the function
\(\mathfrak I_{-,1}^{1-\lambda}(\mathfrak A_{Y,t,h,\alpha}\psi)\) is absolutely continuous on
\([0,1]\), and
\begin{equation}\label{eq:BY-conj}
\frac{1}{\Gamma(\lambda)}
\mathfrak D_{-,1}^{\lambda}\bigl(\mathfrak A_{Y,t,h,\alpha}\psi\bigr)
=
\mathfrak B_{Y,t,h,\alpha}\psi;
\end{equation}
\item there exists \(C<\infty\), depending only on \(t,\alpha\), such that
\begin{equation}\label{eq:BY-close}
\left\|
\mathfrak B_{Y,t,h,\alpha}-L_Y(t,t)I_E
\right\|_{\mathcal L(E)}
\le
Ch^\alpha
\end{equation}
for all sufficiently small \(h>0\).
\end{enumerate}
\end{proposition}

\begin{proof}
For \(\psi\in E\), set
\[
f_h:=\mathfrak A_{Y,t,h,\alpha}\psi.
\]
Using \eqref{eq:Iminus1} with order \(1-\lambda\), Fubini's theorem gives for almost every
\(u\in[0,1]\),
\begin{align*}
(\mathfrak I_{-,1}^{1-\lambda}f_h)(u)
&=
\frac{1}{\Gamma(1-\lambda)}
\int_u^1\int_r^1
(r-u)^{-\lambda}(v-r)^{\lambda-1}
a_{Y,t,h,\alpha}(r,v)\psi(v)\,dv\,dr \\
&=
\int_u^1 F_{Y,t,h,\alpha}(u,v)\psi(v)\,dv,
\end{align*}
where
\[
F_{Y,t,h,\alpha}(u,v)
:=
\frac{1}{\Gamma(1-\lambda)}
\int_0^1 s^{-\lambda}(1-s)^{\lambda-1}
a_{Y,t,h,\alpha}(u+(v-u)s,v)\,ds.
\]
Since \(L_Y\) is \(C^1\), the function \(F_{Y,t,h,\alpha}\) is continuous on \(\Delta_1\), and
\[
\partial_1F_{Y,t,h,\alpha}(u,v)
=
\frac{1}{\Gamma(1-\lambda)}
\int_0^1 s^{-\lambda}(1-s)^\lambda
\partial_1 a_{Y,t,h,\alpha}(u+(v-u)s,v)\,ds
\]
exists and is bounded. Therefore \(u\mapsto (\mathfrak I_{-,1}^{1-\lambda}f_h)(u)\) is
absolutely continuous and differentiation under the integral sign yields
\[
\mathfrak D_{-,1}^{\lambda}f_h(u)
=
\Gamma(\lambda)a_{Y,t,h,\alpha}(u,u)\psi(u)
-\int_u^1 G_{Y,t,h,\alpha}(u,v)\psi(v)\,dv
\]
for almost every \(u\in[0,1]\), where
\[
G_{Y,t,h,\alpha}(u,v)
:=
\frac{\Gamma(\lambda)}{\Gamma(\lambda)\Gamma(1-\lambda)}
\int_0^1 s^{-\lambda}(1-s)^\lambda
\partial_1 a_{Y,t,h,\alpha}(u+(v-u)s,v)\,ds.
\]
Define
\[
\begin{aligned}
(\mathfrak B_{Y,t,h,\alpha}\psi)(u)
&:=
a_{Y,t,h,\alpha}(u,u)\psi(u) \\
&\quad
-\frac{1}{\Gamma(\lambda)\Gamma(1-\lambda)}
\int_u^1\int_0^1 s^{-\lambda}(1-s)^\lambda \\
&\qquad\qquad\qquad \times
\partial_1 a_{Y,t,h,\alpha}(u+(v-u)s,v)\,ds\,\psi(v)\,dv.
\end{aligned}
\]
This proves \eqref{eq:BY-conj}.

To prove \eqref{eq:BY-close}, note first that
\[
|a_{Y,t,h,\alpha}(u,u)-L_Y(t,t)|\le Ch^\alpha
\qquad\text{for }0\le u\le 1,
\]
because \(L_Y\) is \(C^1\). Also,
\[
\partial_1 a_{Y,t,h,\alpha}(x,v)
=
-h^\alpha \partial_1L_Y(t-h^\alpha x,t-h^\alpha v),
\]
so \(|\partial_1 a_{Y,t,h,\alpha}(x,v)|\le Ch^\alpha\) uniformly on \(\Delta_1\). Therefore
\(\mathfrak B_{Y,t,h,\alpha}-a_{Y,t,h,\alpha}(\cdot,\cdot)_{\mathrm{diag}}I_E\) has integral
kernel bounded by \(Ch^\alpha\one_{\{0\le u\le v\le 1\}}\). This kernel has Hilbert--Schmidt norm
\(O(h^\alpha)\), which proves \eqref{eq:BY-close}.
\end{proof}

\begin{corollary}[Uniform invertibility of the conjugated target operator]
\label{cor:BY-invertible}
Under the assumptions of Proposition~\ref{prop:BY}, there exist \(h_0\in(0,t]\) and
\(C_t<\infty\) such that for every \(0<h\le h_0\),
\[
\mathfrak B_{Y,t,h,\alpha}:E\to E
\]
is invertible and
\begin{equation}\label{eq:BY-inv}
\|\mathfrak B_{Y,t,h,\alpha}^{-1}\|_{\mathcal L(E)}\le C_t.
\end{equation}
\end{corollary}

\begin{proof}
By Assumption~\ref{ass:model}, \(L_Y(t,t)>0\). Proposition~\ref{prop:BY} gives
\[
\left\|
\frac{1}{L_Y(t,t)}\mathfrak B_{Y,t,h,\alpha}-I_E
\right\|_{\mathcal L(E)}
\le
\frac{C}{L_Y(t,t)}h^\alpha.
\]
For \(h\) small enough the right-hand side is \(<1/2\), so the Neumann series gives
invertibility and \eqref{eq:BY-inv}.
\end{proof}

\begin{proposition}[Conjugated source asymptotics on the smooth core]
\label{prop:q-smooth}
Assume Assumption~\ref{ass:model}, fix \(t\in(0,T)\) and \(\alpha\in(0,1)\), and suppose
\[
0<H_{XY}<H_Y<\tfrac12.
\]
Set
\[
\lambda:=H_Y+\frac12,
\qquad
\mu:=H_{XY}+\frac12,
\qquad
\beta:=H_Y-H_{XY}.
\]
Let \(\phi\in C_c^\infty((0,1))\). Then there exists \(C_\phi<\infty\) such that for all
sufficiently small \(h>0\),
\begin{equation}\label{eq:q-smooth}
\left\|
\frac{1}{\Gamma(\lambda)}
\mathfrak D_{-,1}^{\lambda}\bigl(\mathfrak A_{XY,t,h,\alpha}\phi\bigr)
-L_Y(t,t)\mathfrak T_t^{\mathrm{fr}}\phi
\right\|_E
\le
C_\phi h^\alpha.
\end{equation}
\end{proposition}

\begin{proof}
Define
\[
a_{XY,t,h,\alpha}(u,v):=
L_{XY}(t-h^\alpha u,t-h^\alpha v),
\qquad
0\le u\le v\le 1.
\]
Using \eqref{eq:Iminus1} and the change of variables \(r=u+(v-u)s\), one obtains for almost
every \(u\in[0,1]\),
\begin{equation}\label{eq:q-representation}
\frac{1}{\Gamma(\lambda)}
\mathfrak D_{-,1}^{\lambda}\bigl(\mathfrak A_{XY,t,h,\alpha}\phi\bigr)(u)
=
-\frac{d}{du}\int_u^1 (v-u)^{-\beta}g_{t,h,\alpha}(u,v)\phi(v)\,dv,
\end{equation}
where
\[
g_{t,h,\alpha}(u,v)
:=
\frac{1}{\Gamma(\lambda)\Gamma(1-\lambda)}
\int_0^1 s^{-\lambda}(1-s)^{\mu-1}
a_{XY,t,h,\alpha}(u+(v-u)s,v)\,ds.
\]
Because \(L_{XY}\) is \(C^1\), one has
\begin{equation}\label{eq:g-close}
\|g_{t,h,\alpha}-g_t^{\mathrm{fr}}\|_{C^1(\Delta_1)}
\le
Ch^\alpha,
\end{equation}
where
\[
g_t^{\mathrm{fr}}
:=
\frac{L_{XY}(t,t)\Gamma(\mu)}{\Gamma(\lambda)\Gamma(1-\beta)}.
\]

Now write
\[
R_{t,h,\alpha,\phi}(u)
:=
\frac{1}{\Gamma(\lambda)}
\mathfrak D_{-,1}^{\lambda}\bigl(\mathfrak A_{XY,t,h,\alpha}\phi\bigr)(u)
-L_Y(t,t)\mathfrak T_t^{\mathrm{fr}}\phi(u).
\]
Using \eqref{eq:q-representation}, the constancy of \(g_t^{\mathrm{fr}}\), and
Proposition~\ref{prop:frozen-fractional}(ii), this becomes
\[
R_{t,h,\alpha,\phi}(u)
=
-\frac{d}{du}\int_u^1 (v-u)^{-\beta}
\bigl(g_{t,h,\alpha}(u,v)-g_t^{\mathrm{fr}}\bigr)\phi(v)\,dv.
\]
Since \(\phi\) is compactly supported in \((0,1)\), the boundary term at \(v=1\) vanishes, and
the substitution \(v=u+r\) yields
\begin{align*}
R_{t,h,\alpha,\phi}(u)
&=
-\int_0^{1-u} r^{-\beta}
\Bigl[
\bigl(\partial_1+\partial_2\bigr)
\bigl(g_{t,h,\alpha}-g_t^{\mathrm{fr}}\bigr)(u,u+r)\,\phi(u+r) \\
&\qquad\qquad\qquad\qquad\qquad
+\bigl(g_{t,h,\alpha}-g_t^{\mathrm{fr}}\bigr)(u,u+r)\,\phi'(u+r)
\Bigr]\,dr.
\end{align*}
By \eqref{eq:g-close}, the bracketed term is \(O_\phi(h^\alpha)\) uniformly in \(u\) and \(r\).
Because \(\beta\in(0,1)\),
\[
\int_0^{1-u} r^{-\beta}\,dr\le \frac{1}{1-\beta}
\qquad\text{for }0\le u\le 1.
\]
Hence \(\|R_{t,h,\alpha,\phi}\|_{L^\infty([0,1])}\le C_\phi h^\alpha\), which implies
\eqref{eq:q-smooth}.
\end{proof}

\begin{theorem}[Strong local convergence on the smooth core]
\label{thm:smooth-core}
Assume Assumptions~\ref{ass:model} and~\ref{ass:obs-source}, let \(\eta>0\), fix \(t\in(0,T)\)
and \(\alpha\in(0,1)\), and suppose
\[
0<H_{XY}<H_Y<\tfrac12.
\]
Then for every \(\phi\in C_c^\infty((0,1))\) there exist \(h_\phi\in(0,t]\) and
\(C_\phi<\infty\) such that for every \(0<h\le h_\phi\),
\[
\phi\in \mathfrak G_{t,h,\alpha}^{\mathrm{near}}
\]
and
\begin{equation}\label{eq:smooth-core-strong}
\|\mathfrak T_{t,h,\alpha}\phi-\mathfrak T_t^{\mathrm{fr}}\phi\|_E
\le
C_\phi h^\alpha.
\end{equation}
\end{theorem}

\begin{proof}
Set
\[
\lambda:=H_Y+\frac12,
\qquad
\beta:=H_Y-H_{XY}.
\]
Fix \(\phi\in C_c^\infty((0,1))\). For \(h>0\) sufficiently small, define
\[
q_{t,h,\alpha,\phi}
:=
\frac{1}{\Gamma(\lambda)}
\mathfrak D_{-,1}^{\lambda}\bigl(\mathfrak A_{XY,t,h,\alpha}\phi\bigr)\in E.
\]
By Corollary~\ref{cor:BY-invertible}, the operator \(\mathfrak B_{Y,t,h,\alpha}\) is invertible
for small \(h\), so we may set
\[
\psi_h:=\mathfrak B_{Y,t,h,\alpha}^{-1}q_{t,h,\alpha,\phi}\in E.
\]
Using Proposition~\ref{prop:q-smooth}, Proposition~\ref{prop:BY}, and
\eqref{eq:BY-inv},
\begin{align*}
\|\psi_h-\mathfrak T_t^{\mathrm{fr}}\phi\|_E
&\le
\|\mathfrak B_{Y,t,h,\alpha}^{-1}\|_{\mathcal L(E)}
\left(
\|q_{t,h,\alpha,\phi}-L_Y(t,t)\mathfrak T_t^{\mathrm{fr}}\phi\|_E\right. \\
&\qquad\qquad\left.
+\|(L_Y(t,t)I_E-\mathfrak B_{Y,t,h,\alpha})\mathfrak T_t^{\mathrm{fr}}\phi\|_E
\right) \\
&\le
C_\phi h^\alpha.
\end{align*}

Let
\[
r_h:=\mathfrak A_{Y,t,h,\alpha}\psi_h-\mathfrak A_{XY,t,h,\alpha}\phi.
\]
By the definition of \(\psi_h\), Proposition~\ref{prop:BY}, and the definition of
\(q_{t,h,\alpha,\phi}\),
\[
\mathfrak D_{-,1}^{\lambda}r_h=0
\qquad\text{in }E.
\]
Therefore \(\mathfrak I_{-,1}^{1-\lambda}r_h\) is absolutely continuous with zero derivative, so
it is constant. Since \((\mathfrak I_{-,1}^{1-\lambda}r_h)(1)=0\), that constant is \(0\). By
injectivity of \(\mathfrak I_{-,1}^{1-\lambda}\), proved in Lemma~\ref{lem:unit-semigroup},
\[
r_h=0.
\]
Thus
\begin{equation}\label{eq:smooth-local-eq}
\mathfrak A_{Y,t,h,\alpha}\psi_h=\mathfrak A_{XY,t,h,\alpha}\phi.
\end{equation}

Now define
\[
w_h:=\mathcal L_{t,h}^{(\alpha)}\phi\in H_{t,h,\alpha}^{\mathrm{near}},
\qquad
u_h:=h^{-\alpha\beta}\mathcal L_{t,h}^{(\alpha)}\psi_h\in H_{t,h,\alpha}^{\mathrm{near}}.
\]
By Proposition~\ref{prop:local-volterra} and \eqref{eq:smooth-local-eq},
\begin{align*}
\mathcal D_{t,h}^{(\alpha)}\mathcal K_{Y,t}u_h
&=
h^{-\alpha\beta}\mathcal D_{t,h}^{(\alpha)}
\mathcal K_{Y,t}\mathcal L_{t,h}^{(\alpha)}\psi_h \\
&=
h^{-\alpha\beta+\alpha H_Y}\mathfrak A_{Y,t,h,\alpha}\psi_h \\
&=
h^{\alpha H_{XY}}\mathfrak A_{XY,t,h,\alpha}\phi \\
&=
\mathcal D_{t,h}^{(\alpha)}\mathcal K_{XY,t}w_h.
\end{align*}
Both \(\mathcal K_{Y,t}u_h\) and \(\mathcal K_{XY,t}w_h\) are supported in
\([t-h^\alpha,t]\), so Lemma~\ref{lem:unit-dilation} implies
\[
\mathcal K_{Y,t}u_h=\mathcal K_{XY,t}w_h.
\]
Hence \(w_h\in\mathcal D(\mathcal S_t)\) and \(\mathcal S_t w_h=u_h\). By
Definition~\ref{def:local-transfer},
\[
\phi=\mathcal D_{t,h}^{(\alpha)}w_h\in \mathfrak G_{t,h,\alpha}^{\mathrm{near}}
\]
and
\[
\mathfrak T_{t,h,\alpha}\phi
=
h^{\alpha\beta}\mathcal D_{t,h}^{(\alpha)}\mathcal S_t w_h
=
h^{\alpha\beta}\mathcal D_{t,h}^{(\alpha)}u_h
=
\psi_h.
\]
Combining this identity with the estimate on \(\psi_h-\mathfrak T_t^{\mathrm{fr}}\phi\) proves
\eqref{eq:smooth-core-strong}.
\end{proof}

\begin{corollary}[Smooth core as a recovery core]
\label{cor:smooth-recovery}
Under the assumptions of Theorem~\ref{thm:smooth-core}, the space \(C_c^\infty((0,1))\)
satisfies items~(ii) and~(iii) of Definition~\ref{def:local-recovery-core}. In particular, if
\(C_c^\infty((0,1))\) is dense in \(\mathcal D(\mathfrak T_t^{\mathrm{fr}})\) with respect to
\(\|\cdot\|_{\mathfrak G_t^{\mathrm{fr}}}\), then it is a local graph recovery core.
\end{corollary}

\begin{proof}
Theorem~\ref{thm:smooth-core} gives items~(ii) and~(iii) directly. If
\(C_c^\infty((0,1))\) is also graph dense, then item~(i) of
Definition~\ref{def:local-recovery-core} holds as well.
\end{proof}

\begin{definition}[Intrinsic exact local graph on the unit interval]
\label{def:local-exact-graph}
Assume Assumptions~\ref{ass:model} and~\ref{ass:obs-source}, let \(\eta>0\), fix
\(t\in(0,T)\), \(\alpha\in(0,1)\), and choose \(h>0\) sufficiently small that
Corollary~\ref{cor:BY-invertible} applies. Define
\[
\Gamma_{t,h,\alpha}^{\mathrm{loc,ex}}
:=
\left\{
(\phi,\psi)\in E\times E:
\mathfrak A_{Y,t,h,\alpha}\psi
=
\mathfrak A_{XY,t,h,\alpha}\phi
\right\}.
\]
\end{definition}

\begin{proposition}[Closed exact local solver on the unit interval]
\label{prop:local-exact-solver}
Assume the hypotheses of Definition~\ref{def:local-exact-graph}. Then:
\begin{enumerate}[label=(\roman*)]
\item \(\Gamma_{t,h,\alpha}^{\mathrm{loc,ex}}\) is a closed linear subspace of \(E\times E\);
\item \(\mathfrak A_{Y,t,h,\alpha}:E\to E\) is injective;
\item there exists a unique closed linear operator
\[
\mathfrak T_{t,h,\alpha}^{\mathrm{loc,ex}}:
\mathcal D(\mathfrak T_{t,h,\alpha}^{\mathrm{loc,ex}})\subset E\to E
\]
whose graph is \(\Gamma_{t,h,\alpha}^{\mathrm{loc,ex}}\).
\end{enumerate}
\end{proposition}

\begin{proof}
Since \(\mathfrak A_{Y,t,h,\alpha}\) and \(\mathfrak A_{XY,t,h,\alpha}\) are bounded on \(E\),
the map
\[
(\phi,\psi)\longmapsto
\mathfrak A_{Y,t,h,\alpha}\psi-\mathfrak A_{XY,t,h,\alpha}\phi
\]
is bounded from \(E\times E\) to \(E\). Therefore
\(\Gamma_{t,h,\alpha}^{\mathrm{loc,ex}}\), being its kernel, is a closed linear subspace. This
proves (i).

To prove injectivity, let \(\psi\in E\) satisfy
\[
\mathfrak A_{Y,t,h,\alpha}\psi=0.
\]
Proposition~\ref{prop:BY} gives
\[
\mathfrak B_{Y,t,h,\alpha}\psi
=
\frac{1}{\Gamma(\lambda)}
\mathfrak D_{-,1}^{\lambda}\bigl(\mathfrak A_{Y,t,h,\alpha}\psi\bigr)
=
0,
\qquad
\lambda:=H_Y+\frac12.
\]
By Corollary~\ref{cor:BY-invertible}, \(\mathfrak B_{Y,t,h,\alpha}\) is invertible, so
\(\psi=0\). This proves (ii).

Part (iii) now follows exactly as in Proposition~\ref{prop:micro-window-solver}: by (ii), for each
\(\phi\in E\) there is at most one \(\psi\in E\) such that \((\phi,\psi)\in
\Gamma_{t,h,\alpha}^{\mathrm{loc,ex}}\). Hence the closed linear subspace
\(\Gamma_{t,h,\alpha}^{\mathrm{loc,ex}}\) is the graph of a unique closed linear operator.
\end{proof}

\begin{proposition}[Intrinsic and transfer-defined local graphs coincide]
\label{prop:local-exact-identification}
Assume the hypotheses of Definition~\ref{def:local-exact-graph}. Then
\begin{equation}\label{eq:local-exact-identification}
\mathcal D(\mathfrak T_{t,h,\alpha}^{\mathrm{loc,ex}})
=
\mathfrak G_{t,h,\alpha}^{\mathrm{near}},
\qquad
\mathfrak T_{t,h,\alpha}^{\mathrm{loc,ex}}\phi
=
\mathfrak T_{t,h,\alpha}\phi
\quad
\text{for every }\phi\in\mathfrak G_{t,h,\alpha}^{\mathrm{near}}.
\end{equation}
In particular,
\[
\Gamma_{t,h,\alpha}^{\mathrm{loc,ex}}
=
\left\{
\bigl(\phi,\mathfrak T_{t,h,\alpha}\phi\bigr):
\phi\in\mathfrak G_{t,h,\alpha}^{\mathrm{near}}
\right\}.
\]
\end{proposition}

\begin{proof}
Let \(\phi\in\mathfrak G_{t,h,\alpha}^{\mathrm{near}}\), and set
\[
\psi:=\mathfrak T_{t,h,\alpha}\phi.
\]
Proposition~\ref{prop:local-transfer}(ii) gives
\[
\mathfrak A_{Y,t,h,\alpha}\psi
=
\mathfrak A_{XY,t,h,\alpha}\phi,
\]
so \((\phi,\psi)\in\Gamma_{t,h,\alpha}^{\mathrm{loc,ex}}\). Hence
\[
\mathfrak G_{t,h,\alpha}^{\mathrm{near}}
\subset
\mathcal D(\mathfrak T_{t,h,\alpha}^{\mathrm{loc,ex}})
\]
and
\[
\mathfrak T_{t,h,\alpha}^{\mathrm{loc,ex}}\phi
=
\mathfrak T_{t,h,\alpha}\phi
\qquad\text{for every }\phi\in\mathfrak G_{t,h,\alpha}^{\mathrm{near}}.
\]

Conversely, let
\[
\phi\in \mathcal D(\mathfrak T_{t,h,\alpha}^{\mathrm{loc,ex}})
\]
and set
\[
\psi:=\mathfrak T_{t,h,\alpha}^{\mathrm{loc,ex}}\phi.
\]
Then
\[
\mathfrak A_{Y,t,h,\alpha}\psi
=
\mathfrak A_{XY,t,h,\alpha}\phi.
\]
Define
\[
w:=\mathcal L_{t,h}^{(\alpha)}\phi\in H_{t,h,\alpha}^{\mathrm{near}},
\qquad
u:=h^{-\alpha(H_Y-H_{XY})}\mathcal L_{t,h}^{(\alpha)}\psi
\in H_{t,h,\alpha}^{\mathrm{near}}.
\]
Applying Proposition~\ref{prop:local-volterra} to both sides of the exact local identity gives
\begin{align*}
\mathcal D_{t,h}^{(\alpha)}\mathcal K_{Y,t}u
&=
h^{-\alpha(H_Y-H_{XY})}
\mathcal D_{t,h}^{(\alpha)}
\mathcal K_{Y,t}\mathcal L_{t,h}^{(\alpha)}\psi \\
&=
h^{-\alpha(H_Y-H_{XY})+\alpha H_Y}
\mathfrak A_{Y,t,h,\alpha}\psi \\
&=
h^{\alpha H_{XY}}
\mathfrak A_{XY,t,h,\alpha}\phi \\
&=
\mathcal D_{t,h}^{(\alpha)}\mathcal K_{XY,t}w.
\end{align*}
Both \(\mathcal K_{Y,t}u\) and \(\mathcal K_{XY,t}w\) are supported in \([t-h^\alpha,t]\), so
Lemma~\ref{lem:unit-dilation} implies
\[
\mathcal K_{Y,t}u=\mathcal K_{XY,t}w.
\]
Hence \(w\in\mathcal D(\mathcal S_t)\) and \(\mathcal S_t w=u\). By
Definition~\ref{def:local-transfer},
\[
\phi=\mathcal D_{t,h}^{(\alpha)}w\in \mathfrak G_{t,h,\alpha}^{\mathrm{near}},
\]
and
\[
\mathfrak T_{t,h,\alpha}\phi
=
h^{\alpha(H_Y-H_{XY})}
\mathcal D_{t,h}^{(\alpha)}\mathcal S_t w
=
h^{\alpha(H_Y-H_{XY})}
\mathcal D_{t,h}^{(\alpha)}u
=
\psi.
\]
This proves \eqref{eq:local-exact-identification}.
\end{proof}

\begin{proposition}[Conjugated local commutator operators on the unit interval]
\label{prop:local-smooth-commutator}
Assume Assumption~\ref{ass:model}, fix \(t\in(0,T)\), \(\alpha\in(0,1)\), and suppose
\[
0<H_{XY}<H_Y<\tfrac12.
\]
Let
\[
\chi\in C^\infty([0,1]),
\qquad
0\le \chi\le 1,
\]
and define
\[
\delta_\chi(u,v):=
\begin{cases}
\dfrac{\chi(v)-\chi(u)}{v-u}, & 0\le u<v\le 1,\\[1ex]
\chi'(u), & 0\le u=v\le 1.
\end{cases}
\]
For \(\Theta\in E\), define
\begin{align}
\bigl(
\mathfrak W_{Y,t,h,\alpha,\chi}^{\mathrm{loc}}
\Theta
\bigr)(u)
&:=
\int_u^1
(v-u)^{\lambda-1}
\bigl(\chi(v)-\chi(u)\bigr)
L_Y(t-h^\alpha u,t-h^\alpha v)\Theta(v)\,dv,
\label{eq:local-com-Y}\\
\bigl(
\mathfrak W_{XY,t,h,\alpha,\chi}^{\mathrm{loc}}
\Theta
\bigr)(u)
&:=
\int_u^1
(v-u)^{\mu-1}
\bigl(\chi(v)-\chi(u)\bigr)
L_{XY}(t-h^\alpha u,t-h^\alpha v)\Theta(v)\,dv,
\label{eq:local-com-XY}
\end{align}
where
\[
\lambda:=H_Y+\frac12,
\qquad
\mu:=H_{XY}+\frac12.
\]
Then there exist bounded operators
\[
\mathfrak Q_{Y,t,h,\alpha,\chi}^{\mathrm{loc}},
\mathfrak Q_{XY,t,h,\alpha,\chi}^{\mathrm{loc}}
\in\mathcal L(E)
\]
such that for every \(\Theta\in E\),
\begin{align}
\frac{1}{\Gamma(\lambda)}
\mathfrak D_{-,1}^{\lambda}
\bigl(
\mathfrak W_{Y,t,h,\alpha,\chi}^{\mathrm{loc}}\Theta
\bigr)
&=
\mathfrak Q_{Y,t,h,\alpha,\chi}^{\mathrm{loc}}\Theta
\qquad\text{in }E,
\label{eq:local-com-conj-Y}\\
\frac{1}{\Gamma(\lambda)}
\mathfrak D_{-,1}^{\lambda}
\bigl(
\mathfrak W_{XY,t,h,\alpha,\chi}^{\mathrm{loc}}\Theta
\bigr)
&=
\mathfrak Q_{XY,t,h,\alpha,\chi}^{\mathrm{loc}}\Theta
\qquad\text{in }E.
\label{eq:local-com-conj-XY}
\end{align}
Moreover, there exists \(C_t<\infty\) such that
\begin{equation}\label{eq:local-com-bound}
\left\|
\mathfrak Q_{Y,t,h,\alpha,\chi}^{\mathrm{loc}}
\right\|_{\mathcal L(E)}
+
\left\|
\mathfrak Q_{XY,t,h,\alpha,\chi}^{\mathrm{loc}}
\right\|_{\mathcal L(E)}
\le
C_t\|\delta_\chi\|_{C^1(\Delta_1)},
\qquad
\Delta_1:=\{0\le u\le v\le 1\},
\end{equation}
for all sufficiently small \(h>0\).
\end{proposition}

\begin{proof}
For the target side, write
\[
\bigl(
\mathfrak W_{Y,t,h,\alpha,\chi}^{\mathrm{loc}}\Theta
\bigr)(u)
=
\int_u^1
(v-u)^{\lambda}
\delta_\chi(u,v)
L_Y(t-h^\alpha u,t-h^\alpha v)\Theta(v)\,dv.
\]
Exactly as in the proof of Proposition~\ref{prop:micro-smooth-commutator}, Fubini's theorem gives
\[
\mathfrak I_{-,1}^{1-\lambda}
\bigl(
\mathfrak W_{Y,t,h,\alpha,\chi}^{\mathrm{loc}}\Theta
\bigr)(u)
=
\int_u^1
(v-u)
F_{Y,t,h,\alpha,\chi}^{\mathrm{loc}}(u,v)\Theta(v)\,dv,
\]
where \(F_{Y,t,h,\alpha,\chi}^{\mathrm{loc}}\) is \(C^1\) on \(\Delta_1\). Differentiating under
the integral sign yields
\[
\frac{1}{\Gamma(\lambda)}
\mathfrak D_{-,1}^{\lambda}
\bigl(
\mathfrak W_{Y,t,h,\alpha,\chi}^{\mathrm{loc}}\Theta
\bigr)(u)
=
F_{Y,t,h,\alpha,\chi}^{\mathrm{loc}}(u,u)\Theta(u)
-\int_u^1
K_{Y,t,h,\alpha,\chi}^{\mathrm{loc}}(u,v)\Theta(v)\,dv
\]
for almost every \(u\in[0,1]\), where
\[
K_{Y,t,h,\alpha,\chi}^{\mathrm{loc}}(u,v)
:=
F_{Y,t,h,\alpha,\chi}^{\mathrm{loc}}(u,v)
+(v-u)\partial_1F_{Y,t,h,\alpha,\chi}^{\mathrm{loc}}(u,v).
\]
This defines the bounded operator \(\mathfrak Q_{Y,t,h,\alpha,\chi}^{\mathrm{loc}}\).

Because \(L_Y\) is \(C^1\) and \(\delta_\chi\in C^1(\Delta_1)\), the same Beta-integral formula
shows
\[
\|F_{Y,t,h,\alpha,\chi}^{\mathrm{loc}}\|_{C^1(\Delta_1)}
\le
C_t\|\delta_\chi\|_{C^1(\Delta_1)}
\]
uniformly for small \(h\). Since \(\Delta_1\) has finite measure, the corresponding multiplication
term and integral term are bounded on \(E\), with norm bounded by
\(C_t\|\delta_\chi\|_{C^1(\Delta_1)}\).

For the source side, write
\[
\bigl(
\mathfrak W_{XY,t,h,\alpha,\chi}^{\mathrm{loc}}\Theta
\bigr)(u)
=
\int_u^1
(v-u)^{\mu}
\delta_\chi(u,v)
L_{XY}(t-h^\alpha u,t-h^\alpha v)\Theta(v)\,dv.
\]
Exactly as in Proposition~\ref{prop:micro-smooth-commutator}, one obtains
\[
\mathfrak I_{-,1}^{1-\lambda}
\bigl(
\mathfrak W_{XY,t,h,\alpha,\chi}^{\mathrm{loc}}\Theta
\bigr)(u)
=
\int_u^1
(v-u)^{1-\beta}
F_{XY,t,h,\alpha,\chi}^{\mathrm{loc}}(u,v)\Theta(v)\,dv,
\]
where \(\beta:=H_Y-H_{XY}\) and \(F_{XY,t,h,\alpha,\chi}^{\mathrm{loc}}\) is \(C^1\) on
\(\Delta_1\) with
\[
\|F_{XY,t,h,\alpha,\chi}^{\mathrm{loc}}\|_{C^1(\Delta_1)}
\le
C_t\|\delta_\chi\|_{C^1(\Delta_1)}.
\]
Since \(1-\beta>0\), differentiation under the integral sign gives
\begin{align*}
\frac{1}{\Gamma(\lambda)}
\mathfrak D_{-,1}^{\lambda}
\bigl(
\mathfrak W_{XY,t,h,\alpha,\chi}^{\mathrm{loc}}\Theta
\bigr)(u)
&=
-(1-\beta)
\int_u^1
(v-u)^{-\beta}
F_{XY,t,h,\alpha,\chi}^{\mathrm{loc}}(u,v)\Theta(v)\,dv\\
&\quad-
\int_u^1
(v-u)^{1-\beta}
\partial_1F_{XY,t,h,\alpha,\chi}^{\mathrm{loc}}(u,v)\Theta(v)\,dv.
\end{align*}
Because \(\beta<\tfrac12\), the kernels
\[
(v-u)^{-\beta}\one_{\Delta_1}(u,v),
\qquad
(v-u)^{1-\beta}\one_{\Delta_1}(u,v)
\]
are Hilbert--Schmidt on \(\Delta_1\). Hence the right-hand side defines a bounded operator
\(\mathfrak Q_{XY,t,h,\alpha,\chi}^{\mathrm{loc}}\in\mathcal L(E)\) with norm bounded by
\(C_t\|\delta_\chi\|_{C^1(\Delta_1)}\). This proves \eqref{eq:local-com-conj-XY} and
\eqref{eq:local-com-bound}.
\end{proof}

\begin{theorem}[Smooth localization of an exact local pair]
\label{thm:local-smooth-localization}
Assume Assumptions~\ref{ass:model} and~\ref{ass:obs-source}, let \(\eta>0\), fix
\(t\in(0,T)\), \(\alpha\in(0,1)\), and suppose
\[
0<H_{XY}<H_Y<\tfrac12.
\]
Let \(h>0\) be sufficiently small that Corollary~\ref{cor:BY-invertible} applies, let
\[
\chi\in C^\infty([0,1]),
\qquad
0\le \chi\le 1,
\]
and let
\[
(\phi,\psi)\in \Gamma_{t,h,\alpha}^{\mathrm{loc,ex}}.
\]
Define the conjugated local defect
\begin{equation}\label{eq:local-smooth-defect}
\mathfrak F_{t,h,\alpha,\chi}^{\mathrm{loc}}(\phi,\psi)
:=
\mathfrak Q_{Y,t,h,\alpha,\chi}^{\mathrm{loc}}\psi
-\mathfrak Q_{XY,t,h,\alpha,\chi}^{\mathrm{loc}}\phi
\in E,
\end{equation}
and the corrected output
\begin{equation}\label{eq:local-smooth-corrected}
\psi_{\chi}^{\mathrm{loc}}
:=
\chi\psi
-\mathfrak B_{Y,t,h,\alpha}^{-1}
\mathfrak F_{t,h,\alpha,\chi}^{\mathrm{loc}}(\phi,\psi).
\end{equation}
Then:
\begin{enumerate}[label=(\roman*)]
\item
\[
(\chi\phi,\psi_{\chi}^{\mathrm{loc}})
\in
\Gamma_{t,h,\alpha}^{\mathrm{loc,ex}};
\]
\item there exists \(C_t<\infty\) such that
\begin{equation}\label{eq:local-smooth-correction}
\|\psi_{\chi}^{\mathrm{loc}}-\chi\psi\|_E
\le
C_t\|\delta_\chi\|_{C^1(\Delta_1)}
\bigl(\|\phi\|_E+\|\psi\|_E\bigr);
\end{equation}
\item consequently,
\begin{equation}\label{eq:local-smooth-graph}
\|\chi\phi\|_E+\frac{\eta}{\nu_Y}\|\psi_{\chi}^{\mathrm{loc}}\|_E
\le
C_t\bigl(1+\|\delta_\chi\|_{C^1(\Delta_1)}\bigr)
\left(
\|\phi\|_E+\frac{\eta}{\nu_Y}\|\psi\|_E
\right).
\end{equation}
\end{enumerate}
\end{theorem}

\begin{proof}
Because \((\phi,\psi)\in\Gamma_{t,h,\alpha}^{\mathrm{loc,ex}}\), one has
\[
\mathfrak A_{Y,t,h,\alpha}\psi=\mathfrak A_{XY,t,h,\alpha}\phi.
\]
Multiplying by \(\chi(u)\) and subtracting gives, for almost every \(u\in[0,1]\),
\[
\mathfrak A_{Y,t,h,\alpha}(\chi\psi)(u)-\mathfrak A_{XY,t,h,\alpha}(\chi\phi)(u)
=
\int_u^1 K_Y(u,v)\,dv-\int_u^1 K_{XY}(u,v)\,dv,
\]
where
\begin{align*}
K_Y(u,v)
&:=
(v-u)^{\lambda-1}\bigl(\chi(v)-\chi(u)\bigr)
L_Y(t-h^\alpha u,t-h^\alpha v)\psi(v),\\
K_{XY}(u,v)
&:=
(v-u)^{\mu-1}\bigl(\chi(v)-\chi(u)\bigr)
L_{XY}(t-h^\alpha u,t-h^\alpha v)\phi(v).
\end{align*}
By the definition of the local commutator operators, this identity can be rewritten as
\[
\mathfrak A_{Y,t,h,\alpha}(\chi\psi)-\mathfrak A_{XY,t,h,\alpha}(\chi\phi)
=
\mathfrak W_{Y,t,h,\alpha,\chi}^{\mathrm{loc}}\psi
-\mathfrak W_{XY,t,h,\alpha,\chi}^{\mathrm{loc}}\phi.
\]
Applying Proposition~\ref{prop:local-smooth-commutator} and then
Proposition~\ref{prop:BY} yields
\[
\mathfrak B_{Y,t,h,\alpha}(\chi\psi)
=
\frac{1}{\Gamma(\lambda)}
\mathfrak D_{-,1}^{\lambda}
\bigl(\mathfrak A_{XY,t,h,\alpha}(\chi\phi)\bigr)
+
\mathfrak F_{t,h,\alpha,\chi}^{\mathrm{loc}}(\phi,\psi)
\qquad\text{in }E.
\]
By the definition of \(\psi_{\chi}^{\mathrm{loc}}\),
\[
\mathfrak B_{Y,t,h,\alpha}\psi_{\chi}^{\mathrm{loc}}
=
\frac{1}{\Gamma(\lambda)}
\mathfrak D_{-,1}^{\lambda}
\bigl(\mathfrak A_{XY,t,h,\alpha}(\chi\phi)\bigr)
\qquad\text{in }E.
\]
Applying Proposition~\ref{prop:BY} again shows that
\[
\mathfrak D_{-,1}^{\lambda}
\bigl(
\mathfrak A_{Y,t,h,\alpha}\psi_{\chi}^{\mathrm{loc}}
-\mathfrak A_{XY,t,h,\alpha}(\chi\phi)
\bigr)
=
0
\qquad\text{in }E.
\]
Exactly as in the proof of Theorem~\ref{thm:smooth-core}, injectivity of
\(\mathfrak I_{-,1}^{1-\lambda}\) implies
\[
\mathfrak A_{Y,t,h,\alpha}\psi_{\chi}^{\mathrm{loc}}
=
\mathfrak A_{XY,t,h,\alpha}(\chi\phi),
\]
so \((\chi\phi,\psi_{\chi}^{\mathrm{loc}})\in \Gamma_{t,h,\alpha}^{\mathrm{loc,ex}}\). This
proves (i).

For (ii), Corollary~\ref{cor:BY-invertible} and \eqref{eq:local-com-bound} give
\begin{align*}
\|\psi_{\chi}^{\mathrm{loc}}-\chi\psi\|_E
&\le
\|\mathfrak B_{Y,t,h,\alpha}^{-1}\|_{\mathcal L(E)}
\left\|
\mathfrak F_{t,h,\alpha,\chi}^{\mathrm{loc}}(\phi,\psi)
\right\|_E\\
&\le
C_t\|\delta_\chi\|_{C^1(\Delta_1)}
\bigl(\|\phi\|_E+\|\psi\|_E\bigr),
\end{align*}
which proves \eqref{eq:local-smooth-correction}.

Finally, \(|\chi|\le1\) and \eqref{eq:local-smooth-correction} imply
\begin{align*}
\|\chi\phi\|_E+\frac{\eta}{\nu_Y}\|\psi_{\chi}^{\mathrm{loc}}\|_E
&\le
\|\phi\|_E+\frac{\eta}{\nu_Y}\|\chi\psi\|_E
+\frac{\eta}{\nu_Y}\|\psi_{\chi}^{\mathrm{loc}}-\chi\psi\|_E\\
&\le
C_t\bigl(1+\|\delta_\chi\|_{C^1(\Delta_1)}\bigr)
\left(
\|\phi\|_E+\frac{\eta}{\nu_Y}\|\psi\|_E
\right),
\end{align*}
which is \eqref{eq:local-smooth-graph}.
\end{proof}

\begin{corollary}[Smooth cutoffs preserve boundary-supported local graph elements]
\label{cor:near-smooth-cutoff}
Assume the hypotheses of Theorem~\ref{thm:local-smooth-localization}. Let
\[
\chi\in C^\infty([0,1]),
\qquad
0\le \chi\le 1,
\]
and define
\[
\chi_{t,h,\alpha}(s)
:=
\chi\!\left(\frac{t-s}{h^\alpha}\right)\one_{[t-h^\alpha,t]}(s),
\qquad 0\le s\le t.
\]
Then for every
\[
w\in \mathfrak G_t\cap H_{t,h,\alpha}^{\mathrm{near}},
\]
one has
\[
\chi_{t,h,\alpha}w\in \mathfrak G_t\cap H_{t,h,\alpha}^{\mathrm{near}},
\]
and there exists \(C_t<\infty\), independent of \(w\) and \(h\), such that
\begin{equation}\label{eq:near-smooth-cutoff}
\|\chi_{t,h,\alpha}w\|_{H_t}
+\frac{\eta}{\nu_Y}h^{\alpha(H_Y-H_{XY})}
\|\mathcal S_t(\chi_{t,h,\alpha}w)\|_{H_t}
\le
C_t\bigl(1+\|\delta_\chi\|_{C^1(\Delta_1)}\bigr)\|w\|_{\mathfrak G_t}.
\end{equation}
\end{corollary}

\begin{proof}
Set
\[
\phi:=\mathcal D_{t,h}^{(\alpha)}w\in \mathfrak G_{t,h,\alpha}^{\mathrm{near}},
\qquad
\psi:=\mathfrak T_{t,h,\alpha}\phi\in E.
\]
By Proposition~\ref{prop:local-exact-identification},
\[
(\phi,\psi)\in \Gamma_{t,h,\alpha}^{\mathrm{loc,ex}}.
\]
Apply Theorem~\ref{thm:local-smooth-localization} to obtain
\[
(\chi\phi,\psi_{\chi}^{\mathrm{loc}})
\in
\Gamma_{t,h,\alpha}^{\mathrm{loc,ex}}
\]
and
\[
\|\chi\phi\|_E+\frac{\eta}{\nu_Y}\|\psi_{\chi}^{\mathrm{loc}}\|_E
\le
C_t\bigl(1+\|\delta_\chi\|_{C^1(\Delta_1)}\bigr)
\left(
\|\phi\|_E+\frac{\eta}{\nu_Y}\|\psi\|_E
\right).
\]
Now define
\[
\widetilde w:=\mathcal L_{t,h}^{(\alpha)}(\chi\phi),
\qquad
\widetilde u:=h^{-\alpha(H_Y-H_{XY})}\mathcal L_{t,h}^{(\alpha)}\psi_{\chi}^{\mathrm{loc}}.
\]
By Proposition~\ref{prop:local-exact-identification}, the exact local pair relation implies
\[
\widetilde w\in \mathfrak G_t\cap H_{t,h,\alpha}^{\mathrm{near}},
\qquad
\mathcal S_t\widetilde w=\widetilde u.
\]
Since
\[
\widetilde w(s)
=
h^{-\alpha/2}\chi\!\left(\frac{t-s}{h^\alpha}\right)\phi\!\left(\frac{t-s}{h^\alpha}\right)
\one_{[t-h^\alpha,t]}(s)
=
\chi_{t,h,\alpha}(s)w(s),
\]
we have \(\widetilde w=\chi_{t,h,\alpha}w\). Finally,
Lemma~\ref{lem:unit-dilation} gives
\[
\|\widetilde w\|_{H_t}=\|\chi\phi\|_E,
\qquad
\|\mathcal S_t\widetilde w\|_{H_t}
=
h^{-\alpha(H_Y-H_{XY})}\|\psi_{\chi}^{\mathrm{loc}}\|_E,
\]
while
\[
\|w\|_{H_t}=\|\phi\|_E,
\qquad
\|\mathcal S_t w\|_{H_t}
=
h^{-\alpha(H_Y-H_{XY})}\|\psi\|_E.
\]
Therefore the preceding bound is exactly \eqref{eq:near-smooth-cutoff}.
\end{proof}

\begin{proposition}[Slow variation of a fixed time-side cutoff on the micro scale]
\label{prop:micro-cutoff-scaling}
Fix \(\alpha\in(0,1)\), let
\[
\chi\in C^\infty([0,1]),
\qquad
0\le \chi\le 1,
\]
and assume that for some \(\delta_*\in(0,1)\),
\[
\chi\equiv 1\ \text{on }[0,1-\delta_*].
\]
For \(h\in(0,1]\), set
\[
\varepsilon_{h,\alpha}:=h^{1-\alpha},
\qquad
\widetilde\chi_{h,\alpha}(r)
:=
\chi\bigl(\varepsilon_{h,\alpha}r\bigr)\one_{[0,\varepsilon_{h,\alpha}^{-1}]}(r),
\qquad r\ge0.
\]
Define, on
\[
\Delta_{\varepsilon_{h,\alpha}^{-1}}
:=
\{0\le r\le q\le \varepsilon_{h,\alpha}^{-1}\},
\]
the quotient kernel
\[
\delta_{\widetilde\chi_{h,\alpha}}(r,q)
:=
\begin{cases}
\dfrac{\widetilde\chi_{h,\alpha}(q)-\widetilde\chi_{h,\alpha}(r)}{q-r},
& 0\le r<q\le \varepsilon_{h,\alpha}^{-1},\\[1ex]
\widetilde\chi_{h,\alpha}'(r), & 0\le r=q\le \varepsilon_{h,\alpha}^{-1}.
\end{cases}
\]
Then:
\begin{enumerate}[label=(\roman*)]
\item
\[
0\le \widetilde\chi_{h,\alpha}\le 1,
\qquad
\widetilde\chi_{h,\alpha}\equiv 1\ \text{on }
\left[0,(1-\delta_*)\varepsilon_{h,\alpha}^{-1}\right],
\qquad
\operatorname{supp}\widetilde\chi_{h,\alpha}\subset
[0,\varepsilon_{h,\alpha}^{-1}];
\]
\item there exists \(C_\chi<\infty\), depending only on \(\chi\), such that
\begin{equation}\label{eq:micro-cutoff-scaling}
\|\delta_{\widetilde\chi_{h,\alpha}}\|_{C^1(\Delta_{\varepsilon_{h,\alpha}^{-1}})}
\le
C_\chi \varepsilon_{h,\alpha}
\qquad\text{for every }h\in(0,1].
\end{equation}
\end{enumerate}
\end{proposition}

\begin{proof}
Item (i) follows from the definition of \(\widetilde\chi_{h,\alpha}\).

For \(0\le r<q\le \varepsilon_{h,\alpha}^{-1}\), set
\[
\varepsilon:=\varepsilon_{h,\alpha}.
\]
Then
\[
\delta_{\widetilde\chi_{h,\alpha}}(r,q)
=
\frac{\chi(\varepsilon q)-\chi(\varepsilon r)}{q-r}
=
\varepsilon
\frac{\chi(\varepsilon q)-\chi(\varepsilon r)}{\varepsilon q-\varepsilon r}.
\]
Hence, if
\[
\delta_\chi(x,y)
:=
\begin{cases}
\dfrac{\chi(y)-\chi(x)}{y-x}, & 0\le x<y\le 1,\\[1ex]
\chi'(x), & 0\le x=y\le 1,
\end{cases}
\]
then
\[
\delta_{\widetilde\chi_{h,\alpha}}(r,q)
=
\varepsilon\,\delta_\chi(\varepsilon r,\varepsilon q)
\qquad\text{on }\Delta_{\varepsilon^{-1}}.
\]
Since \(\chi\in C^\infty([0,1])\), one has \(\delta_\chi\in C^1(\Delta_1)\). Therefore
\[
\|\delta_{\widetilde\chi_{h,\alpha}}\|_{L^\infty(\Delta_{\varepsilon^{-1}})}
\le
\varepsilon \|\delta_\chi\|_{L^\infty(\Delta_1)}.
\]
Differentiating the identity above yields
\[
\partial_r\delta_{\widetilde\chi_{h,\alpha}}(r,q)
=
\varepsilon^2 (\partial_1\delta_\chi)(\varepsilon r,\varepsilon q),
\qquad
\partial_q\delta_{\widetilde\chi_{h,\alpha}}(r,q)
=
\varepsilon^2 (\partial_2\delta_\chi)(\varepsilon r,\varepsilon q),
\]
and thus
\[
\|\nabla \delta_{\widetilde\chi_{h,\alpha}}\|_{L^\infty(\Delta_{\varepsilon^{-1}})}
\le
\varepsilon^2 \|\nabla \delta_\chi\|_{L^\infty(\Delta_1)}
\le
\varepsilon \|\nabla \delta_\chi\|_{L^\infty(\Delta_1)},
\]
because \(\varepsilon\le1\). This proves \eqref{eq:micro-cutoff-scaling}.
\end{proof}

\begin{proposition}[Frozen domain as a fractional-integral range]
\label{prop:frozen-range}
Assume
\[
0<H_{XY}<H_Y<\tfrac12,
\qquad
\beta:=H_Y-H_{XY},
\]
and fix \(t\in(0,T)\). Then
\begin{equation}\label{eq:frozen-range}
\mathcal D(\mathfrak T_t^{\mathrm{fr}})
=
\mathfrak I_{-,1}^{\beta}(E).
\end{equation}
The range identity has the following two forms.
\begin{enumerate}[label=(\roman*)]
\item if \(\phi\in\mathcal D(\mathfrak T_t^{\mathrm{fr}})\), then
\begin{equation}\label{eq:frozen-range-forward}
\phi
=
\mathbf c_t^{-1}\mathfrak I_{-,1}^{\beta}\bigl(\mathfrak T_t^{\mathrm{fr}}\phi\bigr);
\end{equation}
\item if \(g\in E\) and
\[
\phi:=\mathbf c_t^{-1}\mathfrak I_{-,1}^{\beta}g,
\]
then \(\phi\in\mathcal D(\mathfrak T_t^{\mathrm{fr}})\) and
\begin{equation}\label{eq:frozen-range-backward}
\mathfrak T_t^{\mathrm{fr}}\phi=g.
\end{equation}
\end{enumerate}
\end{proposition}

\begin{proof}
Set
\[
\lambda:=H_Y+\frac12,
\qquad
\mu:=H_{XY}+\frac12,
\]
so that \(\lambda=\mu+\beta\). Let \(\phi\in\mathcal D(\mathfrak T_t^{\mathrm{fr}})\), and write
\[
g:=\mathfrak T_t^{\mathrm{fr}}\phi.
\]
By Proposition~\ref{prop:frozen-fractional}(i) and Lemma~\ref{lem:unit-semigroup},
\begin{align*}
L_Y(t,t)\Gamma(\lambda)\mathfrak I_{-,1}^{\lambda}g
&=
L_{XY}(t,t)\Gamma(\mu)\mathfrak I_{-,1}^{\mu}\phi \\
&=
L_Y(t,t)\Gamma(\lambda)\mathbf c_t\,\mathfrak I_{-,1}^{\mu}\phi.
\end{align*}
Hence
\[
\mathfrak I_{-,1}^{\mu}\bigl(\phi-\mathbf c_t^{-1}\mathfrak I_{-,1}^{\beta}g\bigr)=0,
\]
because \(\mathfrak I_{-,1}^{\mu}\mathfrak I_{-,1}^{\beta}=\mathfrak I_{-,1}^{\lambda}\). Since
\(\mathfrak I_{-,1}^{\mu}\) is injective by the same argument as in
Lemma~\ref{lem:unit-semigroup}, this proves \eqref{eq:frozen-range-forward}.

Conversely, let \(g\in E\) and set \(\phi:=\mathbf c_t^{-1}\mathfrak I_{-,1}^{\beta}g\). Then
\[
\mathfrak I_{-,1}^{\mu}\phi
=
\mathbf c_t^{-1}\mathfrak I_{-,1}^{\lambda}g.
\]
Multiplying by the Gamma factors built into \(\mathbf c_t\) gives
\[
L_{XY}(t,t)\Gamma(\mu)\mathfrak I_{-,1}^{\mu}\phi
=
L_Y(t,t)\Gamma(\lambda)\mathfrak I_{-,1}^{\lambda}g,
\]
that is,
\[
L_{XY}(t,t)\mathfrak J_{H_{XY}}\phi
=
L_Y(t,t)\mathfrak J_{H_Y}g.
\]
Therefore \(\mathfrak J_{H_{XY}}\phi\in\mathrm{Ran}(\mathfrak J_{H_Y})\), so
\(\phi\in\mathcal D(\mathfrak T_t^{\mathrm{fr}})\) and
\eqref{eq:frozen-range-backward} follows from Proposition~\ref{prop:frozen-fractional}(i) and
injectivity of \(\mathfrak J_{H_Y}\).
\end{proof}

\begin{definition}[Finite-interval fractional Sobolev space]\label{def:Hbeta-interval}
Let \(\beta\in(0,\tfrac12)\).  The notation \(H^\beta((0,1))\) denotes the Bessel potential space on the
interval \((0,1)\), endowed with its Hilbert norm.
\end{definition}

\begin{theorem}[Density of the smooth core in the frozen graph]
\label{thm:frozen-density}
Assume
\[
0<H_{XY}<H_Y<\tfrac12,
\qquad
\beta:=H_Y-H_{XY},
\]
and fix \(t\in(0,T)\). Then:
\begin{enumerate}[label=(\roman*)]
\item \(\mathcal D(\mathfrak T_t^{\mathrm{fr}})=H^\beta((0,1))\) as sets;
\item the frozen graph norm is equivalent to the \(H^\beta((0,1))\) norm;
\item \(C_c^\infty((0,1))\) is dense in \(\mathcal D(\mathfrak T_t^{\mathrm{fr}})\) with respect
to the frozen graph norm.
\end{enumerate}
\end{theorem}

\begin{proof}
By Proposition~\ref{prop:frozen-range},
\[
\mathcal D(\mathfrak T_t^{\mathrm{fr}})
=
\mathfrak I_{-,1}^{\beta}(E).
\]
For \(\beta\in(0,\tfrac12)\), the classical finite-interval Riemann--Liouville isomorphism
theorem identifies \(\mathfrak I_{-,1}^{\beta}(E)\) with the right-sided fractional Sobolev space
\(H_0^\beta((0,1))\), and the inverse is the right-sided derivative
\(\mathfrak D_{-,1}^{\beta}\) in the \(L^2\) sense. Since \(\beta<\tfrac12\), one has
\[
H_0^\beta((0,1))=H^\beta((0,1)),
\]
so item~(i) follows. This finite-interval fractional-Sobolev identification is used here
as an analytic input; see \cite[Ch.~13]{SKM93} and
\cite[Ch.~2]{Triebel83}.

Let \(\phi\in\mathcal D(\mathfrak T_t^{\mathrm{fr}})\), and set \(g:=\mathfrak T_t^{\mathrm{fr}}\phi\).
Proposition~\ref{prop:frozen-range} gives
\[
\phi=\mathbf c_t^{-1}\mathfrak I_{-,1}^{\beta}g.
\]
Applying the inverse \(\mathfrak D_{-,1}^{\beta}\) on the range
\(\mathfrak I_{-,1}^{\beta}(E)\) therefore yields
\[
\mathfrak D_{-,1}^{\beta}\phi=\mathbf c_t^{-1}g.
\]
By the same classical isomorphism, there exist constants \(c_1,c_2>0\), depending only on
\(\beta\), such that
\[
c_1\bigl(\|\phi\|_E+\|\mathfrak D_{-,1}^{\beta}\phi\|_E\bigr)
\le
\|\phi\|_{H^\beta((0,1))}
\le
c_2\bigl(\|\phi\|_E+\|\mathfrak D_{-,1}^{\beta}\phi\|_E\bigr).
\]
Since \(\mathfrak D_{-,1}^{\beta}\phi=\mathbf c_t^{-1}\mathfrak T_t^{\mathrm{fr}}\phi\), this is
equivalent to
\[
\|\phi\|_{H^\beta((0,1))}
\asymp
\|\phi\|_E+\|\mathfrak T_t^{\mathrm{fr}}\phi\|_E
\asymp
\|\phi\|_{\mathfrak G_t^{\mathrm{fr}}},
\]
where the implicit constants depend only on \(\beta\), \(\eta/\nu_Y\), and \(|\mathbf c_t|\). This
proves item~(ii).

Finally, because \(\beta<\tfrac12\), the density theorem for Bessel potential spaces on
bounded intervals yields
\[
\overline{C_c^\infty((0,1))}^{\,H^\beta((0,1))}
=
H^\beta((0,1)).
\]
See \cite[Ch.~2]{Triebel83}. By item~(ii), this is equivalent to density in the frozen graph
norm, proving item~(iii).
\end{proof}

\begin{corollary}[Smooth core as a local graph recovery core]
\label{cor:smooth-core-recovery}
Assume Assumptions~\ref{ass:model} and~\ref{ass:obs-source}, let \(\eta>0\), fix \(t\in(0,T)\)
and \(\alpha\in(0,1)\), and suppose
\[
0<H_{XY}<H_Y<\tfrac12.
\]
Then \(C_c^\infty((0,1))\) is a local graph recovery core in the sense of
Definition~\ref{def:local-recovery-core}.
\end{corollary}

\begin{proof}
By Theorem~\ref{thm:frozen-density}, \(C_c^\infty((0,1))\) is dense in
\(\mathcal D(\mathfrak T_t^{\mathrm{fr}})\) with respect to the frozen graph norm. By
Theorem~\ref{thm:smooth-core}, it also satisfies items~(ii) and~(iii) of
Definition~\ref{def:local-recovery-core}. Hence it is a local graph recovery core.
\end{proof}

\begin{corollary}[Smooth-core local recovery and Mosco convergence]
\label{cor:unconditional-local}
Under the assumptions of Corollary~\ref{cor:smooth-core-recovery}, the conclusions of
Theorem~\ref{thm:local-recovery} and Corollary~\ref{cor:local-mosco} hold without any additional
recovery-core hypothesis.
\end{corollary}

\begin{proof}
Apply Corollary~\ref{cor:smooth-core-recovery} inside
Theorem~\ref{thm:local-recovery} and Corollary~\ref{cor:local-mosco}.
\end{proof}

\section{Frozen buffered localization}\label{sec:frozen-buffer}

The frozen right-sided fractional transfer has exact locality.  Its action inside a boundary layer
depends only on the input on a slightly larger boundary layer.  On the frozen side this localization
is exact, not asymptotic.

\begin{proposition}[Global fractional representation of the frozen transfer]
\label{prop:frozen-global-derivative}
Assume
\[
0<H_{XY}<H_Y<\tfrac12,
\qquad
\beta:=H_Y-H_{XY},
\]
and fix \(t\in(0,T)\). Then for every
\[
\phi\in\mathcal D(\mathfrak T_t^{\mathrm{fr}}),
\]
one has
\begin{equation}\label{eq:frozen-global-derivative}
\mathfrak T_t^{\mathrm{fr}}\phi
=
\mathbf c_t\,\mathfrak D_{-,1}^{\beta}\phi
\qquad\text{in }E.
\end{equation}
\end{proposition}

\begin{proof}
Let
\[
g:=\mathfrak T_t^{\mathrm{fr}}\phi.
\]
By Proposition~\ref{prop:frozen-range},
\[
\phi=\mathbf c_t^{-1}\mathfrak I_{-,1}^{\beta}g.
\]
Applying \(\mathfrak D_{-,1}^{\beta}\), which is the inverse of \(\mathfrak I_{-,1}^{\beta}\) on its
range by the finite-interval Riemann--Liouville isomorphism used in
Theorem~\ref{thm:frozen-density}, gives
\[
\mathfrak D_{-,1}^{\beta}\phi=\mathbf c_t^{-1}g.
\]
Multiplying by \(\mathbf c_t\) proves \eqref{eq:frozen-global-derivative}.
\end{proof}

\begin{definition}[Buffered cutoff pair on the unit interval]\label{def:buffer-pair}
Let \(0<\rho_{\mathrm{in}}<\rho_{\mathrm{out}}<1\). A pair
\[
(\chi_{\mathrm{in}},\chi_{\mathrm{out}})\in C^\infty([0,1])\times C^\infty([0,1])
\]
is called a \emph{buffered cutoff pair} with inner radius \(\rho_{\mathrm{in}}\) and outer radius
\(\rho_{\mathrm{out}}\) if
\begin{align}
0\le \chi_{\mathrm{in}},\chi_{\mathrm{out}} \le 1
\qquad&\text{on }[0,1], \label{eq:buffer-cutoff-bounds}\\
\operatorname{supp}\chi_{\mathrm{in}} \subset [1-\rho_{\mathrm{in}},1],
\qquad&
\operatorname{supp}\chi_{\mathrm{out}} \subset [1-\rho_{\mathrm{out}},1], \label{eq:buffer-cutoff-support}\\
\chi_{\mathrm{out}}\equiv 1
\qquad&\text{on }[1-\rho_{\mathrm{in}},1]. \label{eq:buffer-cutoff-one}
\end{align}
\end{definition}

\begin{proposition}[Smooth cutoffs preserve the frozen graph space]
\label{prop:frozen-cutoff-bounded}
Assume
\[
0<H_{XY}<H_Y<\tfrac12,
\qquad
\beta:=H_Y-H_{XY},
\]
and fix \(t\in(0,T)\). Let \(\chi\in C^\infty([0,1])\). Then multiplication by \(\chi\) maps
\[
\mathcal D(\mathfrak T_t^{\mathrm{fr}})
\]
into itself, and there exists \(C_\chi<\infty\) such that
\begin{equation}\label{eq:frozen-cutoff-bounded}
\|\chi\phi\|_{\mathfrak G_t^{\mathrm{fr}}}
\le
C_\chi\|\phi\|_{\mathfrak G_t^{\mathrm{fr}}}
\qquad\text{for every }\phi\in\mathcal D(\mathfrak T_t^{\mathrm{fr}}).
\end{equation}
\end{proposition}

\begin{proof}
By Theorem~\ref{thm:frozen-density}, one has
\[
\mathcal D(\mathfrak T_t^{\mathrm{fr}})=H^\beta((0,1)),
\]
and the frozen graph norm is equivalent to the \(H^\beta((0,1))\) norm. Since \(\chi\in
C^\infty([0,1])\), the smooth-multiplier theorem on Bessel potential spaces gives
\[
\|\chi\phi\|_{H^\beta((0,1))}
\le
C_\chi'\|\phi\|_{H^\beta((0,1))}
\qquad\text{for every }\phi\in H^\beta((0,1)).
\]
See \cite[Ch.~2]{Triebel83}. 
Returning to the frozen graph norm via Theorem~\ref{thm:frozen-density}(ii) proves
\eqref{eq:frozen-cutoff-bounded}.
\end{proof}

\begin{lemma}[Terminal vanishing propagates to the frozen transfer]
\label{lem:frozen-terminal-zero}
Assume
\[
0<H_{XY}<H_Y<\tfrac12,
\qquad
\beta:=H_Y-H_{XY},
\]
and fix \(t\in(0,T)\). Let \(\phi\in\mathcal D(\mathfrak T_t^{\mathrm{fr}})\), and suppose that for
some \(a\in(0,1)\),
\[
\phi=0
\qquad\text{a.e. on }[a,1].
\]
Then
\begin{equation}\label{eq:frozen-terminal-zero}
\mathfrak T_t^{\mathrm{fr}}\phi=0
\qquad\text{a.e. on }[a,1].
\end{equation}
\end{lemma}

\begin{proof}
Set
\[
g:=\mathbf c_t^{-1}\mathfrak T_t^{\mathrm{fr}}\phi.
\]
By Proposition~\ref{prop:frozen-range},
\[
\phi=\mathfrak I_{-,1}^{\beta}g.
\]
Applying the semigroup identity gives
\[
\mathfrak I_{-,1}^{1-\beta}\phi
=
\mathfrak I_{-,1}^{1}g.
\]
Now fix \(u\in[a,1]\). Since \(\phi(v)=0\) for almost every \(v\in[u,1]\), the integral formula for
\(\mathfrak I_{-,1}^{1-\beta}\phi\) yields
\[
\bigl(\mathfrak I_{-,1}^{1-\beta}\phi\bigr)(u)=0.
\]
Hence
\[
\bigl(\mathfrak I_{-,1}^{1}g\bigr)(u)=\int_u^1 g(v)\,dv=0
\qquad\text{for every }u\in[a,1].
\]
The function \(u\mapsto \mathfrak I_{-,1}^{1}g(u)\) is absolutely continuous, so differentiating
almost everywhere on \([a,1]\) gives \(g=0\) almost everywhere on that interval. Multiplying by
\(\mathbf c_t\) proves \eqref{eq:frozen-terminal-zero}.
\end{proof}

\begin{theorem}[Exact buffered localization for the frozen transfer]
\label{thm:frozen-buffered}
Assume
\[
0<H_{XY}<H_Y<\tfrac12,
\qquad
\beta:=H_Y-H_{XY},
\]
fix \(t\in(0,T)\), and let \((\chi_{\mathrm{in}},\chi_{\mathrm{out}})\) be a buffered cutoff pair in
the sense of Definition~\ref{def:buffer-pair}. Then:
\begin{enumerate}[label=(\roman*)]
\item for every \(\phi\in\mathcal D(\mathfrak T_t^{\mathrm{fr}})\),
\[
\chi_{\mathrm{out}}\phi\in\mathcal D(\mathfrak T_t^{\mathrm{fr}});
\]
\item there exists \(C_{\mathrm{out}}<\infty\) such that
\begin{equation}\label{eq:frozen-buffered-stability}
\|\chi_{\mathrm{out}}\phi\|_{\mathfrak G_t^{\mathrm{fr}}}
\le
C_{\mathrm{out}}\|\phi\|_{\mathfrak G_t^{\mathrm{fr}}}
\qquad\text{for every }\phi\in\mathcal D(\mathfrak T_t^{\mathrm{fr}});
\end{equation}
\item for every \(\phi\in\mathcal D(\mathfrak T_t^{\mathrm{fr}})\),
\begin{equation}\label{eq:frozen-buffered-exact}
\chi_{\mathrm{in}}\mathfrak T_t^{\mathrm{fr}}\phi
=
\chi_{\mathrm{in}}\mathfrak T_t^{\mathrm{fr}}(\chi_{\mathrm{out}}\phi)
\qquad\text{in }E.
\end{equation}
\end{enumerate}
\end{theorem}

\begin{proof}
Parts (i) and (ii) are exactly Proposition~\ref{prop:frozen-cutoff-bounded} with
\(\chi=\chi_{\mathrm{out}}\).

For part (iii), set
\[
\psi:=(1-\chi_{\mathrm{out}})\phi.
\]
By \eqref{eq:buffer-cutoff-one}, one has \(\psi=0\) on \([1-\rho_{\mathrm{in}},1]\). Since
\(\chi_{\mathrm{out}}\phi\in\mathcal D(\mathfrak T_t^{\mathrm{fr}})\) by part (i), also
\[
\psi=\phi-\chi_{\mathrm{out}}\phi\in\mathcal D(\mathfrak T_t^{\mathrm{fr}}).
\]
Lemma~\ref{lem:frozen-terminal-zero} with \(a=1-\rho_{\mathrm{in}}\) therefore gives
\[
\mathfrak T_t^{\mathrm{fr}}\psi=0
\qquad\text{a.e. on }[1-\rho_{\mathrm{in}},1].
\]
Because \(\operatorname{supp}\chi_{\mathrm{in}}\subset [1-\rho_{\mathrm{in}},1]\), this implies
\[
\chi_{\mathrm{in}}\mathfrak T_t^{\mathrm{fr}}\psi=0.
\]
Since
\[
\psi=\phi-\chi_{\mathrm{out}}\phi,
\]
we conclude that
\[
\chi_{\mathrm{in}}\mathfrak T_t^{\mathrm{fr}}\phi
-\chi_{\mathrm{in}}\mathfrak T_t^{\mathrm{fr}}(\chi_{\mathrm{out}}\phi)
=
\chi_{\mathrm{in}}\mathfrak T_t^{\mathrm{fr}}\psi
=0,
\]
which is exactly \eqref{eq:frozen-buffered-exact}.
\end{proof}

\begin{corollary}[Inner-supported frozen graph functionals reduce to buffered profiles]
\label{cor:frozen-buffered-functional}
Assume the hypotheses of Theorem~\ref{thm:frozen-buffered}. Let \(a,b\in E\) satisfy
\[
\operatorname{supp}a\subset \operatorname{supp}\chi_{\mathrm{in}},
\qquad
\operatorname{supp}b\subset \operatorname{supp}\chi_{\mathrm{in}},
\]
and define
\[
\Lambda(\phi)
:=
\langle a,\phi\rangle_E+\langle b,\mathfrak T_t^{\mathrm{fr}}\phi\rangle_E,
\qquad
\phi\in\mathcal D(\mathfrak T_t^{\mathrm{fr}}).
\]
Then \(\Lambda\) is continuous on \(\mathfrak G_t^{\mathrm{fr}}\), and for every
\(\phi\in\mathcal D(\mathfrak T_t^{\mathrm{fr}})\),
\begin{equation}\label{eq:frozen-buffered-functional}
\Lambda(\phi)=\Lambda(\chi_{\mathrm{out}}\phi).
\end{equation}
Consequently,
\begin{equation}\label{eq:frozen-buffered-dual}
\|\Lambda\|_{(\mathfrak G_t^{\mathrm{fr}})^*}
\le
C_{\mathrm{out}}
\sup\left\{
|\Lambda(\psi)|:
\psi\in\mathcal D(\mathfrak T_t^{\mathrm{fr}}),\
\operatorname{supp}\psi\subset \operatorname{supp}\chi_{\mathrm{out}},\
\|\psi\|_{\mathfrak G_t^{\mathrm{fr}}}\le 1
\right\},
\end{equation}
where \(C_{\mathrm{out}}\) is the constant from \eqref{eq:frozen-buffered-stability}.
\end{corollary}

\begin{proof}
Continuity follows from Cauchy--Schwarz and Definition~\ref{def:local-graph-norm}. Since \(a\) and
\(b\) are supported in \(\operatorname{supp}\chi_{\mathrm{in}}\), one has
\[
\langle a,\phi\rangle_E=\langle a,\chi_{\mathrm{out}}\phi\rangle_E,
\]
because \(\chi_{\mathrm{out}}\equiv1\) on \(\operatorname{supp}\chi_{\mathrm{in}}\), and
Theorem~\ref{thm:frozen-buffered}(iii) gives
\[
\chi_{\mathrm{in}}\mathfrak T_t^{\mathrm{fr}}\phi
=
\chi_{\mathrm{in}}\mathfrak T_t^{\mathrm{fr}}(\chi_{\mathrm{out}}\phi).
\]
Therefore
\[
\langle b,\mathfrak T_t^{\mathrm{fr}}\phi\rangle_E
=
\langle b,\mathfrak T_t^{\mathrm{fr}}(\chi_{\mathrm{out}}\phi)\rangle_E,
\]
because \(b\) is supported in \(\operatorname{supp}\chi_{\mathrm{in}}\). Hence
\eqref{eq:frozen-buffered-functional} holds. For \(\phi\in\mathcal D(\mathfrak T_t^{\mathrm{fr}})\)
with \(\|\phi\|_{\mathfrak G_t^{\mathrm{fr}}}\le1\), set \(\psi:=\chi_{\mathrm{out}}\phi\). Then
\(\operatorname{supp}\psi\subset \operatorname{supp}\chi_{\mathrm{out}}\), and by Theorem~\ref{thm:frozen-buffered}(ii),
\[
\|\psi\|_{\mathfrak G_t^{\mathrm{fr}}}\le C_{\mathrm{out}}.
\]
Hence
\[
|\Lambda(\phi)|
=
|\Lambda(\psi)|
\le
C_{\mathrm{out}}
\sup\left\{
|\Lambda(\vartheta)|:
\vartheta\in\mathcal D(\mathfrak T_t^{\mathrm{fr}}),\
\operatorname{supp}\vartheta\subset \operatorname{supp}\chi_{\mathrm{out}},\
\|\vartheta\|_{\mathfrak G_t^{\mathrm{fr}}}\le 1
\right\},
\]
and taking the supremum over graph-unit \(\phi\) proves \eqref{eq:frozen-buffered-dual}.
\end{proof}

\section{Two-scale factorization of the singular functional}\label{sec:two-scale}

The local transfer family on the unit interval captures graph geometry at the coarse scale
\(h^\alpha\).  The gain functional is more singular.  After restriction to that boundary layer, the
increment profiles still concentrate on the smaller micro-scale
\[
\varepsilon_{h,\alpha}:=h^{1-\alpha}
\]
near the endpoint \(u=0\).  A second dilation identifies the exact micro-local operator induced by
the local transfer family and recovers the same universal half-line profiles as the original
\(h\)-dilation.

\begin{definition}[Micro-dilation and micro-lift on the unit interval]
\label{def:micro-dilation}
For \(\varepsilon\in(0,1]\), define
\[
\mathfrak d_\varepsilon:E\to L^2(\R_+),
\qquad
(\mathfrak d_\varepsilon\phi)(r)
:=
\varepsilon^{1/2}\phi(\varepsilon r)\one_{[0,\varepsilon^{-1}]}(r),
\]
for \(\phi\in E=L^2([0,1])\). Also define
\[
\mathfrak l_\varepsilon:L^2(\R_+)\to E,
\qquad
(\mathfrak l_\varepsilon\Phi)(u)
:=
\varepsilon^{-1/2}\Phi\!\left(\frac{u}{\varepsilon}\right),
\qquad 0\le u\le 1,
\]
for \(\Phi\in L^2(\R_+)\). Finally, let
\[
\mathfrak p_\varepsilon:L^2(\R_+)\to L^2(\R_+),
\qquad
(\mathfrak p_\varepsilon\Phi)(r):=\Phi(r)\one_{[0,\varepsilon^{-1}]}(r),
\]
denote truncation to the interval \([0,\varepsilon^{-1}]\).
\end{definition}

\begin{lemma}[Micro-dilation isometry and truncation identities]
\label{lem:micro-dilation}
Let \(\varepsilon\in(0,1]\). Then:
\begin{enumerate}[label=(\roman*)]
\item \(\mathfrak d_\varepsilon:E\to L^2(\R_+)\) is a linear isometric embedding;
\item for every \(\phi\in E\) and \(\Phi\in L^2(\R_+)\),
\begin{equation}\label{eq:micro-adjoint}
\langle \mathfrak d_\varepsilon\phi,\Phi\rangle_{L^2(\R_+)}
=
\langle \phi,\mathfrak l_\varepsilon\Phi\rangle_E;
\end{equation}
\item
\begin{equation}\label{eq:micro-inverse}
\mathfrak l_\varepsilon\mathfrak d_\varepsilon=I_E,
\qquad
\mathfrak d_\varepsilon\mathfrak l_\varepsilon=\mathfrak p_\varepsilon.
\end{equation}
\end{enumerate}
\end{lemma}

\begin{proof}
For \(\phi\in E\),
\[
\|\mathfrak d_\varepsilon\phi\|_{L^2(\R_+)}^2
=
\int_0^{\varepsilon^{-1}}\varepsilon |\phi(\varepsilon r)|^2\,dr
=
\int_0^1 |\phi(u)|^2\,du
=
\|\phi\|_E^2,
\]
which proves (i). Likewise, for \(\phi\in E\) and \(\Phi\in L^2(\R_+)\), the change of variables
\(u=\varepsilon r\) gives
\[
\langle \mathfrak d_\varepsilon\phi,\Phi\rangle_{L^2(\R_+)}
=
\int_0^{\varepsilon^{-1}} \varepsilon^{1/2}\phi(\varepsilon r)\Phi(r)\,dr
=
\int_0^1 \phi(u)\varepsilon^{-1/2}\Phi\!\left(\frac{u}{\varepsilon}\right)\,du
=
\langle \phi,\mathfrak l_\varepsilon\Phi\rangle_E,
\]
which is \eqref{eq:micro-adjoint}. Finally,
\[
(\mathfrak l_\varepsilon\mathfrak d_\varepsilon\phi)(u)
=
\varepsilon^{-1/2}\varepsilon^{1/2}\phi(u)
=
\phi(u),
\qquad 0\le u\le 1,
\]
while
\[
(\mathfrak d_\varepsilon\mathfrak l_\varepsilon\Phi)(r)
=
\varepsilon^{1/2}\varepsilon^{-1/2}\Phi(r)\one_{[0,\varepsilon^{-1}]}(r)
=
(\mathfrak p_\varepsilon\Phi)(r).
\]
This proves \eqref{eq:micro-inverse}.
\end{proof}

\begin{proposition}[Second dilation inside the \(h^\alpha\) boundary layer]
\label{prop:two-scale-factorization}
Fix \(t\in(0,T)\), \(\alpha\in(0,1)\), and \(h\in(0,t]\), and set
\[
\varepsilon_{h,\alpha}:=h^{1-\alpha}.
\]
Then:
\begin{enumerate}[label=(\roman*)]
\item for every \(f\in H_{t,h,\alpha}^{\mathrm{near}}\),
\begin{equation}\label{eq:two-scale-D}
\mathcal D_{t,h}f
=
\mathfrak d_{\varepsilon_{h,\alpha}}\bigl(\mathcal D_{t,h}^{(\alpha)}f\bigr);
\end{equation}
\item for every \(\Phi\in L^2(\R_+)\),
\begin{equation}\label{eq:two-scale-L}
\mathcal L_{t,h}^{(\alpha)}\mathfrak l_{\varepsilon_{h,\alpha}}\Phi
=
\mathcal L_{t,h}\mathfrak p_{\varepsilon_{h,\alpha}}\Phi;
\end{equation}
\item if \(\Phi\in L^2(\R_+)\) satisfies
\[
\mathfrak l_{\varepsilon_{h,\alpha}}\Phi\in \mathfrak G_{t,h,\alpha}^{\mathrm{near}},
\]
then
\begin{equation}\label{eq:two-scale-transfer}
\varepsilon_{h,\alpha}^{\beta}
\mathfrak d_{\varepsilon_{h,\alpha}}
\Bigl(
\mathfrak T_{t,h,\alpha}\mathfrak l_{\varepsilon_{h,\alpha}}\Phi
\Bigr)
=
h^\beta
\mathcal D_{t,h}\mathcal S_t\mathcal L_{t,h}\mathfrak p_{\varepsilon_{h,\alpha}}\Phi,
\end{equation}
where \(\beta:=H_Y-H_{XY}\).
\end{enumerate}
\end{proposition}

\begin{proof}
For \(f\in H_{t,h,\alpha}^{\mathrm{near}}\) and \(r\ge0\),
\begin{align*}
\bigl(\mathfrak d_{\varepsilon_{h,\alpha}}\mathcal D_{t,h}^{(\alpha)}f\bigr)(r)
&=
\varepsilon_{h,\alpha}^{1/2}
h^{\alpha/2}
f\bigl(t-h^\alpha \varepsilon_{h,\alpha}r\bigr)\one_{[0,\varepsilon_{h,\alpha}^{-1}]}(r) \\
&=
h^{1/2}f(t-hr)\one_{[0,\varepsilon_{h,\alpha}^{-1}]}(r).
\end{align*}
Because \(f\) is supported in \([t-h^\alpha,t]\), the right-hand side vanishes for
\(r>\varepsilon_{h,\alpha}^{-1}=h^{\alpha-1}\). Hence it equals \(\mathcal D_{t,h}f(r)\), proving
\eqref{eq:two-scale-D}.

For \(\Phi\in L^2(\R_+)\) and \(0\le s\le t\),
\begin{align*}
\bigl(\mathcal L_{t,h}^{(\alpha)}\mathfrak l_{\varepsilon_{h,\alpha}}\Phi\bigr)(s)
&=
h^{-\alpha/2}
\varepsilon_{h,\alpha}^{-1/2}
\Phi\!\left(\frac{t-s}{h^\alpha\varepsilon_{h,\alpha}}\right)
\one_{[t-h^\alpha,t]}(s) \\
&=
h^{-1/2}
\Phi\!\left(\frac{t-s}{h}\right)
\one_{[t-h^\alpha,t]}(s).
\end{align*}
On the other hand,
\[
\bigl(\mathcal L_{t,h}\mathfrak p_{\varepsilon_{h,\alpha}}\Phi\bigr)(s)
=
h^{-1/2}
\Phi\!\left(\frac{t-s}{h}\right)
\one_{\{0\le (t-s)/h\le \varepsilon_{h,\alpha}^{-1}\}}(s)\one_{[0,t]}(s).
\]
But
\[
0\le \frac{t-s}{h}\le \varepsilon_{h,\alpha}^{-1}
\iff
t-h^\alpha\le s\le t,
\]
so the two expressions coincide. Therefore \eqref{eq:two-scale-L} holds.

For part (iii), Definition~\ref{def:local-transfer} gives
\[
\mathfrak T_{t,h,\alpha}\mathfrak l_{\varepsilon_{h,\alpha}}\Phi
=
h^{\alpha\beta}
\mathcal D_{t,h}^{(\alpha)}
\mathcal S_t
\mathcal L_{t,h}^{(\alpha)}
\mathfrak l_{\varepsilon_{h,\alpha}}\Phi.
\]
Applying \(\varepsilon_{h,\alpha}^{\beta}\mathfrak d_{\varepsilon_{h,\alpha}}\) and using
\eqref{eq:two-scale-D} and \eqref{eq:two-scale-L} yields
\begin{align*}
\varepsilon_{h,\alpha}^{\beta}
\mathfrak d_{\varepsilon_{h,\alpha}}
\Bigl(
\mathfrak T_{t,h,\alpha}\mathfrak l_{\varepsilon_{h,\alpha}}\Phi
\Bigr)
&=
\varepsilon_{h,\alpha}^{\beta}h^{\alpha\beta}
\mathcal D_{t,h}
\mathcal S_t
\mathcal L_{t,h}
\mathfrak p_{\varepsilon_{h,\alpha}}\Phi \\
&=
h^\beta
\mathcal D_{t,h}\mathcal S_t\mathcal L_{t,h}\mathfrak p_{\varepsilon_{h,\alpha}}\Phi,
\end{align*}
because \(\varepsilon_{h,\alpha}^{\beta}h^{\alpha\beta}=h^\beta\). This proves
\eqref{eq:two-scale-transfer}.
\end{proof}

\begin{definition}[Micro-dilated exact and frozen Volterra operators]
\label{def:micro-volterra}
Fix \(t\in(0,T)\), \(\alpha\in(0,1)\), let \(i\in\{Y,XY\}\), and let \(h>0\) be sufficiently small
that \(h^\alpha\le t\). Set
\[
\varepsilon_{h,\alpha}:=h^{1-\alpha}.
\]
The \emph{micro-dilated exact Volterra operator} is
\[
\mathfrak V_{i,t,h,\alpha}^{\mathrm{mic}}:L^2(\R_+)\to L^2(\R_+),
\]
defined by
\begin{equation}\label{eq:micro-volterra}
\bigl(\mathfrak V_{i,t,h,\alpha}^{\mathrm{mic}}\Phi\bigr)(r)
:=
\int_r^{\varepsilon_{h,\alpha}^{-1}}
(q-r)^{H_i-\frac12}
L_i(t-hr,t-hq)\Phi(q)\,dq,
\qquad r\ge0.
\end{equation}
Also define the \emph{frozen half-line Volterra operator}
\[
\mathfrak V_{i,t}^{\mathrm{hl}}:C_c^\infty((0,\infty))\to L^2(\R_+)
\]
by
\begin{equation}\label{eq:halfline-volterra}
\bigl(\mathfrak V_{i,t}^{\mathrm{hl}}\Phi\bigr)(r)
:=
L_i(t,t)\int_r^\infty (q-r)^{H_i-\frac12}\Phi(q)\,dq,
\qquad r\ge0.
\end{equation}
\end{definition}

\begin{proposition}[Second blow-up of the localized Volterra operators]
\label{prop:micro-volterra}
Assume Assumption~\ref{ass:model}. Fix \(t\in(0,T)\), \(\alpha\in(0,1)\), and let
\(i\in\{Y,XY\}\). For every sufficiently small \(h>0\) and every \(\Phi\in L^2(\R_+)\),
\begin{equation}\label{eq:micro-volterra-exact}
\mathfrak d_{\varepsilon_{h,\alpha}}
\mathfrak A_{i,t,h,\alpha}
\mathfrak l_{\varepsilon_{h,\alpha}}\Phi
=
\varepsilon_{h,\alpha}^{H_i}
\mathfrak V_{i,t,h,\alpha}^{\mathrm{mic}}\Phi
\qquad\text{in }L^2(\R_+).
\end{equation}
In particular,
\begin{equation}\label{eq:micro-volterra-support}
\operatorname{supp}\bigl(\mathfrak V_{i,t,h,\alpha}^{\mathrm{mic}}\Phi\bigr)
\subset
[0,\varepsilon_{h,\alpha}^{-1}]
\qquad\text{for every }\Phi\in L^2(\R_+).
\end{equation}
\end{proposition}

\begin{proof}
Fix \(h>0\) sufficiently small and abbreviate \(\varepsilon:=\varepsilon_{h,\alpha}\). For \(r\ge0\),
\begin{align*}
\bigl(
\mathfrak d_{\varepsilon}\mathfrak A_{i,t,h,\alpha}\mathfrak l_{\varepsilon}\Phi
\bigr)(r)
&=
\varepsilon^{1/2}
\bigl(
\mathfrak A_{i,t,h,\alpha}\mathfrak l_{\varepsilon}\Phi
\bigr)(\varepsilon r)
\one_{[0,\varepsilon^{-1}]}(r) \\
&=
\varepsilon^{1/2}
\int_{\varepsilon r}^1
(v-\varepsilon r)^{H_i-\frac12}
L_i(t-h^\alpha\varepsilon r,t-h^\alpha v)\,
\varepsilon^{-1/2}\Phi(v/\varepsilon)\,dv\,
\one_{[0,\varepsilon^{-1}]}(r).
\end{align*}
With the change of variables \(v=\varepsilon q\), and using \(h^\alpha\varepsilon=h\), this becomes
\begin{align*}
\bigl(
\mathfrak d_{\varepsilon}\mathfrak A_{i,t,h,\alpha}\mathfrak l_{\varepsilon}\Phi
\bigr)(r)
&=
\varepsilon^{H_i}
\int_r^{\varepsilon^{-1}}
(q-r)^{H_i-\frac12}
L_i(t-hr,t-hq)\Phi(q)\,dq\,
\one_{[0,\varepsilon^{-1}]}(r) \\
&=
\varepsilon^{H_i}
\bigl(\mathfrak V_{i,t,h,\alpha}^{\mathrm{mic}}\Phi\bigr)(r),
\end{align*}
which proves \eqref{eq:micro-volterra-exact}. The support statement
\eqref{eq:micro-volterra-support} follows from \eqref{eq:micro-volterra}.
\end{proof}

\begin{corollary}[Exact micro-transfer identity through the micro-Volterra pair]
\label{cor:micro-volterra-transfer}
Assume Assumptions~\ref{ass:model} and~\ref{ass:obs-source}, let \(\eta>0\), fix \(t\in(0,T)\),
and let \(\alpha\in(0,1)\). Then for every sufficiently small \(h>0\) and every
\[
\Phi\in\mathcal D(\mathfrak T_{t,h,\alpha}^{\mathrm{mic}}),
\]
one has
\begin{equation}\label{eq:micro-volterra-transfer}
\mathfrak V_{Y,t,h,\alpha}^{\mathrm{mic}}
\bigl(
\mathfrak T_{t,h,\alpha}^{\mathrm{mic}}\Phi
\bigr)
=
\mathfrak V_{XY,t,h,\alpha}^{\mathrm{mic}}\Phi.
\end{equation}
\end{corollary}

\begin{proof}
Set \(\varepsilon:=\varepsilon_{h,\alpha}\). By Definition~\ref{def:micro-local},
\[
\mathfrak T_{t,h,\alpha}^{\mathrm{mic}}\Phi
=
\varepsilon^\beta
\mathfrak d_{\varepsilon}
\mathfrak T_{t,h,\alpha}
\mathfrak l_{\varepsilon}\Phi.
\]
Hence, using \(\mathfrak l_{\varepsilon}\mathfrak d_{\varepsilon}=I_E\) from
Lemma~\ref{lem:micro-dilation},
\begin{align*}
\mathfrak V_{Y,t,h,\alpha}^{\mathrm{mic}}
\bigl(
\mathfrak T_{t,h,\alpha}^{\mathrm{mic}}\Phi
\bigr)
&=
\varepsilon^{-H_Y}
\mathfrak d_{\varepsilon}
\mathfrak A_{Y,t,h,\alpha}
\mathfrak l_{\varepsilon}
\bigl(
\varepsilon^\beta \mathfrak d_{\varepsilon}\mathfrak T_{t,h,\alpha}\mathfrak l_{\varepsilon}\Phi
\bigr) \\
&=
\varepsilon^{-H_Y+\beta}
\mathfrak d_{\varepsilon}
\mathfrak A_{Y,t,h,\alpha}
\mathfrak T_{t,h,\alpha}
\mathfrak l_{\varepsilon}\Phi.
\end{align*}
By the exact local transfer identity \eqref{eq:local-transfer-eq},
\[
\mathfrak A_{Y,t,h,\alpha}\mathfrak T_{t,h,\alpha}
\mathfrak l_{\varepsilon}\Phi
=
\mathfrak A_{XY,t,h,\alpha}\mathfrak l_{\varepsilon}\Phi.
\]
Since \(\beta=H_Y-H_{XY}\), it follows that
\begin{align*}
\mathfrak V_{Y,t,h,\alpha}^{\mathrm{mic}}
\bigl(
\mathfrak T_{t,h,\alpha}^{\mathrm{mic}}\Phi
\bigr)
&=
\varepsilon^{-H_{XY}}
\mathfrak d_{\varepsilon}
\mathfrak A_{XY,t,h,\alpha}
\mathfrak l_{\varepsilon}\Phi \\
&=
\mathfrak V_{XY,t,h,\alpha}^{\mathrm{mic}}\Phi,
\end{align*}
where the last step is Proposition~\ref{prop:micro-volterra}.
\end{proof}

\begin{proposition}[Frozen half-line limit on compactly supported profiles]
\label{prop:micro-volterra-limit}
Assume Assumption~\ref{ass:model}. Fix \(t\in(0,T)\), \(\alpha\in(0,1)\), let \(i\in\{Y,XY\}\),
and let \(\Phi\in C_c^\infty((0,\infty))\). Then
\begin{equation}\label{eq:micro-volterra-limit}
\mathfrak V_{i,t,h,\alpha}^{\mathrm{mic}}\Phi
\longrightarrow
\mathfrak V_{i,t}^{\mathrm{hl}}\Phi
\qquad\text{in }L^2(\R_+)
\end{equation}
as \(h\downarrow0\).
\end{proposition}

\begin{proof}
Choose \(R>0\) such that \(\operatorname{supp}\Phi\subset[0,R]\). For \(h\) sufficiently small one
has \(\varepsilon_{h,\alpha}^{-1}>R\) and \(hR<t\), so for \(0\le r\le R\),
\[
\bigl(\mathfrak V_{i,t,h,\alpha}^{\mathrm{mic}}\Phi\bigr)(r)
=
\int_r^R
(q-r)^{H_i-\frac12}
L_i(t-hr,t-hq)\Phi(q)\,dq,
\]
while \(\mathfrak V_{i,t,h,\alpha}^{\mathrm{mic}}\Phi(r)=0\) for \(r>R\). Likewise,
\[
\bigl(\mathfrak V_{i,t}^{\mathrm{hl}}\Phi\bigr)(r)
=
L_i(t,t)\int_r^R (q-r)^{H_i-\frac12}\Phi(q)\,dq,
\qquad r\ge0.
\]
Since \(L_i\) is continuous near \((t,t)\), one has pointwise convergence
\[
\bigl(\mathfrak V_{i,t,h,\alpha}^{\mathrm{mic}}\Phi\bigr)(r)
\longrightarrow
\bigl(\mathfrak V_{i,t}^{\mathrm{hl}}\Phi\bigr)(r)
\qquad\text{for every }r\ge0.
\]
Moreover, for \(0\le r\le R\),
\[
\left|
\bigl(\mathfrak V_{i,t,h,\alpha}^{\mathrm{mic}}\Phi\bigr)(r)
\right|
\le
C\int_r^R (q-r)^{H_i-\frac12}|\Phi(q)|\,dq,
\]
with \(C\) independent of \(h\). The right-hand side belongs to \(L^2([0,R])\), because the kernel
\((q-r)^{H_i-\frac12}\one_{\{0\le r\le q\le R\}}\) is Hilbert--Schmidt. Therefore dominated
convergence gives \eqref{eq:micro-volterra-limit}.
\end{proof}

\begin{definition}[Micro-dilated local profiles and transfer family]
\label{def:micro-local}
Fix \(t\in(0,T)\), \(\alpha\in(0,1)\), and \(h\in(0,t]\), and set
\[
\varepsilon_{h,\alpha}:=h^{1-\alpha}.
\]
For \(i\in\{Y,XY\}\), define the micro-dilated local increment profile
\begin{equation}\label{eq:micro-profile}
\mathfrak g_{i,t,h,\alpha}^{\mathrm{mic}}
:=
\varepsilon_{h,\alpha}^{-H_i}
\mathfrak d_{\varepsilon_{h,\alpha}}
\left(
h^{-\alpha H_i}
\mathcal D_{t,h}^{(\alpha)}
\Pi_{t,h}^{\mathrm{near},\alpha}g_{i,t,h}
\right)
\in L^2(\R_+).
\end{equation}
Also define the micro-dilated local transfer operator by
\begin{equation}\label{eq:micro-transfer}
\mathfrak T_{t,h,\alpha}^{\mathrm{mic}}:
\mathcal D(\mathfrak T_{t,h,\alpha}^{\mathrm{mic}})\subset L^2(\R_+)\to L^2(\R_+),
\qquad
\mathfrak T_{t,h,\alpha}^{\mathrm{mic}}\Phi
:=
\varepsilon_{h,\alpha}^{\beta}
\mathfrak d_{\varepsilon_{h,\alpha}}
\Bigl(
\mathfrak T_{t,h,\alpha}\mathfrak l_{\varepsilon_{h,\alpha}}\Phi
\Bigr),
\end{equation}
with domain
\[
\mathcal D(\mathfrak T_{t,h,\alpha}^{\mathrm{mic}})
:=
\left\{
\Phi\in L^2(\R_+):
\mathfrak l_{\varepsilon_{h,\alpha}}\Phi\in
\mathfrak G_{t,h,\alpha}^{\mathrm{near}}
\right\}.
\]
\end{definition}

\begin{proposition}[Exact micro-local identities]
\label{prop:micro-local}
Assume Assumptions~\ref{ass:model} and~\ref{ass:obs-source}, let \(\eta>0\), fix \(t\in(0,T)\),
and let \(\alpha\in(0,1)\). Then for every sufficiently small \(h>0\):
\begin{enumerate}[label=(\roman*)]
\item for \(i\in\{Y,XY\}\),
\begin{equation}\label{eq:micro-profile-trunc}
\mathfrak g_{i,t,h,\alpha}^{\mathrm{mic}}
=
\widehat g_{i,t,h}\,\one_{[0,\varepsilon_{h,\alpha}^{-1}]}
\qquad\text{in }L^2(\R_+),
\end{equation}
where \(\widehat g_{i,t,h}\) is defined in \eqref{eq:profile-gh};
 \item for every \(\Phi\in\mathcal D(\mathfrak T_{t,h,\alpha}^{\mathrm{mic}})\),
\begin{equation}\label{eq:micro-transfer-exact}
\mathfrak T_{t,h,\alpha}^{\mathrm{mic}}\Phi
=
h^\beta
\mathcal D_{t,h}\mathcal S_t\mathcal L_{t,h}\mathfrak p_{\varepsilon_{h,\alpha}}\Phi;
\end{equation}
\item for every \(\Phi\in\mathcal D(\mathfrak T_{t,h,\alpha}^{\mathrm{mic}})\),
\begin{equation}\label{eq:micro-functional}
h^{-H_{XY}}
\mathcal F_{t,h}\bigl(\mathcal L_{t,h}\mathfrak p_{\varepsilon_{h,\alpha}}\Phi\bigr)
=
\left\langle
\mathfrak g_{XY,t,h,\alpha}^{\mathrm{mic}},\Phi
\right\rangle_{L^2(\R_+)}
-
\left\langle
\mathfrak g_{Y,t,h,\alpha}^{\mathrm{mic}},
\mathfrak T_{t,h,\alpha}^{\mathrm{mic}}\Phi
\right\rangle_{L^2(\R_+)}.
\end{equation}
\end{enumerate}
\end{proposition}

\begin{proof}
For \(i\in\{Y,XY\}\), \(\Pi_{t,h}^{\mathrm{near},\alpha}g_{i,t,h}\) is supported in
\([t-h^\alpha,t]\), so Proposition~\ref{prop:two-scale-factorization}(i) gives
\[
\mathfrak d_{\varepsilon_{h,\alpha}}
\Bigl(
\mathcal D_{t,h}^{(\alpha)}
\Pi_{t,h}^{\mathrm{near},\alpha}g_{i,t,h}
\Bigr)
=
\mathcal D_{t,h}\Pi_{t,h}^{\mathrm{near},\alpha}g_{i,t,h}.
\]
Multiplying by \(\varepsilon_{h,\alpha}^{-H_i}h^{-\alpha H_i}=h^{-H_i}\) yields
\[
\mathfrak g_{i,t,h,\alpha}^{\mathrm{mic}}
=
h^{-H_i}
\mathcal D_{t,h}\Pi_{t,h}^{\mathrm{near},\alpha}g_{i,t,h}.
\]
Because
\[
\operatorname{supp}\bigl(\Pi_{t,h}^{\mathrm{near},\alpha}g_{i,t,h}\bigr)
\subset [t-h^\alpha,t],
\]
its \(h\)-dilation is supported in \([0,\varepsilon_{h,\alpha}^{-1}]\). Therefore
\[
h^{-H_i}
\mathcal D_{t,h}\Pi_{t,h}^{\mathrm{near},\alpha}g_{i,t,h}
=
\widehat g_{i,t,h}\,\one_{[0,\varepsilon_{h,\alpha}^{-1}]},
\]
which proves \eqref{eq:micro-profile-trunc}.

Part (ii) is exactly Proposition~\ref{prop:two-scale-factorization}(iii).

For part (iii), set
\[
v_h:=\mathcal L_{t,h}\mathfrak p_{\varepsilon_{h,\alpha}}\Phi.
\]
By \eqref{eq:two-scale-L}, one also has
\[
v_h=\mathcal L_{t,h}^{(\alpha)}\mathfrak l_{\varepsilon_{h,\alpha}}\Phi.
\]
Since \(\Phi\in\mathcal D(\mathfrak T_{t,h,\alpha}^{\mathrm{mic}})\), the right-hand side belongs to
\(\mathfrak G_t\), so \(\mathcal F_{t,h}(v_h)\) is well-defined. By
Definition~\ref{def:graphspace}, Lemma~\ref{lem:lift-duality},
Lemma~\ref{lem:dilation-isometry}, and part (ii),
\begin{align*}
h^{-H_{XY}}\mathcal F_{t,h}(v_h)
&=
h^{-H_{XY}}
\langle g_{XY,t,h},v_h\rangle_{H_t}
-
h^{-H_{XY}}\langle g_{Y,t,h},\mathcal S_t v_h\rangle_{H_t} \\
&=
\langle \widehat g_{XY,t,h},\mathfrak p_{\varepsilon_{h,\alpha}}\Phi\rangle_{L^2(\R_+)}
-
\langle \widehat g_{Y,t,h},\mathfrak T_{t,h,\alpha}^{\mathrm{mic}}\Phi\rangle_{L^2(\R_+)}.
\end{align*}
The first term is \eqref{eq:lift-gh} applied to
\(\mathfrak p_{\varepsilon_{h,\alpha}}\Phi\). For the second term,
\[
h^{-H_{XY}}\langle g_{Y,t,h},\mathcal S_t v_h\rangle_{H_t}
=
\left\langle
h^{-H_Y}\mathcal D_{t,h}g_{Y,t,h},
h^\beta \mathcal D_{t,h}\mathcal S_t v_h
\right\rangle_{L^2(\R_+)},
\]
and then applies \eqref{eq:micro-transfer-exact}.
Because \(\widehat g_{XY,t,h}\) and \(\mathfrak g_{XY,t,h,\alpha}^{\mathrm{mic}}\) agree on
\([0,\varepsilon_{h,\alpha}^{-1}]\) by \eqref{eq:micro-profile-trunc},
\[
\langle \widehat g_{XY,t,h},\mathfrak p_{\varepsilon_{h,\alpha}}\Phi\rangle_{L^2(\R_+)}
=
\left\langle
\mathfrak g_{XY,t,h,\alpha}^{\mathrm{mic}},\Phi
\right\rangle_{L^2(\R_+)},
\]
and the same argument with \(i=Y\) gives
\[
\langle \widehat g_{Y,t,h},\mathfrak T_{t,h,\alpha}^{\mathrm{mic}}\Phi\rangle_{L^2(\R_+)}
=
\left\langle
\mathfrak g_{Y,t,h,\alpha}^{\mathrm{mic}},
\mathfrak T_{t,h,\alpha}^{\mathrm{mic}}\Phi
\right\rangle_{L^2(\R_+)}.
\]
This proves \eqref{eq:micro-functional}.
\end{proof}

\begin{theorem}[Micro-profile limit inside every \(h^\alpha\) layer]
\label{thm:micro-profile}
Assume Assumption~\ref{ass:model}. Fix \(t\in(0,T)\), let \(\alpha\in(0,1)\), and let
\(i\in\{Y,XY\}\). Then
\begin{equation}\label{eq:micro-profile-limit}
\mathfrak g_{i,t,h,\alpha}^{\mathrm{mic}}
\longrightarrow
L_i(t,t)\psi_{H_i}
\qquad\text{in }L^2(\R_+)
\end{equation}
as \(h\downarrow0\).
\end{theorem}

\begin{proof}
By Proposition~\ref{prop:micro-local}(i),
\[
\mathfrak g_{i,t,h,\alpha}^{\mathrm{mic}}
=
\widehat g_{i,t,h}\one_{[0,\varepsilon_{h,\alpha}^{-1}]}.
\]
Hence
\begin{align*}
\left\|
\mathfrak g_{i,t,h,\alpha}^{\mathrm{mic}}-L_i(t,t)\psi_{H_i}
\right\|_{L^2(\R_+)}
&\le
\left\|
\widehat g_{i,t,h}-L_i(t,t)\psi_{H_i}
\right\|_{L^2(\R_+)} \\
&\quad+
|L_i(t,t)|
\left\|
\psi_{H_i}\one_{(\varepsilon_{h,\alpha}^{-1},\infty)}
\right\|_{L^2(\R_+)}.
\end{align*}
The first term tends to \(0\) by Proposition~\ref{prop:profile}. Since
\(\varepsilon_{h,\alpha}^{-1}=h^{\alpha-1}\to\infty\) and \(\psi_{H_i}\in L^2(\R_+)\), the second
term also tends to \(0\). This proves \eqref{eq:micro-profile-limit}.
\end{proof}

\begin{remark}[Interpretation of the two-scale analysis]
The exact local transfer family on the unit interval is an intermediate asymptotic object for the
gain functional.  After the first localization to a boundary layer of width \(h^\alpha\), the
singular coefficients still collapse onto the smaller internal scale
\(\varepsilon_{h,\alpha}=h^{1-\alpha}\) near \(u=0\).  Proposition~\ref{prop:micro-local} identifies the
micro-dilated family \(\mathfrak T_{t,h,\alpha}^{\mathrm{mic}}\), rather than
\(\mathfrak T_{t,h,\alpha}\) alone, as the operator governing the singular gain.
Corollary~\ref{cor:micro-volterra-transfer} identifies that family with the micro-dilated
Volterra pair \(\mathfrak V_{Y,t,h,\alpha}^{\mathrm{mic}},
\mathfrak V_{XY,t,h,\alpha}^{\mathrm{mic}}\).  The argument then compares
\(\mathfrak T_{t,h,\alpha}^{\mathrm{mic}}\) with the half-line fractional model
\(\mathcal T_{t,\beta}^{\mathrm{mod}}\) on graph-bounded families and lifts that micro-local
control to the global supremum in Proposition~\ref{prop:graph-dual}.
\end{remark}

\section{Strong micro-local convergence on the smooth core}\label{sec:micro-smooth}

The second blow-up produces the half-line tangent problem.  That tangent problem controls the exact
micro-transfer family on fixed smooth compactly supported profiles.  This is the half-line analogue
of the unit-interval smooth-core theorem from Section~\ref{sec:smooth-core}.

\begin{definition}[Half-line right fractional integral on compact support]
\label{def:halfline-fractional}
For \(\lambda\in(0,1)\) and \(f\in L^2(\R_+)\) with compact support, define
\begin{equation}\label{eq:Iminus-halfline}
(\mathfrak I_-^\lambda f)(r)
:=
\frac{1}{\Gamma(\lambda)}
\int_r^\infty (q-r)^{\lambda-1}f(q)\,dq,
\qquad r\ge0.
\end{equation}
\end{definition}

\begin{lemma}[Half-line semigroup and injectivity on compact supports]
\label{lem:halfline-semigroup}
Let \(\lambda,\mu\in(0,1)\) with \(\lambda+\mu\le 1\), and let \(f\in L^2(\R_+)\) have compact
support. Then:
\begin{enumerate}[label=(\roman*)]
\item
\begin{equation}\label{eq:halfline-semigroup}
\mathfrak I_-^\lambda \mathfrak I_-^\mu f
=
\mathfrak I_-^{\lambda+\mu}f;
\end{equation}
\item if \(\mathfrak I_-^\lambda f=0\), then \(f=0\) almost everywhere.
\end{enumerate}
\end{lemma}

\begin{proof}
For almost every \(r\ge0\), Fubini's theorem gives
\begin{align*}
(\mathfrak I_-^\lambda \mathfrak I_-^\mu f)(r)
&=
\frac{1}{\Gamma(\lambda)\Gamma(\mu)}
\int_r^\infty \int_s^\infty
(s-r)^{\lambda-1}(q-s)^{\mu-1}f(q)\,dq\,ds \\
&=
\frac{1}{\Gamma(\lambda)\Gamma(\mu)}
\int_r^\infty
\left(
\int_r^q (s-r)^{\lambda-1}(q-s)^{\mu-1}\,ds
\right)f(q)\,dq.
\end{align*}
With \(s=r+(q-r)\theta\), the inner integral is
\[
(q-r)^{\lambda+\mu-1}
\frac{\Gamma(\lambda)\Gamma(\mu)}{\Gamma(\lambda+\mu)},
\]
which proves \eqref{eq:halfline-semigroup}.

For injectivity, assume \(\mathfrak I_-^\lambda f=0\). Applying
\eqref{eq:halfline-semigroup} with \(\mu=1-\lambda\) yields
\[
\mathfrak I_-^1 f=0.
\]
Since
\[
(\mathfrak I_-^1 f)(r)=\int_r^\infty f(q)\,dq,
\]
the function \(\mathfrak I_-^1 f\) is absolutely continuous with derivative \(-f\). Hence
\(f=0\) almost everywhere.
\end{proof}

\begin{proposition}[Conjugated target-side micro operator on fixed compact supports]
\label{prop:BY-micro}
Assume Assumption~\ref{ass:model}, fix \(t\in(0,T)\), \(\alpha\in(0,1)\), and \(R>0\), and set
\[
\lambda:=H_Y+\frac12\in(0,1),
\qquad
E_R:=L^2([0,R]),
\]
identified with its zero extension to \(L^2(\R_+)\). Then for all sufficiently small \(h>0\)
satisfying
\[
hR<t,
\qquad
R<\varepsilon_{h,\alpha}^{-1},
\]
there exists a bounded operator
\[
\mathfrak B_{Y,t,h,\alpha,R}^{\mathrm{mic}}\in \mathcal L(E_R)
\]
such that:
\begin{enumerate}[label=(\roman*)]
\item for every \(\Psi\in E_R\), the function
\(\mathfrak I_-^{1-\lambda}(\mathfrak V_{Y,t,h,\alpha}^{\mathrm{mic}}\Psi)\) is absolutely
continuous on \([0,\infty)\), and
\begin{equation}\label{eq:BY-micro-conj}
\frac{1}{\Gamma(\lambda)}
\mathfrak D_-^\lambda
\bigl(
\mathfrak V_{Y,t,h,\alpha}^{\mathrm{mic}}\Psi
\bigr)
=
\mathfrak B_{Y,t,h,\alpha,R}^{\mathrm{mic}}\Psi
\qquad\text{in }E_R;
\end{equation}
\item there exists \(C_{t,R}<\infty\) such that
\begin{equation}\label{eq:BY-micro-close}
\left\|
\mathfrak B_{Y,t,h,\alpha,R}^{\mathrm{mic}}-L_Y(t,t)I_{E_R}
\right\|_{\mathcal L(E_R)}
\le
C_{t,R}h
\end{equation}
for all sufficiently small \(h>0\).
\end{enumerate}
\end{proposition}

\begin{proof}
For \(0\le r\le q\le R\), define
\[
a_{Y,t,h}(r,q):=L_Y(t-hr,t-hq).
\]
Let \(\Psi\in E_R\). Because \(\Psi\) is supported in \([0,R]\) and \(R<\varepsilon_{h,\alpha}^{-1}\),
\eqref{eq:micro-volterra} gives
\[
\bigl(
\mathfrak V_{Y,t,h,\alpha}^{\mathrm{mic}}\Psi
\bigr)(r)
=
\int_r^R (q-r)^{\lambda-1}a_{Y,t,h}(r,q)\Psi(q)\,dq,
\qquad r\ge0,
\]
with the understanding that the right-hand side is \(0\) for \(r>R\).

Set
\[
f_h:=\mathfrak V_{Y,t,h,\alpha}^{\mathrm{mic}}\Psi.
\]
For \(0\le r\le R\), Definition~\ref{def:halfline-fractional} and Fubini's theorem give
\begin{align*}
(\mathfrak I_-^{1-\lambda}f_h)(r)
&=
\frac{1}{\Gamma(1-\lambda)}
\int_r^R\int_s^R
(s-r)^{-\lambda}(q-s)^{\lambda-1}
a_{Y,t,h}(s,q)\Psi(q)\,dq\,ds \\
&=
\int_r^R F_{Y,t,h,R}^{\mathrm{mic}}(r,q)\Psi(q)\,dq,
\end{align*}
where
\[
F_{Y,t,h,R}^{\mathrm{mic}}(r,q)
:=
\frac{1}{\Gamma(1-\lambda)}
\int_0^1
s^{-\lambda}(1-s)^{\lambda-1}
a_{Y,t,h}(r+(q-r)s,q)\,ds.
\]
Since \(L_Y\) is \(C^1\), the function \(F_{Y,t,h,R}^{\mathrm{mic}}\) is continuous on
\(\Delta_R:=\{0\le r\le q\le R\}\), and its first derivative in \(r\) is bounded. Therefore
\(\mathfrak I_-^{1-\lambda}f_h\) is absolutely continuous on \([0,R]\), vanishes on \([R,\infty)\),
and differentiation under the integral sign yields for almost every \(r\in[0,R]\),
\[
\mathfrak D_-^\lambda f_h(r)
=
\Gamma(\lambda)a_{Y,t,h}(r,r)\Psi(r)
-\int_r^R G_{Y,t,h,R}^{\mathrm{mic}}(r,q)\Psi(q)\,dq,
\]
where
\[
G_{Y,t,h,R}^{\mathrm{mic}}(r,q)
:=
\frac{\Gamma(\lambda)}{\Gamma(\lambda)\Gamma(1-\lambda)}
\int_0^1
s^{-\lambda}(1-s)^\lambda
\partial_1 a_{Y,t,h}(r+(q-r)s,q)\,ds.
\]
Define
\begin{align*}
\bigl(
\mathfrak B_{Y,t,h,\alpha,R}^{\mathrm{mic}}\Psi
\bigr)(r)
:={}&
a_{Y,t,h}(r,r)\Psi(r) \\
&-
\frac{1}{\Gamma(\lambda)\Gamma(1-\lambda)}
\int_r^R\int_0^1
s^{-\lambda}(1-s)^\lambda
\partial_1 a_{Y,t,h}(r+(q-r)s,q)\,ds\,\Psi(q)\,dq
\end{align*}
for \(0\le r\le R\). This proves \eqref{eq:BY-micro-conj}.

To prove \eqref{eq:BY-micro-close}, note first that
\[
|a_{Y,t,h}(r,r)-L_Y(t,t)|\le C_{t,R}h
\qquad\text{for }0\le r\le R,
\]
because \(L_Y\) is \(C^1\). Also,
\[
\partial_1 a_{Y,t,h}(x,q)
=
-h\,\partial_1L_Y(t-hx,t-hq),
\]
so \(|\partial_1 a_{Y,t,h}(x,q)|\le C_{t,R}h\) uniformly on \(\Delta_R\). Set
\[
D_{t,h,R}:=
\mathfrak B_{Y,t,h,\alpha,R}^{\mathrm{mic}}-a_{Y,t,h}(\cdot,\cdot)_{\mathrm{diag}}I_{E_R}.
\]
Then \(D_{t,h,R}\) has integral kernel bounded by \(C_{t,R}h\one_{\Delta_R}\). Its
Hilbert--Schmidt norm is \(O_{t,R}(h)\), which proves \eqref{eq:BY-micro-close}.
\end{proof}

\begin{corollary}[Uniform invertibility of the conjugated micro target operator]
\label{cor:BY-micro-invertible}
Under the assumptions of Proposition~\ref{prop:BY-micro}, there exist \(h_0\in(0,t]\) and
\(C_{t,R}<\infty\) such that for every \(0<h\le h_0\),
\[
\mathfrak B_{Y,t,h,\alpha,R}^{\mathrm{mic}}:E_R\to E_R
\]
is invertible and
\begin{equation}\label{eq:BY-micro-inv}
\left\|
\bigl(
\mathfrak B_{Y,t,h,\alpha,R}^{\mathrm{mic}}
\bigr)^{-1}
\right\|_{\mathcal L(E_R)}
\le
C_{t,R}.
\end{equation}
\end{corollary}

\begin{proof}
By Assumption~\ref{ass:model}, \(L_Y(t,t)>0\). Proposition~\ref{prop:BY-micro} gives
\[
\left\|
\frac{1}{L_Y(t,t)}\mathfrak B_{Y,t,h,\alpha,R}^{\mathrm{mic}}-I_{E_R}
\right\|_{\mathcal L(E_R)}
\le
\frac{C_{t,R}}{L_Y(t,t)}h.
\]
For \(h\) small enough the right-hand side is \(<1/2\), so the Neumann series gives
invertibility and \eqref{eq:BY-micro-inv}.
\end{proof}

\begin{proposition}[Conjugated micro source asymptotics on the half-line smooth core]
\label{prop:q-micro}
Assume Assumption~\ref{ass:model}, fix \(t\in(0,T)\), \(\alpha\in(0,1)\), and suppose
\[
0<H_{XY}<H_Y<\tfrac12.
\]
Set
\[
\lambda:=H_Y+\frac12,
\qquad
\mu:=H_{XY}+\frac12,
\qquad
\beta:=H_Y-H_{XY}.
\]
Let \(\Phi\in C_c^\infty((0,\infty))\), and choose \(R>0\) such that
\[
\operatorname{supp}\Phi\subset [0,R).
\]
Then there exists \(C_{\Phi,t,R}<\infty\) such that for all sufficiently small \(h>0\) with
\(hR<t\) and \(R<\varepsilon_{h,\alpha}^{-1}\),
\begin{equation}\label{eq:q-micro}
\left\|
\frac{1}{\Gamma(\lambda)}
\mathfrak D_-^\lambda
\bigl(
\mathfrak V_{XY,t,h,\alpha}^{\mathrm{mic}}\Phi
\bigr)
-L_Y(t,t)\mathcal T_{t,\beta}^{\mathrm{mod}}\Phi
\right\|_{L^2(\R_+)}
\le
C_{\Phi,t,R}h.
\end{equation}
\end{proposition}

\begin{proof}
Define
\[
a_{XY,t,h}(r,q):=L_{XY}(t-hr,t-hq),
\qquad
0\le r\le q\le R.
\]
For \(0\le r\le R\),
\[
\bigl(
\mathfrak V_{XY,t,h,\alpha}^{\mathrm{mic}}\Phi
\bigr)(r)
=
\int_r^R (q-r)^{\mu-1}a_{XY,t,h}(r,q)\Phi(q)\,dq,
\]
and the left-hand side vanishes for \(r>R\).

Using Definition~\ref{def:halfline-fractional} and the change of variables
\(s=r+(q-r)\theta\), one obtains for almost every \(r\in[0,R]\),
\begin{equation}\label{eq:q-micro-representation}
\frac{1}{\Gamma(\lambda)}
\mathfrak D_-^\lambda
\bigl(
\mathfrak V_{XY,t,h,\alpha}^{\mathrm{mic}}\Phi
\bigr)(r)
=
-\frac{d}{dr}
\int_r^R (q-r)^{-\beta}g_{t,h,R}^{\mathrm{mic}}(r,q)\Phi(q)\,dq,
\end{equation}
where
\[
g_{t,h,R}^{\mathrm{mic}}(r,q)
:=
\frac{1}{\Gamma(\lambda)\Gamma(1-\lambda)}
\int_0^1
\theta^{-\lambda}(1-\theta)^{\mu-1}
a_{XY,t,h}(r+(q-r)\theta,q)\,d\theta.
\]
Because \(L_{XY}\) is \(C^1\), one has
\begin{equation}\label{eq:g-micro-close}
\bigl\|
g_{t,h,R}^{\mathrm{mic}}-g_{t}^{\mathrm{hl}}
\bigr\|_{C^1(\Delta_R)}
\le
C_{t,R}h,
\end{equation}
where
\[
g_t^{\mathrm{hl}}
:=
\frac{L_{XY}(t,t)\Gamma(\mu)}{\Gamma(\lambda)\Gamma(1-\beta)}
=
\frac{L_Y(t,t)\mathbf c_t}{\Gamma(1-\beta)}.
\]

Now write
\[
Q_{t,h,\Phi}^{\mathrm{mic}}(r)
:=
\frac{1}{\Gamma(\lambda)}
\mathfrak D_-^\lambda
\bigl(
\mathfrak V_{XY,t,h,\alpha}^{\mathrm{mic}}\Phi
\bigr)(r)
-L_Y(t,t)\mathcal T_{t,\beta}^{\mathrm{mod}}\Phi(r).
\]
Using \eqref{eq:q-micro-representation}, the constancy of \(g_t^{\mathrm{hl}}\), and
\eqref{eq:model-transfer}, this becomes
\[
Q_{t,h,\Phi}^{\mathrm{mic}}(r)
=
-\frac{d}{dr}
\int_r^R
(q-r)^{-\beta}
\bigl(
g_{t,h,R}^{\mathrm{mic}}(r,q)-g_t^{\mathrm{hl}}
\bigr)\Phi(q)\,dq.
\]
Since \(\operatorname{supp}\Phi\subset [0,R)\), there is no boundary term at \(q=R\). With \(q=r+u\),
\begin{align*}
Q_{t,h,\Phi}^{\mathrm{mic}}(r)
&=
-\int_0^{R-r} u^{-\beta}
\Bigl[
\bigl(\partial_1+\partial_2\bigr)
\bigl(
g_{t,h,R}^{\mathrm{mic}}-g_t^{\mathrm{hl}}
\bigr)(r,r+u)\,\Phi(r+u) \\
&\qquad\qquad\qquad\qquad\qquad
+\bigl(
g_{t,h,R}^{\mathrm{mic}}-g_t^{\mathrm{hl}}
\bigr)(r,r+u)\,\Phi'(r+u)
\Bigr]\,du.
\end{align*}
By \eqref{eq:g-micro-close}, the bracketed term is \(O_{\Phi,t,R}(h)\) uniformly in \(r\) and
\(u\). Since \(\beta\in(0,1)\),
\[
\int_0^{R-r} u^{-\beta}\,du\le \frac{R^{1-\beta}}{1-\beta}
\qquad\text{for }0\le r\le R.
\]
Hence \(\|Q_{t,h,\Phi}^{\mathrm{mic}}\|_{L^\infty([0,R])}\le C_{\Phi,t,R}h\), which proves
\eqref{eq:q-micro}.
\end{proof}

\begin{theorem}[Strong micro-local convergence on the half-line smooth core]
\label{thm:micro-smooth-core}
Assume Assumptions~\ref{ass:model} and~\ref{ass:obs-source}, let \(\eta>0\), fix
\(t\in(0,T)\), \(\alpha\in(0,1)\), and suppose
\[
0<H_{XY}<H_Y<\tfrac12.
\]
Then for every \(\Phi\in C_c^\infty((0,\infty))\) there exist \(R_\Phi<\infty\),
\(h_\Phi\in(0,t]\), and \(C_\Phi<\infty\) such that for every \(0<h\le h_\Phi\),
\[
\operatorname{supp}\Phi\subset [0,R_\Phi),
\qquad
\Phi\in \mathcal D(\mathfrak T_{t,h,\alpha}^{\mathrm{mic}}),
\]
\(\mathfrak T_{t,h,\alpha}^{\mathrm{mic}}\Phi\) is supported in \([0,R_\Phi]\), and
\begin{equation}\label{eq:micro-smooth-core}
\left\|
\mathfrak T_{t,h,\alpha}^{\mathrm{mic}}\Phi
-\mathcal T_{t,\beta}^{\mathrm{mod}}\Phi
\right\|_{L^2(\R_+)}
\le
C_\Phi h.
\end{equation}
\end{theorem}

\begin{proof}
Set
\[
\lambda:=H_Y+\frac12,
\qquad
\beta:=H_Y-H_{XY}.
\]
Fix \(\Phi\in C_c^\infty((0,\infty))\), and choose \(R_\Phi<\infty\) such that
\[
\operatorname{supp}\Phi\subset [0,R_\Phi).
\]
For \(h>0\) sufficiently small, define
\[
q_{t,h,\Phi}^{\mathrm{mic}}
:=
\frac{1}{\Gamma(\lambda)}
\mathfrak D_-^\lambda
\bigl(
\mathfrak V_{XY,t,h,\alpha}^{\mathrm{mic}}\Phi
\bigr)
\in E_{R_\Phi}.
\]
By Corollary~\ref{cor:BY-micro-invertible}, the operator
\(\mathfrak B_{Y,t,h,\alpha,R_\Phi}^{\mathrm{mic}}\) is invertible for small \(h\), so we may set
\[
\Psi_h
:=
\bigl(
\mathfrak B_{Y,t,h,\alpha,R_\Phi}^{\mathrm{mic}}
\bigr)^{-1}
q_{t,h,\Phi}^{\mathrm{mic}}
\in E_{R_\Phi},
\]
again identified with its zero extension to \(L^2(\R_+)\).

Using Proposition~\ref{prop:q-micro}, Proposition~\ref{prop:BY-micro}, and
\eqref{eq:BY-micro-inv},
\begin{align*}
\left\|
\Psi_h-\mathcal T_{t,\beta}^{\mathrm{mod}}\Phi
\right\|_{E_{R_\Phi}}
&\le
\left\|
\bigl(
\mathfrak B_{Y,t,h,\alpha,R_\Phi}^{\mathrm{mic}}
\bigr)^{-1}
\right\|_{\mathcal L(E_{R_\Phi})}
\Bigl(
\left\|
q_{t,h,\Phi}^{\mathrm{mic}}
-L_Y(t,t)\mathcal T_{t,\beta}^{\mathrm{mod}}\Phi
\right\|_{E_{R_\Phi}} \\
&\qquad\qquad\qquad\qquad
+\left\|
\bigl(
L_Y(t,t)I_{E_{R_\Phi}}
-\mathfrak B_{Y,t,h,\alpha,R_\Phi}^{\mathrm{mic}}
\bigr)
\mathcal T_{t,\beta}^{\mathrm{mod}}\Phi
\right\|_{E_{R_\Phi}}
\Bigr) \\
&\le
C_\Phi h.
\end{align*}
Since \(\mathcal T_{t,\beta}^{\mathrm{mod}}\Phi\) is supported in \([0,R_\Phi]\), this already
gives
\begin{equation}\label{eq:micro-psi-est}
\left\|
\Psi_h-\mathcal T_{t,\beta}^{\mathrm{mod}}\Phi
\right\|_{L^2(\R_+)}
\le
C_\Phi h.
\end{equation}

Now set
\[
r_h
:=
\mathfrak V_{Y,t,h,\alpha}^{\mathrm{mic}}\Psi_h
-\mathfrak V_{XY,t,h,\alpha}^{\mathrm{mic}}\Phi.
\]
Because both \(\Psi_h\) and \(\Phi\) are supported in \([0,R_\Phi]\), so is \(r_h\). By the
definitions of \(\Psi_h\) and \(q_{t,h,\Phi}^{\mathrm{mic}}\), together with
Proposition~\ref{prop:BY-micro},
\[
\frac{1}{\Gamma(\lambda)}
\mathfrak D_-^\lambda r_h
=
\mathfrak B_{Y,t,h,\alpha,R_\Phi}^{\mathrm{mic}}\Psi_h
-q_{t,h,\Phi}^{\mathrm{mic}}
=
0
\qquad\text{in }E_{R_\Phi}.
\]
Hence, by Proposition~\ref{prop:BY-micro} for the first term and the proof of
Proposition~\ref{prop:q-micro} for the second, \(\mathfrak I_-^{1-\lambda}r_h\) is absolutely
continuous with zero derivative, so it is constant on \([0,\infty)\). Since \(r_h\) is compactly supported, one has
\[
\bigl(
\mathfrak I_-^{1-\lambda}r_h
\bigr)(r)=0
\qquad\text{for }r\ge R_\Phi,
\]
and therefore \(\mathfrak I_-^{1-\lambda}r_h\equiv0\). By
Lemma~\ref{lem:halfline-semigroup}, this implies
\[
r_h=0,
\]
that is,
\begin{equation}\label{eq:micro-YX-eq}
\mathfrak V_{Y,t,h,\alpha}^{\mathrm{mic}}\Psi_h
=
\mathfrak V_{XY,t,h,\alpha}^{\mathrm{mic}}\Phi.
\end{equation}

Set \(\varepsilon:=\varepsilon_{h,\alpha}\). Since \(R_\Phi<\varepsilon^{-1}\) for \(h\) small,
Proposition~\ref{prop:micro-volterra} and \eqref{eq:micro-YX-eq} give
\[
\mathfrak d_{\varepsilon}
\mathfrak A_{Y,t,h,\alpha}
\mathfrak l_{\varepsilon}\Psi_h
=
\varepsilon^{H_Y}
\mathfrak V_{Y,t,h,\alpha}^{\mathrm{mic}}\Psi_h
=
\varepsilon^\beta
\mathfrak d_{\varepsilon}
\mathfrak A_{XY,t,h,\alpha}
\mathfrak l_{\varepsilon}\Phi.
\]
Applying \(\mathfrak l_{\varepsilon}\) and using
\(\mathfrak l_{\varepsilon}\mathfrak d_{\varepsilon}=I_E\) from
Lemma~\ref{lem:micro-dilation} yields
\begin{equation}\label{eq:micro-local-eq}
\mathfrak A_{Y,t,h,\alpha}
\bigl(
\varepsilon^{-\beta}\mathfrak l_{\varepsilon}\Psi_h
\bigr)
=
\mathfrak A_{XY,t,h,\alpha}\mathfrak l_{\varepsilon}\Phi.
\end{equation}

Now define
\[
w_h:=\mathcal L_{t,h}^{(\alpha)}\mathfrak l_{\varepsilon}\Phi
\in H_{t,h,\alpha}^{\mathrm{near}},
\qquad
u_h:=h^{-\alpha\beta}
\mathcal L_{t,h}^{(\alpha)}
\bigl(
\varepsilon^{-\beta}\mathfrak l_{\varepsilon}\Psi_h
\bigr)
\in H_{t,h,\alpha}^{\mathrm{near}}.
\]
Using Proposition~\ref{prop:local-volterra} and \eqref{eq:micro-local-eq},
\begin{align*}
\mathcal D_{t,h}^{(\alpha)}\mathcal K_{Y,t}u_h
&=
h^{-\alpha\beta+\alpha H_Y}
\mathfrak A_{Y,t,h,\alpha}
\bigl(
\varepsilon^{-\beta}\mathfrak l_{\varepsilon}\Psi_h
\bigr) \\
&=
h^{\alpha H_{XY}}
\mathfrak A_{XY,t,h,\alpha}\mathfrak l_{\varepsilon}\Phi \\
&=
\mathcal D_{t,h}^{(\alpha)}\mathcal K_{XY,t}w_h.
\end{align*}
Both \(\mathcal K_{Y,t}u_h\) and \(\mathcal K_{XY,t}w_h\) are supported in \([t-h^\alpha,t]\), so
Lemma~\ref{lem:unit-dilation} implies
\[
\mathcal K_{Y,t}u_h=\mathcal K_{XY,t}w_h.
\]
Hence \(w_h\in\mathcal D(\mathcal S_t)\), \(\mathcal S_t w_h=u_h\), and
\[
\mathfrak l_{\varepsilon}\Phi
=
\mathcal D_{t,h}^{(\alpha)}w_h
\in
\mathfrak G_{t,h,\alpha}^{\mathrm{near}}.
\]
Therefore \(\Phi\in\mathcal D(\mathfrak T_{t,h,\alpha}^{\mathrm{mic}})\), and by
Definitions~\ref{def:local-transfer} and~\ref{def:micro-local},
\begin{align*}
\mathfrak T_{t,h,\alpha}^{\mathrm{mic}}\Phi
&=
\varepsilon^\beta
\mathfrak d_{\varepsilon}
\mathfrak T_{t,h,\alpha}\mathfrak l_{\varepsilon}\Phi \\
&=
\varepsilon^\beta
\mathfrak d_{\varepsilon}
\bigl(
\varepsilon^{-\beta}\mathfrak l_{\varepsilon}\Psi_h
\bigr) \\
&=
\mathfrak d_{\varepsilon}\mathfrak l_{\varepsilon}\Psi_h
=
\mathfrak p_{\varepsilon}\Psi_h
=
\Psi_h,
\end{align*}
because \(\operatorname{supp}\Psi_h\subset[0,R_\Phi]\subset[0,\varepsilon^{-1}]\). Combining this
identity with \eqref{eq:micro-psi-est} proves \eqref{eq:micro-smooth-core}.
\end{proof}

\section{Fixed-window micro graph theory}\label{sec:micro-window}

The half-line micro family remains globally nonlocal because the right-sided transfer law depends
on the full future tail.  On each fixed window \([0,R]\), the exact micro family closes: inputs
supported in \([0,R]\) produce outputs supported in \([0,R]\).  The problem becomes a
finite-interval graph convergence problem.  Tail leakage is the only nonlocal obstruction to a
global theorem.

\begin{corollary}[General support preservation near the terminal time]
\label{cor:general-support-preserve}
Assume Assumptions~\ref{ass:model} and~\ref{ass:obs-source}, let \(\eta>0\), fix \(t\in(0,T)\),
and let \(a\in(0,t)\). If \(w\in\mathfrak G_t\) satisfies
\[
\operatorname{supp}(w)\subset[a,t],
\]
then
\[
\operatorname{supp}(\mathcal S_t w)\subset[a,t].
\]
\end{corollary}

\begin{proof}
The argument is identical to that of Proposition~\ref{prop:support-preserve}, applied to the
truncated operator \(\mathcal K_{Y,a}\).
\end{proof}

\begin{definition}[Fixed-window half-line Volterra model and graph norm]
\label{def:micro-window-model}
Assume
\[
0<H_{XY}<H_Y<\tfrac12,
\qquad
\beta:=H_Y-H_{XY},
\]
fix \(t\in(0,T)\), and let \(R>0\). Set
\[
E_R:=L^2([0,R]),
\]
identified with its zero extension to \(L^2(\R_+)\). For \(H\in(0,\tfrac12)\), define
\[
\mathfrak J_{H,R}:E_R\to E_R,
\qquad
(\mathfrak J_{H,R}\Phi)(r):=\int_r^R (q-r)^{H-\frac12}\Phi(q)\,dq,
\qquad 0\le r\le R.
\]
Define the fixed-window model transfer domain
\[
\mathcal D(\mathfrak T_{t,R}^{\mathrm{hl}})
:=
\left\{
\Phi\in E_R:
\mathfrak J_{H_{XY},R}\Phi\in \operatorname{Ran}(\mathfrak J_{H_Y,R})
\right\},
\]
and the corresponding transfer operator
\[
\mathfrak T_{t,R}^{\mathrm{hl}}:
\mathcal D(\mathfrak T_{t,R}^{\mathrm{hl}})\subset E_R\to E_R,
\qquad
\mathfrak T_{t,R}^{\mathrm{hl}}\Phi
:=
\frac{L_{XY}(t,t)}{L_Y(t,t)}
\mathfrak J_{H_Y,R}^{-1}\mathfrak J_{H_{XY},R}\Phi.
\]
For \(\Phi\in\mathcal D(\mathfrak T_{t,R}^{\mathrm{hl}})\), define the fixed-window model graph norm
\begin{equation}\label{eq:micro-window-model-norm}
\|\Phi\|_{\mathfrak G_{t,R}^{\mathrm{hl}}}^2
:=
\|\Phi\|_{E_R}^2
+\frac{\eta^2}{\nu_Y^2}\|\mathfrak T_{t,R}^{\mathrm{hl}}\Phi\|_{E_R}^2.
\end{equation}
\end{definition}

\begin{theorem}[Fractional-derivative representation of the fixed-window half-line model]
\label{thm:micro-window-model}
Assume the hypotheses of Definition~\ref{def:micro-window-model}. Then:
\begin{enumerate}[label=(\roman*)]
\item for every \(\Phi\in C_c^\infty((0,R))\),
\begin{equation}\label{eq:micro-window-model-smooth}
\mathfrak T_{t,R}^{\mathrm{hl}}\Phi
=
\mathcal T_{t,\beta}^{\mathrm{mod}}\Phi
\qquad\text{in }E_R;
\end{equation}
\item \(\mathcal D(\mathfrak T_{t,R}^{\mathrm{hl}})=H^\beta((0,R))\) as sets;
\item the graph norm \(\|\cdot\|_{\mathfrak G_{t,R}^{\mathrm{hl}}}\) is equivalent to the
\(H^\beta((0,R))\) norm;
\item \(C_c^\infty((0,R))\) is dense in \(\mathcal D(\mathfrak T_{t,R}^{\mathrm{hl}})\) with
respect to \(\|\cdot\|_{\mathfrak G_{t,R}^{\mathrm{hl}}}\).
\end{enumerate}
\end{theorem}

\begin{proof}
Define the unitary scaling map
\[
U_R:E_R\to E=L^2([0,1]),
\qquad
(U_R\Phi)(u):=R^{1/2}\Phi(Ru),
\qquad 0\le u\le 1.
\]
For \(H\in(0,\tfrac12)\) and \(\Phi\in E_R\), a change of variables gives
\begin{equation}\label{eq:window-scaling-J}
U_R\mathfrak J_{H,R}U_R^{-1}
=
R^{H+\frac12}\mathfrak J_H.
\end{equation}
In particular, \(\mathfrak J_{H,R}\) is injective because \(\mathfrak J_H\) is injective by
Lemma~\ref{lem:unit-semigroup}.
Hence \(\Phi\in\mathcal D(\mathfrak T_{t,R}^{\mathrm{hl}})\) if and only if \(U_R\Phi\in
\mathcal D(\mathfrak T_t^{\mathrm{fr}})\), and in that case
\begin{equation}\label{eq:window-scaling-T}
U_R\mathfrak T_{t,R}^{\mathrm{hl}}U_R^{-1}
=
R^{-\beta}\mathfrak T_t^{\mathrm{fr}}.
\end{equation}

Now let \(\Phi\in C_c^\infty((0,R))\). Then \(U_R\Phi\in C_c^\infty((0,1))\), so
Proposition~\ref{prop:frozen-fractional}(ii) gives
\[
\mathfrak T_t^{\mathrm{fr}}(U_R\Phi)
=
\mathbf c_t\,\mathfrak D_{-,1}^{\beta}(U_R\Phi).
\]
Because \(\Phi\) vanishes near \(R\), a direct change of variables in the definition of the
right-sided derivative yields
\[
\mathfrak D_{-,1}^{\beta}(U_R\Phi)
=
R^{\beta+\frac12}U_R(\mathfrak D_-^\beta \Phi).
\]
Combining this with \eqref{eq:window-scaling-T} gives
\[
U_R\mathfrak T_{t,R}^{\mathrm{hl}}\Phi
=
R^{-\beta}\mathbf c_t \mathfrak D_{-,1}^{\beta}(U_R\Phi)
=
\mathbf c_t U_R(\mathfrak D_-^\beta\Phi)
=
U_R(\mathcal T_{t,\beta}^{\mathrm{mod}}\Phi),
\]
which proves \eqref{eq:micro-window-model-smooth}.

Items~(ii)--(iv) follow from \eqref{eq:window-scaling-T}, Theorem~\ref{thm:frozen-density}, and the
fact that \(U_R\) is a topological isomorphism between \(H^\beta((0,R))\) and \(H^\beta((0,1))\).
\end{proof}

\begin{corollary}[Fixed-window model functional vanishes on the full graph domain]
\label{cor:micro-window-model-functional}
Assume the hypotheses of Theorem~\ref{thm:micro-window-model}. Define
\[
\Lambda_{t,R}^{\mathrm{hl}}(\Phi)
:=
L_{XY}(t,t)\langle \psi_{H_{XY}},\Phi\rangle_{L^2(\R_+)}
-L_Y(t,t)
\left\langle
\psi_{H_Y},
\mathfrak T_{t,R}^{\mathrm{hl}}\Phi
\right\rangle_{L^2(\R_+)},
\qquad
\Phi\in\mathcal D(\mathfrak T_{t,R}^{\mathrm{hl}}).
\]
Then \(\Lambda_{t,R}^{\mathrm{hl}}\) is continuous on \(\mathfrak G_{t,R}^{\mathrm{hl}}\), and
\[
\Lambda_{t,R}^{\mathrm{hl}}(\Phi)=0
\qquad\text{for every }\Phi\in \mathcal D(\mathfrak T_{t,R}^{\mathrm{hl}}).
\]
\end{corollary}

\begin{proof}
Continuity follows from Cauchy--Schwarz, \(\psi_{H_i}\in L^2(\R_+)\), and
\eqref{eq:micro-window-model-norm}. By Theorem~\ref{thm:micro-window-model}(iv), choose
\(\Phi_n\in C_c^\infty((0,R))\) such that
\[
\|\Phi_n-\Phi\|_{\mathfrak G_{t,R}^{\mathrm{hl}}}\to 0.
\]
By Theorem~\ref{thm:micro-window-model}(i), one has
\[
\mathfrak T_{t,R}^{\mathrm{hl}}\Phi_n=\mathcal T_{t,\beta}^{\mathrm{mod}}\Phi_n,
\]
so Corollary~\ref{cor:model-functional} gives \(\Lambda_{t,R}^{\mathrm{hl}}(\Phi_n)=0\) for every
\(n\). Passing to the limit by continuity proves the claim.
\end{proof}

\begin{definition}[Exact micro graph space on a fixed window]
\label{def:micro-window-exact}
Assume Assumptions~\ref{ass:model} and~\ref{ass:obs-source}, let \(\eta>0\), fix
\(t\in(0,T)\), \(\alpha\in(0,1)\), \(R>0\), and \(h>0\) such that
\[
hR<t,
\qquad
R<\varepsilon_{h,\alpha}^{-1}.
\]
Define the exact micro graph space
\[
\mathfrak G_{t,h,\alpha,R}^{\mathrm{mic}}
:=
\left\{
\Phi\in \mathcal D(\mathfrak T_{t,h,\alpha}^{\mathrm{mic}}):
\operatorname{supp}\Phi\subset [0,R]
\right\}.
\]
For \(\Phi\in\mathfrak G_{t,h,\alpha,R}^{\mathrm{mic}}\), define
\begin{equation}\label{eq:micro-window-graph-norm}
\|\Phi\|_{\mathfrak G_{t,h,\alpha,R}^{\mathrm{mic}}}^2
:=
\|\Phi\|_{E_R}^2
+\frac{\eta^2}{\nu_Y^2}
\|\mathfrak T_{t,h,\alpha}^{\mathrm{mic}}\Phi\|_{E_R}^2.
\end{equation}
For \(i\in\{Y,XY\}\), define the fixed-window exact micro Volterra operator
\[
\mathfrak V_{i,t,h,\alpha,R}^{\mathrm{mic}}:E_R\to E_R
\]
by
\[
\bigl(
\mathfrak V_{i,t,h,\alpha,R}^{\mathrm{mic}}\Phi
\bigr)(r)
:=
\int_r^R
(q-r)^{H_i-\frac12}
L_i(t-hr,t-hq)\Phi(q)\,dq,
\qquad 0\le r\le R.
\]
\end{definition}

\begin{proposition}[Exact micro transfer closes on every fixed window]
\label{prop:micro-window-exact}
Assume the hypotheses of Definition~\ref{def:micro-window-exact}. Then for every
\(\Phi\in\mathfrak G_{t,h,\alpha,R}^{\mathrm{mic}}\):
\begin{enumerate}[label=(\roman*)]
\item
\begin{equation}\label{eq:micro-window-support}
\operatorname{supp}\bigl(\mathfrak T_{t,h,\alpha}^{\mathrm{mic}}\Phi\bigr)\subset [0,R];
\end{equation}
\item
\begin{equation}\label{eq:micro-window-identity}
\mathfrak V_{Y,t,h,\alpha,R}^{\mathrm{mic}}
\bigl(
\mathfrak T_{t,h,\alpha}^{\mathrm{mic}}\Phi
\bigr)
=
\mathfrak V_{XY,t,h,\alpha,R}^{\mathrm{mic}}\Phi
\qquad\text{in }E_R.
\end{equation}
\end{enumerate}
\end{proposition}

\begin{proof}
Set \(\varepsilon:=\varepsilon_{h,\alpha}\), and define
\[
w_h:=\mathcal L_{t,h}\Phi.
\]
Because \(\operatorname{supp}\Phi\subset [0,R]\) and \(R<\varepsilon^{-1}\), one has
\[
\mathfrak p_{\varepsilon}\Phi=\Phi
\]
and therefore, by \eqref{eq:two-scale-L},
\[
w_h
=
\mathcal L_{t,h}\Phi
=
\mathcal L_{t,h}^{(\alpha)}\mathfrak l_{\varepsilon}\Phi.
\]
In particular, \(w_h\in\mathfrak G_t\), and
\[
\operatorname{supp}(w_h)\subset [t-hR,t].
\]
Corollary~\ref{cor:general-support-preserve} with \(a=t-hR\) therefore gives
\[
\operatorname{supp}(\mathcal S_t w_h)\subset [t-hR,t].
\]
By Proposition~\ref{prop:micro-local}(ii),
\[
\mathfrak T_{t,h,\alpha}^{\mathrm{mic}}\Phi
=
h^\beta \mathcal D_{t,h}\mathcal S_t w_h.
\]
Since \(\mathcal S_t w_h\) is supported in \([t-hR,t]\), its \(h\)-dilation is supported in
\([0,R]\), proving \eqref{eq:micro-window-support}.

Corollary~\ref{cor:micro-volterra-transfer} gives
\[
\mathfrak V_{Y,t,h,\alpha}^{\mathrm{mic}}
\bigl(
\mathfrak T_{t,h,\alpha}^{\mathrm{mic}}\Phi
\bigr)
=
\mathfrak V_{XY,t,h,\alpha}^{\mathrm{mic}}\Phi
\qquad\text{in }L^2(\R_+).
\]
Because both \(\Phi\) and \(\mathfrak T_{t,h,\alpha}^{\mathrm{mic}}\Phi\) are supported in
\([0,R]\), the operators on both sides coincide with their fixed-window restrictions on \(E_R\),
which proves \eqref{eq:micro-window-identity}.
\end{proof}

\begin{proposition}[Operator-norm freeze of the exact micro Volterra pair on fixed windows]
\label{prop:micro-window-freeze}
Assume Assumption~\ref{ass:model}, fix \(t\in(0,T)\), \(\alpha\in(0,1)\), and \(R>0\), and let
\(i\in\{Y,XY\}\). Then there exists \(C_{t,R}<\infty\) such that for all sufficiently small \(h>0\)
with \(hR<t\),
\begin{equation}\label{eq:micro-window-freeze}
\left\|
\mathfrak V_{i,t,h,\alpha,R}^{\mathrm{mic}}
-L_i(t,t)\mathfrak J_{H_i,R}
\right\|_{\mathcal L(E_R)}
\le
C_{t,R}h.
\end{equation}
\end{proposition}

\begin{proof}
The difference kernel is
\[
K_{i,t,h,R}^{\mathrm{err}}(r,q)
:=
(q-r)^{H_i-\frac12}
\bigl(
L_i(t-hr,t-hq)-L_i(t,t)
\bigr)\one_{\{0\le r\le q\le R\}}.
\]
Because \(L_i\) is \(C^1\) near \((t,t)\), one has
\[
|L_i(t-hr,t-hq)-L_i(t,t)|\le C_{t,R}h
\qquad\text{for }0\le r\le q\le R
\]
when \(hR<t\). Hence
\[
|K_{i,t,h,R}^{\mathrm{err}}(r,q)|
\le
C_{t,R}h (q-r)^{H_i-\frac12}\one_{\{0\le r\le q\le R\}}.
\]
Since this majorant is Hilbert--Schmidt on \(\Delta_R\), the Hilbert--Schmidt norm of the
difference is \(O_{t,R}(h)\), which proves \eqref{eq:micro-window-freeze}.
\end{proof}

\begin{proposition}[Weak graph compactness of the exact micro family on fixed windows]
\label{prop:micro-window-compact}
Assume Assumptions~\ref{ass:model} and~\ref{ass:obs-source}, let \(\eta>0\), fix
\(t\in(0,T)\), \(\alpha\in(0,1)\), \(R>0\), and suppose
\[
0<H_{XY}<H_Y<\tfrac12.
\]
Let \(h_n\downarrow0\) with \(h_nR<t\) and \(R<\varepsilon_{h_n,\alpha}^{-1}\) for every \(n\), and
let
\[
\Phi_n\in \mathfrak G_{t,h_n,\alpha,R}^{\mathrm{mic}}
\]
satisfy
\[
\sup_{n\ge1}\|\Phi_n\|_{\mathfrak G_{t,h_n,\alpha,R}^{\mathrm{mic}}}<\infty.
\]
Then there exist a subsequence, still denoted \((\Phi_n)\), and some
\[
\Phi\in \mathcal D(\mathfrak T_{t,R}^{\mathrm{hl}})
\]
such that
\begin{align*}
\Phi_n &\rightharpoonup \Phi
\qquad\text{weakly in }E_R, \\
\mathfrak T_{t,h_n,\alpha}^{\mathrm{mic}}\Phi_n
&\rightharpoonup
\mathfrak T_{t,R}^{\mathrm{hl}}\Phi
\qquad\text{weakly in }E_R.
\end{align*}
Moreover,
\[
\|\Phi\|_{\mathfrak G_{t,R}^{\mathrm{hl}}}
\le
\liminf_{n\to\infty}
\|\Phi_n\|_{\mathfrak G_{t,h_n,\alpha,R}^{\mathrm{mic}}}.
\]
\end{proposition}

\begin{proof}
By Definition~\ref{def:micro-window-exact} and Proposition~\ref{prop:micro-window-exact}(i), both
\((\Phi_n)\) and \((\mathfrak T_{t,h_n,\alpha}^{\mathrm{mic}}\Phi_n)\) are bounded in \(E_R\).
Passing to a subsequence, assume
\[
\Phi_n\rightharpoonup \Phi,
\qquad
\mathfrak T_{t,h_n,\alpha}^{\mathrm{mic}}\Phi_n\rightharpoonup \Psi
\qquad\text{weakly in }E_R
\]
for some \(\Phi,\Psi\in E_R\).

By Proposition~\ref{prop:micro-window-exact}(ii),
\[
\mathfrak V_{Y,t,h_n,\alpha,R}^{\mathrm{mic}}
\bigl(
\mathfrak T_{t,h_n,\alpha}^{\mathrm{mic}}\Phi_n
\bigr)
=
\mathfrak V_{XY,t,h_n,\alpha,R}^{\mathrm{mic}}\Phi_n.
\]
Using Proposition~\ref{prop:micro-window-freeze} exactly as in the proof of
Proposition~\ref{prop:local-graph-compact}, one obtains
\[
\mathfrak V_{XY,t,h_n,\alpha,R}^{\mathrm{mic}}\Phi_n
\rightharpoonup
L_{XY}(t,t)\mathfrak J_{H_{XY},R}\Phi
\qquad\text{weakly in }E_R
\]
and
\[
\mathfrak V_{Y,t,h_n,\alpha,R}^{\mathrm{mic}}
\bigl(
\mathfrak T_{t,h_n,\alpha}^{\mathrm{mic}}\Phi_n
\bigr)
\rightharpoonup
L_Y(t,t)\mathfrak J_{H_Y,R}\Psi
\qquad\text{weakly in }E_R.
\]
Passing to the weak limit in the exact identity gives
\[
L_Y(t,t)\mathfrak J_{H_Y,R}\Psi
=
L_{XY}(t,t)\mathfrak J_{H_{XY},R}\Phi.
\]
By Definition~\ref{def:micro-window-model}, this means \(\Phi\in\mathcal D(\mathfrak T_{t,R}^{\mathrm{hl}})\),
and by injectivity of \(\mathfrak J_{H_Y,R}\),
\[
\Psi=\mathfrak T_{t,R}^{\mathrm{hl}}\Phi.
\]
Weak lower semicontinuity of the norm in \(E_R\) then yields the graph-norm liminf exactly as in
Proposition~\ref{prop:local-graph-compact}.
\end{proof}

\begin{corollary}[Fixed-window micro Mosco convergence]
\label{cor:micro-window-mosco}
Assume the hypotheses of Proposition~\ref{prop:micro-window-compact}. Let \(h_n\downarrow0\) with
\(h_nR<t\) and \(R<\varepsilon_{h_n,\alpha}^{-1}\) for every \(n\). Define
\[
\Gamma_{t,h_n,\alpha,R}^{\mathrm{mic}}
:=
\left\{
\bigl(
\Phi,\mathfrak T_{t,h_n,\alpha}^{\mathrm{mic}}\Phi
\bigr):
\Phi\in\mathfrak G_{t,h_n,\alpha,R}^{\mathrm{mic}}
\right\}
\subset E_R\times E_R,
\]
and define the fixed-window model graph
\[
\Gamma_{t,R}^{\mathrm{hl}}
:=
\left\{
\bigl(
\Phi,\mathfrak T_{t,R}^{\mathrm{hl}}\Phi
\bigr):
\Phi\in \mathcal D(\mathfrak T_{t,R}^{\mathrm{hl}})
\right\}.
\]
Then \(\Gamma_{t,h_n,\alpha,R}^{\mathrm{mic}}\to \Gamma_{t,R}^{\mathrm{hl}}\) in the sequential Mosco
sense.
\end{corollary}

\begin{proof}
The weak-compactness direction is Proposition~\ref{prop:micro-window-compact}. For the recovery
direction, Theorem~\ref{thm:micro-window-model}(iv) gives density of \(C_c^\infty((0,R))\) in the
fixed-window model graph norm, while Theorem~\ref{thm:micro-smooth-core} gives strong convergence
of \(\mathfrak T_{t,h,\alpha}^{\mathrm{mic}}\) to \(\mathcal T_{t,\beta}^{\mathrm{mod}}\) on that
smooth core, which coincides with \(\mathfrak T_{t,R}^{\mathrm{hl}}\) by
Theorem~\ref{thm:micro-window-model}(i). The same diagonal argument as in
Theorem~\ref{thm:local-recovery} and Corollary~\ref{cor:local-mosco} therefore yields the Mosco
recovery property on \(E_R\times E_R\).
\end{proof}

\section{Buffered micro tail forcing}\label{sec:micro-tail}

Fixed-window theory isolates the remaining nonlocal obstruction.  Truncating a global micro pair at
an outer radius \(R\) breaks the exact Volterra identity on an inner window \([0,R_0]\) only through
the tail beyond \(R\).  The defect is written exactly.  After the same right-sided fractional
conjugation used in Section~\ref{sec:micro-smooth}, the resulting forcing is quantitatively small.

\begin{definition}[Buffered micro tail operators]
\label{def:micro-tail}
Assume Assumption~\ref{ass:model}, fix \(t\in(0,T)\), \(\alpha\in(0,1)\), and suppose
\[
0<H_{XY}<H_Y<\tfrac12.
\]
Set
\[
\lambda:=H_Y+\frac12,
\qquad
\mu:=H_{XY}+\frac12,
\qquad
\beta:=H_Y-H_{XY},
\]
and let \(0<R_0<R\). For \(h>0\) sufficiently small that \(hR<t\) and
\(R<\varepsilon_{h,\alpha}^{-1}\), define
\[
F_{h,R}:=L^2([R,\varepsilon_{h,\alpha}^{-1}]),
\]
identified with its zero extension to \(L^2(\R_+)\). For \(\Theta\in F_{h,R}\), define the
target-side and source-side buffered micro tail operators on \(L^2(\R_+)\) by
\begin{align}
\bigl(
\mathfrak W_{Y,t,h,\alpha,R_0,R}^{\mathrm{tail}}\Theta
\bigr)(r)
&:=
\int_{\max\{R,r\}}^{\varepsilon_{h,\alpha}^{-1}}
(q-r)^{\lambda-1}
L_Y(t-hr,t-hq)\Theta(q)\,dq,
\label{eq:micro-tail-Y}\\
\bigl(
\mathfrak W_{XY,t,h,\alpha,R_0,R}^{\mathrm{tail}}\Theta
\bigr)(r)
&:=
\int_{\max\{R,r\}}^{\varepsilon_{h,\alpha}^{-1}}
(q-r)^{\mu-1}
L_{XY}(t-hr,t-hq)\Theta(q)\,dq
\label{eq:micro-tail-XY}
\end{align}
for \(r\ge0\). On the inner window \([0,R_0]\), these reduce to \eqref{eq:micro-tail-Y} and
\eqref{eq:micro-tail-XY} with lower limit \(R\), since \(R_0<R\).
\end{definition}

\begin{proposition}[Conjugated target-side buffered tail forcing]
\label{prop:micro-tail-Y}
Assume the hypotheses of Definition~\ref{def:micro-tail}. Then for every sufficiently small \(h>0\)
there exists a bounded operator
\[
\mathfrak Q_{Y,t,h,\alpha,R_0,R}^{\mathrm{tail}}
\in \mathcal L(F_{h,R},E_{R_0})
\]
such that for every \(\Theta\in F_{h,R}\),
\begin{equation}\label{eq:micro-tail-Y-conj}
\frac{1}{\Gamma(\lambda)}
\mathfrak D_-^\lambda
\bigl(
\mathfrak W_{Y,t,h,\alpha,R_0,R}^{\mathrm{tail}}\Theta
\bigr)
=
\mathfrak Q_{Y,t,h,\alpha,R_0,R}^{\mathrm{tail}}\Theta
\qquad\text{in }E_{R_0}.
\end{equation}
Moreover, there exists \(C_{t,R_0}<\infty\) such that
\begin{equation}\label{eq:micro-tail-Y-bound}
\left\|
\mathfrak Q_{Y,t,h,\alpha,R_0,R}^{\mathrm{tail}}
\right\|_{\mathcal L(F_{h,R},E_{R_0})}
\le
C_{t,R_0}h\sqrt{\varepsilon_{h,\alpha}^{-1}-R}
\le
C_{t,R_0}h^{\frac{1+\alpha}{2}}
\end{equation}
for all sufficiently small \(h>0\).
\end{proposition}

\begin{proof}
For \(0\le r\le R_0\) and \(q\in[R,\varepsilon_{h,\alpha}^{-1}]\), define
\[
a_{Y,t,h}(r,q):=L_Y(t-hr,t-hq).
\]
Since \(q-r\ge R-R_0>0\), the function
\[
k_{Y,t,h,q}(r):=(q-r)^{\lambda-1}a_{Y,t,h}(r,q)
\]
is \(C^1\) in \(r\) on \([0,R_0]\). The same computation as in the proof of
Proposition~\ref{prop:BY-micro}, now with fixed upper endpoint \(q\), gives
\[
\frac{1}{\Gamma(\lambda)}
\mathfrak D_-^\lambda k_{Y,t,h,q}(r)
=
K_{Y,t,h}^{\mathrm{tail}}(r,q)
\qquad\text{for }0\le r\le R_0,
\]
where
\[
K_{Y,t,h}^{\mathrm{tail}}(r,q)
:=
-
\frac{1}{\Gamma(\lambda)\Gamma(1-\lambda)}
\int_0^1
s^{-\lambda}(1-s)^\lambda
\partial_1 a_{Y,t,h}(r+(q-r)s,q)\,ds.
\]
Hence Fubini's theorem yields
\[
\frac{1}{\Gamma(\lambda)}
\mathfrak D_-^\lambda
\bigl(
\mathfrak W_{Y,t,h,\alpha,R_0,R}^{\mathrm{tail}}\Theta
\bigr)(r)
=
\int_R^{\varepsilon_{h,\alpha}^{-1}}
K_{Y,t,h}^{\mathrm{tail}}(r,q)\Theta(q)\,dq
\]
for almost every \(r\in[0,R_0]\). This defines the bounded operator
\[
\bigl(
\mathfrak Q_{Y,t,h,\alpha,R_0,R}^{\mathrm{tail}}\Theta
\bigr)(r)
:=
\int_R^{\varepsilon_{h,\alpha}^{-1}}
K_{Y,t,h}^{\mathrm{tail}}(r,q)\Theta(q)\,dq,
\qquad 0\le r\le R_0,
\]
and proves \eqref{eq:micro-tail-Y-conj}.

Because \(L_Y\) is \(C^1\),
\[
\partial_1 a_{Y,t,h}(x,q)
=
-h\,\partial_1L_Y(t-hx,t-hq),
\]
so \(|\partial_1 a_{Y,t,h}(x,q)|\le C_t h\) uniformly on
\(\{0\le x\le q\le \varepsilon_{h,\alpha}^{-1}\}\). Therefore
\[
|K_{Y,t,h}^{\mathrm{tail}}(r,q)|\le C_t h
\qquad\text{for }0\le r\le R_0,\quad R\le q\le \varepsilon_{h,\alpha}^{-1}.
\]
The kernel is Hilbert--Schmidt, with
\[
\left\|
K_{Y,t,h}^{\mathrm{tail}}
\right\|_{L^2([0,R_0]\times [R,\varepsilon_{h,\alpha}^{-1}])}
\le
C_{t,R_0}h\sqrt{\varepsilon_{h,\alpha}^{-1}-R},
\]
which proves the first bound in \eqref{eq:micro-tail-Y-bound}. Since
\(\varepsilon_{h,\alpha}^{-1}=h^{\alpha-1}\), one also has
\[
h\sqrt{\varepsilon_{h,\alpha}^{-1}-R}
\le
h\sqrt{\varepsilon_{h,\alpha}^{-1}}
=
h^{\frac{1+\alpha}{2}},
\]
which proves the second bound.
\end{proof}

\begin{proposition}[Conjugated source-side buffered tail forcing]
\label{prop:micro-tail-XY}
Assume the hypotheses of Definition~\ref{def:micro-tail}. Then for every sufficiently small \(h>0\)
there exists a bounded operator
\[
\mathfrak Q_{XY,t,h,\alpha,R_0,R}^{\mathrm{tail}}
\in \mathcal L(F_{h,R},E_{R_0})
\]
such that for every \(\Theta\in F_{h,R}\),
\begin{equation}\label{eq:micro-tail-XY-conj}
\frac{1}{\Gamma(\lambda)}
\mathfrak D_-^\lambda
\bigl(
\mathfrak W_{XY,t,h,\alpha,R_0,R}^{\mathrm{tail}}\Theta
\bigr)
=
\mathfrak Q_{XY,t,h,\alpha,R_0,R}^{\mathrm{tail}}\Theta
\qquad\text{in }E_{R_0}.
\end{equation}
Moreover, there exist \(C_{t,R_0}<\infty\) and \(C_{t,R_0,\alpha}<\infty\) such that
\begin{equation}\label{eq:micro-tail-XY-bound}
\left\|
\mathfrak Q_{XY,t,h,\alpha,R_0,R}^{\mathrm{tail}}
\right\|_{\mathcal L(F_{h,R},E_{R_0})}
\le
C_{t,R_0}(R-R_0)^{-\beta}
+C_{t,R_0,\alpha}h^{\frac{1+\alpha}{2}}
\end{equation}
for all sufficiently small \(h>0\).
\end{proposition}

\begin{proof}
For \(0\le r\le R_0\) and \(q\in[R,\varepsilon_{h,\alpha}^{-1}]\), define
\[
a_{XY,t,h}(r,q):=L_{XY}(t-hr,t-hq).
\]
Exactly as in the proof of Proposition~\ref{prop:q-micro}, but with fixed upper endpoint \(q\),
one obtains
\[
\frac{1}{\Gamma(\lambda)}
\mathfrak D_-^\lambda
\Bigl(
(q-\cdot)^{\mu-1}a_{XY,t,h}(\cdot,q)
\Bigr)(r)
=
K_{XY,t,h}^{\mathrm{tail}}(r,q),
\qquad 0\le r\le R_0,
\]
where
\[
K_{XY,t,h}^{\mathrm{tail}}(r,q)
:=
\beta (q-r)^{-\beta-1}g_{t,h}^{\mathrm{tail}}(r,q)
-
(q-r)^{-\beta}\partial_1 g_{t,h}^{\mathrm{tail}}(r,q)
\]
and
\[
g_{t,h}^{\mathrm{tail}}(r,q)
:=
\frac{1}{\Gamma(\lambda)\Gamma(1-\lambda)}
\int_0^1
\theta^{-\lambda}(1-\theta)^{\mu-1}
a_{XY,t,h}(r+(q-r)\theta,q)\,d\theta.
\]
Therefore Fubini's theorem yields
\[
\frac{1}{\Gamma(\lambda)}
\mathfrak D_-^\lambda
\bigl(
\mathfrak W_{XY,t,h,\alpha,R_0,R}^{\mathrm{tail}}\Theta
\bigr)(r)
=
\int_R^{\varepsilon_{h,\alpha}^{-1}}
K_{XY,t,h}^{\mathrm{tail}}(r,q)\Theta(q)\,dq
\]
for almost every \(r\in[0,R_0]\). This defines the bounded operator
\[
\bigl(
\mathfrak Q_{XY,t,h,\alpha,R_0,R}^{\mathrm{tail}}\Theta
\bigr)(r)
:=
\int_R^{\varepsilon_{h,\alpha}^{-1}}
K_{XY,t,h}^{\mathrm{tail}}(r,q)\Theta(q)\,dq,
\qquad 0\le r\le R_0,
\]
and proves \eqref{eq:micro-tail-XY-conj}.

Since \(L_{XY}\) is \(C^1\), one has
\[
\bigl|g_{t,h}^{\mathrm{tail}}(r,q)\bigr|\le C_t,
\qquad
\bigl|\partial_1 g_{t,h}^{\mathrm{tail}}(r,q)\bigr|\le C_t h
\]
uniformly for \(0\le r\le R_0\) and \(R\le q\le \varepsilon_{h,\alpha}^{-1}\). Accordingly, write
\[
K_{XY,t,h}^{\mathrm{tail}}=K_{XY,t,h}^{(1)}+K_{XY,t,h}^{(2)},
\]
with
\[
\bigl|K_{XY,t,h}^{(1)}(r,q)\bigr|
\le
C_t (q-r)^{-\beta-1},
\qquad
\bigl|K_{XY,t,h}^{(2)}(r,q)\bigr|
\le
C_t h (q-r)^{-\beta}.
\]

For \(K_{XY,t,h}^{(1)}\), Schur's test gives
\[
\sup_{0\le r\le R_0}\int_R^{\varepsilon_{h,\alpha}^{-1}} (q-r)^{-\beta-1}\,dq
\le
\frac{1}{\beta}(R-R_0)^{-\beta}
\]
and
\[
\sup_{R\le q\le \varepsilon_{h,\alpha}^{-1}}
\int_0^{R_0} (q-r)^{-\beta-1}\,dr
\le
\frac{1}{\beta}(R-R_0)^{-\beta}.
\]
Hence
\begin{equation}\label{eq:micro-tail-XY-main}
\left\|
K_{XY,t,h}^{(1)}
\right\|_{\mathcal L(F_{h,R},E_{R_0})}
\le
C_{t,R_0}(R-R_0)^{-\beta}.
\end{equation}

For \(K_{XY,t,h}^{(2)}\), use the Hilbert--Schmidt bound:
\begin{align*}
\left\|
K_{XY,t,h}^{(2)}
\right\|_{L^2([0,R_0]\times [R,\varepsilon_{h,\alpha}^{-1}])}^2
&\le
C_{t,R_0} h^2
\int_0^{R_0}\int_R^{\varepsilon_{h,\alpha}^{-1}}
(q-r)^{-2\beta}\,dq\,dr \\
&\le
C_{t,R_0,\alpha} h^2\varepsilon_{h,\alpha}^{-(1-2\beta)},
\end{align*}
because \(2\beta<1\). Since \(\varepsilon_{h,\alpha}=h^{1-\alpha}\),
\[
h\,\varepsilon_{h,\alpha}^{-\frac{1-2\beta}{2}}
=
h^{1-\frac{(1-\alpha)(1-2\beta)}{2}}
\le
h^{\frac{1+\alpha}{2}}
\qquad\text{for }0<h\le1.
\]
Therefore
\begin{equation}\label{eq:micro-tail-XY-error}
\left\|
K_{XY,t,h}^{(2)}
\right\|_{\mathcal L(F_{h,R},E_{R_0})}
\le
C_{t,R_0,\alpha}h^{\frac{1+\alpha}{2}}.
\end{equation}
Combining \eqref{eq:micro-tail-XY-main} and \eqref{eq:micro-tail-XY-error} proves
\eqref{eq:micro-tail-XY-bound}.
\end{proof}

\begin{theorem}[Buffered truncation defect for a global exact micro pair]
\label{thm:micro-tail-defect}
Assume Assumptions~\ref{ass:model} and~\ref{ass:obs-source}, let \(\eta>0\), fix \(t\in(0,T)\),
\(\alpha\in(0,1)\), and suppose
\[
0<H_{XY}<H_Y<\tfrac12.
\]
Let \(0<R_0<R\), choose \(h>0\) sufficiently small that \(hR<t\) and
\(R<\varepsilon_{h,\alpha}^{-1}\), and let
\[
\Phi\in\mathcal D(\mathfrak T_{t,h,\alpha}^{\mathrm{mic}}),
\qquad
\Psi:=\mathfrak T_{t,h,\alpha}^{\mathrm{mic}}\Phi.
\]
Define
\[
\Phi^{\le R}:=\Phi\one_{[0,R]},
\qquad
\Phi^{>R}:=\Phi\one_{(R,\varepsilon_{h,\alpha}^{-1}]},
\qquad
\Psi^{\le R}:=\Psi\one_{[0,R]},
\qquad
\Psi^{>R}:=\Psi\one_{(R,\varepsilon_{h,\alpha}^{-1}]}.
\]
Then:
\begin{enumerate}[label=(\roman*)]
\item the windowed truncation defect satisfies
\begin{equation}\label{eq:micro-tail-defect}
\frac{1}{\Gamma(\lambda)}
\mathfrak D_-^\lambda
\left(
\mathfrak V_{Y,t,h,\alpha,R}^{\mathrm{mic}}\Psi^{\le R}
-
\mathfrak V_{XY,t,h,\alpha,R}^{\mathrm{mic}}\Phi^{\le R}
\right)
=
\mathfrak Q_{XY,t,h,\alpha,R_0,R}^{\mathrm{tail}}\Phi^{>R}
-
\mathfrak Q_{Y,t,h,\alpha,R_0,R}^{\mathrm{tail}}\Psi^{>R}
\end{equation}
in \(E_{R_0}\);
\item there exists \(C_{t,R_0,\alpha}<\infty\) such that
\begin{align}
\left\|
\frac{1}{\Gamma(\lambda)}
\mathfrak D_-^\lambda
\left(
\mathfrak V_{Y,t,h,\alpha,R}^{\mathrm{mic}}\Psi^{\le R}
-
\mathfrak V_{XY,t,h,\alpha,R}^{\mathrm{mic}}\Phi^{\le R}
\right)
\right\|_{E_{R_0}}
\notag\\
\le
C_{t,R_0,\alpha}
\left(
(R-R_0)^{-\beta}\|\Phi^{>R}\|_{L^2(\R_+)}
+h^{\frac{1+\alpha}{2}}
\bigl(
\|\Phi^{>R}\|_{L^2(\R_+)}+\|\Psi^{>R}\|_{L^2(\R_+)}
\bigr)
\right).
\label{eq:micro-tail-defect-bound}
\end{align}
\end{enumerate}
\end{theorem}

\begin{proof}
Because \(\Psi=\mathfrak T_{t,h,\alpha}^{\mathrm{mic}}\Phi\), Corollary~\ref{cor:micro-volterra-transfer}
gives
\[
\mathfrak V_{Y,t,h,\alpha}^{\mathrm{mic}}\Psi
=
\mathfrak V_{XY,t,h,\alpha}^{\mathrm{mic}}\Phi
\qquad\text{in }L^2(\R_+).
\]
Set \(\varepsilon:=\varepsilon_{h,\alpha}\). By Proposition~\ref{prop:micro-local}(ii),
\[
\Psi
=
h^\beta
\mathcal D_{t,h}\mathcal S_t\mathcal L_{t,h}\mathfrak p_{\varepsilon}\Phi.
\]
Since \(\mathfrak p_{\varepsilon}\Phi\) is supported in \([0,\varepsilon^{-1}]\), the lift
\(\mathcal L_{t,h}\mathfrak p_{\varepsilon}\Phi\) is supported in
\([t-h\varepsilon^{-1},t]=[t-h^\alpha,t]\). Corollary~\ref{cor:general-support-preserve} therefore
implies
\[
\operatorname{supp}\bigl(\mathcal S_t\mathcal L_{t,h}\mathfrak p_{\varepsilon}\Phi\bigr)
\subset [t-h^\alpha,t],
\]
and hence
\[
\operatorname{supp}(\Psi)\subset [0,\varepsilon^{-1}].
\]
Splitting the exact Volterra identity at \(q=R\) therefore gives, in \(E_{R_0}\),
\[
\mathfrak V_{Y,t,h,\alpha,R}^{\mathrm{mic}}\Psi^{\le R}
-
\mathfrak V_{XY,t,h,\alpha,R}^{\mathrm{mic}}\Phi^{\le R}
=
\mathfrak W_{XY,t,h,\alpha,R_0,R}^{\mathrm{tail}}\Phi^{>R}
-
\mathfrak W_{Y,t,h,\alpha,R_0,R}^{\mathrm{tail}}\Psi^{>R}.
\]
Applying Propositions~\ref{prop:micro-tail-Y} and~\ref{prop:micro-tail-XY} proves
\eqref{eq:micro-tail-defect}. The norm estimate \eqref{eq:micro-tail-defect-bound} follows by applying
\eqref{eq:micro-tail-Y-bound} and \eqref{eq:micro-tail-XY-bound} to the two tail terms and then
using the triangle inequality.
\end{proof}

\section{Intrinsic exact fixed-window solver}\label{sec:micro-solver}

The fixed-window graph spaces from Section~\ref{sec:micro-window} were defined by restricting the
global micro transfer to compactly supported inputs.  The exact finite-window problem can also be
formulated intrinsically.  The windowed Volterra relation determines a unique closed solver, and
that intrinsic solver coincides with the supported-global graph.

\begin{definition}[Intrinsic exact fixed-window graph]
\label{def:micro-window-solver}
Assume Assumptions~\ref{ass:model} and~\ref{ass:obs-source}, let \(\eta>0\), fix
\(t\in(0,T)\), \(\alpha\in(0,1)\), \(R>0\), and choose \(h>0\) such that
\[
hR<t,
\qquad
R<\varepsilon_{h,\alpha}^{-1}.
\]
Define the intrinsic exact fixed-window graph by
\[
\Gamma_{t,h,\alpha,R}^{\mathrm{ex}}
:=
\left\{
(\Phi,\Psi)\in E_R\times E_R:
\mathfrak V_{Y,t,h,\alpha,R}^{\mathrm{mic}}\Psi
=
\mathfrak V_{XY,t,h,\alpha,R}^{\mathrm{mic}}\Phi
\right\}.
\]
\end{definition}

\begin{proposition}[Closed exact fixed-window solver]
\label{prop:micro-window-solver}
Assume the hypotheses of Definition~\ref{def:micro-window-solver}, and suppose in addition that
\(h>0\) is small enough for Corollary~\ref{cor:BY-micro-invertible} to apply on the window
\([0,R]\). Then:
\begin{enumerate}[label=(\roman*)]
\item \(\Gamma_{t,h,\alpha,R}^{\mathrm{ex}}\) is a closed linear subspace of \(E_R\times E_R\);
\item the operator \(\mathfrak V_{Y,t,h,\alpha,R}^{\mathrm{mic}}:E_R\to E_R\) is injective;
\item there exists a unique closed linear operator
\[
\mathfrak T_{t,h,\alpha,R}^{\mathrm{ex}}:
\mathcal D(\mathfrak T_{t,h,\alpha,R}^{\mathrm{ex}})\subset E_R\to E_R
\]
whose graph is \(\Gamma_{t,h,\alpha,R}^{\mathrm{ex}}\).
\end{enumerate}
\end{proposition}

\begin{proof}
The map
\[
(\Phi,\Psi)\longmapsto
\mathfrak V_{Y,t,h,\alpha,R}^{\mathrm{mic}}\Psi
-
\mathfrak V_{XY,t,h,\alpha,R}^{\mathrm{mic}}\Phi
\]
is bounded from \(E_R\times E_R\) to \(E_R\), because both fixed-window Volterra operators are
bounded. Therefore \(\Gamma_{t,h,\alpha,R}^{\mathrm{ex}}\), being its kernel, is a closed linear
subspace. This proves (i).

To prove injectivity, let \(\Psi\in E_R\) satisfy
\[
\mathfrak V_{Y,t,h,\alpha,R}^{\mathrm{mic}}\Psi=0.
\]
Since \(\Psi\) is supported in \([0,R]\), Proposition~\ref{prop:BY-micro} gives
\[
\mathfrak B_{Y,t,h,\alpha,R}^{\mathrm{mic}}\Psi
=
\frac{1}{\Gamma(\lambda)}
\mathfrak D_-^\lambda
\bigl(
\mathfrak V_{Y,t,h,\alpha}^{\mathrm{mic}}\Psi
\bigr)
=
0
\qquad\text{in }E_R,
\]
where \(\lambda=H_Y+\frac12\). Corollary~\ref{cor:BY-micro-invertible} implies
\(\Psi=0\). This proves (ii).

Part (iii) follows from (ii): for each \(\Phi\in E_R\) there is at most one
\(\Psi\in E_R\) such that \((\Phi,\Psi)\in\Gamma_{t,h,\alpha,R}^{\mathrm{ex}}\). Hence the closed
subspace \(\Gamma_{t,h,\alpha,R}^{\mathrm{ex}}\) is the graph of a unique closed linear operator.
\end{proof}

\begin{proposition}[Intrinsic and supported-global exact graphs coincide]
\label{prop:micro-window-identification}
Assume the hypotheses of Proposition~\ref{prop:micro-window-solver}. Then
\begin{equation}\label{eq:micro-window-identification}
\mathcal D(\mathfrak T_{t,h,\alpha,R}^{\mathrm{ex}})
=
\mathfrak G_{t,h,\alpha,R}^{\mathrm{mic}},
\qquad
\mathfrak T_{t,h,\alpha,R}^{\mathrm{ex}}\Phi
=
\mathfrak T_{t,h,\alpha}^{\mathrm{mic}}\Phi
\quad
\text{for every }\Phi\in \mathfrak G_{t,h,\alpha,R}^{\mathrm{mic}}.
\end{equation}
In particular, the graph of \(\mathfrak T_{t,h,\alpha,R}^{\mathrm{ex}}\) is exactly
\[
\Gamma_{t,h,\alpha,R}^{\mathrm{mic}}
:=
\left\{
\bigl(
\Phi,\mathfrak T_{t,h,\alpha}^{\mathrm{mic}}\Phi
\bigr):
\Phi\in\mathfrak G_{t,h,\alpha,R}^{\mathrm{mic}}
\right\}
\subset E_R\times E_R.
\]
\end{proposition}

\begin{proof}
Let \(\Phi\in\mathfrak G_{t,h,\alpha,R}^{\mathrm{mic}}\), and set
\[
\Psi:=\mathfrak T_{t,h,\alpha}^{\mathrm{mic}}\Phi.
\]
By Proposition~\ref{prop:micro-window-exact}(i), \(\Psi\in E_R\). By
Proposition~\ref{prop:micro-window-exact}(ii),
\[
\mathfrak V_{Y,t,h,\alpha,R}^{\mathrm{mic}}\Psi
=
\mathfrak V_{XY,t,h,\alpha,R}^{\mathrm{mic}}\Phi
\qquad\text{in }E_R.
\]
Hence \((\Phi,\Psi)\in \Gamma_{t,h,\alpha,R}^{\mathrm{ex}}\), so
\[
\mathfrak G_{t,h,\alpha,R}^{\mathrm{mic}}
\subset
\mathcal D(\mathfrak T_{t,h,\alpha,R}^{\mathrm{ex}})
\]
and
\[
\mathfrak T_{t,h,\alpha,R}^{\mathrm{ex}}\Phi
=
\mathfrak T_{t,h,\alpha}^{\mathrm{mic}}\Phi
\qquad\text{for every }\Phi\in \mathfrak G_{t,h,\alpha,R}^{\mathrm{mic}}.
\]

Conversely, let
\[
\Phi\in\mathcal D(\mathfrak T_{t,h,\alpha,R}^{\mathrm{ex}})
\]
and set
\[
\Psi:=\mathfrak T_{t,h,\alpha,R}^{\mathrm{ex}}\Phi\in E_R.
\]
Then
\begin{equation}\label{eq:intrinsic-window-identity}
\mathfrak V_{Y,t,h,\alpha,R}^{\mathrm{mic}}\Psi
=
\mathfrak V_{XY,t,h,\alpha,R}^{\mathrm{mic}}\Phi
\qquad\text{in }E_R.
\end{equation}
Since \(\Phi\) and \(\Psi\) are supported in \([0,R]\subset[0,\varepsilon_{h,\alpha}^{-1}]\), one
has
\[
\mathfrak p_{\varepsilon_{h,\alpha}}\Phi=\Phi,
\qquad
\mathfrak p_{\varepsilon_{h,\alpha}}\Psi=\Psi.
\]
Set \(\varepsilon:=\varepsilon_{h,\alpha}\). By Proposition~\ref{prop:micro-volterra},
\[
\mathfrak d_{\varepsilon}
\mathfrak A_{Y,t,h,\alpha}
\mathfrak l_{\varepsilon}\Psi
=
\varepsilon^{H_Y}
\mathfrak V_{Y,t,h,\alpha}^{\mathrm{mic}}\Psi,
\qquad
\mathfrak d_{\varepsilon}
\mathfrak A_{XY,t,h,\alpha}
\mathfrak l_{\varepsilon}\Phi
=
\varepsilon^{H_{XY}}
\mathfrak V_{XY,t,h,\alpha}^{\mathrm{mic}}\Phi.
\]
Multiplying \eqref{eq:intrinsic-window-identity} by \(\varepsilon^{H_Y}\), using
\(\beta=H_Y-H_{XY}\), and then applying \(\mathfrak l_{\varepsilon}\) together with
\(\mathfrak l_{\varepsilon}\mathfrak d_{\varepsilon}=I_E\) from
Lemma~\ref{lem:micro-dilation}, yields
\begin{equation}\label{eq:intrinsic-window-local}
\mathfrak A_{Y,t,h,\alpha}
\bigl(
\varepsilon^{-\beta}\mathfrak l_{\varepsilon}\Psi
\bigr)
=
\mathfrak A_{XY,t,h,\alpha}\mathfrak l_{\varepsilon}\Phi.
\end{equation}

Now define
\[
w_h:=\mathcal L_{t,h}^{(\alpha)}\mathfrak l_{\varepsilon}\Phi
\in H_{t,h,\alpha}^{\mathrm{near}},
\qquad
u_h:=
h^{-\alpha\beta}
\mathcal L_{t,h}^{(\alpha)}
\bigl(
\varepsilon^{-\beta}\mathfrak l_{\varepsilon}\Psi
\bigr)
\in H_{t,h,\alpha}^{\mathrm{near}}.
\]
Using Proposition~\ref{prop:local-volterra} and \eqref{eq:intrinsic-window-local},
\begin{align*}
\mathcal D_{t,h}^{(\alpha)}\mathcal K_{Y,t}u_h
&=
h^{-\alpha\beta+\alpha H_Y}
\mathfrak A_{Y,t,h,\alpha}
\bigl(
\varepsilon^{-\beta}\mathfrak l_{\varepsilon}\Psi
\bigr) \\
&=
h^{\alpha H_{XY}}
\mathfrak A_{XY,t,h,\alpha}\mathfrak l_{\varepsilon}\Phi \\
&=
\mathcal D_{t,h}^{(\alpha)}\mathcal K_{XY,t}w_h.
\end{align*}
Both \(\mathcal K_{Y,t}u_h\) and \(\mathcal K_{XY,t}w_h\) are supported in
\([t-h^\alpha,t]\), so Lemma~\ref{lem:unit-dilation} implies
\[
\mathcal K_{Y,t}u_h=\mathcal K_{XY,t}w_h.
\]
Hence \(w_h\in\mathcal D(\mathcal S_t)\), \(\mathcal S_t w_h=u_h\), and
\[
\mathfrak l_{\varepsilon}\Phi
=
\mathcal D_{t,h}^{(\alpha)}w_h
\in
\mathfrak G_{t,h,\alpha}^{\mathrm{near}}.
\]
Therefore \(\Phi\in\mathfrak G_{t,h,\alpha,R}^{\mathrm{mic}}\). Moreover, by
Definitions~\ref{def:local-transfer} and~\ref{def:micro-local},
\begin{align*}
\mathfrak T_{t,h,\alpha}^{\mathrm{mic}}\Phi
&=
\varepsilon^\beta
\mathfrak d_{\varepsilon}
\mathfrak T_{t,h,\alpha}\mathfrak l_{\varepsilon}\Phi \\
&=
\varepsilon^\beta
\mathfrak d_{\varepsilon}
\bigl(
\varepsilon^{-\beta}\mathfrak l_{\varepsilon}\Psi
\bigr) \\
&=
\mathfrak d_{\varepsilon}\mathfrak l_{\varepsilon}\Psi
=
\mathfrak p_{\varepsilon}\Psi
=
\Psi,
\end{align*}
because \(\operatorname{supp}\Psi\subset [0,R]\subset [0,\varepsilon^{-1}]\). This proves the
reverse inclusion and completes \eqref{eq:micro-window-identification}.
\end{proof}

\begin{corollary}[Intrinsic exact fixed-window graph norm]
\label{cor:micro-window-exact-norm}
Under the hypotheses of Proposition~\ref{prop:micro-window-identification}, define
\[
\|\Phi\|_{\mathfrak G_{t,h,\alpha,R}^{\mathrm{ex}}}^2
:=
\|\Phi\|_{E_R}^2
+
\frac{\eta^2}{\nu_Y^2}
\|\mathfrak T_{t,h,\alpha,R}^{\mathrm{ex}}\Phi\|_{E_R}^2,
\qquad
\Phi\in\mathcal D(\mathfrak T_{t,h,\alpha,R}^{\mathrm{ex}}).
\]
Then \(\mathcal D(\mathfrak T_{t,h,\alpha,R}^{\mathrm{ex}})\), equipped with this norm, is a
Hilbert space, and the identification in \eqref{eq:micro-window-identification} is isometric:
\[
\|\Phi\|_{\mathfrak G_{t,h,\alpha,R}^{\mathrm{ex}}}
=
\|\Phi\|_{\mathfrak G_{t,h,\alpha,R}^{\mathrm{mic}}}
\qquad\text{for every }\Phi\in\mathfrak G_{t,h,\alpha,R}^{\mathrm{mic}}.
\]
\end{corollary}

\begin{proof}
Since \(\mathfrak T_{t,h,\alpha,R}^{\mathrm{ex}}\) is closed by
Proposition~\ref{prop:micro-window-solver}, its graph norm defines a Hilbert space. The isometric
identification is exactly Proposition~\ref{prop:micro-window-identification} together with
Definition~\ref{def:micro-window-exact}.
\end{proof}

\section{Exact fixed-window regularity}\label{sec:micro-regularity}

The intrinsic windowed solver requires a domain description.  This is the only point where the
argument uses an additional derivative of the slowly varying amplitudes.  With that extra
regularity, the source-side conjugated operator splits into a diagonal multiplier times the
fractional model derivative plus a bounded lower-order remainder.

\begin{definition}[Extra amplitude regularity on a fixed time horizon]
\label{def:model-c2}
We say that Assumption~\ref{ass:model} holds with \emph{second-order amplitude regularity} if each
\[
L_i,\qquad i\in\{Y,XY\},
\]
extends to a \(C^2\) function on a neighbourhood of \(\Delta_T\).
\end{definition}

\begin{lemma}[Distributional model derivative on a fixed window]
\label{lem:window-distribution}
Assume
\[
0<H_{XY}<H_Y<\tfrac12,
\qquad
\beta:=H_Y-H_{XY},
\]
fix \(t\in(0,T)\), and let \(R>0\). For \(\Phi\in E_R=L^2([0,R])\), define
\[
(\mathfrak J_{\beta,R}\Phi)(r):=\int_r^R (q-r)^{-\beta}\Phi(q)\,dq,
\qquad 0\le r\le R,
\]
and let
\[
\mathfrak D_{\beta,R}^{\mathrm{dist}}\Phi
:=
-\frac{d}{dr}\mathfrak J_{\beta,R}\Phi
\]
in \(\mathcal D'((0,R))\). Then:
\begin{enumerate}[label=(\roman*)]
\item \(\Phi\in H^\beta((0,R))\) if and only if
\(\mathfrak D_{\beta,R}^{\mathrm{dist}}\Phi\in E_R\);
\item in that case,
\begin{equation}\label{eq:window-distribution}
\mathfrak T_{t,R}^{\mathrm{hl}}\Phi
=
\frac{\mathbf c_t}{\Gamma(1-\beta)}
\mathfrak D_{\beta,R}^{\mathrm{dist}}\Phi
\qquad\text{in }E_R,
\end{equation}
where \(\mathbf c_t\) is defined in \eqref{eq:model-coeff}.
\end{enumerate}
\end{lemma}

\begin{proof}
The equivalence in item~(i) is exactly the classical finite-interval Riemann--Liouville
characterization of \(H^\beta((0,R))\) by the right-sided derivative, i.e. the rescaled version of
the analytic input invoked in Theorem~\ref{thm:frozen-density}. In particular,
\[
H^\beta((0,R))
=
\left\{
\Phi\in E_R:
\mathfrak D_{\beta,R}^{\mathrm{dist}}\Phi\in E_R
\right\},
\]
and the map \(\Phi\mapsto \mathfrak D_{\beta,R}^{\mathrm{dist}}\Phi\) is bounded from
\(H^\beta((0,R))\) to \(E_R\).

For \(\Phi\in C_c^\infty((0,R))\), Definition~\ref{def:LW} and the proof of
Theorem~\ref{thm:micro-window-model} give
\[
\mathfrak T_{t,R}^{\mathrm{hl}}\Phi
=
\mathbf c_t \mathfrak D_-^\beta \Phi
=
-\frac{\mathbf c_t}{\Gamma(1-\beta)}
\frac{d}{dr}
\int_r^R (q-r)^{-\beta}\Phi(q)\,dq,
\]
which is exactly \eqref{eq:window-distribution} on the smooth core. Since
\(C_c^\infty((0,R))\) is dense in \(H^\beta((0,R))\) by Theorem~\ref{thm:micro-window-model}(iv),
\(\mathfrak T_{t,R}^{\mathrm{hl}}\) is closed with domain \(H^\beta((0,R))\), and
\(\Phi\mapsto \mathfrak D_{\beta,R}^{\mathrm{dist}}\Phi\) is bounded on \(H^\beta((0,R))\), the
identity extends to every \(\Phi\in H^\beta((0,R))\). This proves item~(ii).
\end{proof}

\begin{proposition}[Source-side conjugated decomposition on a fixed window]
\label{prop:micro-window-decompose}
Assume Assumption~\ref{ass:model} with the extra regularity from Definition~\ref{def:model-c2}.
Fix \(t\in(0,T)\), \(\alpha\in(0,1)\), \(R>0\), and suppose
\[
0<H_{XY}<H_Y<\tfrac12.
\]
Set
\[
\lambda:=H_Y+\frac12,
\qquad
\mu:=H_{XY}+\frac12,
\qquad
\beta:=H_Y-H_{XY},
\]
and define on \(\Delta_R:=\{0\le r\le q\le R\}\)
\[
g_{t,h,R}^{\mathrm{win}}(r,q)
:=
\frac{1}{\Gamma(\lambda)\Gamma(1-\lambda)}
\int_0^1
\theta^{-\lambda}(1-\theta)^{\mu-1}
L_{XY}(t-h(r+(q-r)\theta),t-hq)\,d\theta.
\]
Also define
\[
g_t^{\mathrm{hl}}
:=
\frac{L_{XY}(t,t)\Gamma(\mu)}{\Gamma(\lambda)\Gamma(1-\beta)}
=
\frac{L_Y(t,t)\mathbf c_t}{\Gamma(1-\beta)}.
\]
Then there exist \(h_0\in(0,t]\), \(C_{t,R}<\infty\), a multiplier
\[
\mathfrak M_{t,h,\alpha,R}^{\mathrm{win}}\in \mathcal L(E_R),
\qquad
(\mathfrak M_{t,h,\alpha,R}^{\mathrm{win}}\Phi)(r)
:=
m_{t,h,\alpha,R}(r)\Phi(r),
\]
and a bounded remainder operator
\[
\mathfrak K_{t,h,\alpha,R}^{\mathrm{win}}\in\mathcal L(E_R),
\]
such that for every \(0<h\le h_0\):
\begin{enumerate}[label=(\roman*)]
\item
\begin{equation}\label{eq:window-mult-close}
\|m_{t,h,\alpha,R}-1\|_{C^1([0,R])}\le C_{t,R}h;
\end{equation}
\item
\begin{equation}\label{eq:window-remainder-bound}
\|\mathfrak K_{t,h,\alpha,R}^{\mathrm{win}}\|_{\mathcal L(E_R)}
\le
C_{t,R}h;
\end{equation}
\item for every \(\Phi\in C_c^\infty((0,R))\),
\begin{equation}\label{eq:window-source-decompose}
\frac{1}{\Gamma(\lambda)}
\mathfrak D_-^\lambda
\bigl(
\mathfrak V_{XY,t,h,\alpha,R}^{\mathrm{mic}}\Phi
\bigr)
=
L_Y(t,t)\mathfrak M_{t,h,\alpha,R}^{\mathrm{win}}
\mathfrak T_{t,R}^{\mathrm{hl}}\Phi
+
\mathfrak K_{t,h,\alpha,R}^{\mathrm{win}}\Phi
\qquad\text{in }E_R.
\end{equation}
\item for every \(\Phi\in E_R\),
\begin{equation}\label{eq:window-source-decompose-dist}
\frac{1}{\Gamma(\lambda)}
\mathfrak D_-^\lambda
\bigl(
\mathfrak V_{XY,t,h,\alpha,R}^{\mathrm{mic}}\Phi
\bigr)
=
\frac{L_Y(t,t)\mathbf c_t}{\Gamma(1-\beta)}
\mathfrak M_{t,h,\alpha,R}^{\mathrm{win}}
\mathfrak D_{\beta,R}^{\mathrm{dist}}\Phi
+
\mathfrak K_{t,h,\alpha,R}^{\mathrm{win}}\Phi
\qquad\text{in }\mathcal D'((0,R)),
\end{equation}
Here the left-hand side denotes the distributional derivative of
\(\mathfrak I_-^{1-\lambda}(\mathfrak V_{XY,t,h,\alpha,R}^{\mathrm{mic}}\Phi)\).
\end{enumerate}
\end{proposition}

\begin{proof}
Because \(L_{XY}\) is \(C^2\), the function \(g_{t,h,R}^{\mathrm{win}}\) is \(C^2\) on
\(\Delta_R\), and
\begin{equation}\label{eq:window-g-c2}
\|g_{t,h,R}^{\mathrm{win}}-g_t^{\mathrm{hl}}\|_{C^2(\Delta_R)}
\le
C_{t,R}h
\end{equation}
for all sufficiently small \(h>0\). Define
\[
m_{t,h,\alpha,R}(r):=
\frac{g_{t,h,R}^{\mathrm{win}}(r,r)}{g_t^{\mathrm{hl}}},
\qquad 0\le r\le R.
\]
Then \eqref{eq:window-mult-close} follows from \eqref{eq:window-g-c2}.

For \(0\le r<q\le R\), set
\[
b_{t,h,\alpha,R}(r,q)
:=
\frac{g_{t,h,R}^{\mathrm{win}}(r,q)-g_{t,h,R}^{\mathrm{win}}(r,r)}{q-r},
\]
and extend continuously to the diagonal by
\[
b_{t,h,\alpha,R}(r,r):=\partial_2 g_{t,h,R}^{\mathrm{win}}(r,r).
\]
By \eqref{eq:window-g-c2}, one has
\begin{equation}\label{eq:b-window}
\|b_{t,h,\alpha,R}\|_{C^1(\Delta_R)}\le C_{t,R}h.
\end{equation}

Now let \(\Phi\in C_c^\infty((0,R))\). The proof of Proposition~\ref{prop:q-micro} gives
\[
\frac{1}{\Gamma(\lambda)}
\mathfrak D_-^\lambda
\bigl(
\mathfrak V_{XY,t,h,\alpha,R}^{\mathrm{mic}}\Phi
\bigr)(r)
=
-\frac{d}{dr}
\int_r^R
(q-r)^{-\beta}
g_{t,h,R}^{\mathrm{win}}(r,q)\Phi(q)\,dq
\]
for almost every \(r\in[0,R]\). Decomposing
\[
g_{t,h,R}^{\mathrm{win}}(r,q)
=
g_{t,h,R}^{\mathrm{win}}(r,r)
+
(q-r)b_{t,h,\alpha,R}(r,q),
\]
we obtain
\begin{align*}
\frac{1}{\Gamma(\lambda)}
\mathfrak D_-^\lambda
\bigl(
\mathfrak V_{XY,t,h,\alpha,R}^{\mathrm{mic}}\Phi
\bigr)(r)
&=
-g_{t,h,R}^{\mathrm{win}}(r,r)\frac{d}{dr}
\int_r^R (q-r)^{-\beta}\Phi(q)\,dq \\
&\quad
-\bigl(\partial_1 g_{t,h,R}^{\mathrm{win}}\bigr)(r,r)
\int_r^R (q-r)^{-\beta}\Phi(q)\,dq \\
&\quad
-\frac{d}{dr}
\int_r^R (q-r)^{1-\beta}
b_{t,h,\alpha,R}(r,q)\Phi(q)\,dq.
\end{align*}
Since
\[
g_{t,h,R}^{\mathrm{win}}(r,r)
=
g_t^{\mathrm{hl}}m_{t,h,\alpha,R}(r)
=
\frac{L_Y(t,t)\mathbf c_t}{\Gamma(1-\beta)}m_{t,h,\alpha,R}(r),
\]
Lemma~\ref{lem:window-distribution} yields
\[
-g_{t,h,R}^{\mathrm{win}}(r,r)\frac{d}{dr}
\int_r^R (q-r)^{-\beta}\Phi(q)\,dq
=
L_Y(t,t)\mathfrak M_{t,h,\alpha,R}^{\mathrm{win}}
\mathfrak T_{t,R}^{\mathrm{hl}}\Phi(r).
\]
Define \(\mathfrak K_{t,h,\alpha,R}^{\mathrm{win}}\) by
\begin{align*}
\bigl(
\mathfrak K_{t,h,\alpha,R}^{\mathrm{win}}\Phi
\bigr)(r)
:={}&
-\bigl(\partial_1 g_{t,h,R}^{\mathrm{win}}\bigr)(r,r)
\int_r^R (q-r)^{-\beta}\Phi(q)\,dq \\
&-
\frac{d}{dr}
\int_r^R (q-r)^{1-\beta}
b_{t,h,\alpha,R}(r,q)\Phi(q)\,dq.
\end{align*}
This proves \eqref{eq:window-source-decompose} on the smooth core.

It remains to show that \(\mathfrak K_{t,h,\alpha,R}^{\mathrm{win}}\) extends boundedly to \(E_R\)
with norm \(O(h)\). For the first term, \eqref{eq:window-g-c2} and the Hilbert--Schmidt property
of the kernel \((q-r)^{-\beta}\one_{\Delta_R}\) give an \(O(h)\) operator bound. For the second
term, the lower endpoint contributes no boundary term because \(1-\beta>0\), so differentiation
under the integral sign yields
\[
-\frac{d}{dr}
\int_r^R (q-r)^{1-\beta}
b_{t,h,\alpha,R}(r,q)\Phi(q)\,dq
=
\int_r^R
K_{t,h,\alpha,R}^{\mathrm{rem}}(r,q)\Phi(q)\,dq,
\]
where
\[
K_{t,h,\alpha,R}^{\mathrm{rem}}(r,q)
=
-(1-\beta)(q-r)^{-\beta}b_{t,h,\alpha,R}(r,q)
-
(q-r)^{1-\beta}\partial_1 b_{t,h,\alpha,R}(r,q).
\]
By \eqref{eq:b-window},
\[
|K_{t,h,\alpha,R}^{\mathrm{rem}}(r,q)|
\le
C_{t,R}h
\bigl(
(q-r)^{-\beta}+(q-r)^{1-\beta}
\bigr)\one_{\Delta_R}(r,q),
\]
and the right-hand side is Hilbert--Schmidt on \(\Delta_R\) because \(\beta<\tfrac12\). Hence
\eqref{eq:window-remainder-bound} follows.

To prove \eqref{eq:window-source-decompose-dist}, first note that the same Fubini computation as
above shows that for every \(\Phi\in E_R\),
\[
\bigl(
\mathfrak I_-^{1-\lambda}
\mathfrak V_{XY,t,h,\alpha,R}^{\mathrm{mic}}\Phi
\bigr)(r)
=
\int_r^R
(q-r)^{-\beta}
g_{t,h,R}^{\mathrm{win}}(r,q)\Phi(q)\,dq,
\qquad 0\le r\le R.
\]
Because the kernel \((q-r)^{-\beta}g_{t,h,R}^{\mathrm{win}}(r,q)\one_{\Delta_R}(r,q)\) is
square-integrable on \(\Delta_R\), the map
\[
\Phi\longmapsto
\mathfrak I_-^{1-\lambda}
\mathfrak V_{XY,t,h,\alpha,R}^{\mathrm{mic}}\Phi
\]
is bounded from \(E_R\) to \(E_R\), hence the left-hand side of
\eqref{eq:window-source-decompose-dist} defines a continuous linear map \(E_R\to\mathcal D'((0,R))\)
after distributional differentiation. By the boundedness of \(\mathfrak J_{\beta,R}\) on \(E_R\),
the definition of \(\mathfrak D_{\beta,R}^{\mathrm{dist}}\), and
\eqref{eq:window-remainder-bound}, the right-hand side of
\eqref{eq:window-source-decompose-dist} is also continuous from \(E_R\) to \(\mathcal D'((0,R))\).
Since \(C_c^\infty((0,R))\) is dense in \(E_R\) and \eqref{eq:window-source-decompose-dist}
coincides with \eqref{eq:window-source-decompose} on the smooth core by
Lemma~\ref{lem:window-distribution}, the distributional identity extends to every \(\Phi\in E_R\).
\end{proof}

\begin{theorem}[Fractional Sobolev characterization of the exact fixed-window solver]
\label{thm:micro-window-regularity}
Assume Assumptions~\ref{ass:model} and~\ref{ass:obs-source}, and suppose that
Definition~\ref{def:model-c2} holds. Let \(\eta>0\), fix \(t\in(0,T)\), \(\alpha\in(0,1)\), and
\(R>0\), and suppose
\[
0<H_{XY}<H_Y<\tfrac12.
\]
Then there exist \(h_0\in(0,t]\) and \(C_{t,R}<\infty\) such that for every \(0<h\le h_0\) with
\[
hR<t,
\qquad
R<\varepsilon_{h,\alpha}^{-1},
\]
the following hold:
\begin{enumerate}[label=(\roman*)]
\item
\begin{equation}\label{eq:window-domain-regularity}
\mathcal D(\mathfrak T_{t,h,\alpha,R}^{\mathrm{ex}})
=
H^\beta((0,R));
\end{equation}
\item there exist bounded operators
\[
\mathfrak A_{t,h,\alpha,R}^{\mathrm{win}},
\mathfrak R_{t,h,\alpha,R}^{\mathrm{win}}
\in \mathcal L(E_R)
\]
such that \(\mathfrak A_{t,h,\alpha,R}^{\mathrm{win}}\) is invertible,
\begin{equation}\label{eq:window-A-close}
\left\|
\mathfrak A_{t,h,\alpha,R}^{\mathrm{win}}-I_{E_R}
\right\|_{\mathcal L(E_R)}
+
\left\|
\mathfrak R_{t,h,\alpha,R}^{\mathrm{win}}
\right\|_{\mathcal L(E_R)}
\le
C_{t,R}h,
\end{equation}
and
\begin{equation}\label{eq:window-solver-representation}
\mathfrak T_{t,h,\alpha,R}^{\mathrm{ex}}\Phi
=
\mathfrak A_{t,h,\alpha,R}^{\mathrm{win}}
\mathfrak T_{t,R}^{\mathrm{hl}}\Phi
+
\mathfrak R_{t,h,\alpha,R}^{\mathrm{win}}\Phi
\qquad
\text{for every }\Phi\in H^\beta((0,R));
\end{equation}
\item the exact graph norm is uniformly equivalent to the \(H^\beta\) norm:
\begin{equation}\label{eq:window-graph-equivalence}
C_{t,R}^{-1}\|\Phi\|_{H^\beta((0,R))}
\le
\|\Phi\|_{\mathfrak G_{t,h,\alpha,R}^{\mathrm{ex}}}
\le
C_{t,R}\|\Phi\|_{H^\beta((0,R))}
\qquad
\text{for every }\Phi\in H^\beta((0,R)).
\end{equation}
\end{enumerate}
\end{theorem}

\begin{proof}
Let \(h_0\) be small enough that Proposition~\ref{prop:micro-window-decompose} and
Corollary~\ref{cor:BY-micro-invertible} both hold for every \(0<h\le h_0\). For such \(h\), define
\[
\mathfrak A_{t,h,\alpha,R}^{\mathrm{win}}
:=
\bigl(
\mathfrak B_{Y,t,h,\alpha,R}^{\mathrm{mic}}
\bigr)^{-1}
L_Y(t,t)\mathfrak M_{t,h,\alpha,R}^{\mathrm{win}},
\qquad
\mathfrak R_{t,h,\alpha,R}^{\mathrm{win}}
:=
\bigl(
\mathfrak B_{Y,t,h,\alpha,R}^{\mathrm{mic}}
\bigr)^{-1}
\mathfrak K_{t,h,\alpha,R}^{\mathrm{win}}.
\]
By \eqref{eq:BY-micro-close}, \eqref{eq:BY-micro-inv}, \eqref{eq:window-mult-close}, and
\eqref{eq:window-remainder-bound}, one has
\[
\left\|
\mathfrak A_{t,h,\alpha,R}^{\mathrm{win}}-I_{E_R}
\right\|_{\mathcal L(E_R)}
\le
C_{t,R}h,
\qquad
\left\|
\mathfrak R_{t,h,\alpha,R}^{\mathrm{win}}
\right\|_{\mathcal L(E_R)}
\le
C_{t,R}h,
\]
after possibly shrinking \(h_0\). In particular \(\mathfrak A_{t,h,\alpha,R}^{\mathrm{win}}\) is
invertible with uniformly bounded inverse, which proves \eqref{eq:window-A-close}.

Now let \(\Phi\in C_c^\infty((0,R))\), and define
\[
\Psi_h
:=
\mathfrak A_{t,h,\alpha,R}^{\mathrm{win}}
\mathfrak T_{t,R}^{\mathrm{hl}}\Phi
+
\mathfrak R_{t,h,\alpha,R}^{\mathrm{win}}\Phi
\in E_R.
\]
By construction and Proposition~\ref{prop:micro-window-decompose},
\[
\mathfrak B_{Y,t,h,\alpha,R}^{\mathrm{mic}}\Psi_h
=
L_Y(t,t)\mathfrak M_{t,h,\alpha,R}^{\mathrm{win}}
\mathfrak T_{t,R}^{\mathrm{hl}}\Phi
+
\mathfrak K_{t,h,\alpha,R}^{\mathrm{win}}\Phi
=
\frac{1}{\Gamma(\lambda)}
\mathfrak D_-^\lambda
\bigl(
\mathfrak V_{XY,t,h,\alpha,R}^{\mathrm{mic}}\Phi
\bigr)
\]
in \(E_R\). Exactly as in the proof of Theorem~\ref{thm:micro-smooth-core}, both
\[
\mathfrak V_{Y,t,h,\alpha,R}^{\mathrm{mic}}\Psi_h
\quad\text{and}\quad
\mathfrak V_{XY,t,h,\alpha,R}^{\mathrm{mic}}\Phi
\]
are compactly supported in \([0,R]\), have the same conjugated \(\lambda\)-derivative, and
therefore coincide. Hence
\[
(\Phi,\Psi_h)\in \Gamma_{t,h,\alpha,R}^{\mathrm{ex}},
\]
so by Proposition~\ref{prop:micro-window-solver},
\[
\mathfrak T_{t,h,\alpha,R}^{\mathrm{ex}}\Phi=\Psi_h
\qquad\text{for every }\Phi\in C_c^\infty((0,R)).
\]
Thus \eqref{eq:window-solver-representation} holds on the smooth core.

Since \(C_c^\infty((0,R))\) is dense in \(H^\beta((0,R))\) and
\(\mathfrak T_{t,R}^{\mathrm{hl}}\) is closed with domain \(H^\beta((0,R))\),
\eqref{eq:window-solver-representation} defines a closed operator on \(H^\beta((0,R))\). Passing
to the limit along smooth approximations and using the boundedness of
\(\mathfrak V_{Y,t,h,\alpha,R}^{\mathrm{mic}}\) and \(\mathfrak V_{XY,t,h,\alpha,R}^{\mathrm{mic}}\),
we conclude that
\[
H^\beta((0,R))
\subset
\mathcal D(\mathfrak T_{t,h,\alpha,R}^{\mathrm{ex}})
\]
and \eqref{eq:window-solver-representation} holds for every \(\Phi\in H^\beta((0,R))\).

  For the reverse inclusion, let
  \[
  \Phi\in\mathcal D(\mathfrak T_{t,h,\alpha,R}^{\mathrm{ex}}),
  \qquad
  \Psi:=\mathfrak T_{t,h,\alpha,R}^{\mathrm{ex}}\Phi\in E_R.
\]
Then \((\Phi,\Psi)\in\Gamma_{t,h,\alpha,R}^{\mathrm{ex}}\), so
  \[
  \mathfrak V_{Y,t,h,\alpha,R}^{\mathrm{mic}}\Psi
  =
  \mathfrak V_{XY,t,h,\alpha,R}^{\mathrm{mic}}\Phi
  \qquad\text{in }E_R.
  \]
  Proposition~\ref{prop:BY-micro}(i) and
  Proposition~\ref{prop:micro-window-decompose}(iv) therefore give
  \[
  \mathfrak B_{Y,t,h,\alpha,R}^{\mathrm{mic}}\Psi
  =
  \frac{L_Y(t,t)\mathbf c_t}{\Gamma(1-\beta)}
  \mathfrak M_{t,h,\alpha,R}^{\mathrm{win}}
  \mathfrak D_{\beta,R}^{\mathrm{dist}}\Phi
  +
  \mathfrak K_{t,h,\alpha,R}^{\mathrm{win}}\Phi
  \]
  in \(\mathcal D'((0,R))\). Since
  \(\mathfrak B_{Y,t,h,\alpha,R}^{\mathrm{mic}}\Psi\in E_R\) and
  \(\mathfrak K_{t,h,\alpha,R}^{\mathrm{win}}\Phi\in E_R\), the multiplier
  \(\mathfrak M_{t,h,\alpha,R}^{\mathrm{win}}\), being boundedly invertible for small \(h\), shows
  that \(\mathfrak D_{\beta,R}^{\mathrm{dist}}\Phi\in E_R\). Lemma~\ref{lem:window-distribution}
  then yields \(\Phi\in H^\beta((0,R))\). This proves \eqref{eq:window-domain-regularity}, and
  \eqref{eq:window-solver-representation} follows from the uniqueness part of
  Proposition~\ref{prop:micro-window-solver}.
  
  Finally, \eqref{eq:window-solver-representation}, \eqref{eq:window-A-close}, and the boundedness
  of \(\bigl(\mathfrak A_{t,h,\alpha,R}^{\mathrm{win}}\bigr)^{-1}\) imply
  \begin{align*}
  \|\mathfrak T_{t,h,\alpha,R}^{\mathrm{ex}}\Phi\|_{E_R}
  &\le
  C_{t,R}\bigl(
  \|\mathfrak T_{t,R}^{\mathrm{hl}}\Phi\|_{E_R}+\|\Phi\|_{E_R}
  \bigr),\\
  \|\mathfrak T_{t,R}^{\mathrm{hl}}\Phi\|_{E_R}
  &\le
  C_{t,R}\bigl(
  \|\mathfrak T_{t,h,\alpha,R}^{\mathrm{ex}}\Phi\|_{E_R}+\|\Phi\|_{E_R}
  \bigr)
  \end{align*}
  for every \(\Phi\in H^\beta((0,R))\). Combining these bounds with
  Theorem~\ref{thm:micro-window-model}(iii), which identifies the model graph norm with the
  \(H^\beta((0,R))\) norm, proves \eqref{eq:window-graph-equivalence}.
\end{proof}

\begin{proposition}[Conjugated characterization of the exact fixed-window graph]
\label{prop:micro-window-conjugated}
Assume the hypotheses of Theorem~\ref{thm:micro-window-regularity}. Then for every
\[
\Phi\in H^\beta((0,R)),
\qquad
\Psi\in E_R,
\]
the following are equivalent:
\begin{enumerate}[label=(\roman*)]
\item \((\Phi,\Psi)\in \Gamma_{t,h,\alpha,R}^{\mathrm{ex}}\);
\item
\begin{equation}\label{eq:window-conjugated-characterization}
\mathfrak B_{Y,t,h,\alpha,R}^{\mathrm{mic}}\Psi
=
L_Y(t,t)\mathfrak M_{t,h,\alpha,R}^{\mathrm{win}}
\mathfrak T_{t,R}^{\mathrm{hl}}\Phi
+
\mathfrak K_{t,h,\alpha,R}^{\mathrm{win}}\Phi
\qquad\text{in }E_R.
\end{equation}
\end{enumerate}
\end{proposition}

\begin{proof}
If \((\Phi,\Psi)\in \Gamma_{t,h,\alpha,R}^{\mathrm{ex}}\), then
\[
\Psi=\mathfrak T_{t,h,\alpha,R}^{\mathrm{ex}}\Phi
\]
by Proposition~\ref{prop:micro-window-solver}. Applying
\eqref{eq:window-solver-representation} from Theorem~\ref{thm:micro-window-regularity} and then the
definitions of \(\mathfrak A_{t,h,\alpha,R}^{\mathrm{win}}\) and
\(\mathfrak R_{t,h,\alpha,R}^{\mathrm{win}}\) gives
\eqref{eq:window-conjugated-characterization}.

Conversely, suppose \eqref{eq:window-conjugated-characterization} holds. Since
\(\Phi\in H^\beta((0,R))\), Theorem~\ref{thm:micro-window-regularity} yields
\[
\Psi_*:=\mathfrak T_{t,h,\alpha,R}^{\mathrm{ex}}\Phi\in E_R
\]
and
\[
\mathfrak B_{Y,t,h,\alpha,R}^{\mathrm{mic}}\Psi_*
=
L_Y(t,t)\mathfrak M_{t,h,\alpha,R}^{\mathrm{win}}
\mathfrak T_{t,R}^{\mathrm{hl}}\Phi
+
\mathfrak K_{t,h,\alpha,R}^{\mathrm{win}}\Phi
\]
in \(E_R\). Subtracting from \eqref{eq:window-conjugated-characterization} gives
\[
\mathfrak B_{Y,t,h,\alpha,R}^{\mathrm{mic}}(\Psi-\Psi_*)=0.
\]
Corollary~\ref{cor:BY-micro-invertible} therefore implies \(\Psi=\Psi_*\). Hence
\[
(\Phi,\Psi)\in\Gamma_{t,h,\alpha,R}^{\mathrm{ex}},
\]
which proves (i).
\end{proof}

\begin{theorem}[Stable correction of a full-window conjugated defect]
\label{thm:micro-window-correction}
Assume the hypotheses of Theorem~\ref{thm:micro-window-regularity}. Then there exists
\(C_{t,R}<\infty\) such that for every
\[
\Phi\in H^\beta((0,R)),
\qquad
\Psi\in E_R,
\]
the conjugated defect
\begin{equation}\label{eq:window-defect}
\mathfrak F_{t,h,\alpha,R}^{\mathrm{ex}}(\Phi,\Psi)
:=
\mathfrak B_{Y,t,h,\alpha,R}^{\mathrm{mic}}\Psi
-
L_Y(t,t)\mathfrak M_{t,h,\alpha,R}^{\mathrm{win}}
\mathfrak T_{t,R}^{\mathrm{hl}}\Phi
-
\mathfrak K_{t,h,\alpha,R}^{\mathrm{win}}\Phi
\in E_R
\end{equation}
satisfies the exact correction formula
\begin{equation}\label{eq:window-correction-formula}
\mathfrak T_{t,h,\alpha,R}^{\mathrm{ex}}\Phi
=
\Psi
-
\bigl(
\mathfrak B_{Y,t,h,\alpha,R}^{\mathrm{mic}}
\bigr)^{-1}
\mathfrak F_{t,h,\alpha,R}^{\mathrm{ex}}(\Phi,\Psi),
\end{equation}
and the stability estimate
\begin{equation}\label{eq:window-correction-bound}
\left\|
\Psi-\mathfrak T_{t,h,\alpha,R}^{\mathrm{ex}}\Phi
\right\|_{E_R}
\le
C_{t,R}
\left\|
\mathfrak F_{t,h,\alpha,R}^{\mathrm{ex}}(\Phi,\Psi)
\right\|_{E_R}.
\end{equation}
\end{theorem}

\begin{proof}
Let
\[
f:=\mathfrak F_{t,h,\alpha,R}^{\mathrm{ex}}(\Phi,\Psi)\in E_R,
\]
and define
\[
\widetilde\Psi
:=
\Psi
-
\bigl(
\mathfrak B_{Y,t,h,\alpha,R}^{\mathrm{mic}}
\bigr)^{-1}f
\in E_R.
\]
By construction,
\[
\mathfrak B_{Y,t,h,\alpha,R}^{\mathrm{mic}}\widetilde\Psi
=
L_Y(t,t)\mathfrak M_{t,h,\alpha,R}^{\mathrm{win}}
\mathfrak T_{t,R}^{\mathrm{hl}}\Phi
+
\mathfrak K_{t,h,\alpha,R}^{\mathrm{win}}\Phi
\qquad\text{in }E_R.
\]
Proposition~\ref{prop:micro-window-conjugated} therefore implies
\[
(\Phi,\widetilde\Psi)\in\Gamma_{t,h,\alpha,R}^{\mathrm{ex}},
\]
hence
\[
\widetilde\Psi=\mathfrak T_{t,h,\alpha,R}^{\mathrm{ex}}\Phi.
\]
This proves \eqref{eq:window-correction-formula}. The bound
\eqref{eq:window-correction-bound} then follows from
\eqref{eq:BY-micro-inv}.
\end{proof}

\begin{definition}[Smooth buffered cutoff on the micro half-line]
\label{def:micro-smooth-cutoff}
Let \(0<R_0<R_1<R_2<\infty\). A function
\[
\chi\in C^\infty([0,\infty))
\]
is called a \emph{smooth buffered cutoff} with radii \((R_0,R_1,R_2)\) if
\[
0\le \chi\le 1
\quad\text{on }[0,\infty),
\qquad
\chi\equiv1\ \text{on }[0,R_1],
\qquad
\operatorname{supp}\chi\subset [0,R_2].
\]
\end{definition}

\begin{proposition}[Full-window commutator forcing for a smooth cutoff]
\label{prop:micro-smooth-commutator}
Assume Assumption~\ref{ass:model}, fix \(t\in(0,T)\), \(\alpha\in(0,1)\), and suppose
\[
0<H_{XY}<H_Y<\tfrac12.
\]
Set
\[
\lambda:=H_Y+\frac12,
\qquad
\mu:=H_{XY}+\frac12,
\qquad
\beta:=H_Y-H_{XY}.
\]
Let \(0<R_1<R_2\), let
\[
\chi\in C^\infty([0,\infty)),
\qquad
0\le\chi\le1,
\qquad
\operatorname{supp}\chi\subset[0,R_2],
\]
and choose
\(h>0\) sufficiently small that
\[
hR_2<t,
\qquad
R_2<\varepsilon_{h,\alpha}^{-1}.
\]
For \(\Theta\in L^2(\R_+)\), define on \(E_{R_2}=L^2([0,R_2])\) the commutator Volterra operators
\begin{align}
\bigl(
\mathfrak W_{Y,t,h,\alpha,\chi}^{\mathrm{com}}\Theta
\bigr)(r)
&:=
\int_r^{\varepsilon_{h,\alpha}^{-1}}
(q-r)^{\lambda-1}
\bigl(\chi(q)-\chi(r)\bigr)
L_Y(t-hr,t-hq)\Theta(q)\,dq,
\label{eq:micro-smooth-com-Y}\\
\bigl(
\mathfrak W_{XY,t,h,\alpha,\chi}^{\mathrm{com}}\Theta
\bigr)(r)
&:=
\int_r^{\varepsilon_{h,\alpha}^{-1}}
(q-r)^{\mu-1}
\bigl(\chi(q)-\chi(r)\bigr)
L_{XY}(t-hr,t-hq)\Theta(q)\,dq
\label{eq:micro-smooth-com-XY}
\end{align}
for \(0\le r\le R_2\). Then there exist bounded operators
\[
\mathfrak Q_{Y,t,h,\alpha,\chi}^{\mathrm{com}},
\mathfrak Q_{XY,t,h,\alpha,\chi}^{\mathrm{com}}
\in \mathcal L(L^2(\R_+),E_{R_2})
\]
such that for every \(\Theta\in L^2(\R_+)\),
\begin{align}
\frac{1}{\Gamma(\lambda)}
\mathfrak D_-^\lambda
\bigl(
\mathfrak W_{Y,t,h,\alpha,\chi}^{\mathrm{com}}\Theta
\bigr)
&=
\mathfrak Q_{Y,t,h,\alpha,\chi}^{\mathrm{com}}\Theta
\qquad\text{in }E_{R_2},
\label{eq:micro-smooth-com-conj-Y}\\
\frac{1}{\Gamma(\lambda)}
\mathfrak D_-^\lambda
\bigl(
\mathfrak W_{XY,t,h,\alpha,\chi}^{\mathrm{com}}\Theta
\bigr)
&=
\mathfrak Q_{XY,t,h,\alpha,\chi}^{\mathrm{com}}\Theta
\qquad\text{in }E_{R_2}.
\label{eq:micro-smooth-com-conj-XY}
\end{align}
\end{proposition}

\begin{proof}
Set \(\varepsilon:=\varepsilon_{h,\alpha}\). For \(0\le r<q\le \varepsilon^{-1}\), define
\[
\delta_\chi(r,q):=\frac{\chi(q)-\chi(r)}{q-r},
\]
and extend continuously to the diagonal by \(\delta_\chi(r,r):=\chi'(r)\). Since \(\chi\) is
smooth and compactly supported in \([0,R_2]\), the function \(\delta_\chi\) is \(C^1\) on
\[
\Delta_{R_2,\varepsilon}:=\{0\le r\le q\le \varepsilon^{-1}\},
\]
with finite \(C^1\) norm.

For the target-side commutator, write
\[
\bigl(
\mathfrak W_{Y,t,h,\alpha,\chi}^{\mathrm{com}}\Theta
\bigr)(r)
=
\int_r^{\varepsilon^{-1}}
(q-r)^\lambda
\delta_\chi(r,q)L_Y(t-hr,t-hq)\Theta(q)\,dq.
\]
Exactly the same Fubini computation as in Proposition~\ref{prop:micro-tail-Y}, now with exponent
\(\lambda\) instead of \(\lambda-1\), shows that
\[
\mathfrak I_-^{1-\lambda}
\bigl(
\mathfrak W_{Y,t,h,\alpha,\chi}^{\mathrm{com}}\Theta
\bigr)(r)
=
\int_r^{\varepsilon^{-1}}
(q-r)F_{Y,t,h,\chi}(r,q)\Theta(q)\,dq,
\]
where \(F_{Y,t,h,\chi}\) is \(C^1\) on \(\Delta_{R_2,\varepsilon}\). Differentiating under the
integral sign therefore yields
\[
\frac{1}{\Gamma(\lambda)}
\mathfrak D_-^\lambda
\bigl(
\mathfrak W_{Y,t,h,\alpha,\chi}^{\mathrm{com}}\Theta
\bigr)(r)
=
F_{Y,t,h,\chi}(r,r)\Theta(r)
-
\int_r^{\varepsilon^{-1}}
K_{Y,t,h,\chi}^{\mathrm{com}}(r,q)\Theta(q)\,dq
\]
for almost every \(r\in[0,R_2]\), with bounded kernel
\[
K_{Y,t,h,\chi}^{\mathrm{com}}(r,q)
:=
F_{Y,t,h,\chi}(r,q)
+
(q-r)\partial_1F_{Y,t,h,\chi}(r,q).
\]
Since \(\Delta_{R_2,\varepsilon}\) has finite measure, this defines a bounded operator
\(\mathfrak Q_{Y,t,h,\alpha,\chi}^{\mathrm{com}}\in\mathcal L(L^2(\R_+),E_{R_2})\), proving
\eqref{eq:micro-smooth-com-conj-Y}.

For the source-side commutator, write
\[
\bigl(
\mathfrak W_{XY,t,h,\alpha,\chi}^{\mathrm{com}}\Theta
\bigr)(r)
=
\int_r^{\varepsilon^{-1}}
(q-r)^\mu
\delta_\chi(r,q)L_{XY}(t-hr,t-hq)\Theta(q)\,dq.
\]
Exactly as in Proposition~\ref{prop:q-micro}, the same Beta-function computation gives
\[
\mathfrak I_-^{1-\lambda}
\bigl(
\mathfrak W_{XY,t,h,\alpha,\chi}^{\mathrm{com}}\Theta
\bigr)(r)
=
\int_r^{\varepsilon^{-1}}
(q-r)^{1-\beta}
F_{XY,t,h,\chi}(r,q)\Theta(q)\,dq,
\]
where \(F_{XY,t,h,\chi}\) is \(C^1\) on \(\Delta_{R_2,\varepsilon}\). Because \(1-\beta>0\),
differentiation under the integral sign yields
\begin{align*}
\frac{1}{\Gamma(\lambda)}
\mathfrak D_-^\lambda
\bigl(
\mathfrak W_{XY,t,h,\alpha,\chi}^{\mathrm{com}}\Theta
\bigr)(r)
&=
-(1-\beta)
\int_r^{\varepsilon^{-1}}
(q-r)^{-\beta}
F_{XY,t,h,\chi}(r,q)\Theta(q)\,dq\\
&\quad-
\int_r^{\varepsilon^{-1}}
(q-r)^{1-\beta}
\partial_1F_{XY,t,h,\chi}(r,q)\Theta(q)\,dq
\end{align*}
for almost every \(r\in[0,R_2]\). Since \(\beta<\tfrac12\) and
\(\Delta_{R_2,\varepsilon}\) has finite measure, both kernels are square-integrable on
\(\Delta_{R_2,\varepsilon}\). Hence this defines a bounded operator
\(\mathfrak Q_{XY,t,h,\alpha,\chi}^{\mathrm{com}}\in\mathcal L(L^2(\R_+),E_{R_2})\), proving
\eqref{eq:micro-smooth-com-conj-XY}.
\end{proof}

\begin{theorem}[Smooth buffered localization of a global exact micro pair]
\label{thm:micro-smooth-buffer}
Assume Assumptions~\ref{ass:model} and~\ref{ass:obs-source}, let \(\eta>0\), fix \(t\in(0,T)\),
\(\alpha\in(0,1)\), and suppose
\[
0<H_{XY}<H_Y<\tfrac12.
\]
Let \(0<R_0<R_1<R_2\), let \(\chi\) be a smooth buffered cutoff with radii
\((R_0,R_1,R_2)\), and choose \(h>0\) sufficiently small that
\[
hR_2<t,
\qquad
R_2<\varepsilon_{h,\alpha}^{-1}.
\]
Let
\[
\Phi\in\mathcal D(\mathfrak T_{t,h,\alpha}^{\mathrm{mic}}),
\qquad
\Psi:=\mathfrak T_{t,h,\alpha}^{\mathrm{mic}}\Phi.
\]
Then:
\begin{enumerate}[label=(\roman*)]
\item \(\chi\Phi\in H^\beta((0,R_2))\), and the localized pair satisfies the full-window defect
identity
\begin{equation}\label{eq:micro-smooth-buffer-full}
\mathfrak F_{t,h,\alpha,R_2}^{\mathrm{ex}}(\chi\Phi,\chi\Psi)
=
\mathfrak Q_{Y,t,h,\alpha,\chi}^{\mathrm{com}}\Psi
-
\mathfrak Q_{XY,t,h,\alpha,\chi}^{\mathrm{com}}\Phi
\qquad\text{in }E_{R_2};
\end{equation}
\item on the inner window \([0,R_0]\), the full-window defect reduces exactly to the buffered tail
forcing:
\begin{equation}\label{eq:micro-smooth-buffer-inner}
\one_{[0,R_0]}
\mathfrak F_{t,h,\alpha,R_2}^{\mathrm{ex}}(\chi\Phi,\chi\Psi)
=
\mathfrak Q_{XY,t,h,\alpha,R_0,R_1}^{\mathrm{tail}}
\bigl((1-\chi)\Phi\bigr)
-
\mathfrak Q_{Y,t,h,\alpha,R_0,R_1}^{\mathrm{tail}}
\bigl((1-\chi)\Psi\bigr)
\qquad\text{in }E_{R_0};
\end{equation}
\item there exists \(C_{t,R_0,\alpha}<\infty\) such that
\begin{align}
\left\|
\one_{[0,R_0]}
\mathfrak F_{t,h,\alpha,R_2}^{\mathrm{ex}}(\chi\Phi,\chi\Psi)
\right\|_{E_{R_0}}
\notag\\
\le
C_{t,R_0,\alpha}
\bigl((R_1-R_0)^{-\beta}+h^{\frac{1+\alpha}{2}}\bigr)
\bigl(\|(1-\chi)\Phi\|_{L^2(\R_+)}+\|(1-\chi)\Psi\|_{L^2(\R_+)}\bigr);
\label{eq:micro-smooth-buffer-bound}
\end{align}
\item defining
\begin{equation}\label{eq:micro-smooth-buffer-corrected}
\Psi_\chi^{\mathrm{corr}}
:=
\chi\Psi
-
\bigl(
\mathfrak B_{Y,t,h,\alpha,R_2}^{\mathrm{mic}}
\bigr)^{-1}
\mathfrak F_{t,h,\alpha,R_2}^{\mathrm{ex}}(\chi\Phi,\chi\Psi),
\end{equation}
one has
\[
(\chi\Phi,\Psi_\chi^{\mathrm{corr}})
\in
\Gamma_{t,h,\alpha,R_2}^{\mathrm{ex}}.
\]
\end{enumerate}
\end{theorem}

\begin{proof}
Set \(\varepsilon:=\varepsilon_{h,\alpha}\). Since \(\chi\) is supported in \([0,R_2]\), one has
\[
\chi\Phi,\chi\Psi\in E_{R_2}.
\]
By Corollary~\ref{cor:micro-volterra-transfer},
\[
\mathfrak V_{Y,t,h,\alpha}^{\mathrm{mic}}\Psi
=
\mathfrak V_{XY,t,h,\alpha}^{\mathrm{mic}}\Phi
\qquad\text{in }L^2(\R_+).
\]
Multiplying by \(\chi(r)\) and subtracting from the Volterra operators applied to \(\chi\Phi\) and
\(\chi\Psi\) gives, on \([0,R_2]\),
\[
\mathfrak V_{Y,t,h,\alpha,R_2}^{\mathrm{mic}}(\chi\Psi)
-
\mathfrak V_{XY,t,h,\alpha,R_2}^{\mathrm{mic}}(\chi\Phi)
=
\mathfrak W_{Y,t,h,\alpha,\chi}^{\mathrm{com}}\Psi
-
\mathfrak W_{XY,t,h,\alpha,\chi}^{\mathrm{com}}\Phi.
\]
Applying Proposition~\ref{prop:micro-smooth-commutator} and the distributional source-side
decomposition \eqref{eq:window-source-decompose-dist} from
Proposition~\ref{prop:micro-window-decompose} therefore yields
\begin{align*}
\mathfrak B_{Y,t,h,\alpha,R_2}^{\mathrm{mic}}(\chi\Psi)
&=
\frac{L_Y(t,t)\mathbf c_t}{\Gamma(1-\beta)}
\mathfrak M_{t,h,\alpha,R_2}^{\mathrm{win}}
\mathfrak D_{\beta,R_2}^{\mathrm{dist}}(\chi\Phi)\\
&\quad+
\mathfrak K_{t,h,\alpha,R_2}^{\mathrm{win}}(\chi\Phi)
+
\mathfrak Q_{Y,t,h,\alpha,\chi}^{\mathrm{com}}\Psi
-
\mathfrak Q_{XY,t,h,\alpha,\chi}^{\mathrm{com}}\Phi
\end{align*}
in \(\mathcal D'((0,R_2))\). Since the rightmost two terms belong to \(E_{R_2}\), Lemma~\ref{lem:window-distribution}
implies \(\chi\Phi\in H^\beta((0,R_2))\). This proves the regularity claim in (i), and
\eqref{eq:micro-smooth-buffer-full} is now exactly the definition of the defect
\(\mathfrak F_{t,h,\alpha,R_2}^{\mathrm{ex}}(\chi\Phi,\chi\Psi)\).

For \(0\le r\le R_0\), one has \(\chi(r)=1\). Since \(\chi\equiv1\) on \([0,R_1]\),
\((1-\chi)\Theta\) vanishes on \([0,R_1]\) for every \(\Theta\). Therefore
\[
\bigl(
\mathfrak W_{Y,t,h,\alpha,\chi}^{\mathrm{com}}\Psi
\bigr)(r)
=
-
\int_{R_1}^{\varepsilon^{-1}}
(q-r)^{\lambda-1}
L_Y(t-hr,t-hq)(1-\chi(q))\Psi(q)\,dq
\]
and similarly
\[
\bigl(
\mathfrak W_{XY,t,h,\alpha,\chi}^{\mathrm{com}}\Phi
\bigr)(r)
=
-
\int_{R_1}^{\varepsilon^{-1}}
(q-r)^{\mu-1}
L_{XY}(t-hr,t-hq)(1-\chi(q))\Phi(q)\,dq.
\]
Because \(\Phi\) and \(\Psi\) are supported in \([0,\varepsilon^{-1}]\), the functions
\((1-\chi)\Phi\) and \((1-\chi)\Psi\) belong to \(F_{h,R_1}\). Thus, on \([0,R_0]\),
\[
\mathfrak W_{Y,t,h,\alpha,\chi}^{\mathrm{com}}\Psi
=
-\mathfrak W_{Y,t,h,\alpha,R_0,R_1}^{\mathrm{tail}}
\bigl((1-\chi)\Psi\bigr),
\]
and likewise for the source side. Applying Propositions~\ref{prop:micro-tail-Y}
and~\ref{prop:micro-tail-XY} proves \eqref{eq:micro-smooth-buffer-inner}. The bound
\eqref{eq:micro-smooth-buffer-bound} follows from
\eqref{eq:micro-tail-Y-bound}, \eqref{eq:micro-tail-XY-bound}, and the triangle inequality.

Finally, part (iv) is exactly Theorem~\ref{thm:micro-window-correction} applied with
\(\Phi\) replaced by \(\chi\Phi\) and \(\Psi\) replaced by \(\chi\Psi\).
\end{proof}

\section{Quantitative fixed-window localization}\label{sec:micro-window-buffer}

Localization needs a reusable fixed-window estimate.  The smooth-buffer theorem starts from a
global exact micro pair and produces an exact corrected pair on the localized window.  On a finite
window, the commutator created by a smooth cutoff is controlled quantitatively by the width of the
transition layer.  Exact fixed-window graph elements can therefore be localized with an explicit
loss.

\begin{lemma}[Quantitative smooth buffered cutoffs]
\label{lem:quant-buffer}
There exists a universal constant \(C_{\chi}<\infty\) such that for every
\[
0<R_1<R_2<\infty
\]
there exists a smooth buffered cutoff
\[
\chi_{R_1,R_2}\in C^\infty([0,\infty))
\]
with
\[
0\le \chi_{R_1,R_2}\le 1,
\qquad
\chi_{R_1,R_2}\equiv1\ \text{on }[0,R_1],
\qquad
\operatorname{supp}\chi_{R_1,R_2}\subset[0,R_2],
\]
and
\begin{equation}\label{eq:quant-buffer-derivatives}
\|\chi_{R_1,R_2}'\|_{L^\infty(\R_+)}
\le
\frac{C_{\chi}}{R_2-R_1},
\qquad
\|\chi_{R_1,R_2}''\|_{L^\infty(\R_+)}
\le
\frac{C_{\chi}}{(R_2-R_1)^2}.
\end{equation}
If
\[
\delta_{R_1,R_2}(r,q)
:=
\frac{\chi_{R_1,R_2}(q)-\chi_{R_1,R_2}(r)}{q-r},
\qquad 0\le r<q\le R_2,
\]
with continuous extension \(\delta_{R_1,R_2}(r,r):=\chi_{R_1,R_2}'(r)\), then
\begin{equation}\label{eq:quant-buffer-delta}
\|\delta_{R_1,R_2}\|_{C^1(\Delta_{R_2})}
\le
C_{\chi}
\left(
\frac{1}{R_2-R_1}
+
\frac{1}{(R_2-R_1)^2}
\right),
\qquad
\Delta_{R_2}:=\{0\le r\le q\le R_2\}.
\end{equation}
\end{lemma}

\begin{proof}
Choose \(\vartheta\in C^\infty(\R)\) such that
\[
0\le\vartheta\le1,
\qquad
\vartheta\equiv1\ \text{on }(-\infty,0],
\qquad
\vartheta\equiv0\ \text{on }[1,\infty).
\]
For \(0<R_1<R_2\), define
\[
\chi_{R_1,R_2}(r)
:=
\vartheta\!\left(\frac{r-R_1}{R_2-R_1}\right),
\qquad r\ge0.
\]
Then \(\chi_{R_1,R_2}\) is a smooth buffered cutoff with the required support and plateau
properties, and \eqref{eq:quant-buffer-derivatives} follows by the chain rule.

For \(0\le r<q\le R_2\), the mean-value theorem gives
\[
|\delta_{R_1,R_2}(r,q)|
\le
\|\chi_{R_1,R_2}'\|_{L^\infty(\R_+)},
\]
so the first term in \eqref{eq:quant-buffer-delta} follows from
\eqref{eq:quant-buffer-derivatives}. Moreover,
\[
\partial_r \delta_{R_1,R_2}(r,q)
=
\frac{\chi_{R_1,R_2}(q)-\chi_{R_1,R_2}(r)-\chi_{R_1,R_2}'(r)(q-r)}{(q-r)^2},
\]
and the analogous formula holds for \(\partial_q\delta_{R_1,R_2}\). Taylor's theorem therefore
implies
\[
|\partial_r\delta_{R_1,R_2}(r,q)|
+
|\partial_q\delta_{R_1,R_2}(r,q)|
\le
\frac12\|\chi_{R_1,R_2}''\|_{L^\infty(\R_+)},
\]
which extends continuously to the diagonal. Combining this with
\eqref{eq:quant-buffer-derivatives} proves \eqref{eq:quant-buffer-delta}.
\end{proof}

\begin{proposition}[Quantitative commutator control on a fixed window]
\label{prop:quant-window-comm}
Assume Assumption~\ref{ass:model}, fix \(t\in(0,T)\), \(\alpha\in(0,1)\), and suppose
\[
0<H_{XY}<H_Y<\tfrac12.
\]
Then for every \(R_2<\infty\) there exists \(C_{t,R_2}<\infty\) such that the following holds.
Let \(\chi\) be a smooth buffered cutoff with support in \([0,R_2]\), define
\[
\delta_{\chi}(r,q):=\frac{\chi(q)-\chi(r)}{q-r},
\qquad
\Lambda_{\chi,R_2}:=\|\delta_{\chi}\|_{C^1(\Delta_{R_2})},
\]
and choose \(h>0\) sufficiently small that
\[
hR_2<t,
\qquad
R_2<\varepsilon_{h,\alpha}^{-1}.
\]
Then for every \(\Theta\in E_{R_2}=L^2([0,R_2])\),
\begin{equation}\label{eq:quant-window-comm}
\left\|
\mathfrak Q_{Y,t,h,\alpha,\chi}^{\mathrm{com}}\Theta
\right\|_{E_{R_2}}
+
\left\|
\mathfrak Q_{XY,t,h,\alpha,\chi}^{\mathrm{com}}\Theta
\right\|_{E_{R_2}}
\le
C_{t,R_2}\Lambda_{\chi,R_2}\|\Theta\|_{E_{R_2}}.
\end{equation}
\end{proposition}

\begin{proof}
Let \(\Theta\in E_{R_2}\). Since \(\Theta\) is supported in \([0,R_2]\), the formulas in the proof
of Proposition~\ref{prop:micro-smooth-commutator} reduce to integrals over the fixed triangle
\(\Delta_{R_2}\).

On the target side, the kernel \(F_{Y,t,h,\chi}\) appearing there depends linearly on
\(\delta_{\chi}\) and its first derivatives, multiplied by bounded \(C^1\) coefficients coming from
\(L_Y(t-hr,t-hq)\). Hence
\[
\|F_{Y,t,h,\chi}\|_{C^1(\Delta_{R_2})}\le C_{t,R_2}\Lambda_{\chi,R_2},
\]
and therefore the kernel \(K_{Y,t,h,\chi}^{\mathrm{com}}\) from
\eqref{eq:micro-smooth-com-conj-Y} satisfies
\[
\|K_{Y,t,h,\chi}^{\mathrm{com}}\|_{L^\infty(\Delta_{R_2})}
\le
C_{t,R_2}\Lambda_{\chi,R_2}.
\]
Since \(\Delta_{R_2}\) has finite measure, the corresponding integral operator is bounded from
\(E_{R_2}\) to itself with norm at most \(C_{t,R_2}\Lambda_{\chi,R_2}\).

On the source side, the proof of Proposition~\ref{prop:micro-smooth-commutator} gives a conjugated
kernel of the form
\[
(q-r)^{-\beta}F_{XY,t,h,\chi}(r,q)
+
(q-r)^{1-\beta}\partial_1F_{XY,t,h,\chi}(r,q),
\qquad (r,q)\in\Delta_{R_2},
\]
with
\[
\|F_{XY,t,h,\chi}\|_{C^1(\Delta_{R_2})}
\le
C_{t,R_2}\Lambda_{\chi,R_2}.
\]
Because \(\beta<\tfrac12\), the kernels
\[
(q-r)^{-\beta}\one_{\Delta_{R_2}}(r,q),
\qquad
(q-r)^{1-\beta}\one_{\Delta_{R_2}}(r,q)
\]
belong to \(L^2(\Delta_{R_2})\). Hence the source-side conjugated commutator operator is
Hilbert--Schmidt on \(E_{R_2}\), with norm bounded by \(C_{t,R_2}\Lambda_{\chi,R_2}\). Combining
the two bounds proves \eqref{eq:quant-window-comm}.
\end{proof}

\begin{theorem}[Quantitative smooth localization of an exact fixed-window pair]
\label{thm:quant-window-buffer}
Assume the hypotheses of Theorem~\ref{thm:micro-window-regularity}. Then for every \(R_2<\infty\)
there exists \(C_{t,R_2}<\infty\) such that the following holds whenever \(h>0\) is sufficiently
small and
\[
0<R_1<R_2<\varepsilon_{h,\alpha}^{-1},\qquad hR_2<t.
\]
Let \(\chi\) be a smooth buffered cutoff with support in \([0,R_2]\), and let
\[
(\Phi,\Psi)\in\Gamma_{t,h,\alpha,R_2}^{\mathrm{ex}}.
\]
Define
\[
\delta_{\chi}(r,q):=\frac{\chi(q)-\chi(r)}{q-r},
\qquad
\Lambda_{\chi,R_2}:=\|\delta_{\chi}\|_{C^1(\Delta_{R_2})}.
\]
Then:
\begin{enumerate}[label=(\roman*)]
\item \(\chi\Phi\in H^\beta((0,R_2))\), and the localized pair satisfies
\begin{equation}\label{eq:quant-window-defect}
\mathfrak F_{t,h,\alpha,R_2}^{\mathrm{ex}}(\chi\Phi,\chi\Psi)
=
\mathfrak Q_{Y,t,h,\alpha,\chi}^{\mathrm{com}}\Psi
-
\mathfrak Q_{XY,t,h,\alpha,\chi}^{\mathrm{com}}\Phi
\qquad\text{in }E_{R_2};
\end{equation}
\item the corrected output
\begin{equation}\label{eq:quant-window-corrected}
\Psi_{\chi}^{\mathrm{ex}}
:=
\chi\Psi
-
\bigl(\mathfrak B_{Y,t,h,\alpha,R_2}^{\mathrm{mic}}\bigr)^{-1}
\mathfrak F_{t,h,\alpha,R_2}^{\mathrm{ex}}(\chi\Phi,\chi\Psi)
\end{equation}
satisfies
\[
(\chi\Phi,\Psi_{\chi}^{\mathrm{ex}})
\in
\Gamma_{t,h,\alpha,R_2}^{\mathrm{ex}};
\]
\item one has the quantitative correction bound
\begin{equation}\label{eq:quant-window-correction}
\left\|
\Psi_{\chi}^{\mathrm{ex}}-\chi\Psi
\right\|_{E_{R_2}}
\le
C_{t,R_2}\Lambda_{\chi,R_2}
\bigl(
\|\Phi\|_{E_{R_2}}+\|\Psi\|_{E_{R_2}}
\bigr);
\end{equation}
\item consequently,
\begin{equation}\label{eq:quant-window-graph}
\|\chi\Phi\|_{\mathfrak G_{t,h,\alpha,R_2}^{\mathrm{ex}}}
\le
C_{t,R_2}
\bigl(1+\Lambda_{\chi,R_2}\bigr)
\|\Phi\|_{\mathfrak G_{t,h,\alpha,R_2}^{\mathrm{ex}}}.
\end{equation}
\end{enumerate}
If \(\chi=\chi_{R_1,R_2}\) is chosen as in Lemma~\ref{lem:quant-buffer}, then
\begin{equation}\label{eq:quant-window-correction-width}
\left\|
\Psi_{\chi_{R_1,R_2}}^{\mathrm{ex}}-\chi_{R_1,R_2}\Psi
\right\|_{E_{R_2}}
\le
C_{t,R_2}
\left(
\frac{1}{R_2-R_1}
+
\frac{1}{(R_2-R_1)^2}
\right)
\|\Phi\|_{\mathfrak G_{t,h,\alpha,R_2}^{\mathrm{ex}}}.
\end{equation}
\end{theorem}

\begin{proof}
By Proposition~\ref{prop:micro-window-identification},
\[
\Phi\in\mathfrak G_{t,h,\alpha,R_2}^{\mathrm{mic}},
\qquad
\Psi=\mathfrak T_{t,h,\alpha}^{\mathrm{mic}}\Phi,
\]
with both \(\Phi\) and \(\Psi\) supported in \([0,R_2]\). Repeating the proof of
Theorem~\ref{thm:micro-smooth-buffer}(i), but now with no outer tail contribution because
\(\Phi,\Psi\in E_{R_2}\), yields \(\chi\Phi\in H^\beta((0,R_2))\) together with the defect identity
\eqref{eq:quant-window-defect}.

Part (ii) is exactly Theorem~\ref{thm:micro-window-correction} applied to the pair
\((\chi\Phi,\chi\Psi)\). For part (iii), combine the correction estimate
\eqref{eq:window-correction-bound} with \eqref{eq:quant-window-defect} and then invoke
Proposition~\ref{prop:quant-window-comm}:
\begin{align*}
\left\|
\Psi_{\chi}^{\mathrm{ex}}-\chi\Psi
\right\|_{E_{R_2}}
&\le
C_{t,R_2}
\left\|
\mathfrak F_{t,h,\alpha,R_2}^{\mathrm{ex}}(\chi\Phi,\chi\Psi)
\right\|_{E_{R_2}}\\
&\le
C_{t,R_2}\Lambda_{\chi,R_2}
\bigl(
\|\Phi\|_{E_{R_2}}+\|\Psi\|_{E_{R_2}}
\bigr).
\end{align*}
This proves \eqref{eq:quant-window-correction}.

Finally,
\[
\mathfrak T_{t,h,\alpha,R_2}^{\mathrm{ex}}(\chi\Phi)=\Psi_{\chi}^{\mathrm{ex}}
\]
by part (ii), so using \(|\chi|\le1\) and \eqref{eq:quant-window-correction} gives
\begin{align*}
\|\chi\Phi\|_{\mathfrak G_{t,h,\alpha,R_2}^{\mathrm{ex}}}
&\le
\|\Phi\|_{E_{R_2}}
+
\frac{\eta}{\nu_Y}
\left(
\|\chi\Psi\|_{E_{R_2}}
+
\|\Psi_{\chi}^{\mathrm{ex}}-\chi\Psi\|_{E_{R_2}}
\right)\\
&\le
C_{t,R_2}
\bigl(1+\Lambda_{\chi,R_2}\bigr)
\left(
\|\Phi\|_{E_{R_2}}
+
\frac{\eta}{\nu_Y}\|\Psi\|_{E_{R_2}}
\right),
\end{align*}
which is exactly \eqref{eq:quant-window-graph}.

If \(\chi=\chi_{R_1,R_2}\) is the cutoff from Lemma~\ref{lem:quant-buffer}, then
\eqref{eq:quant-window-correction-width} follows from
\eqref{eq:quant-window-correction}, \eqref{eq:quant-buffer-delta}, and the elementary bound
\[
\|\Phi\|_{E_{R_2}}+\|\Psi\|_{E_{R_2}}
\le
C_{\eta,\nu_Y}
\|\Phi\|_{\mathfrak G_{t,h,\alpha,R_2}^{\mathrm{ex}}}.
\qedhere
\]
\end{proof}

\section{Fixed-window dual collapse and abstract global reduction}
\label{sec:micro-dual}

The local estimates reduce global collapse to a localization passage from the global micro graph to
exact fixed windows.  Two consequences have already been proved:
\begin{enumerate}[label=(\roman*)]
\item on every fixed window, the exact micro dual norm collapses to \(0\) as \(h\downarrow0\);
\item any global localization principle reducing graph-unit micro sequences to fixed windows forces
the full observable micro dual norm to collapse as well.
\end{enumerate}

\begin{definition}[Global and fixed-window exact micro functionals]
\label{def:micro-dual}
Assume Assumptions~\ref{ass:model} and~\ref{ass:obs-source}, let \(\eta>0\), fix \(t\in(0,T)\),
\(\alpha\in(0,1)\), and suppose
\[
0<H_{XY}<H_Y<\tfrac12.
\]
For
\[
\Phi\in\mathcal D(\mathfrak T_{t,h,\alpha}^{\mathrm{mic}}),
\]
define the global micro graph norm
\begin{equation}\label{eq:micro-global-graph-norm}
\|\Phi\|_{\mathfrak G_{t,h,\alpha}^{\mathrm{mic}}}^2
:=
\|\Phi\|_{L^2(\R_+)}^2
+
\frac{\eta^2}{\nu_Y^2}
\left\|
\mathfrak T_{t,h,\alpha}^{\mathrm{mic}}\Phi
\right\|_{L^2(\R_+)}^2,
\end{equation}
and the global exact micro functional
\begin{equation}\label{eq:micro-global-functional}
\Lambda_{t,h,\alpha}^{\mathrm{mic}}(\Phi)
:=
\left\langle
\mathfrak g_{XY,t,h,\alpha}^{\mathrm{mic}},\Phi
\right\rangle_{L^2(\R_+)}
-
\left\langle
\mathfrak g_{Y,t,h,\alpha}^{\mathrm{mic}},
\mathfrak T_{t,h,\alpha}^{\mathrm{mic}}\Phi
\right\rangle_{L^2(\R_+)}.
\end{equation}
For \(R>0\) with \(hR<t\) and \(R<\varepsilon_{h,\alpha}^{-1}\), and
\[
\Phi\in\mathcal D(\mathfrak T_{t,h,\alpha,R}^{\mathrm{ex}}),
\]
define the fixed-window exact functional
\begin{equation}\label{eq:micro-window-functional-ex}
\Lambda_{t,h,\alpha,R}^{\mathrm{ex}}(\Phi)
:=
\left\langle
\mathfrak g_{XY,t,h,\alpha}^{\mathrm{mic}},\Phi
\right\rangle_{L^2(\R_+)}
-
\left\langle
\mathfrak g_{Y,t,h,\alpha}^{\mathrm{mic}},
\mathfrak T_{t,h,\alpha,R}^{\mathrm{ex}}\Phi
\right\rangle_{L^2(\R_+)}.
\end{equation}
Finally, define the corresponding dual norms
\begin{align}
\mathfrak D_{t,h,\alpha}^{\mathrm{mic}}
&:=
\sup\left\{
\left|
\Lambda_{t,h,\alpha}^{\mathrm{mic}}(\Phi)
\right|:
\Phi\in\mathcal D(\mathfrak T_{t,h,\alpha}^{\mathrm{mic}}),
\ \|\Phi\|_{\mathfrak G_{t,h,\alpha}^{\mathrm{mic}}}\le1
\right\},
\label{eq:micro-global-dual}\\
\mathfrak D_{t,h,\alpha,R}^{\mathrm{ex}}
&:=
\sup\left\{
\left|
\Lambda_{t,h,\alpha,R}^{\mathrm{ex}}(\Phi)
\right|:
\Phi\in\mathcal D(\mathfrak T_{t,h,\alpha,R}^{\mathrm{ex}}),
\ \|\Phi\|_{\mathfrak G_{t,h,\alpha,R}^{\mathrm{ex}}}\le1
\right\}.
\label{eq:micro-window-dual}
\end{align}
\end{definition}

\begin{proposition}[Continuity of the global and fixed-window exact functionals]
\label{prop:micro-dual-cont}
Assume the hypotheses of Definition~\ref{def:micro-dual}. Then:
\begin{enumerate}[label=(\roman*)]
\item \(\Lambda_{t,h,\alpha}^{\mathrm{mic}}\) is continuous on
\(\mathfrak G_{t,h,\alpha}^{\mathrm{mic}}\);
\item if \(R>0\) with \(hR<t\) and \(R<\varepsilon_{h,\alpha}^{-1}\), then
\(\Lambda_{t,h,\alpha,R}^{\mathrm{ex}}\) is continuous on
\(\mathfrak G_{t,h,\alpha,R}^{\mathrm{ex}}\).
\end{enumerate}
\end{proposition}

\begin{proof}
Since \(\mathfrak g_{XY,t,h,\alpha}^{\mathrm{mic}},\mathfrak g_{Y,t,h,\alpha}^{\mathrm{mic}}\in
L^2(\R_+)\), Cauchy--Schwarz gives
\[
\left|
\Lambda_{t,h,\alpha}^{\mathrm{mic}}(\Phi)
\right|
\le
\|\mathfrak g_{XY,t,h,\alpha}^{\mathrm{mic}}\|_{L^2}\|\Phi\|_{L^2}
+
\|\mathfrak g_{Y,t,h,\alpha}^{\mathrm{mic}}\|_{L^2}
\left\|
\mathfrak T_{t,h,\alpha}^{\mathrm{mic}}\Phi
\right\|_{L^2},
\]
which proves item~(i) via \eqref{eq:micro-global-graph-norm}. The fixed-window statement is
identical, using \eqref{eq:micro-window-functional-ex} and
Corollary~\ref{cor:micro-window-exact-norm}.
\end{proof}

\begin{theorem}[Fixed-window exact dual norms collapse]
\label{thm:micro-window-dual-collapse}
Assume Assumptions~\ref{ass:model} and~\ref{ass:obs-source}, let \(\eta>0\), fix \(t\in(0,T)\),
\(\alpha\in(0,1)\), \(R>0\), and suppose
\[
0<H_{XY}<H_Y<\tfrac12.
\]
Then
\begin{equation}\label{eq:micro-window-dual-collapse}
\mathfrak D_{t,h,\alpha,R}^{\mathrm{ex}}
\longrightarrow0
\qquad\text{as }h\downarrow0
\end{equation}
along every sequence with \(hR<t\) and \(R<\varepsilon_{h,\alpha}^{-1}\).
\end{theorem}

\begin{proof}
Suppose the claim fails. Then there exist \(\epsilon_0>0\), a sequence \(h_n\downarrow0\) with
\[
h_nR<t,
\qquad
R<\varepsilon_{h_n,\alpha}^{-1},
\]
and elements
\[
\Phi_n\in\mathcal D(\mathfrak T_{t,h_n,\alpha,R}^{\mathrm{ex}})
\]
such that
\[
\|\Phi_n\|_{\mathfrak G_{t,h_n,\alpha,R}^{\mathrm{ex}}}\le1,
\qquad
\left|
\Lambda_{t,h_n,\alpha,R}^{\mathrm{ex}}(\Phi_n)
\right|
\ge
\epsilon_0.
\]
By Corollary~\ref{cor:micro-window-exact-norm} and Proposition~\ref{prop:micro-window-identification},
we may regard \(\Phi_n\) as elements of \(\mathfrak G_{t,h_n,\alpha,R}^{\mathrm{mic}}\) with the same
norm bound. Proposition~\ref{prop:micro-window-compact} therefore yields a subsequence, still
denoted \((\Phi_n)\), and some
\[
\Phi\in \mathcal D(\mathfrak T_{t,R}^{\mathrm{hl}})
\]
such that
\begin{align}
\Phi_n &\rightharpoonup \Phi
\qquad\text{weakly in }E_R,
\label{eq:dual-collapse-weak-1}\\
\mathfrak T_{t,h_n,\alpha,R}^{\mathrm{ex}}\Phi_n
&\rightharpoonup
\mathfrak T_{t,R}^{\mathrm{hl}}\Phi
\qquad\text{weakly in }E_R.
\label{eq:dual-collapse-weak-2}
\end{align}
By Theorem~\ref{thm:micro-profile},
\[
\mathfrak g_{i,t,h_n,\alpha}^{\mathrm{mic}}
\longrightarrow
L_i(t,t)\psi_{H_i}
\qquad\text{strongly in }L^2(\R_+)
\]
for \(i\in\{Y,XY\}\). Since \(\Phi_n\) and
\(\mathfrak T_{t,h_n,\alpha,R}^{\mathrm{ex}}\Phi_n\) are supported in \([0,R]\), the pairings in
\eqref{eq:micro-window-functional-ex} may be taken over \(E_R\). Combining the strong profile
convergence with \eqref{eq:dual-collapse-weak-1}--\eqref{eq:dual-collapse-weak-2} gives
\[
\Lambda_{t,h_n,\alpha,R}^{\mathrm{ex}}(\Phi_n)
\longrightarrow
L_{XY}(t,t)\langle \psi_{H_{XY}},\Phi\rangle_{L^2(\R_+)}
-
L_Y(t,t)
\left\langle
\psi_{H_Y},
\mathfrak T_{t,R}^{\mathrm{hl}}\Phi
\right\rangle_{L^2(\R_+)}.
\]
The right-hand side is exactly \(\Lambda_{t,R}^{\mathrm{hl}}(\Phi)\), which vanishes by
Corollary~\ref{cor:micro-window-model-functional}. Hence
\[
\Lambda_{t,h_n,\alpha,R}^{\mathrm{ex}}(\Phi_n)\to0,
\]
contradicting \(|\Lambda_{t,h_n,\alpha,R}^{\mathrm{ex}}(\Phi_n)|\ge\epsilon_0\).
\end{proof}

\begin{theorem}[Abstract global reduction from exact fixed-window localization]
\label{thm:abstract-global-reduction}
Assume the hypotheses of Definition~\ref{def:micro-dual}. Suppose the following localization
property holds:

\smallskip
\noindent
for every sequence \(h_n\downarrow0\), every sequence
\[
\Phi_n\in\mathcal D(\mathfrak T_{t,h_n,\alpha}^{\mathrm{mic}})
\qquad\text{with}\qquad
\|\Phi_n\|_{\mathfrak G_{t,h_n,\alpha}^{\mathrm{mic}}}\le1,
\]
and every \(\varepsilon>0\), there exist \(R<\infty\), \(n_0\in\mathbb N\), and elements
\[
\widetilde\Phi_n\in \mathcal D(\mathfrak T_{t,h_n,\alpha,R}^{\mathrm{ex}})
\qquad (n\ge n_0)
\]
such that
\begin{align}
\|\widetilde\Phi_n\|_{\mathfrak G_{t,h_n,\alpha,R}^{\mathrm{ex}}}
&\le
1+\varepsilon,
\label{eq:abstract-localize-norm}\\
\left|
\Lambda_{t,h_n,\alpha}^{\mathrm{mic}}(\Phi_n)
-
\Lambda_{t,h_n,\alpha,R}^{\mathrm{ex}}(\widetilde\Phi_n)
\right|
&\le
\varepsilon
\qquad\text{for every }n\ge n_0.
\label{eq:abstract-localize-functional}
\end{align}

\smallskip
\noindent
Then
\begin{equation}\label{eq:abstract-global-collapse}
\mathfrak D_{t,h,\alpha}^{\mathrm{mic}}
\longrightarrow0
\qquad\text{as }h\downarrow0.
\end{equation}
\end{theorem}

\begin{proof}
Suppose \eqref{eq:abstract-global-collapse} fails. Then there exist \(\epsilon_0>0\) and
\(h_n\downarrow0\) such that
\[
\mathfrak D_{t,h_n,\alpha}^{\mathrm{mic}}\ge 4\epsilon_0
\qquad\text{for every }n.
\]
By Definition~\ref{def:micro-dual}, choose
\[
\Phi_n\in\mathcal D(\mathfrak T_{t,h_n,\alpha}^{\mathrm{mic}})
\]
with
\[
\|\Phi_n\|_{\mathfrak G_{t,h_n,\alpha}^{\mathrm{mic}}}\le1,
\qquad
\left|
\Lambda_{t,h_n,\alpha}^{\mathrm{mic}}(\Phi_n)
\right|
\ge
2\epsilon_0.
\]
Apply the localization property with \(\varepsilon=\epsilon_0\). Then there exist \(R<\infty\),
\(n_0\), and
\[
\widetilde\Phi_n\in \mathcal D(\mathfrak T_{t,h_n,\alpha,R}^{\mathrm{ex}})
\qquad(n\ge n_0)
\]
such that
\[
\|\widetilde\Phi_n\|_{\mathfrak G_{t,h_n,\alpha,R}^{\mathrm{ex}}}\le1+\epsilon_0,
\]
and
\[
\left|
\Lambda_{t,h_n,\alpha}^{\mathrm{mic}}(\Phi_n)
-
\Lambda_{t,h_n,\alpha,R}^{\mathrm{ex}}(\widetilde\Phi_n)
\right|
\le
\epsilon_0.
\]
Hence, for every \(n\ge n_0\),
\[
\left|
\Lambda_{t,h_n,\alpha,R}^{\mathrm{ex}}(\widetilde\Phi_n)
\right|
\ge
\epsilon_0.
\]
After normalizing by \(1+\epsilon_0\), this implies
\[
\mathfrak D_{t,h_n,\alpha,R}^{\mathrm{ex}}
\ge
\frac{\epsilon_0}{1+\epsilon_0}
\qquad\text{for every }n\ge n_0,
\]
which contradicts Theorem~\ref{thm:micro-window-dual-collapse}.
\end{proof}

\section{Uniform profile tails and buffer selection}
\label{sec:good-buffers}

The unresolved global step is localization of graph-unit micro sequences on one common finite
window without losing the functional.  Two preparatory facts follow from the existing theory:
\begin{enumerate}[label=(\roman*)]
\item the micro profiles are uniformly tight in \(L^2(\R_+)\) as \(h\downarrow0\);
\item every graph-unit sequence admits, after subselection, one common annulus on which both the
input and output masses are arbitrarily small.
\end{enumerate}

\begin{proposition}[Uniform tail tightness of the micro profiles]
\label{prop:micro-profile-tails}
Assume Assumption~\ref{ass:model}, fix \(t\in(0,T)\) and \(\alpha\in(0,1)\), and let
\[
i\in\{Y,XY\}.
\]
For every \(\varepsilon>0\), there exist \(R_\varepsilon<\infty\) and \(h_\varepsilon>0\) such that
\begin{equation}\label{eq:micro-profile-tail}
\left\|
\one_{(R_\varepsilon,\infty)}
\mathfrak g_{i,t,h,\alpha}^{\mathrm{mic}}
\right\|_{L^2(\R_+)}
\le
\varepsilon
\qquad\text{for every }0<h\le h_\varepsilon.
\end{equation}
Moreover, \(R_\varepsilon\) and \(h_\varepsilon\) may be chosen simultaneously for \(i=Y\) and
\(i=XY\).
\end{proposition}

\begin{proof}
Fix \(\varepsilon>0\). By Theorem~\ref{thm:micro-profile},
\[
\mathfrak g_{i,t,h,\alpha}^{\mathrm{mic}}
\longrightarrow
L_i(t,t)\psi_{H_i}
\qquad\text{in }L^2(\R_+)
\]
as \(h\downarrow0\). Since \(L_i(t,t)\psi_{H_i}\in L^2(\R_+)\), there exists \(R_\varepsilon<\infty\)
such that
\[
\left\|
\one_{(R_\varepsilon,\infty)}L_i(t,t)\psi_{H_i}
\right\|_{L^2(\R_+)}
\le
\frac{\varepsilon}{2}.
\]
Then choose \(h_\varepsilon>0\) such that
\[
\left\|
\mathfrak g_{i,t,h,\alpha}^{\mathrm{mic}}-L_i(t,t)\psi_{H_i}
\right\|_{L^2(\R_+)}
\le
\frac{\varepsilon}{2}
\qquad\text{for every }0<h\le h_\varepsilon.
\]
Therefore, for \(0<h\le h_\varepsilon\),
\begin{align*}
\left\|
\one_{(R_\varepsilon,\infty)}
\mathfrak g_{i,t,h,\alpha}^{\mathrm{mic}}
\right\|_{L^2(\R_+)}
&\le
\left\|
\mathfrak g_{i,t,h,\alpha}^{\mathrm{mic}}-L_i(t,t)\psi_{H_i}
\right\|_{L^2(\R_+)}\\
&\quad+
\left\|
\one_{(R_\varepsilon,\infty)}L_i(t,t)\psi_{H_i}
\right\|_{L^2(\R_+)}
\le
\varepsilon.
\end{align*}
This proves \eqref{eq:micro-profile-tail}. The simultaneous statement for \(i=Y\) and \(i=XY\)
follows by taking the larger of the two radii and the smaller of the two \(h\)-thresholds.
\end{proof}

\begin{proposition}[Annular low-mass buffer selection for graph-unit vectors]
\label{prop:good-buffer}
Assume the hypotheses of Definition~\ref{def:micro-dual}, let
\[
\Phi\in\mathcal D(\mathfrak T_{t,h,\alpha}^{\mathrm{mic}}),
\qquad
\Psi:=\mathfrak T_{t,h,\alpha}^{\mathrm{mic}}\Phi,
\]
and suppose
\[
\|\Phi\|_{\mathfrak G_{t,h,\alpha}^{\mathrm{mic}}}\le1.
\]
Fix \(A\ge0\), \(\ell>0\), and \(M\in\mathbb N\), and define the disjoint annuli
\[
I_j:=\bigl[A+(j-1)\ell,\ A+j\ell\bigr],
\qquad j=1,\dots,M.
\]
Then there exists \(j_*\in\{1,\dots,M\}\) such that
\begin{equation}\label{eq:good-buffer}
\|\Phi\|_{L^2(I_{j_*})}^2
+
\frac{\eta^2}{\nu_Y^2}
\|\Psi\|_{L^2(I_{j_*})}^2
\le
\frac{1}{M}.
\end{equation}
\end{proposition}

\begin{proof}
Summing over \(j=1,\dots,M\) and using the disjointness of the \(I_j\) gives
\[
\sum_{j=1}^M
\left(
\|\Phi\|_{L^2(I_j)}^2
+
\frac{\eta^2}{\nu_Y^2}
\|\Psi\|_{L^2(I_j)}^2
\right)
\le
\|\Phi\|_{L^2(\R_+)}^2
+
\frac{\eta^2}{\nu_Y^2}
\|\Psi\|_{L^2(\R_+)}^2
=
\|\Phi\|_{\mathfrak G_{t,h,\alpha}^{\mathrm{mic}}}^2
\le
1.
\]
Hence at least one summand is at most \(1/M\), which is exactly \eqref{eq:good-buffer}.
\end{proof}

\begin{corollary}[Common low-mass buffers along graph-unit sequences]
\label{cor:good-buffer-sequence}
Assume the hypotheses of Definition~\ref{def:micro-dual}. Let \(h_n\downarrow0\), and for each
\(n\) let
\[
\Phi_n\in\mathcal D(\mathfrak T_{t,h_n,\alpha}^{\mathrm{mic}}),
\qquad
\Psi_n:=\mathfrak T_{t,h_n,\alpha}^{\mathrm{mic}}\Phi_n,
\qquad
\|\Phi_n\|_{\mathfrak G_{t,h_n,\alpha}^{\mathrm{mic}}}\le1.
\]
Fix \(A\ge0\), \(\ell>0\), and \(M\in\mathbb N\), and define
\[
I_j:=\bigl[A+(j-1)\ell,\ A+j\ell\bigr],
\qquad j=1,\dots,M.
\]
Then there exist an index \(j_*\in\{1,\dots,M\}\) and a subsequence \((n_k)\) such that
\begin{equation}\label{eq:good-buffer-subseq}
\|\Phi_{n_k}\|_{L^2(I_{j_*})}^2
+
\frac{\eta^2}{\nu_Y^2}
\|\Psi_{n_k}\|_{L^2(I_{j_*})}^2
\le
\frac{1}{M}
\qquad\text{for every }k.
\end{equation}
\end{corollary}

\begin{proof}
For each \(n\), Proposition~\ref{prop:good-buffer} gives an index
\[
j_n\in\{1,\dots,M\}
\]
such that
\[
\|\Phi_n\|_{L^2(I_{j_n})}^2
+
\frac{\eta^2}{\nu_Y^2}
\|\Psi_n\|_{L^2(I_{j_n})}^2
\le
\frac{1}{M}.
\]
Since \(\{1,\dots,M\}\) is finite, one index \(j_*\) occurs infinitely often. Passing to that
subsequence gives \eqref{eq:good-buffer-subseq}.
\end{proof}

\begin{corollary}[Profile-tight common buffers]
\label{cor:profile-tight-buffer}
Assume the hypotheses of Definition~\ref{def:micro-dual}. Let \(h_n\downarrow0\), and let
\[
\Phi_n\in\mathcal D(\mathfrak T_{t,h_n,\alpha}^{\mathrm{mic}}),
\qquad
\Psi_n:=\mathfrak T_{t,h_n,\alpha}^{\mathrm{mic}}\Phi_n,
\qquad
\|\Phi_n\|_{\mathfrak G_{t,h_n,\alpha}^{\mathrm{mic}}}\le1.
\]
Fix \(\varepsilon>0\), \(\ell>0\), and \(M\in\mathbb N\). Then, after passing to a subsequence,
there exist finite radii
\[
0<R_0<R_1<R_2
\]
such that
\[
R_1-R_0\ge \ell,
\qquad
R_2-R_1=\ell,
\]
and:
\begin{enumerate}[label=(\roman*)]
\item for \(i\in\{Y,XY\}\),
\begin{equation}\label{eq:profile-tight-tail}
\left\|
\one_{(R_0,\infty)}
\mathfrak g_{i,t,h_n,\alpha}^{\mathrm{mic}}
\right\|_{L^2(\R_+)}
\le
\varepsilon
\qquad\text{for every }n;
\end{equation}
\item the common transition annulus satisfies
\begin{equation}\label{eq:profile-tight-annulus}
\|\Phi_n\|_{L^2([R_1,R_2])}^2
+
\frac{\eta^2}{\nu_Y^2}
\|\Psi_n\|_{L^2([R_1,R_2])}^2
\le
\frac{1}{M}
\qquad\text{for every }n.
\end{equation}
\end{enumerate}
\end{corollary}

\begin{proof}
By Proposition~\ref{prop:micro-profile-tails}, after discarding finitely many indices there exists
\(R_0<\infty\) such that \eqref{eq:profile-tight-tail} holds for both \(i=Y\) and \(i=XY\) for every
remaining \(n\). Set
\[
A:=R_0+\ell.
\]
Applying Corollary~\ref{cor:good-buffer-sequence} with this \(A\), the same \(\ell\), and the same
\(M\), and then relabeling the resulting subsequence, yields an index \(j_*\in\{1,\dots,M\}\) such
that
\[
\|\Phi_n\|_{L^2(I_{j_*})}^2
+
\frac{\eta^2}{\nu_Y^2}
\|\Psi_n\|_{L^2(I_{j_*})}^2
\le
\frac{1}{M}
\qquad\text{for every }n,
\]
where
\[
I_{j_*}
=
\bigl[A+(j_*-1)\ell,\ A+j_*\ell\bigr].
\]
Now define
\[
R_1:=A+(j_*-1)\ell,
\qquad
R_2:=A+j_*\ell.
\]
Then \(R_2-R_1=\ell\), \(R_1\ge A=R_0+\ell\), and therefore \(R_1-R_0\ge\ell\). The annulus
estimate \eqref{eq:profile-tight-annulus} is exactly the displayed bound above.
\end{proof}

\section{Non-escaping sequences and functional concentration}
\label{sec:nonescaping}

Common low-mass buffers do not by themselves localize every sequence.  If a graph-unit micro
sequence escapes to infinity, the micro functional vanishes because the profiles are uniformly
tight.  The remaining regime is the non-escaping one, where a fixed amount of graph mass stays on
some common finite window.

\begin{proposition}[Escape to infinity forces functional collapse]
\label{prop:escape-functional}
Assume the hypotheses of Definition~\ref{def:micro-dual}. Let \(h_n\downarrow0\), and let
\[
\Phi_n\in\mathcal D(\mathfrak T_{t,h_n,\alpha}^{\mathrm{mic}}),
\qquad
\Psi_n:=\mathfrak T_{t,h_n,\alpha}^{\mathrm{mic}}\Phi_n,
\qquad
\|\Phi_n\|_{\mathfrak G_{t,h_n,\alpha}^{\mathrm{mic}}}\le1.
\]
Assume moreover that for every \(R<\infty\),
\begin{equation}\label{eq:escape-head}
\|\Phi_n\|_{L^2([0,R])}
+
\|\Psi_n\|_{L^2([0,R])}
\longrightarrow0
\qquad\text{as }n\to\infty.
\end{equation}
Then
\begin{equation}\label{eq:escape-functional}
\Lambda_{t,h_n,\alpha}^{\mathrm{mic}}(\Phi_n)\longrightarrow0
\qquad\text{as }n\to\infty.
\end{equation}
\end{proposition}

\begin{proof}
Fix \(\varepsilon>0\). By Proposition~\ref{prop:micro-profile-tails}, there exist \(R_\varepsilon\)
and \(n_\varepsilon\) such that for every \(n\ge n_\varepsilon\),
\begin{equation}\label{eq:escape-profile-tail}
\left\|
\one_{(R_\varepsilon,\infty)}
\mathfrak g_{i,t,h_n,\alpha}^{\mathrm{mic}}
\right\|_{L^2(\R_+)}
\le
\varepsilon
\qquad\text{for }i\in\{Y,XY\}.
\end{equation}
By \eqref{eq:escape-head}, after increasing \(n_\varepsilon\) if necessary, we may also assume that
for every \(n\ge n_\varepsilon\),
\begin{equation}\label{eq:escape-small-head}
\|\Phi_n\|_{L^2([0,R_\varepsilon])}
+
\|\Psi_n\|_{L^2([0,R_\varepsilon])}
\le
\varepsilon.
\end{equation}
Using \eqref{eq:micro-global-functional}, split both pairings at \(R_\varepsilon\). For \(n\ge
n_\varepsilon\),
\begin{align*}
\left|
\Lambda_{t,h_n,\alpha}^{\mathrm{mic}}(\Phi_n)
\right|
&\le
\left\|
\one_{[0,R_\varepsilon]}
\mathfrak g_{XY,t,h_n,\alpha}^{\mathrm{mic}}
\right\|_{L^2}
\|\Phi_n\|_{L^2([0,R_\varepsilon])}\\
&\quad+
\left\|
\one_{(R_\varepsilon,\infty)}
\mathfrak g_{XY,t,h_n,\alpha}^{\mathrm{mic}}
\right\|_{L^2}
\|\Phi_n\|_{L^2(\R_+)}\\
&\quad+
\left\|
\one_{[0,R_\varepsilon]}
\mathfrak g_{Y,t,h_n,\alpha}^{\mathrm{mic}}
\right\|_{L^2}
\|\Psi_n\|_{L^2([0,R_\varepsilon])}\\
&\quad+
\left\|
\one_{(R_\varepsilon,\infty)}
\mathfrak g_{Y,t,h_n,\alpha}^{\mathrm{mic}}
\right\|_{L^2}
\|\Psi_n\|_{L^2(\R_+)}.
\end{align*}
The full \(L^2\) norms of the micro profiles are uniformly bounded for large \(n\) by
Theorem~\ref{thm:micro-profile}. Also,
\[
\|\Phi_n\|_{L^2(\R_+)}\le1,
\qquad
\|\Psi_n\|_{L^2(\R_+)}\le \frac{\nu_Y}{\eta},
\]
because \(\|\Phi_n\|_{\mathfrak G_{t,h_n,\alpha}^{\mathrm{mic}}}\le1\). Combining these bounds with
\eqref{eq:escape-profile-tail} and \eqref{eq:escape-small-head} yields
\[
\limsup_{n\to\infty}
\left|
\Lambda_{t,h_n,\alpha}^{\mathrm{mic}}(\Phi_n)
\right|
\le
C_{t,\eta,\nu_Y}\varepsilon.
\]
Since \(\varepsilon>0\) was arbitrary, \eqref{eq:escape-functional} follows.
\end{proof}

\begin{corollary}[Non-escaping alternative for nonvanishing functional sequences]
\label{cor:nonescaping}
Assume the hypotheses of Definition~\ref{def:micro-dual}. Let \(h_n\downarrow0\), and let
\[
\Phi_n\in\mathcal D(\mathfrak T_{t,h_n,\alpha}^{\mathrm{mic}}),
\qquad
\Psi_n:=\mathfrak T_{t,h_n,\alpha}^{\mathrm{mic}}\Phi_n,
\qquad
\|\Phi_n\|_{\mathfrak G_{t,h_n,\alpha}^{\mathrm{mic}}}\le1.
\]
Suppose that
\begin{equation}\label{eq:nonescaping-assumption}
\limsup_{n\to\infty}
\left|
\Lambda_{t,h_n,\alpha}^{\mathrm{mic}}(\Phi_n)
\right|
>0.
\end{equation}
Then there exist \(R_*<\infty\), \(\delta_*>0\), and a subsequence, not relabeled, such that
\begin{equation}\label{eq:nonescaping-head}
\|\Phi_n\|_{L^2([0,R_*])}
+
\|\Psi_n\|_{L^2([0,R_*])}
\ge
\delta_*
\qquad\text{for every }n.
\end{equation}
\end{corollary}

\begin{proof}
If no such \(R_*\) and \(\delta_*\) existed, then after passing to a subsequence one would have
\eqref{eq:escape-head} for every \(R<\infty\). Proposition~\ref{prop:escape-functional} would then
imply \(\Lambda_{t,h_n,\alpha}^{\mathrm{mic}}(\Phi_n)\to0\), contradicting
\eqref{eq:nonescaping-assumption}.
\end{proof}

\begin{corollary}[Structured reduction of the global localization problem]
\label{cor:structured-reduction}
Assume the hypotheses of Definition~\ref{def:micro-dual}. To verify the localization hypothesis in
Theorem~\ref{thm:abstract-global-reduction}, it is enough to treat graph-unit sequences
\[
\Phi_n\in\mathcal D(\mathfrak T_{t,h_n,\alpha}^{\mathrm{mic}}),
\qquad
\Psi_n:=\mathfrak T_{t,h_n,\alpha}^{\mathrm{mic}}\Phi_n,
\qquad
\|\Phi_n\|_{\mathfrak G_{t,h_n,\alpha}^{\mathrm{mic}}}\le1,
\]
for which, after passage to a subsequence, there exist:
\begin{enumerate}[label=(\roman*)]
\item a finite radius \(R_*<\infty\) and a constant \(\delta_*>0\) such that
\[
\|\Phi_n\|_{L^2([0,R_*])}
+
\|\Psi_n\|_{L^2([0,R_*])}
\ge
\delta_*
\qquad\text{for every }n;
\]
\item for every prescribed \(\varepsilon>0\), \(\ell>0\), and \(M\in\mathbb N\), finite radii
\[
0<R_0<R_1<R_2
\]
such that
\[
R_1-R_0\ge\ell,
\qquad
R_2-R_1=\ell,
\]
and the profile-tail and good-buffer bounds
\eqref{eq:profile-tight-tail} and \eqref{eq:profile-tight-annulus} hold.
\end{enumerate}
\end{corollary}

\begin{proof}
If \(\Lambda_{t,h_n,\alpha}^{\mathrm{mic}}(\Phi_n)\to0\) already, then the localization hypothesis in
Theorem~\ref{thm:abstract-global-reduction} is already satisfied: for any \(\varepsilon>0\), taking
\(\widetilde\Phi_n:=0\) on any fixed window gives
\[
\|\widetilde\Phi_n\|_{\mathfrak G_{t,h_n,\alpha,R}^{\mathrm{ex}}}=0\le1+\varepsilon
\]
and, for large \(n\),
\[
\left|
\Lambda_{t,h_n,\alpha}^{\mathrm{mic}}(\Phi_n)
-
\Lambda_{t,h_n,\alpha,R}^{\mathrm{ex}}(\widetilde\Phi_n)
\right|
=
\left|
\Lambda_{t,h_n,\alpha}^{\mathrm{mic}}(\Phi_n)
\right|
\le
\varepsilon.
\]
Thus only sequences satisfying \eqref{eq:nonescaping-assumption} need to be considered. For such
sequences, Corollary~\ref{cor:nonescaping} gives the non-escaping head bound, and
Corollary~\ref{cor:profile-tight-buffer} then gives the common profile-tight buffers after
passing to a further subsequence.
\end{proof}

\section{Localization from full-window defect control}
\label{sec:defect-localization}

The final global reduction has a definite form.  A smoothly truncated global exact micro pair must
have a small \emph{full-window} conjugated defect on one common finite window.  The correction
theorem and the profile-tightness estimates then upgrade the truncated pair to an exact fixed-window
approximant with small norm loss and small functional loss.

\begin{theorem}[Exact fixed-window localization from a small full-window defect]
\label{thm:defect-localization}
Assume the hypotheses of Definition~\ref{def:micro-dual}. Let \(h_n\downarrow0\), and let
\[
\Phi_n\in\mathcal D(\mathfrak T_{t,h_n,\alpha}^{\mathrm{mic}}),
\qquad
\Psi_n:=\mathfrak T_{t,h_n,\alpha}^{\mathrm{mic}}\Phi_n,
\qquad
\|\Phi_n\|_{\mathfrak G_{t,h_n,\alpha}^{\mathrm{mic}}}\le1.
\]
Fix radii
\[
0<R_0<R_1<R_2<\infty,
\]
and let \(\chi\) be a smooth buffered cutoff with radii \((R_0,R_1,R_2)\). Assume that for some
\(\varepsilon>0\) and all sufficiently large \(n\),
\begin{align}
\left\|
\one_{(R_1,\infty)}
\mathfrak g_{i,t,h_n,\alpha}^{\mathrm{mic}}
\right\|_{L^2(\R_+)}
&\le
\varepsilon
\qquad\text{for }i\in\{Y,XY\},
\label{eq:defect-localize-profile}\\
\left\|
\mathfrak F_{t,h_n,\alpha,R_2}^{\mathrm{ex}}(\chi\Phi_n,\chi\Psi_n)
\right\|_{E_{R_2}}
&\le
\varepsilon.
\label{eq:defect-localize-defect}
\end{align}
Define
\begin{equation}\label{eq:defect-localize-corrected}
\widetilde\Psi_n
:=
\chi\Psi_n
-
\bigl(
\mathfrak B_{Y,t,h_n,\alpha,R_2}^{\mathrm{mic}}
\bigr)^{-1}
\mathfrak F_{t,h_n,\alpha,R_2}^{\mathrm{ex}}(\chi\Phi_n,\chi\Psi_n),
\end{equation}
and set
\[
\widetilde\Phi_n:=\chi\Phi_n.
\]
Then there exists \(C_{t,R_2,\eta,\nu_Y}<\infty\) such that, for all sufficiently large \(n\),
\begin{enumerate}[label=(\roman*)]
\item
\[
(\widetilde\Phi_n,\widetilde\Psi_n)\in \Gamma_{t,h_n,\alpha,R_2}^{\mathrm{ex}};
\]
\item
\begin{equation}\label{eq:defect-localize-norm}
\|\widetilde\Phi_n\|_{\mathfrak G_{t,h_n,\alpha,R_2}^{\mathrm{ex}}}
\le
1+C_{t,R_2,\eta,\nu_Y}\varepsilon;
\end{equation}
\item
\begin{equation}\label{eq:defect-localize-functional}
\left|
\Lambda_{t,h_n,\alpha}^{\mathrm{mic}}(\Phi_n)
-
\Lambda_{t,h_n,\alpha,R_2}^{\mathrm{ex}}(\widetilde\Phi_n)
\right|
\le
C_{t,R_2,\eta,\nu_Y}\varepsilon.
\end{equation}
\end{enumerate}
\end{theorem}

\begin{proof}
Part (i) is exactly Theorem~\ref{thm:micro-window-correction} applied to the pair
\((\chi\Phi_n,\chi\Psi_n)\).

For part (ii), \(|\chi|\le1\) gives
\[
\|\widetilde\Phi_n\|_{E_{R_2}}
\le
\|\Phi_n\|_{L^2(\R_+)}
\le
1.
\]
By Theorem~\ref{thm:micro-window-correction} and \eqref{eq:defect-localize-defect},
\[
\|\widetilde\Psi_n-\chi\Psi_n\|_{E_{R_2}}
\le
C_{t,R_2}\varepsilon.
\]
Since \(|\chi|\le1\) and \(\|\Phi_n\|_{\mathfrak G_{t,h_n,\alpha}^{\mathrm{mic}}}\le1\),
\[
\|\chi\Psi_n\|_{E_{R_2}}
\le
\|\Psi_n\|_{L^2(\R_+)}
\le
\frac{\nu_Y}{\eta}.
\]
Therefore
using \(\|\Phi_n\|_{\mathfrak G_{t,h_n,\alpha}^{\mathrm{mic}}}^2\le1\),
\begin{align*}
\|\widetilde\Phi_n\|_{\mathfrak G_{t,h_n,\alpha,R_2}^{\mathrm{ex}}}^2
&=
\|\widetilde\Phi_n\|_{E_{R_2}}^2
+
\frac{\eta^2}{\nu_Y^2}\|\widetilde\Psi_n\|_{E_{R_2}}^2\\
&\le
\|\Phi_n\|_{L^2(\R_+)}^2
+
\frac{\eta^2}{\nu_Y^2}
\bigl(
\|\Psi_n\|_{L^2(\R_+)}+C_{t,R_2}\varepsilon
\bigr)^2\\
&\le
\|\Phi_n\|_{L^2(\R_+)}^2
+
\frac{\eta^2}{\nu_Y^2}\|\Psi_n\|_{L^2(\R_+)}^2
+
C_{t,R_2,\eta,\nu_Y}\varepsilon\\
&\le
1+C_{t,R_2,\eta,\nu_Y}\varepsilon.
\end{align*}
After enlarging the constant if necessary,
\[
\|\widetilde\Phi_n\|_{\mathfrak G_{t,h_n,\alpha,R_2}^{\mathrm{ex}}}
\le
1+C_{t,R_2,\eta,\nu_Y}\varepsilon,
\]
which proves \eqref{eq:defect-localize-norm}.

For part (iii), using \eqref{eq:micro-global-functional},
\eqref{eq:micro-window-functional-ex}, and part (i), one has
\begin{align*}
\Lambda_{t,h_n,\alpha}^{\mathrm{mic}}(\Phi_n)
-\Lambda_{t,h_n,\alpha,R_2}^{\mathrm{ex}}(\widetilde\Phi_n)
&=
\left\langle
\mathfrak g_{XY,t,h_n,\alpha}^{\mathrm{mic}},
(1-\chi)\Phi_n
\right\rangle_{L^2(\R_+)}\\
&\quad-
\left\langle
\mathfrak g_{Y,t,h_n,\alpha}^{\mathrm{mic}},
(1-\chi)\Psi_n
\right\rangle_{L^2(\R_+)}\\
&\quad-
\left\langle
\mathfrak g_{Y,t,h_n,\alpha}^{\mathrm{mic}},
\widetilde\Psi_n-\chi\Psi_n
\right\rangle_{L^2(\R_+)}.
\end{align*}
Because \((1-\chi)\) is supported in \([R_1,\infty)\), the first two terms are bounded by
\[
\left\|
\one_{(R_1,\infty)}
\mathfrak g_{XY,t,h_n,\alpha}^{\mathrm{mic}}
\right\|_{L^2}
\|\Phi_n\|_{L^2(\R_+)}
+
\left\|
\one_{(R_1,\infty)}
\mathfrak g_{Y,t,h_n,\alpha}^{\mathrm{mic}}
\right\|_{L^2}
\|\Psi_n\|_{L^2(\R_+)}.
\]
Using \eqref{eq:defect-localize-profile} and
\[
\|\Phi_n\|_{L^2(\R_+)}\le1,
\qquad
\|\Psi_n\|_{L^2(\R_+)}\le \frac{\nu_Y}{\eta},
\]
these terms are at most \(C_{\eta,\nu_Y}\varepsilon\). For the correction term, the full
\(L^2\)-norm of \(\mathfrak g_{Y,t,h_n,\alpha}^{\mathrm{mic}}\) is uniformly bounded for large \(n\)
by Theorem~\ref{thm:micro-profile}, while Theorem~\ref{thm:micro-window-correction} and
\eqref{eq:defect-localize-defect} give
\[
\|\widetilde\Psi_n-\chi\Psi_n\|_{E_{R_2}}
\le
C_{t,R_2}\varepsilon.
\]
Hence the third term is also bounded by \(C_{t,R_2,\eta,\nu_Y}\varepsilon\), which proves
\eqref{eq:defect-localize-functional}.
\end{proof}

\begin{corollary}[Reduction of abstract localization to full-window defect control]
\label{cor:defect-localization}
Assume the hypotheses of Definition~\ref{def:micro-dual}. To verify the localization hypothesis in
Theorem~\ref{thm:abstract-global-reduction}, it is enough to prove the following statement:

\smallskip
\noindent
for every sequence \(h_n\downarrow0\) and every graph-unit sequence
\[
\Phi_n\in\mathcal D(\mathfrak T_{t,h_n,\alpha}^{\mathrm{mic}}),
\qquad
\Psi_n:=\mathfrak T_{t,h_n,\alpha}^{\mathrm{mic}}\Phi_n,
\qquad
\|\Phi_n\|_{\mathfrak G_{t,h_n,\alpha}^{\mathrm{mic}}}\le1,
\]
which satisfies the structured reduction of Corollary~\ref{cor:structured-reduction}, and for every
\(\varepsilon>0\), there exist common radii
\[
0<R_0<R_1<R_2<\infty
\]
and one smooth buffered cutoff \(\chi\) with radii \((R_0,R_1,R_2)\) such that, after passing to a
subsequence if needed,
\eqref{eq:defect-localize-profile} and \eqref{eq:defect-localize-defect} hold.
\end{corollary}

\begin{proof}
By Corollary~\ref{cor:structured-reduction}, it is enough to treat non-escaping sequences with a
common head window and common profile-tight buffers. For such sequences, if
\eqref{eq:defect-localize-profile} and \eqref{eq:defect-localize-defect} hold, then
Theorem~\ref{thm:defect-localization} produces exact fixed-window approximants on the common window
\([0,R_2]\) with arbitrarily small losses in both graph norm and functional value. This gives
the localization property required in Theorem~\ref{thm:abstract-global-reduction}.
\end{proof}

\section{Annulus-outer splitting of the full-window defect}
\label{sec:defect-splitting}

The full-window defect has two sources: the transition annulus and the outer tail.  Because the
cutoff is constant on \([0,R_1]\), mass below \(R_1\) does not enter the commutator.  The annulus
\([R_1,R_2]\) is therefore the only local contribution, and its size is reduced to an explicit
quantitative bound.  The complementary contribution is the outer tail \((R_2,\infty)\).

\begin{proposition}[Annulus-outer decomposition of the full-window defect]
\label{prop:defect-split}
Assume the hypotheses of Definition~\ref{def:micro-dual}. Let
\[
0<R_0<R_1<R_2<\infty,
\]
let \(\chi\) be a smooth buffered cutoff with radii \((R_0,R_1,R_2)\), and let
\[
\Phi\in\mathcal D(\mathfrak T_{t,h,\alpha}^{\mathrm{mic}}),
\qquad
\Psi:=\mathfrak T_{t,h,\alpha}^{\mathrm{mic}}\Phi.
\]
Assume in addition that
\[
hR_2<t,
\qquad
R_2<\varepsilon_{h,\alpha}^{-1}.
\]
Define
\[
\Phi^{\mathrm{ann}}:=\Phi\one_{[R_1,R_2]},
\qquad
\Phi^{\mathrm{out}}:=\Phi\one_{(R_2,\infty)},
\qquad
\Psi^{\mathrm{ann}}:=\Psi\one_{[R_1,R_2]},
\qquad
\Psi^{\mathrm{out}}:=\Psi\one_{(R_2,\infty)}.
\]
Then
\begin{equation}\label{eq:defect-split}
\mathfrak F_{t,h,\alpha,R_2}^{\mathrm{ex}}(\chi\Phi,\chi\Psi)
=
\mathfrak A_{t,h,\alpha,\chi}^{\mathrm{ann}}(\Phi,\Psi)
+
\mathfrak A_{t,h,\alpha,\chi}^{\mathrm{out}}(\Phi,\Psi)
\qquad\text{in }E_{R_2},
\end{equation}
where
\begin{align}
\mathfrak A_{t,h,\alpha,\chi}^{\mathrm{ann}}(\Phi,\Psi)
&:=
\mathfrak Q_{Y,t,h,\alpha,\chi}^{\mathrm{com}}\Psi^{\mathrm{ann}}
-
\mathfrak Q_{XY,t,h,\alpha,\chi}^{\mathrm{com}}\Phi^{\mathrm{ann}},
\label{eq:defect-annulus}\\
\mathfrak A_{t,h,\alpha,\chi}^{\mathrm{out}}(\Phi,\Psi)
&:=
\mathfrak Q_{Y,t,h,\alpha,\chi}^{\mathrm{com}}\Psi^{\mathrm{out}}
-
\mathfrak Q_{XY,t,h,\alpha,\chi}^{\mathrm{com}}\Phi^{\mathrm{out}}.
\label{eq:defect-outer}
\end{align}
Moreover, there exists \(C_{t,R_2}<\infty\) such that
\begin{equation}\label{eq:defect-annulus-bound}
\left\|
\mathfrak A_{t,h,\alpha,\chi}^{\mathrm{ann}}(\Phi,\Psi)
\right\|_{E_{R_2}}
\le
C_{t,R_2}\Lambda_{\chi,R_2}
\bigl(
\|\Phi\|_{L^2([R_1,R_2])}
+
\|\Psi\|_{L^2([R_1,R_2])}
\bigr),
\end{equation}
where
\[
\Lambda_{\chi,R_2}:=
\left\|
\frac{\chi(q)-\chi(r)}{q-r}
\right\|_{C^1(\Delta_{R_2})},
\qquad
\Delta_{R_2}:=\{0\le r\le q\le R_2\}.
\]
If \(\chi=\chi_{R_1,R_2}\) is the canonical cutoff from Lemma~\ref{lem:quant-buffer}, then
\begin{equation}\label{eq:defect-annulus-width}
\left\|
\mathfrak A_{t,h,\alpha,\chi_{R_1,R_2}}^{\mathrm{ann}}(\Phi,\Psi)
\right\|_{E_{R_2}}
\le
C_{t,R_2}
\left(
\frac{1}{R_2-R_1}
+
\frac{1}{(R_2-R_1)^2}
\right)
\bigl(
\|\Phi\|_{L^2([R_1,R_2])}
+
\|\Psi\|_{L^2([R_1,R_2])}
\bigr).
\end{equation}
\end{proposition}

\begin{proof}
By Theorem~\ref{thm:micro-smooth-buffer}(i),
\[
\mathfrak F_{t,h,\alpha,R_2}^{\mathrm{ex}}(\chi\Phi,\chi\Psi)
=
\mathfrak Q_{Y,t,h,\alpha,\chi}^{\mathrm{com}}\Psi
-
\mathfrak Q_{XY,t,h,\alpha,\chi}^{\mathrm{com}}\Phi
\qquad\text{in }E_{R_2}.
\]
Since
\[
\Phi=\Phi\one_{[0,R_1]}+\Phi^{\mathrm{ann}}+\Phi^{\mathrm{out}},
\qquad
\Psi=\Psi\one_{[0,R_1]}+\Psi^{\mathrm{ann}}+\Psi^{\mathrm{out}},
\]
linearity gives
\begin{align*}
\mathfrak F_{t,h,\alpha,R_2}^{\mathrm{ex}}(\chi\Phi,\chi\Psi)
&=
\Bigl(
\mathfrak Q_{Y,t,h,\alpha,\chi}^{\mathrm{com}}(\Psi\one_{[0,R_1]})
-
\mathfrak Q_{XY,t,h,\alpha,\chi}^{\mathrm{com}}(\Phi\one_{[0,R_1]})
\Bigr)\\
&\quad+
\mathfrak A_{t,h,\alpha,\chi}^{\mathrm{ann}}(\Phi,\Psi)
+
\mathfrak A_{t,h,\alpha,\chi}^{\mathrm{out}}(\Phi,\Psi).
\end{align*}
It therefore remains to show that the first bracket vanishes.

Let \(\Theta\in L^2(\R_+)\) be supported in \([0,R_1]\). If the integrand in
\eqref{eq:micro-smooth-com-Y} or \eqref{eq:micro-smooth-com-XY} is nonzero, then
\[
r\le q\le R_1.
\]
Because \(\chi\equiv1\) on \([0,R_1]\), one has \(\chi(q)-\chi(r)=0\). Therefore
\[
\mathfrak W_{Y,t,h,\alpha,\chi}^{\mathrm{com}}\Theta
=
\mathfrak W_{XY,t,h,\alpha,\chi}^{\mathrm{com}}\Theta
=
0
\qquad\text{on }[0,R_2].
\]
Applying \eqref{eq:micro-smooth-com-conj-Y} and \eqref{eq:micro-smooth-com-conj-XY} gives
\[
\mathfrak Q_{Y,t,h,\alpha,\chi}^{\mathrm{com}}\Theta
=
\mathfrak Q_{XY,t,h,\alpha,\chi}^{\mathrm{com}}\Theta
=
0
\qquad\text{in }E_{R_2},
\]
which proves \eqref{eq:defect-split}.

For \eqref{eq:defect-annulus-bound}, note that
\(\Phi^{\mathrm{ann}},\Psi^{\mathrm{ann}}\in E_{R_2}\), so Proposition~\ref{prop:quant-window-comm}
applies:
\begin{align*}
\left\|
\mathfrak A_{t,h,\alpha,\chi}^{\mathrm{ann}}(\Phi,\Psi)
\right\|_{E_{R_2}}
&\le
\left\|
\mathfrak Q_{Y,t,h,\alpha,\chi}^{\mathrm{com}}\Psi^{\mathrm{ann}}
\right\|_{E_{R_2}}
+
\left\|
\mathfrak Q_{XY,t,h,\alpha,\chi}^{\mathrm{com}}\Phi^{\mathrm{ann}}
\right\|_{E_{R_2}}\\
&\le
C_{t,R_2}\Lambda_{\chi,R_2}
\bigl(
\|\Psi^{\mathrm{ann}}\|_{E_{R_2}}
+
\|\Phi^{\mathrm{ann}}\|_{E_{R_2}}
\bigr),
\end{align*}
which is exactly \eqref{eq:defect-annulus-bound}. The width estimate
\eqref{eq:defect-annulus-width} follows from \eqref{eq:defect-annulus-bound} and
Lemma~\ref{lem:quant-buffer}.
\end{proof}

\begin{proposition}[Uniform control of the annulus commutator]
\label{prop:annulus-control}
Assume Assumption~\ref{ass:model}, fix \(t\in(0,T)\), \(\alpha\in(0,1)\), and suppose
\[
0<H_{XY}<H_Y<\tfrac12.
\]
Fix \(\ell>0\). Then there exists \(C_{t,\ell}<\infty\) such that the following holds. Let
\[
0<R_1<R_2<\infty,
\qquad
R_2-R_1=\ell,
\]
let \(\chi:=\chi_{R_1,R_2}\) be the canonical cutoff from Lemma~\ref{lem:quant-buffer}, and choose
\(h>0\) sufficiently small that
\[
hR_2<t,
\qquad
R_2<\varepsilon_{h,\alpha}^{-1}.
\]
Then for every \(\Theta\in L^2([R_1,R_2])\),
\begin{equation}\label{eq:annulus-control-single}
\left\|
\mathfrak Q_{Y,t,h,\alpha,\chi}^{\mathrm{com}}\Theta
\right\|_{E_{R_2}}
+
\left\|
\mathfrak Q_{XY,t,h,\alpha,\chi}^{\mathrm{com}}\Theta
\right\|_{E_{R_2}}
\le
C_{t,\ell}\|\Theta\|_{L^2([R_1,R_2])}.
\end{equation}
Consequently, for every
\[
\Phi\in\mathcal D(\mathfrak T_{t,h,\alpha}^{\mathrm{mic}}),
\qquad
\Psi:=\mathfrak T_{t,h,\alpha}^{\mathrm{mic}}\Phi,
\]
one has
\begin{equation}\label{eq:annulus-control-pair}
\left\|
\mathfrak A_{t,h,\alpha,\chi}^{\mathrm{ann}}(\Phi,\Psi)
\right\|_{E_{R_2}}
\le
C_{t,\ell}
\bigl(
\|\Phi\|_{L^2([R_1,R_2])}
+
\|\Psi\|_{L^2([R_1,R_2])}
\bigr).
\end{equation}
\end{proposition}

\begin{proof}
Let
\[
S_0:=(R_1-\ell)_+,
\qquad
J:=[S_0,R_2].
\]
Then \(|J|\le 2\ell\), and the support of every \(\Theta\in L^2([R_1,R_2])\) lies inside \(J\).

We first control the output on the local strip \(J\). Since \(\Theta\) is supported in \([R_1,R_2]\)
and \(\chi=\chi_{R_1,R_2}\) satisfies the derivative bounds
\eqref{eq:quant-buffer-derivatives}, the proof of Proposition~\ref{prop:quant-window-comm},
restricted to the strip \(J\), gives
\begin{equation}\label{eq:annulus-local-strip}
\left\|
\one_J\mathfrak Q_{Y,t,h,\alpha,\chi}^{\mathrm{com}}\Theta
\right\|_{E_{R_2}}
+
\left\|
\one_J\mathfrak Q_{XY,t,h,\alpha,\chi}^{\mathrm{com}}\Theta
\right\|_{E_{R_2}}
\le
C_{t,\ell}\|\Theta\|_{L^2([R_1,R_2])},
\end{equation}
because the output interval has length at most \(2\ell\), the input interval has length \(\ell\),
and the singular integrals in the source-side estimate are therefore taken over a set whose size is
controlled solely by \(\ell\).

It remains to control the far-left part \([0,S_0]\). If \(S_0=0\), there is nothing to prove, so
assume \(S_0>0\), equivalently \(R_1>\ell\). For \(0\le r\le S_0\) and \(q\in[R_1,R_2]\), one
has \(\chi(r)=1\), hence
\[
\chi(q)-\chi(r)=-(1-\chi(q)).
\]
Therefore, for every \(\Theta\in L^2([R_1,R_2])\), defining
\[
\widetilde\Theta:=(1-\chi)\Theta\in L^2([R_1,R_2]),
\]
the commutator Volterra terms reduce on \([0,S_0]\) to
\begin{align*}
\mathfrak W_{Y,t,h,\alpha,\chi}^{\mathrm{com}}\Theta
&=
-\mathfrak W_{Y,t,h,\alpha,S_0,R_1}^{\mathrm{tail}}\widetilde\Theta,\\
\mathfrak W_{XY,t,h,\alpha,\chi}^{\mathrm{com}}\Theta
&=
-\mathfrak W_{XY,t,h,\alpha,S_0,R_1}^{\mathrm{tail}}\widetilde\Theta.
\end{align*}
The proofs of Propositions~\ref{prop:micro-tail-Y} and~\ref{prop:micro-tail-XY} now apply with the
additional information that \(\widetilde\Theta\) is supported in the interval \([R_1,R_2]\) of
length \(\ell\). On the target side, the conjugated kernel is uniformly \(O(h)\), so
\[
\left\|
\one_{[0,S_0]}
\mathfrak Q_{Y,t,h,\alpha,\chi}^{\mathrm{com}}\Theta
\right\|_{E_{R_2}}
\le
C_t h\sqrt{S_0\ell}\,\|\Theta\|_{L^2([R_1,R_2])}.
\]
Since \(S_0\le R_2<\varepsilon_{h,\alpha}^{-1}=h^{\alpha-1}\), this gives
\begin{equation}\label{eq:annulus-inner-Y}
\left\|
\one_{[0,S_0]}
\mathfrak Q_{Y,t,h,\alpha,\chi}^{\mathrm{com}}\Theta
\right\|_{E_{R_2}}
\le
C_{t,\ell}\|\Theta\|_{L^2([R_1,R_2])}.
\end{equation}

On the source side, the proof of Proposition~\ref{prop:micro-tail-XY} yields a decomposition of the
conjugated kernel into a main term bounded by \(C_t(q-r)^{-\beta-1}\) and an error bounded by
\(C_t h(q-r)^{-\beta}\). Since \(r\le S_0=R_1-\ell\) and \(q\in[R_1,R_2]\), one has
\[
q-r\ge \ell.
\]
Schur's test therefore bounds the contribution of the main kernel by
\[
C_{t,\ell}\|\Theta\|_{L^2([R_1,R_2])}.
\]
For the error kernel, the Hilbert--Schmidt estimate gives
\[
C_t h\sqrt{S_0}\,\ell^{\frac12-\beta}\|\Theta\|_{L^2([R_1,R_2])}
\le
C_{t,\ell}\|\Theta\|_{L^2([R_1,R_2])}.
\]
Hence
\begin{equation}\label{eq:annulus-inner-XY}
\left\|
\one_{[0,S_0]}
\mathfrak Q_{XY,t,h,\alpha,\chi}^{\mathrm{com}}\Theta
\right\|_{E_{R_2}}
\le
C_{t,\ell}\|\Theta\|_{L^2([R_1,R_2])}.
\end{equation}

Combining \eqref{eq:annulus-local-strip}, \eqref{eq:annulus-inner-Y}, and
\eqref{eq:annulus-inner-XY} proves \eqref{eq:annulus-control-single}. The pair estimate
\eqref{eq:annulus-control-pair} follows from the definition \eqref{eq:defect-annulus} and the
triangle inequality.
\end{proof}

\begin{corollary}[Good buffers imply annulus control]
\label{cor:annulus-control}
Assume the hypotheses of Definition~\ref{def:micro-dual}. For every \(\varepsilon>0\) and every
\(\ell>0\), after
passing to a subsequence, there exist common radii
\[
0<R_0<R_1<R_2<\infty,
\qquad
R_2-R_1=\ell,
\]
such that:
\begin{enumerate}[label=(\roman*)]
\item for \(i\in\{Y,XY\}\),
\[
\left\|
\one_{(R_1,\infty)}
\mathfrak g_{i,t,h_n,\alpha}^{\mathrm{mic}}
\right\|_{L^2(\R_+)}
\le
\varepsilon
\qquad\text{for every }n;
\]
\item for the canonical cutoff \(\chi_{R_1,R_2}\),
\begin{equation}\label{eq:annulus-control-seq}
\left\|
\mathfrak A_{t,h_n,\alpha,\chi_{R_1,R_2}}^{\mathrm{ann}}(\Phi_n,\Psi_n)
\right\|_{E_{R_2}}
\le
\varepsilon
\qquad\text{for every }n.
\end{equation}
\end{enumerate}
\end{corollary}

\begin{proof}
Let \(C_{t,\ell}\) be the constant from Proposition~\ref{prop:annulus-control}, and
choose \(M\in\mathbb N\) so large that
\[
C_{t,\ell}\left(1+\frac{\nu_Y}{\eta}\right)M^{-1/2}\le \varepsilon.
\]
Applying Corollary~\ref{cor:profile-tight-buffer} with this \(M\), with this \(\ell\), and with the
same \(\varepsilon\), and then relabeling the resulting subsequence, yields common radii
\[
0<R_0<R_1<R_2<\infty,
\qquad
R_1-R_0\ge \ell,
\qquad
R_2-R_1=\ell,
\]
such that the profile-tail estimate beyond \(R_0\) holds and
\[
\|\Phi_n\|_{L^2([R_1,R_2])}^2
+
\frac{\eta^2}{\nu_Y^2}
\|\Psi_n\|_{L^2([R_1,R_2])}^2
\le
\frac{1}{M}
\qquad\text{for every }n.
\]
Since \(R_1>R_0\), the same profile-tail bound holds beyond \(R_1\), which proves (i). The
annulus mass bound implies
\[
\|\Phi_n\|_{L^2([R_1,R_2])}
+
\|\Psi_n\|_{L^2([R_1,R_2])}
\le
\left(1+\frac{\nu_Y}{\eta}\right)M^{-1/2},
\]
and after discarding finitely many terms we may also assume
\[
h_nR_2<t,
\qquad
R_2<\varepsilon_{h_n,\alpha}^{-1}
\qquad\text{for every }n.
\]
Proposition~\ref{prop:annulus-control} then gives \eqref{eq:annulus-control-seq}.
\end{proof}

\begin{lemma}[Uniform separated buffered tail bounds]
\label{lem:tail-uniform}
Assume the hypotheses of Definition~\ref{def:micro-tail}. Then there exist constants
\[
C_t<\infty,
\qquad
C_{t,\alpha}<\infty
\]
such that for every sufficiently small \(h>0\),
\begin{align}
\left\|
\mathfrak Q_{Y,t,h,\alpha,R_0,R}^{\mathrm{tail}}
\right\|_{\mathcal L(F_{h,R},E_{R_0})}
&\le
C_t h^\alpha,
\label{eq:tail-uniform-Y}\\
\left\|
\mathfrak Q_{XY,t,h,\alpha,R_0,R}^{\mathrm{tail}}
\right\|_{\mathcal L(F_{h,R},E_{R_0})}
&\le
C_t(R-R_0)^{-\beta}
+
C_{t,\alpha}h^{\alpha+\beta-\alpha\beta}.
\label{eq:tail-uniform-XY}
\end{align}
\end{lemma}

\begin{proof}
The target-side estimate follows from \eqref{eq:micro-tail-Y-bound} and the proof of
Proposition~\ref{prop:micro-tail-Y}: since \(R_0<R<\varepsilon_{h,\alpha}^{-1}\),
\[
h\sqrt{\varepsilon_{h,\alpha}^{-1}-R}
\le
h\sqrt{\varepsilon_{h,\alpha}^{-1}}
=
h^{\frac{1+\alpha}{2}}
\le
h^\alpha
\qquad(0<h\le1),
\]
which proves \eqref{eq:tail-uniform-Y}.

For the source side, the main term is exactly \eqref{eq:micro-tail-XY-main}, hence bounded by
\(C_t(R-R_0)^{-\beta}\). For the error term, the proof of
Proposition~\ref{prop:micro-tail-XY} gives
\[
\left\|
K_{XY,t,h}^{(2)}
\right\|_{L^2([0,R_0]\times[R,\varepsilon_{h,\alpha}^{-1}])}^2
\le
C_t h^2 R_0\,\varepsilon_{h,\alpha}^{-(1-2\beta)}.
\]
Since \(R_0<\varepsilon_{h,\alpha}^{-1}\), one obtains
\[
\left\|
K_{XY,t,h}^{(2)}
\right\|_{\mathcal L(F_{h,R},E_{R_0})}
\le
C_t h\,\varepsilon_{h,\alpha}^{-(1-\beta)}
=
C_{t,\alpha}h^{\alpha+\beta-\alpha\beta},
\]
which proves \eqref{eq:tail-uniform-XY}.
\end{proof}

\begin{proposition}[Separated part of the outer-tail commutator]
\label{prop:outer-tail-separated}
Assume the hypotheses of Definition~\ref{def:micro-dual}. Let
\[
0<R_1<R_2<\infty,
\qquad
\ell:=R_2-R_1,
\]
let \(\chi:=\chi_{R_1,R_2}\) be the canonical cutoff from Lemma~\ref{lem:quant-buffer}, and let
\[
\Phi\in\mathcal D(\mathfrak T_{t,h,\alpha}^{\mathrm{mic}}),
\qquad
\Psi:=\mathfrak T_{t,h,\alpha}^{\mathrm{mic}}\Phi.
\]
Assume
\[
hR_2<t,
\qquad
R_2<\varepsilon_{h,\alpha}^{-1}.
\]
Then
\begin{equation}\label{eq:outer-tail-separated-identity}
\one_{[0,R_1]}\mathfrak A_{t,h,\alpha,\chi}^{\mathrm{out}}(\Phi,\Psi)
=
\mathfrak Q_{XY,t,h,\alpha,R_1,R_2}^{\mathrm{tail}}\Phi^{\mathrm{out}}
-
\mathfrak Q_{Y,t,h,\alpha,R_1,R_2}^{\mathrm{tail}}\Psi^{\mathrm{out}}
\qquad\text{in }E_{R_1},
\end{equation}
and there exist constants \(C_t<\infty\) and \(C_{t,\alpha,\eta,\nu_Y}<\infty\) such that
\begin{equation}\label{eq:outer-tail-separated-bound}
\left\|
\one_{[0,R_1]}\mathfrak A_{t,h,\alpha,\chi}^{\mathrm{out}}(\Phi,\Psi)
\right\|_{E_{R_2}}
\le
C_t\ell^{-\beta}
+
C_{t,\alpha,\eta,\nu_Y}h^\alpha
\end{equation}
for all sufficiently small \(h>0\).
\end{proposition}

\begin{proof}
For \(0\le r\le R_1\) and \(q\ge R_2\), one has \(\chi(r)=1\) and \(\chi(q)=0\). Hence
\[
\chi(q)-\chi(r)=-1.
\]
Therefore, on \([0,R_1]\),
\[
\mathfrak W_{Y,t,h,\alpha,\chi}^{\mathrm{com}}\Psi^{\mathrm{out}}
=
-\mathfrak W_{Y,t,h,\alpha,R_1,R_2}^{\mathrm{tail}}\Psi^{\mathrm{out}},
\]
and similarly
\[
\mathfrak W_{XY,t,h,\alpha,\chi}^{\mathrm{com}}\Phi^{\mathrm{out}}
=
-\mathfrak W_{XY,t,h,\alpha,R_1,R_2}^{\mathrm{tail}}\Phi^{\mathrm{out}}.
\]
Applying \eqref{eq:micro-smooth-com-conj-Y}, \eqref{eq:micro-smooth-com-conj-XY},
\eqref{eq:micro-tail-Y-conj}, and \eqref{eq:micro-tail-XY-conj} yields
\eqref{eq:outer-tail-separated-identity}.

Now use Lemma~\ref{lem:tail-uniform} with \(R_0=R_1\) and \(R=R_2\):
\begin{align*}
\left\|
\one_{[0,R_1]}\mathfrak A_{t,h,\alpha,\chi}^{\mathrm{out}}(\Phi,\Psi)
\right\|_{E_{R_2}}
&\le
\left\|
\mathfrak Q_{XY,t,h,\alpha,R_1,R_2}^{\mathrm{tail}}\Phi^{\mathrm{out}}
\right\|_{E_{R_1}}
+
\left\|
\mathfrak Q_{Y,t,h,\alpha,R_1,R_2}^{\mathrm{tail}}\Psi^{\mathrm{out}}
\right\|_{E_{R_1}}\\
&\le
\bigl(
C_t\ell^{-\beta}
+
C_{t,\alpha}h^{\alpha+\beta-\alpha\beta}
\bigr)\|\Phi^{\mathrm{out}}\|_{L^2(\R_+)}
+
C_th^\alpha\|\Psi^{\mathrm{out}}\|_{L^2(\R_+)}.
\end{align*}
Since
\[
\|\Phi^{\mathrm{out}}\|_{L^2(\R_+)}
\le
\|\Phi\|_{L^2(\R_+)}
\le
1,
\qquad
\|\Psi^{\mathrm{out}}\|_{L^2(\R_+)}
\le
\|\Psi\|_{L^2(\R_+)}
\le
\frac{\nu_Y}{\eta},
\]
and \(h^{\alpha+\beta-\alpha\beta}\le h^\alpha\) for \(0<h\le1\), this proves
\eqref{eq:outer-tail-separated-bound}.
\end{proof}

\begin{corollary}[Localization reduces to boundary-strip outer-tail control]
\label{cor:outer-tail-reduction}
Assume the hypotheses of Definition~\ref{def:micro-dual}. To verify the localization hypothesis in
Theorem~\ref{thm:abstract-global-reduction}, it is enough to prove the following statement:

\smallskip
\noindent
for every sequence \(h_n\downarrow0\), every graph-unit sequence
\[
\Phi_n\in\mathcal D(\mathfrak T_{t,h_n,\alpha}^{\mathrm{mic}}),
\qquad
\Psi_n:=\mathfrak T_{t,h_n,\alpha}^{\mathrm{mic}}\Phi_n,
\qquad
\|\Phi_n\|_{\mathfrak G_{t,h_n,\alpha}^{\mathrm{mic}}}\le1,
\]
which satisfies the structured reduction of Corollary~\ref{cor:structured-reduction}, and every
\(\varepsilon>0\), there exists \(\ell>0\) such that after passing to a subsequence and fixing the
common radii
\[
0<R_0<R_1<R_2<\infty,
\qquad
R_2-R_1=\ell,
\]
and the canonical cutoff \(\chi=\chi_{R_1,R_2}\) supplied by Corollary~\ref{cor:annulus-control},
one has
\begin{equation}\label{eq:outer-tail-strip}
\left\|
\one_{[R_1,R_2]}
\mathfrak A_{t,h_n,\alpha,\chi}^{\mathrm{out}}(\Phi_n,\Psi_n)
\right\|_{E_{R_2}}
\le
\varepsilon
\qquad\text{for every }n.
\end{equation}
\end{corollary}

\begin{proof}
Choose \(\ell>0\) so large that
\[
C_t\ell^{-\beta}\le \varepsilon,
\]
where \(C_t\) is the constant from Proposition~\ref{prop:outer-tail-separated}. By
Corollary~\ref{cor:annulus-control}, after passing to a subsequence we may assume that the common
profile-tail bound holds beyond \(R_1\) and that
\[
\left\|
\mathfrak A_{t,h_n,\alpha,\chi}^{\mathrm{ann}}(\Phi_n,\Psi_n)
\right\|_{E_{R_2}}
\le
\varepsilon
\qquad\text{for every }n.
\]
After discarding finitely many terms one also has
\[
h_nR_2<t,
\qquad
R_2<\varepsilon_{h_n,\alpha}^{-1}
\qquad\text{for every }n,
\]
and
\[
C_{t,\alpha,\eta,\nu_Y}h_n^\alpha\le \varepsilon
\qquad\text{for every }n.
\]
Hence Proposition~\ref{prop:outer-tail-separated} gives
\[
\left\|
\one_{[0,R_1]}
\mathfrak A_{t,h_n,\alpha,\chi}^{\mathrm{out}}(\Phi_n,\Psi_n)
\right\|_{E_{R_2}}
\le
2\varepsilon
\qquad\text{for every }n.
\]
If \eqref{eq:outer-tail-strip} also holds, then
\[
\left\|
\mathfrak A_{t,h_n,\alpha,\chi}^{\mathrm{out}}(\Phi_n,\Psi_n)
\right\|_{E_{R_2}}
\le
3\varepsilon
\qquad\text{for every }n.
\]
Together with the annulus bound and Proposition~\ref{prop:defect-split}, this gives
\[
\left\|
\mathfrak F_{t,h_n,\alpha,R_2}^{\mathrm{ex}}(\chi\Phi_n,\chi\Psi_n)
\right\|_{E_{R_2}}
\le
4\varepsilon
\qquad\text{for every }n.
\]
The profile-tail bound beyond \(R_1\) and Theorem~\ref{thm:defect-localization} therefore give
exact fixed-window approximants with arbitrarily small losses in norm and functional value. By
Corollary~\ref{cor:defect-localization}, this verifies the localization hypothesis in
Theorem~\ref{thm:abstract-global-reduction}.
\end{proof}

\begin{remark}[Explicit form of the annulus term]
\label{rem:annulus-explicit}
Proposition~\ref{prop:defect-split} removes the annulus contribution from the list of hidden
unknowns.  It is the explicit operator \(\mathfrak A_{t,h,\alpha,\chi}^{\mathrm{ann}}\), and
Proposition~\ref{prop:annulus-control} and Corollary~\ref{cor:annulus-control} now show that this
annulus term can be made uniformly small along structured non-escaping sequences. Thus the
remaining global work is reduced to the strip contribution
\(\one_{[R_1,R_2]}\mathfrak A_{t,h,\alpha,\chi}^{\mathrm{out}}\). As
Section~\ref{sec:strip-annulus} shows, this is no longer a tail-driven quantity once the
exact micro transfer identity is used.
\end{remark}

\section{Flat outer-edge structure of the canonical cutoff}
\label{sec:flat-cutoff}

The boundary-strip term is created where the canonical cutoff meets the outer edge \(R_2\).  The
needed property is stronger than \(C^2\).  Because the profile used in Lemma~\ref{lem:quant-buffer}
is identically zero on \([1,\infty)\), the rescaled cutoff is flat to every order at \(R_2\).  This
flatness softens the conjugated commutator kernels in the strip analysis.

\begin{lemma}[Arbitrary-order flatness at the outer edge]
\label{lem:quant-buffer-flat}
For every pair of integers \(m\ge0\) and \(k\in\{0,1,2\}\), there exists a constant
\[
C_{\chi,m,k}<\infty
\]
such that for every \(0<R_1<R_2<\infty\), the canonical cutoff \(\chi_{R_1,R_2}\) from
Lemma~\ref{lem:quant-buffer} satisfies
\begin{equation}\label{eq:quant-buffer-flat}
\left|
\chi_{R_1,R_2}^{(k)}(r)
\right|
\le
C_{\chi,m,k}
\frac{(R_2-r)^m}{(R_2-R_1)^{m+k}}
\qquad\text{for every }r\in[R_1,R_2].
\end{equation}
\end{lemma}

\begin{proof}
Recall from the proof of Lemma~\ref{lem:quant-buffer} that
\[
\chi_{R_1,R_2}(r)
=
\vartheta\!\left(\frac{r-R_1}{R_2-R_1}\right),
\qquad r\ge0,
\]
for some fixed \(\vartheta\in C^\infty(\R)\) such that
\[
\vartheta\equiv 1\ \text{on }(-\infty,0],
\qquad
\vartheta\equiv 0\ \text{on }[1,\infty).
\]
Hence, for every \(k\in\{0,1,2\}\), the derivative \(\vartheta^{(k)}\) vanishes identically on
\([1,\infty)\), and therefore all one-sided derivatives of \(\vartheta^{(k)}\) at \(1\) vanish.
Since \(\vartheta^{(k)}\) is smooth on a neighborhood of \(1\), Taylor's theorem at the point
\(1\) shows that for every \(m\ge0\) there exists \(C_{\chi,m,k}<\infty\) such that
\[
\left|
\vartheta^{(k)}(u)
\right|
\le
C_{\chi,m,k}(1-u)^m
\qquad\text{for every }u\in[0,1].
\]
Now let
\[
u:=\frac{r-R_1}{R_2-R_1}\in[0,1].
\]
Then
\[
\chi_{R_1,R_2}^{(k)}(r)
=
\frac{1}{(R_2-R_1)^k}
\vartheta^{(k)}(u),
\]
and
\[
1-u
=
\frac{R_2-r}{R_2-R_1}.
\]
Substituting these identities into the previous bound yields \eqref{eq:quant-buffer-flat}.
\end{proof}

\begin{corollary}[Boundary-strip decay of the cutoff quotient]
\label{cor:quant-buffer-flat}
Let \(0<R_1<R_2<\infty\), set \(\ell:=R_2-R_1\), and let \(\chi_{R_1,R_2}\) be the canonical
cutoff from Lemma~\ref{lem:quant-buffer}. For \(r\in[R_1,R_2]\) and \(q\ge R_2\), define
\[
\delta_{R_1,R_2}(r,q)
:=
\frac{\chi_{R_1,R_2}(q)-\chi_{R_1,R_2}(r)}{q-r}
=
-\frac{\chi_{R_1,R_2}(r)}{q-r}.
\]
Then, for every integer \(m\ge0\), there exists \(C_{\chi,m}<\infty\) such that
\begin{align}
\left|
\delta_{R_1,R_2}(r,q)
\right|
&\le
C_{\chi,m}
\frac{(R_2-r)^m}{\ell^m(q-r)},
\label{eq:quant-buffer-flat-delta}\\
\left|
\partial_r\delta_{R_1,R_2}(r,q)
\right|
+\left|
\partial_q\delta_{R_1,R_2}(r,q)
\right|
&\le
C_{\chi,m}
\left(
\frac{(R_2-r)^m}{\ell^{m+1}(q-r)}
+
\frac{(R_2-r)^m}{\ell^m(q-r)^2}
\right).
\label{eq:quant-buffer-flat-delta-deriv}
\end{align}
\end{corollary}

\begin{proof}
The identity
\[
\delta_{R_1,R_2}(r,q)
=
-\frac{\chi_{R_1,R_2}(r)}{q-r}
\]
follows because \(\chi_{R_1,R_2}(q)=0\) for \(q\ge R_2\). The bound
\eqref{eq:quant-buffer-flat-delta} then follows from \eqref{eq:quant-buffer-flat} with \(k=0\).

For the derivatives, direct differentiation gives
\[
\partial_r\delta_{R_1,R_2}(r,q)
=
-\frac{\chi_{R_1,R_2}'(r)}{q-r}
-\frac{\chi_{R_1,R_2}(r)}{(q-r)^2},
\qquad
\partial_q\delta_{R_1,R_2}(r,q)
=
\frac{\chi_{R_1,R_2}(r)}{(q-r)^2}.
\]
Applying \eqref{eq:quant-buffer-flat} with \(k=0\) and \(k=1\) yields
\eqref{eq:quant-buffer-flat-delta-deriv}.
\end{proof}

\section{Crossing-segment geometry at the outer edge}
\label{sec:edge-geometry}

The remaining boundary-strip term comes from Volterra segments that start at
\(r\in[R_1,R_2]\) in the transition annulus and end at \(q\ge R_2\).  The canonical cutoff vanishes
once the segment crosses \(R_2\), so each weighted Beta-average samples only the initial fraction
\[
\rho_{R_1,R_2}(r,q):=\frac{R_2-r}{q-r}\in[0,1].
\]
The strip analysis combines this truncation geometry with the exact conjugated Volterra formulas.

\begin{lemma}[Weighted segment truncation across the outer edge]
\label{lem:edge-segment}
Let \(a<1\), \(b>-1\), \(m\in\mathbb N_0\), and \(k\in\{0,1\}\). Then there exists a constant
\[
C_{a,b,m,k}<\infty
\]
such that for every \(0<R_1<R_2<\infty\), with \(\ell:=R_2-R_1\), the canonical cutoff
\(\chi_{R_1,R_2}\) from Lemma~\ref{lem:quant-buffer} satisfies
\begin{equation}\label{eq:edge-segment}
\int_0^1
\theta^{-a}(1-\theta)^b
\left|
\chi_{R_1,R_2}^{(k)}\!\bigl(r+(q-r)\theta\bigr)
\right|
d\theta
\le
C_{a,b,m,k}\,\ell^{-m-k}(R_2-r)^m\,
\rho_{R_1,R_2}(r,q)^{1-a}
\end{equation}
for every \(r\in[R_1,R_2]\) and every \(q\ge R_2\).
\end{lemma}

\begin{proof}
Fix \(r\in[R_1,R_2]\) and \(q\ge R_2\), and set \(\rho:=\rho_{R_1,R_2}(r,q)\). Then
\[
r+(q-r)\theta\ge R_2
\qquad\text{whenever }\theta\ge \rho.
\]
Because \(\chi_{R_1,R_2}\) vanishes identically on \([R_2,\infty)\), the integrand in
\eqref{eq:edge-segment} is zero for \(\theta\ge \rho\). Hence
\[
\int_0^1
\theta^{-a}(1-\theta)^b
\left|
\chi_{R_1,R_2}^{(k)}\!\bigl(r+(q-r)\theta\bigr)
\right|
d\theta
=
\int_0^\rho
\theta^{-a}(1-\theta)^b
\left|
\chi_{R_1,R_2}^{(k)}\!\bigl(r+(q-r)\theta\bigr)
\right|
d\theta.
\]

For \(0\le\theta\le \rho\), the point
\[
x:=r+(q-r)\theta
\]
lies in \([R_1,R_2]\), and
\[
R_2-x
=
R_2-r-(q-r)\theta
=
(q-r)(\rho-\theta).
\]
Therefore Lemma~\ref{lem:quant-buffer-flat} gives
\[
\left|
\chi_{R_1,R_2}^{(k)}(x)
\right|
\le
C_{\chi,m,k}
\frac{(R_2-x)^m}{\ell^{m+k}}
=
C_{\chi,m,k}
\frac{(q-r)^m(\rho-\theta)^m}{\ell^{m+k}}
\]
for every \(0\le\theta\le \rho\). Hence
\begin{align*}
\int_0^\rho
\theta^{-a}(1-\theta)^b
\left|
\chi_{R_1,R_2}^{(k)}\!\bigl(r+(q-r)\theta\bigr)
\right|
d\theta
&\le
C_{\chi,m,k}\frac{(q-r)^m}{\ell^{m+k}}
\int_0^\rho
\theta^{-a}(1-\theta)^b(\rho-\theta)^m\,d\theta\\
&=
C_{\chi,m,k}\frac{(q-r)^m\rho^{m+1-a}}{\ell^{m+k}}
\int_0^1
u^{-a}(1-\rho u)^b(1-u)^m\,du.
\end{align*}
If \(b\ge0\), then \((1-\rho u)^b\le1\). If \(-1<b<0\), then \(1-\rho u\ge1-u\), hence
\((1-\rho u)^b\le(1-u)^b\). In both cases
\[
\int_0^1
u^{-a}(1-\rho u)^b(1-u)^m\,du
\le
C_{a,b,m}<\infty.
\]
Since \((q-r)^m\rho^m=(R_2-r)^m\), this proves \eqref{eq:edge-segment}.
\end{proof}

\begin{proposition}[Outer-edge Beta averages for the canonical cutoff]
\label{prop:edge-beta}
Assume Assumption~\ref{ass:model}, fix \(t\in(0,T)\), and suppose
\[
0<H_{XY}<H_Y<\tfrac12.
\]
Set
\[
\lambda:=H_Y+\frac12,
\qquad
\mu:=H_{XY}+\frac12.
\]
Let \(0<R_1<R_2<\infty\), let \(\chi_{R_1,R_2}\) be the canonical cutoff from
Lemma~\ref{lem:quant-buffer}, set \(\ell:=R_2-R_1\), fix \(m\in\mathbb N_0\), and define for
\(r\in[R_1,R_2]\), \(q\ge R_2\),
\begin{align}
g_{Y,t,h,R_1,R_2}^{\mathrm{edge}}(r,q)
&:=
\frac{1}{\Gamma(\lambda)\Gamma(1-\lambda)}
\int_0^1
\theta^{-\lambda}(1-\theta)^\lambda
\chi_{R_1,R_2}\!\bigl(r+(q-r)\theta\bigr)
\notag\\
&\qquad\qquad\qquad\qquad\qquad\cdot
L_Y\!\bigl(t-h(r+(q-r)\theta),t-hq\bigr)\,d\theta,
\label{eq:edge-beta-Y}\\
g_{XY,t,h,R_1,R_2}^{\mathrm{edge}}(r,q)
&:=
\frac{1}{\Gamma(\lambda)\Gamma(1-\lambda)}
\int_0^1
\theta^{-\lambda}(1-\theta)^{\mu-1}
\chi_{R_1,R_2}\!\bigl(r+(q-r)\theta\bigr)
\notag\\
&\qquad\qquad\qquad\qquad\qquad\cdot
L_{XY}\!\bigl(t-h(r+(q-r)\theta),t-hq\bigr)\,d\theta.
\label{eq:edge-beta-XY}
\end{align}
Then there exists \(C_{t,m}<\infty\) such that
\begin{equation}\label{eq:edge-beta-bound}
\left|
g_{Y,t,h,R_1,R_2}^{\mathrm{edge}}(r,q)
\right|
+
\left|
g_{XY,t,h,R_1,R_2}^{\mathrm{edge}}(r,q)
\right|
\le
C_{t,m}\,\ell^{-m}(R_2-r)^m\,\rho_{R_1,R_2}(r,q)^{1-\lambda}
\end{equation}
for every \(r\in[R_1,R_2]\), every \(q\ge R_2\), and every sufficiently small \(h>0\) such that
\(hR_2<t\).
\end{proposition}

\begin{proof}
Because \(L_Y\) and \(L_{XY}\) are continuous on compact subsets of
\(\{0\le s\le u\le T\}\), one has
\[
\left|
L_Y\!\bigl(t-h(r+(q-r)\theta),t-hq\bigr)
\right|
+
\left|
L_{XY}\!\bigl(t-h(r+(q-r)\theta),t-hq\bigr)
\right|
\le
C_t
\]
uniformly for the indicated range of parameters. Apply Lemma~\ref{lem:edge-segment} with
\[
a=\lambda,
\qquad
b=\lambda,
\qquad
m=m,
\qquad
k=0
\]
to \eqref{eq:edge-beta-Y}, and with
\[
a=\lambda,
\qquad
b=\mu-1,
\qquad
m=m,
\qquad
k=0
\]
to \eqref{eq:edge-beta-XY}. Since \(\lambda<1\) and \(\mu-1>-1\), both choices are admissible and
yield \eqref{eq:edge-beta-bound}.
\end{proof}

\begin{proposition}[Derivative bounds for the outer-edge Beta averages]
\label{prop:edge-beta-deriv}
Assume the hypotheses of Proposition~\ref{prop:edge-beta}. Then there exists \(C_{t,m}<\infty\) such
that for every \(r\in[R_1,R_2]\), every \(q\ge R_2\), and every sufficiently small \(h>0\) such
that \(hR_2<t\),
\begin{equation}\label{eq:edge-beta-deriv}
\left|
\partial_r g_{Y,t,h,R_1,R_2}^{\mathrm{edge}}(r,q)
\right|
+
\left|
\partial_r g_{XY,t,h,R_1,R_2}^{\mathrm{edge}}(r,q)
\right|
\le
C_{t,m}(\ell^{-1}+h)\ell^{-m}(R_2-r)^m\rho_{R_1,R_2}(r,q)^{1-\lambda}.
\end{equation}
\end{proposition}

\begin{proof}
Differentiate \eqref{eq:edge-beta-Y} under the integral sign. Since
\[
\partial_r\bigl(r+(q-r)\theta\bigr)=1-\theta,
\]
one obtains
\begin{align*}
\partial_r g_{Y,t,h,R_1,R_2}^{\mathrm{edge}}(r,q)
&=
\frac{1}{\Gamma(\lambda)\Gamma(1-\lambda)}
\int_0^1
\theta^{-\lambda}(1-\theta)^{\lambda+1}
\chi_{R_1,R_2}'\!\bigl(r+(q-r)\theta\bigr)
\notag\\
&\qquad\qquad\qquad\qquad\qquad\cdot
L_Y\!\bigl(t-h(r+(q-r)\theta),t-hq\bigr)\,d\theta
\notag\\
&\quad-
\frac{h}{\Gamma(\lambda)\Gamma(1-\lambda)}
\int_0^1
\theta^{-\lambda}(1-\theta)^{\lambda+1}
\chi_{R_1,R_2}\!\bigl(r+(q-r)\theta\bigr)
\notag\\
&\qquad\qquad\qquad\qquad\qquad\cdot
\partial_1L_Y\!\bigl(t-h(r+(q-r)\theta),t-hq\bigr)\,d\theta.
\end{align*}
Because \(L_Y\) and \(\partial_1L_Y\) are uniformly bounded on the compact set under consideration, applying
Lemma~\ref{lem:edge-segment} with
\[
a=\lambda,\qquad b=\lambda+1,\qquad m=m,\qquad k=1
\]
to the first integral and with
\[
a=\lambda,\qquad b=\lambda+1,\qquad m=m,\qquad k=0
\]
to the second gives
\[
\left|
\partial_r g_{Y,t,h,R_1,R_2}^{\mathrm{edge}}(r,q)
\right|
\le
C_{t,m}(\ell^{-1}+h)\ell^{-m}(R_2-r)^m\rho_{R_1,R_2}(r,q)^{1-\lambda}.
\]

The source-side estimate is identical. Differentiating \eqref{eq:edge-beta-XY} yields
\begin{align*}
\partial_r g_{XY,t,h,R_1,R_2}^{\mathrm{edge}}(r,q)
&=
\frac{1}{\Gamma(\lambda)\Gamma(1-\lambda)}
\int_0^1
\theta^{-\lambda}(1-\theta)^{\mu}
\chi_{R_1,R_2}'\!\bigl(r+(q-r)\theta\bigr)
\notag\\
&\qquad\qquad\qquad\qquad\qquad\cdot
L_{XY}\!\bigl(t-h(r+(q-r)\theta),t-hq\bigr)\,d\theta
\notag\\
&\quad-
\frac{h}{\Gamma(\lambda)\Gamma(1-\lambda)}
\int_0^1
\theta^{-\lambda}(1-\theta)^{\mu}
\chi_{R_1,R_2}\!\bigl(r+(q-r)\theta\bigr)
\notag\\
&\qquad\qquad\qquad\qquad\qquad\cdot
\partial_1L_{XY}\!\bigl(t-h(r+(q-r)\theta),t-hq\bigr)\,d\theta.
\end{align*}
Now apply Lemma~\ref{lem:edge-segment} with
\[
a=\lambda,\qquad b=\mu,\qquad m=m,\qquad k=1
\]
and
\[
a=\lambda,\qquad b=\mu,\qquad m=m,\qquad k=0.
\]
Since \(\mu>0\), both choices are admissible, and they give the same bound for
\(\partial_r g_{XY,t,h,R_1,R_2}^{\mathrm{edge}}\). Combining the two estimates proves
\eqref{eq:edge-beta-deriv}.
\end{proof}

\begin{proposition}[Exact strip kernel representation for the outer-tail term]
\label{prop:edge-strip-kernel}
Assume the hypotheses of Definition~\ref{def:micro-dual}. Let
\[
0<R_1<R_2<\infty,
\qquad
\ell:=R_2-R_1,
\]
let \(\chi:=\chi_{R_1,R_2}\) be the canonical cutoff from Lemma~\ref{lem:quant-buffer}, and let
\[
\Phi\in\mathcal D(\mathfrak T_{t,h,\alpha}^{\mathrm{mic}}),
\qquad
\Psi:=\mathfrak T_{t,h,\alpha}^{\mathrm{mic}}\Phi.
\]
Fix \(m\in\mathbb N_0\). Assume
\[
hR_2<t,
\qquad
R_2<\varepsilon_{h,\alpha}^{-1}.
\]
Then there exist integral operators
\[
\mathfrak Q_{Y,t,h,R_1,R_2}^{\mathrm{edge}},
\mathfrak Q_{XY,t,h,R_1,R_2}^{\mathrm{edge}}
\in \mathcal L(F_{h,R_2},E_{R_2}),
\qquad
F_{h,R_2}:=L^2([R_2,\varepsilon_{h,\alpha}^{-1}]),
\]
such that
\begin{equation}\label{eq:edge-strip-identity}
\one_{[R_1,R_2]}
\mathfrak A_{t,h,\alpha,\chi}^{\mathrm{out}}(\Phi,\Psi)
=
\mathfrak Q_{Y,t,h,R_1,R_2}^{\mathrm{edge}}\Psi^{\mathrm{out}}
-
\mathfrak Q_{XY,t,h,R_1,R_2}^{\mathrm{edge}}\Phi^{\mathrm{out}}
\qquad\text{in }E_{R_2}.
\end{equation}
Moreover, these operators admit kernels \(K_{Y,t,h,R_1,R_2}^{\mathrm{edge}}\) and
\(K_{XY,t,h,R_1,R_2}^{\mathrm{edge}}\), supported in
\[
\{(r,q): R_1\le r\le R_2,\ q\ge R_2\},
\]
such that for \(r\in[R_1,R_2]\) and \(q\ge R_2\),
\begin{align}
\left|
K_{Y,t,h,R_1,R_2}^{\mathrm{edge}}(r,q)
\right|
&\le
C_{t,m}(\ell^{-1}+h)\ell^{-m}(R_2-r)^m\rho_{R_1,R_2}(r,q)^{1-\lambda},
\label{eq:edge-strip-Y}\\
\left|
K_{XY,t,h,R_1,R_2}^{\mathrm{edge}}(r,q)
\right|
&\le
C_{t,m}\ell^{-m}(R_2-r)^m(q-r)^{-\beta-1}\rho_{R_1,R_2}(r,q)^{1-\lambda}
\notag\\
&\quad+
C_{t,m}(\ell^{-1}+h)\ell^{-m}(R_2-r)^m(q-r)^{-\beta}\rho_{R_1,R_2}(r,q)^{1-\lambda}.
\label{eq:edge-strip-XY}
\end{align}
\end{proposition}

\begin{proof}
For \(r\in[R_1,R_2]\) and \(q\ge R_2\), one has \(\chi(q)=0\). Hence
\[
\chi(q)-\chi(r)=-\chi(r).
\]
Therefore, when restricted to the outer tail \(q\ge R_2\), the commutator Volterra operators from
Proposition~\ref{prop:micro-smooth-commutator} are exactly
\begin{align*}
\bigl(
\mathfrak W_{Y,t,h,\alpha,\chi}^{\mathrm{com}}\Psi^{\mathrm{out}}
\bigr)(r)
&=
-
\int_{R_2}^{\varepsilon_{h,\alpha}^{-1}}
(q-r)^{\lambda-1}
\chi(r)L_Y(t-hr,t-hq)\Psi^{\mathrm{out}}(q)\,dq,\\
\bigl(
\mathfrak W_{XY,t,h,\alpha,\chi}^{\mathrm{com}}\Phi^{\mathrm{out}}
\bigr)(r)
&=
-
\int_{R_2}^{\varepsilon_{h,\alpha}^{-1}}
(q-r)^{\mu-1}
\chi(r)L_{XY}(t-hr,t-hq)\Phi^{\mathrm{out}}(q)\,dq.
\end{align*}
Applying \eqref{eq:micro-smooth-com-conj-Y} and \eqref{eq:micro-smooth-com-conj-XY}, and then
restricting to \([R_1,R_2]\), gives \eqref{eq:edge-strip-identity} with the corresponding strip
operators.

For the target-side bound, fix \(q\ge R_2\) and set
\[
a_{Y,t,h,R_1,R_2}^{\mathrm{edge}}(r,q):=-\chi(r)L_Y(t-hr,t-hq).
\]
The same fixed-\(q\) computation as in Proposition~\ref{prop:micro-tail-Y} yields
\[
K_{Y,t,h,R_1,R_2}^{\mathrm{edge}}(r,q)
=
-
\frac{1}{\Gamma(\lambda)\Gamma(1-\lambda)}
\int_0^1
\theta^{-\lambda}(1-\theta)^\lambda
\partial_1 a_{Y,t,h,R_1,R_2}^{\mathrm{edge}}(r+(q-r)\theta,q)\,d\theta.
\]
Since
\[
\partial_1 a_{Y,t,h,R_1,R_2}^{\mathrm{edge}}(x,q)
=
-\chi'(x)L_Y(t-hx,t-hq)
+
h\chi(x)\partial_1L_Y(t-hx,t-hq),
\]
Lemma~\ref{lem:edge-segment}, applied with
\[
a=\lambda,\qquad b=\lambda,\qquad m=m,\qquad k=1
\]
to the first term and with
\[
a=\lambda,\qquad b=\lambda,\qquad m=m,\qquad k=0
\]
to the second, gives \eqref{eq:edge-strip-Y}.

For the source side, the same fixed-\(q\) computation as in Proposition~\ref{prop:micro-tail-XY}
gives
\[
K_{XY,t,h,R_1,R_2}^{\mathrm{edge}}(r,q)
=
\beta(q-r)^{-\beta-1}g_{XY,t,h,R_1,R_2}^{\mathrm{edge}}(r,q)
-
(q-r)^{-\beta}\partial_r g_{XY,t,h,R_1,R_2}^{\mathrm{edge}}(r,q).
\]
Now combine Proposition~\ref{prop:edge-beta} with Proposition~\ref{prop:edge-beta-deriv} to obtain
\eqref{eq:edge-strip-XY}. This completes the proof.
\end{proof}

\begin{remark}[Interpretation of the strip analysis]
\label{rem:edge-geometry}
Proposition~\ref{prop:edge-beta} identifies the first unavoidable geometric factor in every
outer-strip computation:
\[
\rho_{R_1,R_2}(r,q)^{1-\lambda}
=
\left(\frac{R_2-r}{q-r}\right)^{1-\lambda}.
\]
Proposition~\ref{prop:edge-strip-kernel} sharpens this reduction. After retaining arbitrary-order
flatness, every strip kernel carries the extra factor \(\ell^{-m}(R_2-r)^m\), but the source strip
kernel still contains the singular prefactors \((q-r)^{-\beta-1}\) and \((q-r)^{-\beta}\). So the
unresolved boundary-strip estimate is no longer about cutoff regularity alone; it is about how
that boundary vanishing interacts with those source-side singularities and, ultimately, with the
target/source cancellation inside the full defect.
\end{remark}

\section{Exact strip reduction to annulus data}
\label{sec:strip-annulus}

The explicit edge kernels present the strip term as an outer-tail object.  For a global exact micro
pair, the boundary-strip outer-tail contribution is determined entirely by the transition annulus
\([R_1,R_2]\).  The true tail \((R_2,\infty)\) disappears once the exact micro transfer identity is
used.

\begin{definition}[Annulus-driven preconjugated strip discrepancy]
\label{def:strip-annulus}
Assume the hypotheses of Definition~\ref{def:micro-dual}. Let
\[
0<R_1<R_2<\infty,
\]
set
\[
\lambda:=H_Y+\frac12,
\qquad
\mu:=H_{XY}+\frac12.
\]
Let \(\chi:=\chi_{R_1,R_2}\) be the canonical cutoff from Lemma~\ref{lem:quant-buffer}, and let
\[
\Theta_X,\Theta_Y\in L^2([R_1,R_2]).
\]
Define the annulus-driven preconjugated strip discrepancy by
\begin{equation}\label{eq:strip-annulus}
\bigl(
\mathfrak U_{t,h,\alpha,R_1,R_2}^{\mathrm{ann}}(\Theta_X,\Theta_Y)
\bigr)(r)
:=
\chi(r)
\int_r^{R_2}
(q-r)^{\lambda-1}
L_Y(t-hr,t-hq)\Theta_Y(q)\,dq
\end{equation}
for \(r\in[R_1,R_2]\), minus
\[
\chi(r)
\int_r^{R_2}
(q-r)^{\mu-1}
L_{XY}(t-hr,t-hq)\Theta_X(q)\,dq,
\]
and set
\[
\mathfrak U_{t,h,\alpha,R_1,R_2}^{\mathrm{ann}}(\Theta_X,\Theta_Y)(r):=0
\qquad\text{for }r\notin [R_1,R_2].
\]
\end{definition}

\begin{proposition}[Exact strip outer-tail reduction to annulus data]
\label{prop:strip-annulus}
Assume the hypotheses of Definition~\ref{def:micro-dual}. Let
\[
0<R_1<R_2<\infty,
\qquad
\chi:=\chi_{R_1,R_2},
\]
and let
\[
\Phi\in\mathcal D(\mathfrak T_{t,h,\alpha}^{\mathrm{mic}}),
\qquad
\Psi:=\mathfrak T_{t,h,\alpha}^{\mathrm{mic}}\Phi.
\]
Define
\[
\Phi^{\mathrm{ann}}:=\Phi\one_{[R_1,R_2]},
\qquad
\Phi^{\mathrm{out}}:=\Phi\one_{(R_2,\varepsilon_{h,\alpha}^{-1}]},
\]
\[
\Psi^{\mathrm{ann}}:=\Psi\one_{[R_1,R_2]},
\qquad
\Psi^{\mathrm{out}}:=\Psi\one_{(R_2,\varepsilon_{h,\alpha}^{-1}]}.
\]
Assume
\[
hR_2<t,
\qquad
R_2<\varepsilon_{h,\alpha}^{-1}.
\]
Then
\begin{equation}\label{eq:strip-annulus-identity}
\one_{[R_1,R_2]}
\Bigl(
\mathfrak W_{Y,t,h,\alpha,\chi}^{\mathrm{com}}\Psi^{\mathrm{out}}
-
\mathfrak W_{XY,t,h,\alpha,\chi}^{\mathrm{com}}\Phi^{\mathrm{out}}
\Bigr)
=
\mathfrak U_{t,h,\alpha,R_1,R_2}^{\mathrm{ann}}
\bigl(\Phi^{\mathrm{ann}},\Psi^{\mathrm{ann}}\bigr)
\qquad\text{in }E_{R_2}.
\end{equation}
In particular, the boundary-strip outer-tail defect
\[
\one_{[R_1,R_2]}
\mathfrak A_{t,h,\alpha,\chi}^{\mathrm{out}}(\Phi,\Psi)
\]
depends only on the annulus restrictions \(\Phi^{\mathrm{ann}}\) and \(\Psi^{\mathrm{ann}}\).
\end{proposition}

\begin{proof}
Set \(\varepsilon:=\varepsilon_{h,\alpha}\). By Corollary~\ref{cor:micro-volterra-transfer},
\[
\mathfrak V_{Y,t,h,\alpha}^{\mathrm{mic}}\Psi
=
\mathfrak V_{XY,t,h,\alpha}^{\mathrm{mic}}\Phi
\qquad\text{in }L^2(\R_+).
\]
Fix \(r\in[R_1,R_2]\). Since
\[
\Psi=\Psi^{\mathrm{ann}}+\Psi^{\mathrm{out}}
\qquad\text{and}\qquad
\Phi=\Phi^{\mathrm{ann}}+\Phi^{\mathrm{out}}
\]
on \([r,\varepsilon^{-1}]\), and there is no contribution from \([0,R_1]\) because \(q\ge r\ge R_1\),
the exact micro transfer identity gives
\begin{align*}
\int_r^{R_2}
(q-r)^{\lambda-1}
L_Y(t-hr,t-hq)\Psi^{\mathrm{ann}}(q)\,dq
&+
\int_{R_2}^{\varepsilon^{-1}}
(q-r)^{\lambda-1}
L_Y(t-hr,t-hq)\Psi^{\mathrm{out}}(q)\,dq\\
&=\;
\int_r^{R_2}
(q-r)^{\mu-1}
L_{XY}(t-hr,t-hq)\Phi^{\mathrm{ann}}(q)\,dq
\\
&\qquad
+\int_{R_2}^{\varepsilon^{-1}}
(q-r)^{\mu-1}
\\
&\qquad\qquad \times
L_{XY}(t-hr,t-hq)\Phi^{\mathrm{out}}(q)\,dq.
\end{align*}
Rearranging yields
\begin{align*}
&-\chi(r)
\int_{R_2}^{\varepsilon^{-1}}
(q-r)^{\lambda-1}
L_Y(t-hr,t-hq)\Psi^{\mathrm{out}}(q)\,dq\\
&\qquad+
\chi(r)
\int_{R_2}^{\varepsilon^{-1}}
(q-r)^{\mu-1}
L_{XY}(t-hr,t-hq)\Phi^{\mathrm{out}}(q)\,dq\\
&=
\chi(r)
\int_r^{R_2}
(q-r)^{\lambda-1}
L_Y(t-hr,t-hq)\Psi^{\mathrm{ann}}(q)\,dq\\
&\qquad-
\chi(r)
\int_r^{R_2}
(q-r)^{\mu-1}
L_{XY}(t-hr,t-hq)\Phi^{\mathrm{ann}}(q)\,dq.
\end{align*}
Because \(\chi(q)=0\) for \(q\ge R_2\), the left-hand side is exactly
\[
\bigl(
\mathfrak W_{Y,t,h,\alpha,\chi}^{\mathrm{com}}\Psi^{\mathrm{out}}
-
\mathfrak W_{XY,t,h,\alpha,\chi}^{\mathrm{com}}\Phi^{\mathrm{out}}
\bigr)(r),
\]
while the right-hand side is
\[
\bigl(
\mathfrak U_{t,h,\alpha,R_1,R_2}^{\mathrm{ann}}
(\Phi^{\mathrm{ann}},\Psi^{\mathrm{ann}})
\bigr)(r).
\]
This proves \eqref{eq:strip-annulus-identity} for almost every \(r\in[R_1,R_2]\), and both sides
vanish outside \([R_1,R_2]\) by definition.

Finally, \(\one_{[R_1,R_2]}\mathfrak A_{t,h,\alpha,\chi}^{\mathrm{out}}(\Phi,\Psi)\) is obtained
from the left-hand side of \eqref{eq:strip-annulus-identity} by the fixed conjugation formulas
\eqref{eq:micro-smooth-com-conj-Y} and \eqref{eq:micro-smooth-com-conj-XY}. Therefore it is
completely determined by \(\Phi^{\mathrm{ann}}\) and \(\Psi^{\mathrm{ann}}\) as claimed.
\end{proof}

\begin{corollary}[Reduction of the strip problem to annulus-local data]
\label{cor:strip-annulus}
Assume the hypotheses of Corollary~\ref{cor:outer-tail-reduction}. Then the localization hypothesis
in Theorem~\ref{thm:abstract-global-reduction} will follow once one proves the following statement:

\smallskip
\noindent
for every structured non-escaping graph-unit sequence and every \(\varepsilon>0\), there exist
common radii
\[
0<R_0<R_1<R_2<\infty
\]
and the canonical cutoff \(\chi_{R_1,R_2}\) such that, after passage to a subsequence,
\begin{equation}\label{eq:strip-annulus-small}
\left\|
\one_{[R_1,R_2]}
\mathfrak A_{t,h_n,\alpha,\chi_{R_1,R_2}}^{\mathrm{out}}(\Phi_n,\Psi_n)
\right\|_{E_{R_2}}
\le
\varepsilon
\qquad\text{for every }n,
\end{equation}
with the understanding from Proposition~\ref{prop:strip-annulus} that the left-hand side is now a
functional only of the common annulus restrictions
\[
\Phi_n\one_{[R_1,R_2]},
\qquad
\Psi_n\one_{[R_1,R_2]}.
\]
In particular, no genuine outer-tail degree of freedom remains in the final strip estimate.
\end{corollary}

\begin{proof}
Corollary~\ref{cor:outer-tail-reduction}, combined with
Proposition~\ref{prop:strip-annulus}, gives the reduction.  The exact micro transfer identity
changes the status of \eqref{eq:outer-tail-strip}: it is no longer a tail-forcing quantity, but an
annulus-local quantity.
\end{proof}

\begin{definition}[Admissible annulus traces and exact strip operator]
\label{def:annulus-strip-op}
Assume the hypotheses of Definition~\ref{def:micro-dual}. Let
\[
0<R_1<R_2<\infty,
\qquad
\chi:=\chi_{R_1,R_2}.
\]
Define the admissible annulus-trace class
\[
\mathcal A_{t,h,\alpha,R_1,R_2}^{\mathrm{ann}}
\subset
L^2([R_1,R_2])\times L^2([R_1,R_2])
\]
to consist of all pairs \((\Theta_X,\Theta_Y)\) for which there exists
\[
\Phi\in\mathcal D(\mathfrak T_{t,h,\alpha}^{\mathrm{mic}}),
\qquad
\Psi:=\mathfrak T_{t,h,\alpha}^{\mathrm{mic}}\Phi,
\]
such that
\[
\Theta_X=\Phi\one_{[R_1,R_2]},
\qquad
\Theta_Y=\Psi\one_{[R_1,R_2]}.
\]
For every \((\Theta_X,\Theta_Y)\in \mathcal A_{t,h,\alpha,R_1,R_2}^{\mathrm{ann}}\), define
\begin{equation}\label{eq:annulus-strip-op}
\mathfrak S_{t,h,\alpha,R_1,R_2}^{\mathrm{ann}}(\Theta_X,\Theta_Y)
:=
\one_{[R_1,R_2]}
\mathfrak A_{t,h,\alpha,\chi}^{\mathrm{out}}(\Phi,\Psi)
\in E_{R_2},
\end{equation}
where \((\Phi,\Psi)\) is any global exact micro pair realizing the annulus traces
\((\Theta_X,\Theta_Y)\).
\end{definition}

\begin{proposition}[Well-definedness of the exact annulus strip operator]
\label{prop:annulus-strip-op}
Assume the hypotheses of Definition~\ref{def:annulus-strip-op}. Then the definition
\eqref{eq:annulus-strip-op} is independent of the chosen realizing global exact micro pair. In
other words,
\[
\mathfrak S_{t,h,\alpha,R_1,R_2}^{\mathrm{ann}}
:
\mathcal A_{t,h,\alpha,R_1,R_2}^{\mathrm{ann}}
\to
E_{R_2}
\]
is well-defined.
\end{proposition}

\begin{proof}
Let \((\Phi^{(1)},\Psi^{(1)})\) and \((\Phi^{(2)},\Psi^{(2)})\) be two global exact micro pairs
with the same annulus traces:
\[
\Phi^{(1)}\one_{[R_1,R_2]}=\Phi^{(2)}\one_{[R_1,R_2]},
\qquad
\Psi^{(1)}\one_{[R_1,R_2]}=\Psi^{(2)}\one_{[R_1,R_2]}.
\]
Then Proposition~\ref{prop:strip-annulus} applies to both pairs and gives
\begin{align*}
\one_{[R_1,R_2]}
\mathfrak A_{t,h,\alpha,\chi}^{\mathrm{out}}(\Phi^{(1)},\Psi^{(1)})
&=
\mathfrak U_{t,h,\alpha,R_1,R_2}^{\mathrm{ann}}
\bigl(
\Phi^{(1)}\one_{[R_1,R_2]},
\Psi^{(1)}\one_{[R_1,R_2]}
\bigr),\\
\one_{[R_1,R_2]}
\mathfrak A_{t,h,\alpha,\chi}^{\mathrm{out}}(\Phi^{(2)},\Psi^{(2)})
&=
\mathfrak U_{t,h,\alpha,R_1,R_2}^{\mathrm{ann}}
\bigl(
\Phi^{(2)}\one_{[R_1,R_2]},
\Psi^{(2)}\one_{[R_1,R_2]}
\bigr).
\end{align*}
The right-hand sides are equal because the annulus traces agree. Hence the two realizations of
\eqref{eq:annulus-strip-op} coincide.
\end{proof}

\begin{theorem}[Annulus-strip boundedness reduces the global problem]
\label{thm:annulus-strip-reduction}
Assume the hypotheses of Definition~\ref{def:micro-dual}. Suppose that the following annulus-strip
boundedness property holds:

\smallskip
\noindent
for every \(\ell>0\) there exists \(C_{t,\ell}<\infty\) such that for every sufficiently small
\(h>0\), every pair of radii
\[
0<R_1<R_2<\infty,
\qquad
R_2-R_1=\ell,
\]
and every admissible annulus trace
\[
(\Theta_X,\Theta_Y)\in \mathcal A_{t,h,\alpha,R_1,R_2}^{\mathrm{ann}},
\]
one has
\begin{equation}\label{eq:annulus-strip-bounded}
\left\|
\mathfrak S_{t,h,\alpha,R_1,R_2}^{\mathrm{ann}}(\Theta_X,\Theta_Y)
\right\|_{E_{R_2}}
\le
C_{t,\ell}
\bigl(
\|\Theta_X\|_{L^2([R_1,R_2])}
+
\|\Theta_Y\|_{L^2([R_1,R_2])}
\bigr).
\end{equation}

\smallskip
\noindent
Then the localization hypothesis in Theorem~\ref{thm:abstract-global-reduction} holds. In
particular,
\begin{equation}\label{eq:annulus-strip-global-collapse}
\mathfrak D_{t,h,\alpha}^{\mathrm{mic}}\longrightarrow0
\qquad\text{as }h\downarrow0.
\end{equation}
\end{theorem}

\begin{proof}
By Theorem~\ref{thm:abstract-global-reduction}, it is enough to verify the localization hypothesis.
By Corollary~\ref{cor:defect-localization}, it is enough to verify the full-window defect control
required there. By Corollary~\ref{cor:structured-reduction}, it is enough to treat structured
non-escaping graph-unit sequences
\[
\Phi_n\in\mathcal D(\mathfrak T_{t,h_n,\alpha}^{\mathrm{mic}}),
\qquad
\Psi_n:=\mathfrak T_{t,h_n,\alpha}^{\mathrm{mic}}\Phi_n,
\qquad
\|\Phi_n\|_{\mathfrak G_{t,h_n,\alpha}^{\mathrm{mic}}}\le1.
\]

Fix \(\varepsilon>0\). Choose \(\ell>0\) so large that
\[
C_t\ell^{-\beta}\le \varepsilon,
\]
where \(C_t\) is the constant from Proposition~\ref{prop:outer-tail-separated}. Let
\(C_{t,\ell}\) be the constant from the annulus-strip boundedness hypothesis
\eqref{eq:annulus-strip-bounded}, and choose \(M\in\mathbb N\) so large that
\[
C_{t,\ell}\left(1+\frac{\nu_Y}{\eta}\right)M^{-1/2}\le \varepsilon.
\]

Applying Corollary~\ref{cor:profile-tight-buffer} with this \(M\), this \(\ell\), and this
\(\varepsilon\), and then relabeling the resulting subsequence, yields common radii
\[
0<R_0<R_1<R_2<\infty,
\qquad
R_1-R_0\ge \ell,
\qquad
R_2-R_1=\ell,
\]
such that
\[
\left\|
\one_{(R_0,\infty)}
\mathfrak g_{i,t,h_n,\alpha}^{\mathrm{mic}}
\right\|_{L^2(\R_+)}
\le
\varepsilon
\qquad\text{for }i\in\{Y,XY\},
\]
and
\[
\|\Phi_n\|_{L^2([R_1,R_2])}^2
+
\frac{\eta^2}{\nu_Y^2}
\|\Psi_n\|_{L^2([R_1,R_2])}^2
\le
\frac{1}{M}
\qquad\text{for every }n.
\]
Since \(R_1>R_0\), the same profile-tail bound holds beyond \(R_1\). Also,
\[
\|\Phi_n\|_{L^2([R_1,R_2])}
+
\|\Psi_n\|_{L^2([R_1,R_2])}
\le
\left(1+\frac{\nu_Y}{\eta}\right)M^{-1/2}.
\]

After discarding finitely many terms, we may also assume
\[
h_nR_2<t,
\qquad
R_2<\varepsilon_{h_n,\alpha}^{-1},
\qquad
C_{t,\alpha,\eta,\nu_Y}h_n^\alpha\le \varepsilon
\qquad\text{for every }n,
\]
where \(C_{t,\alpha,\eta,\nu_Y}\) is the constant from
Proposition~\ref{prop:outer-tail-separated}.

By Corollary~\ref{cor:annulus-control},
\[
\left\|
\mathfrak A_{t,h_n,\alpha,\chi_{R_1,R_2}}^{\mathrm{ann}}(\Phi_n,\Psi_n)
\right\|_{E_{R_2}}
\le
\varepsilon
\qquad\text{for every }n.
\]
By Proposition~\ref{prop:outer-tail-separated},
\[
\left\|
\one_{[0,R_1]}
\mathfrak A_{t,h_n,\alpha,\chi_{R_1,R_2}}^{\mathrm{out}}(\Phi_n,\Psi_n)
\right\|_{E_{R_2}}
\le
2\varepsilon
\qquad\text{for every }n.
\]
Finally, Definition~\ref{def:annulus-strip-op}, Proposition~\ref{prop:annulus-strip-op}, and the
boundedness hypothesis \eqref{eq:annulus-strip-bounded} give
\begin{align*}
\left\|
\one_{[R_1,R_2]}
\mathfrak A_{t,h_n,\alpha,\chi_{R_1,R_2}}^{\mathrm{out}}(\Phi_n,\Psi_n)
\right\|_{E_{R_2}}
&=
\left\|
\mathfrak S_{t,h_n,\alpha,R_1,R_2}^{\mathrm{ann}}
\bigl(
\Phi_n\one_{[R_1,R_2]},
\Psi_n\one_{[R_1,R_2]}
\bigr)
\right\|_{E_{R_2}}\\
&\le
C_{t,\ell}
\bigl(
\|\Phi_n\|_{L^2([R_1,R_2])}
+
\|\Psi_n\|_{L^2([R_1,R_2])}
\bigr)\\
&\le
\varepsilon
\qquad\text{for every }n.
\end{align*}
Therefore
\[
\left\|
\mathfrak A_{t,h_n,\alpha,\chi_{R_1,R_2}}^{\mathrm{out}}(\Phi_n,\Psi_n)
\right\|_{E_{R_2}}
\le
3\varepsilon
\qquad\text{for every }n,
\]
and Proposition~\ref{prop:defect-split} gives
\[
\left\|
\mathfrak F_{t,h_n,\alpha,R_2}^{\mathrm{ex}}
\bigl(
\chi_{R_1,R_2}\Phi_n,\chi_{R_1,R_2}\Psi_n
\bigr)
\right\|_{E_{R_2}}
\le
4\varepsilon
\qquad\text{for every }n.
\]
By Theorem~\ref{thm:defect-localization}, this yields exact fixed-window approximants with
arbitrarily small loss in norm and functional value. Hence the localization hypothesis in
Theorem~\ref{thm:abstract-global-reduction} holds, and \eqref{eq:annulus-strip-global-collapse}
follows.
\end{proof}

\section{Fixed-width normalization of the annulus problem}
\label{sec:annulus-normalization}

The strip quantity lives on one annulus of fixed width \(\ell=R_2-R_1\).  After translation, it
becomes a single operator family on the fixed window \([0,\ell]\).

\begin{definition}[Translated annulus traces and fixed-width strip operator]
\label{def:annulus-normalized}
Assume the hypotheses of Definition~\ref{def:annulus-strip-op}. Let
\[
0<R_1<R_2<\infty,
\qquad
\ell:=R_2-R_1.
\]
For \(\Theta\in L^2([R_1,R_2])\), define the translation isometry
\[
(\mathfrak T_{R_1}\Theta)(u):=\Theta(R_1+u),
\qquad 0\le u\le \ell.
\]
Define the translated admissible trace class
\[
\widetilde{\mathcal A}_{t,h,\alpha,R_1,\ell}^{\mathrm{ann}}
\subset
L^2([0,\ell])\times L^2([0,\ell])
\]
by
\[
(\widetilde\Theta_X,\widetilde\Theta_Y)\in
\widetilde{\mathcal A}_{t,h,\alpha,R_1,\ell}^{\mathrm{ann}}
\]
if and only if
\[
(\mathfrak T_{R_1}^{-1}\widetilde\Theta_X,\mathfrak T_{R_1}^{-1}\widetilde\Theta_Y)
\in
\mathcal A_{t,h,\alpha,R_1,R_2}^{\mathrm{ann}}.
\]
For such a pair, define
\begin{equation}\label{eq:annulus-normalized-op}
\widetilde{\mathfrak S}_{t,h,\alpha,R_1,\ell}^{\mathrm{ann}}
(\widetilde\Theta_X,\widetilde\Theta_Y)
:=
\mathfrak T_{R_1}
\mathfrak S_{t,h,\alpha,R_1,R_2}^{\mathrm{ann}}
(
\mathfrak T_{R_1}^{-1}\widetilde\Theta_X,
\mathfrak T_{R_1}^{-1}\widetilde\Theta_Y
)
\in
L^2([0,\ell]).
\end{equation}
\end{definition}

\begin{proposition}[Fixed-width normalization of the annulus-strip problem]
\label{prop:annulus-normalized}
Assume the hypotheses of Definition~\ref{def:annulus-normalized}. Then:
\begin{enumerate}[label=(\roman*)]
\item the canonical cutoff translates to the fixed profile
\[
\widetilde\chi_\ell(u):=\chi_{R_1,R_2}(R_1+u)=\vartheta\!\left(\frac{u}{\ell}\right),
\qquad 0\le u\le \ell,
\]
where \(\vartheta\) is the fixed template from Lemma~\ref{lem:quant-buffer};
\item for \(i\in\{Y,XY\}\), the translated coefficients
\[
\widetilde L_{i,t,h,R_1,\ell}(u,v)
:=
L_i\bigl(t-h(R_1+u),t-h(R_1+v)\bigr),
\qquad 0\le u\le v\le \ell,
\]
are uniformly \(C^1\) on \(\Delta_\ell:=\{0\le u\le v\le \ell\}\), with bounds depending only on
\(t\) and \(\ell\), provided \(h(R_1+\ell)<t\);
\item for every admissible translated trace pair
\[
(\widetilde\Theta_X,\widetilde\Theta_Y)\in
\widetilde{\mathcal A}_{t,h,\alpha,R_1,\ell}^{\mathrm{ann}},
\]
the fixed-width strip operator is given by the same change of variables applied to the
annulus-driven discrepancy:
\begin{align}
\bigl(
\widetilde{\mathfrak U}_{t,h,\alpha,R_1,\ell}^{\mathrm{ann}}
(\widetilde\Theta_X,\widetilde\Theta_Y)
\bigr)(u)
&:=
\widetilde\chi_\ell(u)
\int_u^\ell
(v-u)^{\lambda-1}
\widetilde L_{Y,t,h,R_1,\ell}(u,v)\widetilde\Theta_Y(v)\,dv
\notag\\
&\quad-
\widetilde\chi_\ell(u)
\int_u^\ell
(v-u)^{\mu-1}
\widetilde L_{XY,t,h,R_1,\ell}(u,v)\widetilde\Theta_X(v)\,dv,
\label{eq:annulus-normalized-U}
\end{align}
and proving either \eqref{eq:annulus-strip-bounded} or \eqref{eq:strip-annulus-small} is
equivalent to proving the corresponding statement for the family
\(\widetilde{\mathfrak S}_{t,h,\alpha,R_1,\ell}^{\mathrm{ann}}\) on the fixed window \([0,\ell]\).
\end{enumerate}
\end{proposition}

\begin{proof}
Item~(i) follows from the construction in Lemma~\ref{lem:quant-buffer}:
\[
\chi_{R_1,R_2}(R_1+u)
=
\vartheta\!\left(\frac{(R_1+u)-R_1}{R_2-R_1}\right)
=
\vartheta\!\left(\frac{u}{\ell}\right).
\]
For item~(ii), the bounds on \(\widetilde L_{i,t,h,R_1,\ell}\) and their first derivatives follow
from Assumption~\ref{ass:model}, the condition \(h(R_1+\ell)<t\), and the fact that
\((u,v)\in\Delta_\ell\) ranges over a set of fixed diameter \(\ell\). Differentiation in \(u\) or
\(v\) only produces one factor of \(h\), which is uniformly controlled for \(0<h\le1\).

For item~(iii), let
\[
\Theta_X:=\mathfrak T_{R_1}^{-1}\widetilde\Theta_X,
\qquad
\Theta_Y:=\mathfrak T_{R_1}^{-1}\widetilde\Theta_Y.
\]
By Definition~\ref{def:annulus-normalized},
\[
(\Theta_X,\Theta_Y)\in \mathcal A_{t,h,\alpha,R_1,R_2}^{\mathrm{ann}}.
\]
Applying Proposition~\ref{prop:strip-annulus} on the annulus \([R_1,R_2]\), and then changing
variables \(r=R_1+u\), \(q=R_1+v\), gives exactly \eqref{eq:annulus-normalized-U}. The operator
identity \eqref{eq:annulus-normalized-op} is the same translation at the conjugated level.
Since \(\mathfrak T_{R_1}\) is an \(L^2\)-isometry, every estimate on
\(\mathfrak S_{t,h,\alpha,R_1,R_2}^{\mathrm{ann}}\) is equivalent to the corresponding estimate on
\(\widetilde{\mathfrak S}_{t,h,\alpha,R_1,\ell}^{\mathrm{ann}}\), and likewise for the
sequence-level condition \eqref{eq:strip-annulus-small}.
\end{proof}

\section{Frozen limit of the normalized annulus discrepancy}
\label{sec:annulus-frozen}

Fixed-width normalization removes the moving geometry.  With \(\ell\) fixed, only the slowly
varying coefficients \(L_Y\) and \(L_{XY}\) still depend on the global location \(R_1\) and on
\(h\).  On common finite buffers, those coefficients freeze to their diagonal values.  The next
theorem identifies the corresponding constant-coefficient limit of the annulus-driven discrepancy.

\begin{definition}[Frozen normalized annulus discrepancy]
\label{def:annulus-frozen}
Assume the hypotheses of Definition~\ref{def:annulus-normalized}, and fix \(\ell>0\). For
\[
\lambda:=H_Y+\frac12,
\qquad
\mu:=H_{XY}+\frac12,
\]
and
\[
(\widetilde\Theta_X,\widetilde\Theta_Y)\in L^2([0,\ell])\times L^2([0,\ell]),
\]
define the frozen normalized annulus discrepancy by
\begin{align}
\bigl(
\widetilde{\mathfrak U}_{t,\ell}^{\mathrm{fr}}
(\widetilde\Theta_X,\widetilde\Theta_Y)
\bigr)(u)
&:=
\widetilde\chi_\ell(u)\,L_Y(t,t)
\int_u^\ell
(v-u)^{\lambda-1}\widetilde\Theta_Y(v)\,dv
\notag\\
&\quad-
\widetilde\chi_\ell(u)\,L_{XY}(t,t)
\int_u^\ell
(v-u)^{\mu-1}\widetilde\Theta_X(v)\,dv
\label{eq:annulus-frozen}
\end{align}
for \(0\le u\le \ell\), where \(\widetilde\chi_\ell\) is the fixed cutoff profile from
Proposition~\ref{prop:annulus-normalized}(i).
\end{definition}

\begin{theorem}[Strong freezing on each fixed annulus width]
\label{thm:annulus-frozen}
Assume Assumption~\ref{ass:model}, fix \(t\in(0,T)\), \(\alpha\in(0,1)\), and suppose
\[
0<H_{XY}<H_Y<\tfrac12.
\]
Fix \(\ell>0\) and \(R_*<\infty\). Then there exist \(h_0\in(0,1]\) and \(C_{t,\ell,R_*}<\infty\)
such that for every
\[
0\le R_1\le R_*,
\qquad
0<h\le h_0,
\qquad
h(R_1+\ell)<t,
\]
and every
\[
(\widetilde\Theta_X,\widetilde\Theta_Y)\in L^2([0,\ell])\times L^2([0,\ell]),
\]
one has
\begin{equation}\label{eq:annulus-frozen-bound}
\left\|
\widetilde{\mathfrak U}_{t,h,\alpha,R_1,\ell}^{\mathrm{ann}}
(\widetilde\Theta_X,\widetilde\Theta_Y)
-
\widetilde{\mathfrak U}_{t,\ell}^{\mathrm{fr}}
(\widetilde\Theta_X,\widetilde\Theta_Y)
\right\|_{L^2([0,\ell])}
\le
C_{t,\ell,R_*}\,h
\bigl(
\|\widetilde\Theta_X\|_{L^2([0,\ell])}
+
\|\widetilde\Theta_Y\|_{L^2([0,\ell])}
\bigr).
\end{equation}
In particular, for every fixed \((\widetilde\Theta_X,\widetilde\Theta_Y)\in
L^2([0,\ell])\times L^2([0,\ell])\),
\begin{equation}\label{eq:annulus-frozen-conv}
\widetilde{\mathfrak U}_{t,h,\alpha,R_1,\ell}^{\mathrm{ann}}
(\widetilde\Theta_X,\widetilde\Theta_Y)
\longrightarrow
\widetilde{\mathfrak U}_{t,\ell}^{\mathrm{fr}}
(\widetilde\Theta_X,\widetilde\Theta_Y)
\qquad\text{in }L^2([0,\ell])
\end{equation}
uniformly for \(R_1\in[0,R_*]\).
\end{theorem}

\begin{proof}
Choose \(h_0\in(0,1]\) so small that
\[
h_0(R_*+\ell)<t.
\]
By Proposition~\ref{prop:annulus-normalized}(iii),
\begin{align*}
&\widetilde{\mathfrak U}_{t,h,\alpha,R_1,\ell}^{\mathrm{ann}}
(\widetilde\Theta_X,\widetilde\Theta_Y)
-
\widetilde{\mathfrak U}_{t,\ell}^{\mathrm{fr}}
(\widetilde\Theta_X,\widetilde\Theta_Y)\\
&=
\widetilde\chi_\ell(\cdot)
\int_{\cdot}^{\ell}
(v-\cdot)^{\lambda-1}
\bigl(
\widetilde L_{Y,t,h,R_1,\ell}(\cdot,v)-L_Y(t,t)
\bigr)\widetilde\Theta_Y(v)\,dv\\
&\quad-
\widetilde\chi_\ell(\cdot)
\int_{\cdot}^{\ell}
(v-\cdot)^{\mu-1}
\bigl(
\widetilde L_{XY,t,h,R_1,\ell}(\cdot,v)-L_{XY}(t,t)
\bigr)\widetilde\Theta_X(v)\,dv.
\end{align*}
Since \(L_Y\) and \(L_{XY}\) are \(C^1\) on compact subsets of \(\{0\le s\le u\le T\}\), and
\[
t-h(R_1+u)\in [t-h_0(R_*+\ell),t],
\qquad
t-h(R_1+v)\in [t-h_0(R_*+\ell),t],
\]
the mean-value theorem gives
\begin{align}
\sup_{(u,v)\in\Delta_\ell}
\left|
\widetilde L_{Y,t,h,R_1,\ell}(u,v)-L_Y(t,t)
\right|
&\le
C_{t,R_*+\ell}\,h,
\label{eq:annulus-freeze-Y}\\
\sup_{(u,v)\in\Delta_\ell}
\left|
\widetilde L_{XY,t,h,R_1,\ell}(u,v)-L_{XY}(t,t)
\right|
&\le
C_{t,R_*+\ell}\,h.
\label{eq:annulus-freeze-XY}
\end{align}
Define Volterra kernels on \(\Delta_\ell:=\{0\le u\le v\le \ell\}\) by
\[
K_{Y,\ell}(u,v):=\widetilde\chi_\ell(u)(v-u)^{\lambda-1},
\qquad
K_{XY,\ell}(u,v):=\widetilde\chi_\ell(u)(v-u)^{\mu-1}.
\]
Because
\[
\lambda-1=H_Y-\frac12>-\frac12,
\qquad
\mu-1=H_{XY}-\frac12>-\frac12,
\]
both kernels belong to \(L^2(\Delta_\ell)\). Hence the corresponding Volterra operators are
Hilbert--Schmidt on \(L^2([0,\ell])\), and there exists \(C_\ell<\infty\) such that
\begin{align}
\left\|
\widetilde\chi_\ell(\cdot)
\int_{\cdot}^{\ell}
(v-\cdot)^{\lambda-1}f(v)\,dv
\right\|_{L^2([0,\ell])}
&\le
C_\ell\|f\|_{L^2([0,\ell])},
\label{eq:annulus-freeze-op-Y}\\
\left\|
\widetilde\chi_\ell(\cdot)
\int_{\cdot}^{\ell}
(v-\cdot)^{\mu-1}f(v)\,dv
\right\|_{L^2([0,\ell])}
&\le
C_\ell\|f\|_{L^2([0,\ell])}.
\label{eq:annulus-freeze-op-XY}
\end{align}
Combining \eqref{eq:annulus-freeze-Y}, \eqref{eq:annulus-freeze-XY},
\eqref{eq:annulus-freeze-op-Y}, and \eqref{eq:annulus-freeze-op-XY} yields
\eqref{eq:annulus-frozen-bound}. The convergence statement \eqref{eq:annulus-frozen-conv} follows
from the same estimates after letting the coefficient-freezing errors tend to zero.
\end{proof}

\begin{remark}[Role of the frozen strip estimate]
\label{rem:annulus-frozen}
Theorem~\ref{thm:annulus-frozen} identifies the strip contribution with a fixed constant-coefficient
model on every common finite buffer, but only at the preconjugated level. The subsequent
conjugated-strip analysis upgrades this freezing statement to the functional needed for
\eqref{eq:strip-annulus-small}, where the exact admissibility of the annulus traces enforces the
cancellation used in the global collapse argument.
\end{remark}

\section{Conjugated freezing of the normalized strip operator}
\label{sec:annulus-conjugated}

The observable strip term is the right-fractional conjugate of the preconjugated discrepancy from
Section~\ref{sec:annulus-frozen}, using the same conjugation that appears throughout the analysis.
On a fixed-width annulus this conjugated problem is stable: after freezing the slowly varying
coefficients, the exact strip operator differs from a frozen constant-coefficient strip operator by
\(O(h)\).

\begin{definition}[Normalized annulus Volterra pieces and frozen strip operator]
\label{def:annulus-conjugated}
Assume Assumption~\ref{ass:model}, fix \(t\in(0,T)\), and suppose
\[
0<H_{XY}<H_Y<\tfrac12.
\]
Set
\[
\lambda:=H_Y+\frac12,
\qquad
\mu:=H_{XY}+\frac12,
\qquad
\ell>0,
\qquad
E_\ell:=L^2([0,\ell]).
\]
For \(R_1\ge0\) and \(h>0\) satisfying \(h(R_1+\ell)<t\), define on \(E_\ell\)
\begin{align}
\bigl(
\widetilde{\mathfrak V}_{Y,t,h,R_1,\ell}^{\mathrm{ann}}\Theta
\bigr)(u)
&:=
\widetilde\chi_\ell(u)
\int_u^\ell
(v-u)^{\lambda-1}
\widetilde L_{Y,t,h,R_1,\ell}(u,v)\Theta(v)\,dv,
\label{eq:annulus-VY}\\
\bigl(
\widetilde{\mathfrak V}_{XY,t,h,R_1,\ell}^{\mathrm{ann}}\Theta
\bigr)(u)
&:=
\widetilde\chi_\ell(u)
\int_u^\ell
(v-u)^{\mu-1}
\widetilde L_{XY,t,h,R_1,\ell}(u,v)\Theta(v)\,dv,
\label{eq:annulus-VXY}
\end{align}
and the frozen counterparts
\begin{align}
\bigl(
\widetilde{\mathfrak V}_{Y,t,\ell}^{\mathrm{fr}}\Theta
\bigr)(u)
&:=
L_Y(t,t)\widetilde\chi_\ell(u)
\int_u^\ell
(v-u)^{\lambda-1}\Theta(v)\,dv,
\label{eq:annulus-VY-fr}\\
\bigl(
\widetilde{\mathfrak V}_{XY,t,\ell}^{\mathrm{fr}}\Theta
\bigr)(u)
&:=
L_{XY}(t,t)\widetilde\chi_\ell(u)
\int_u^\ell
(v-u)^{\mu-1}\Theta(v)\,dv.
\label{eq:annulus-VXY-fr}
\end{align}
Whenever the bounded operators
\[
\widetilde{\mathfrak Q}_{Y,t,h,R_1,\ell}^{\mathrm{ann}},
\quad
\widetilde{\mathfrak Q}_{XY,t,h,R_1,\ell}^{\mathrm{ann}},
\quad
\widetilde{\mathfrak Q}_{Y,t,\ell}^{\mathrm{fr}},
\quad
\widetilde{\mathfrak Q}_{XY,t,\ell}^{\mathrm{fr}}
\]
are defined by right-fractional conjugation of these Volterra pieces, define the frozen normalized
strip operator by
\begin{equation}\label{eq:annulus-S-fr}
\widetilde{\mathfrak S}_{t,\ell}^{\mathrm{fr}}
(\widetilde\Theta_X,\widetilde\Theta_Y)
:=
\widetilde{\mathfrak Q}_{Y,t,\ell}^{\mathrm{fr}}\widetilde\Theta_Y
-
\widetilde{\mathfrak Q}_{XY,t,\ell}^{\mathrm{fr}}\widetilde\Theta_X.
\end{equation}
\end{definition}

\begin{theorem}[Conjugated freezing on each fixed annulus width]
\label{thm:annulus-conjugated}
Assume Assumption~\ref{ass:model} with the extra regularity from
Definition~\ref{def:model-c2}. Fix \(t\in(0,T)\), \(\alpha\in(0,1)\), and suppose
\[
0<H_{XY}<H_Y<\tfrac12.
\]
Fix \(\ell>0\) and \(R_*<\infty\). Then there exist \(h_0\in(0,1]\) and
\(C_{t,\ell,R_*}<\infty\) such that for every
\[
0\le R_1\le R_*,
\qquad
0<h\le h_0,
\qquad
h(R_1+\ell)<t,
\]
the following hold:
\begin{enumerate}[label=(\roman*)]
\item there exist bounded operators on \(E_\ell\),
\[
\widetilde{\mathfrak Q}_{Y,t,h,R_1,\ell}^{\mathrm{ann}},
\quad
\widetilde{\mathfrak Q}_{XY,t,h,R_1,\ell}^{\mathrm{ann}},
\quad
\widetilde{\mathfrak Q}_{Y,t,\ell}^{\mathrm{fr}},
\quad
\widetilde{\mathfrak Q}_{XY,t,\ell}^{\mathrm{fr}},
\]
such that for every \(\Theta\in E_\ell\),
\begin{align}
\frac{1}{\Gamma(\lambda)}
\mathfrak D_-^\lambda
\bigl(
\widetilde{\mathfrak V}_{Y,t,h,R_1,\ell}^{\mathrm{ann}}\Theta
\bigr)
&=
\widetilde{\mathfrak Q}_{Y,t,h,R_1,\ell}^{\mathrm{ann}}\Theta,
\label{eq:annulus-QY}\\
\frac{1}{\Gamma(\lambda)}
\mathfrak D_-^\lambda
\bigl(
\widetilde{\mathfrak V}_{XY,t,h,R_1,\ell}^{\mathrm{ann}}\Theta
\bigr)
&=
\widetilde{\mathfrak Q}_{XY,t,h,R_1,\ell}^{\mathrm{ann}}\Theta,
\label{eq:annulus-QXY}\\
\frac{1}{\Gamma(\lambda)}
\mathfrak D_-^\lambda
\bigl(
\widetilde{\mathfrak V}_{Y,t,\ell}^{\mathrm{fr}}\Theta
\bigr)
&=
\widetilde{\mathfrak Q}_{Y,t,\ell}^{\mathrm{fr}}\Theta,
\label{eq:annulus-QY-fr}\\
\frac{1}{\Gamma(\lambda)}
\mathfrak D_-^\lambda
\bigl(
\widetilde{\mathfrak V}_{XY,t,\ell}^{\mathrm{fr}}\Theta
\bigr)
&=
\widetilde{\mathfrak Q}_{XY,t,\ell}^{\mathrm{fr}}\Theta;
\label{eq:annulus-QXY-fr}
\end{align}
\item these operators satisfy the uniform freezing bounds
\begin{align}
\left\|
\widetilde{\mathfrak Q}_{Y,t,h,R_1,\ell}^{\mathrm{ann}}
-
\widetilde{\mathfrak Q}_{Y,t,\ell}^{\mathrm{fr}}
\right\|_{\mathcal L(E_\ell)}
&\le
C_{t,\ell,R_*}h,
\label{eq:annulus-QY-close}\\
\left\|
\widetilde{\mathfrak Q}_{XY,t,h,R_1,\ell}^{\mathrm{ann}}
-
\widetilde{\mathfrak Q}_{XY,t,\ell}^{\mathrm{fr}}
\right\|_{\mathcal L(E_\ell)}
&\le
C_{t,\ell,R_*}h;
\label{eq:annulus-QXY-close}
\end{align}
\item for every admissible normalized annulus trace pair
\[
(\widetilde\Theta_X,\widetilde\Theta_Y)\in
\widetilde{\mathcal A}_{t,h,\alpha,R_1,\ell}^{\mathrm{ann}},
\]
one has
\begin{equation}\label{eq:annulus-S-close}
\left\|
\widetilde{\mathfrak S}_{t,h,\alpha,R_1,\ell}^{\mathrm{ann}}
(\widetilde\Theta_X,\widetilde\Theta_Y)
-
\widetilde{\mathfrak S}_{t,\ell}^{\mathrm{fr}}
(\widetilde\Theta_X,\widetilde\Theta_Y)
\right\|_{E_\ell}
\le
C_{t,\ell,R_*}h
\bigl(
\|\widetilde\Theta_X\|_{E_\ell}
+
\|\widetilde\Theta_Y\|_{E_\ell}
\bigr).
\end{equation}
\end{enumerate}
\end{theorem}

\begin{proof}
Choose \(h_0\in(0,1]\) so small that
\[
h_0(R_*+\ell)<t.
\]
Set
\[
b_{Y,t,h,R_1,\ell}(u,v)
:=
\widetilde\chi_\ell(u)\widetilde L_{Y,t,h,R_1,\ell}(u,v),
\qquad
b_{Y,t,\ell}^{\mathrm{fr}}(u,v)
:=
\widetilde\chi_\ell(u)L_Y(t,t),
\]
and
\[
b_{XY,t,h,R_1,\ell}(u,v)
:=
\widetilde\chi_\ell(u)\widetilde L_{XY,t,h,R_1,\ell}(u,v),
\qquad
b_{XY,t,\ell}^{\mathrm{fr}}(u,v)
:=
\widetilde\chi_\ell(u)L_{XY}(t,t).
\]
By Proposition~\ref{prop:annulus-normalized}(ii) and the smoothness of \(\widetilde\chi_\ell\),
there exists \(C_{t,\ell,R_*}<\infty\) such that
\begin{align}
\left\|
b_{Y,t,h,R_1,\ell}-b_{Y,t,\ell}^{\mathrm{fr}}
\right\|_{C^1(\Delta_\ell)}
&\le
C_{t,\ell,R_*}h,
\label{eq:annulus-bY-close}\\
\left\|
b_{XY,t,h,R_1,\ell}-b_{XY,t,\ell}^{\mathrm{fr}}
\right\|_{C^2(\Delta_\ell)}
&\le
C_{t,\ell,R_*}h.
\label{eq:annulus-bXY-close}
\end{align}

Item~(i) for the target side follows by applying the proof of
Proposition~\ref{prop:BY-micro} with \(R=\ell\) and coefficient \(b_{Y,t,h,R_1,\ell}\), and
again with the frozen coefficient \(b_{Y,t,\ell}^{\mathrm{fr}}\). Because the difference of these
coefficients is \(O(h)\) in \(C^1(\Delta_\ell)\), the same proof gives
\eqref{eq:annulus-QY-close}.

For the source side, apply the proof of Proposition~\ref{prop:micro-window-decompose} with
\(R=\ell\), replacing \(L_{XY}(t-h r,t-h q)\) by \(b_{XY,t,h,R_1,\ell}(r,q)\), and again with the
frozen coefficient \(b_{XY,t,\ell}^{\mathrm{fr}}(r,q)\). The \(C^2\) control
\eqref{eq:annulus-bXY-close} yields the operator-norm estimate \eqref{eq:annulus-QXY-close}. This
proves item~(i) and item~(ii).

For item~(iii), let
\[
(\widetilde\Theta_X,\widetilde\Theta_Y)\in
\widetilde{\mathcal A}_{t,h,\alpha,R_1,\ell}^{\mathrm{ann}}.
\]
By Definition~\ref{def:annulus-normalized}, choose the corresponding unshifted traces
\[
\Theta_X:=\mathfrak T_{R_1}^{-1}\widetilde\Theta_X,
\qquad
\Theta_Y:=\mathfrak T_{R_1}^{-1}\widetilde\Theta_Y,
\]
so that
\[
(\Theta_X,\Theta_Y)\in \mathcal A_{t,h,\alpha,R_1,R_2}^{\mathrm{ann}},
\qquad
R_2=R_1+\ell.
\]
Applying Proposition~\ref{prop:strip-annulus} on \([R_1,R_2]\), then translating to
\([0,\ell]\), shows that
\[
\widetilde{\mathfrak S}_{t,h,\alpha,R_1,\ell}^{\mathrm{ann}}
(\widetilde\Theta_X,\widetilde\Theta_Y)
=
\widetilde{\mathfrak Q}_{Y,t,h,R_1,\ell}^{\mathrm{ann}}\widetilde\Theta_Y
-
\widetilde{\mathfrak Q}_{XY,t,h,R_1,\ell}^{\mathrm{ann}}\widetilde\Theta_X.
\]
Subtracting \eqref{eq:annulus-S-fr}, using \eqref{eq:annulus-QY-close} and
\eqref{eq:annulus-QXY-close}, gives \eqref{eq:annulus-S-close}.
\end{proof}

\begin{remark}[Reduction to frozen admissibility]
\label{rem:annulus-conjugated}
Theorem~\ref{thm:annulus-conjugated} reduces the final step from the exact normalized strip operator to
a fixed frozen strip operator on \([0,\ell]\).  Coefficient-freezing is no longer the obstruction;
the admissible trace class of the frozen model is. The subsequent annulus
and localization sections show that these frozen admissibility constraints are exactly what drives
the strip contribution to zero.
\end{remark}

\begin{corollary}[Boundedness of the frozen normalized strip operator]
\label{cor:annulus-frozen-bounded}
Under the assumptions of Definition~\ref{def:annulus-conjugated}, there exists
\(C_{t,\ell}^{\mathrm{fr}}<\infty\) such that
\begin{equation}\label{eq:annulus-frozen-bounded}
\left\|
\widetilde{\mathfrak S}_{t,\ell}^{\mathrm{fr}}
(\widetilde\Theta_X,\widetilde\Theta_Y)
\right\|_{E_\ell}
\le
C_{t,\ell}^{\mathrm{fr}}
\bigl(
\|\widetilde\Theta_X\|_{E_\ell}
+
\|\widetilde\Theta_Y\|_{E_\ell}
\bigr)
\end{equation}
for every
\[
(\widetilde\Theta_X,\widetilde\Theta_Y)\in E_\ell\times E_\ell.
\]
\end{corollary}

\begin{proof}
This follows from Definition~\ref{def:annulus-conjugated} and
Theorem~\ref{thm:annulus-conjugated}(i): the frozen strip operator is the difference of two bounded
operators on \(E_\ell\).
\end{proof}

\begin{theorem}[Bounded-buffer strip smallness]
\label{thm:annulus-strip-small}
Assume Assumptions~\ref{ass:model} and~\ref{ass:obs-source}, and assume in addition the extra
regularity from Definition~\ref{def:model-c2}. Fix \(t\in(0,T)\), \(\alpha\in(0,1)\), and suppose
\[
0<H_{XY}<H_Y<\tfrac12.
\]
Let \(h_n\downarrow0\), and let
\[
\Phi_n\in\mathcal D(\mathfrak T_{t,h_n,\alpha}^{\mathrm{mic}}),
\qquad
\Psi_n:=\mathfrak T_{t,h_n,\alpha}^{\mathrm{mic}}\Phi_n,
\qquad
\|\Phi_n\|_{\mathfrak G_{t,h_n,\alpha}^{\mathrm{mic}}}\le1.
\]
Then for every \(\varepsilon>0\), after passage to a subsequence there exist common radii
\[
0<R_0<R_1<R_2<\infty
\]
and the canonical cutoff \(\chi_{R_1,R_2}\) such that
\[
R_1-R_0\ge R_2-R_1,
\]
and
\begin{equation}\label{eq:annulus-strip-small-thm}
\left\|
\one_{[R_1,R_2]}
\mathfrak A_{t,h_n,\alpha,\chi_{R_1,R_2}}^{\mathrm{out}}(\Phi_n,\Psi_n)
\right\|_{E_{R_2}}
\le
\varepsilon
\qquad\text{for every }n.
\end{equation}
\end{theorem}

\begin{proof}
Fix \(\varepsilon>0\), choose \(\ell>0\) arbitrarily, and set
\[
A:=2\ell.
\]

Let \(C_{t,\ell}^{\mathrm{fr}}\) be the constant from
Corollary~\ref{cor:annulus-frozen-bounded}. Choose \(M\in\mathbb N\) so large that
\[
C_{t,\ell}^{\mathrm{fr}}
\left(1+\frac{\nu_Y}{\eta}\right)M^{-1/2}
\le
\frac{\varepsilon}{4}.
\]

Applying Corollary~\ref{cor:good-buffer-sequence} with this \(A\), this \(\ell\), and this \(M\),
and then relabeling the resulting subsequence, yields one index \(j_*\in\{1,\dots,M\}\) such that
for every \(n\),
\[
\|\Phi_n\|_{L^2(I_{j_*})}^2
+
\frac{\eta^2}{\nu_Y^2}
\|\Psi_n\|_{L^2(I_{j_*})}^2
\le
\frac{1}{M},
\]
where
\[
I_{j_*}
=
\bigl[A+(j_*-1)\ell,\ A+j_*\ell\bigr].
\]
Set
\[
R_1:=A+(j_*-1)\ell,
\qquad
R_2:=A+j_*\ell,
\qquad
R_0:=A-\ell.
\]
Then
\[
0<R_0<R_1<R_2<\infty,
\qquad
R_2-R_1=\ell,
\qquad
R_1-R_0\ge \ell.
\]
Also,
\begin{equation}\label{eq:annulus-strip-trace}
\|\Phi_n\|_{L^2([R_1,R_2])}
+
\|\Psi_n\|_{L^2([R_1,R_2])}
\le
\left(1+\frac{\nu_Y}{\eta}\right)M^{-1/2}
\qquad\text{for every }n.
\end{equation}

Set
\[
R_*:=A+M\ell.
\]
Let \(C_{t,\ell,R_*}\) be the constant from Theorem~\ref{thm:annulus-conjugated}. Since
\(R_2\le R_*\) and \(h_n\downarrow0\), after discarding finitely many more terms we may also assume
for every remaining \(n\) that
\[
h_nR_2<t,
\qquad
R_2<\varepsilon_{h_n,\alpha}^{-1},
\qquad
C_{t,\ell,R_*}h_n
\left(1+\frac{\nu_Y}{\eta}\right)M^{-1/2}
\le
\frac{\varepsilon}{4}.
\]

Define normalized annulus traces by
\[
\widetilde\Theta_{X,n}
:=
\mathfrak T_{R_1}\bigl(\Phi_n\one_{[R_1,R_2]}\bigr),
\qquad
\widetilde\Theta_{Y,n}
:=
\mathfrak T_{R_1}\bigl(\Psi_n\one_{[R_1,R_2]}\bigr)
\in E_\ell.
\]
By Definition~\ref{def:annulus-normalized},
\[
(\widetilde\Theta_{X,n},\widetilde\Theta_{Y,n})
\in
\widetilde{\mathcal A}_{t,h_n,\alpha,R_1,\ell}^{\mathrm{ann}}
\qquad\text{for every }n.
\]
Moreover, \eqref{eq:annulus-strip-trace} gives
\begin{equation}\label{eq:annulus-strip-normalized}
\|\widetilde\Theta_{X,n}\|_{E_\ell}
+
\|\widetilde\Theta_{Y,n}\|_{E_\ell}
\le
\left(1+\frac{\nu_Y}{\eta}\right)M^{-1/2}.
\end{equation}

By Definition~\ref{def:annulus-normalized},
\[
\left\|
\one_{[R_1,R_2]}
\mathfrak A_{t,h_n,\alpha,\chi_{R_1,R_2}}^{\mathrm{out}}(\Phi_n,\Psi_n)
\right\|_{E_{R_2}}
=
\left\|
\widetilde{\mathfrak S}_{t,h_n,\alpha,R_1,\ell}^{\mathrm{ann}}
(\widetilde\Theta_{X,n},\widetilde\Theta_{Y,n})
\right\|_{E_\ell}.
\]
Using Theorem~\ref{thm:annulus-conjugated}(iii),
Corollary~\ref{cor:annulus-frozen-bounded}, and \eqref{eq:annulus-strip-normalized}, we obtain
\begin{align*}
&\left\|
\widetilde{\mathfrak S}_{t,h_n,\alpha,R_1,\ell}^{\mathrm{ann}}
(\widetilde\Theta_{X,n},\widetilde\Theta_{Y,n})
\right\|_{E_\ell}\\
&\le
\left\|
\widetilde{\mathfrak S}_{t,\ell}^{\mathrm{fr}}
(\widetilde\Theta_{X,n},\widetilde\Theta_{Y,n})
\right\|_{E_\ell}
+
\left\|
\widetilde{\mathfrak S}_{t,h_n,\alpha,R_1,\ell}^{\mathrm{ann}}
(\widetilde\Theta_{X,n},\widetilde\Theta_{Y,n})
-
\widetilde{\mathfrak S}_{t,\ell}^{\mathrm{fr}}
(\widetilde\Theta_{X,n},\widetilde\Theta_{Y,n})
\right\|_{E_\ell}\\
&\le
C_{t,\ell}^{\mathrm{fr}}
\bigl(
\|\widetilde\Theta_{X,n}\|_{E_\ell}
+
\|\widetilde\Theta_{Y,n}\|_{E_\ell}
\bigr)
+
C_{t,\ell,R_*}h_n
\bigl(
\|\widetilde\Theta_{X,n}\|_{E_\ell}
+
\|\widetilde\Theta_{Y,n}\|_{E_\ell}
\bigr)\\
&\le
\frac{\varepsilon}{4}
+
\frac{\varepsilon}{4}
\le
\varepsilon.
\end{align*}
This proves \eqref{eq:annulus-strip-small-thm}.
\end{proof}

\begin{corollary}[Global micro dual collapse]
\label{cor:global-collapse}
Under the assumptions of Theorem~\ref{thm:annulus-strip-small}, the localization hypothesis in
Theorem~\ref{thm:abstract-global-reduction} holds. Consequently,
\begin{equation}\label{eq:global-micro-collapse}
\mathfrak D_{t,h,\alpha}^{\mathrm{mic}}\longrightarrow0
\qquad\text{as }h\downarrow0.
\end{equation}
\end{corollary}

\begin{proof}
Theorem~\ref{thm:annulus-strip-small} furnishes the strip-smallness statement
\eqref{eq:strip-annulus-small} required in Corollary~\ref{cor:strip-annulus}. Hence the
localization hypothesis in Theorem~\ref{thm:abstract-global-reduction} follows. The convergence
\eqref{eq:global-micro-collapse} is then the conclusion of
Theorem~\ref{thm:abstract-global-reduction}.
\end{proof}

\section{From micro dual collapse to observable screening}
\label{sec:screening}

Micro dual collapse is a boundary-layer statement.  Translated back to the original observable
functional \(\mathcal F_{t,h}\), it identifies the global micro dual norm with the dual norm of
\(\mathcal F_{t,h}\) restricted to one \(h^\alpha\) boundary layer.  A cutoff-stability assumption
then upgrades the boundary-layer statement to the global screening statement for the observable
gain.

\begin{proposition}[Exact identification of the micro dual norm with the near-layer dual norm]
\label{prop:micro-near-dual}
Assume the hypotheses of Definition~\ref{def:micro-dual}. Fix \(\alpha\in(0,1)\). Then for every
sufficiently small \(h>0\),
\begin{equation}\label{eq:micro-near-dual}
\mathfrak D_{t,h,\alpha}^{\mathrm{mic}}
=
\sup\left\{
h^{-H_{XY}}\left|\mathcal F_{t,h}(v)\right|:
v\in\mathfrak G_t\cap H_{t,h,\alpha}^{\mathrm{near}},
\ \|v\|_{\mathfrak G_t}\le1
\right\}.
\end{equation}
\end{proposition}

\begin{proof}
Fix \(h>0\) sufficiently small.

Because
\[
\mathfrak l_{\varepsilon_{h,\alpha}}\mathfrak p_{\varepsilon_{h,\alpha}}
=
\mathfrak l_{\varepsilon_{h,\alpha}},
\]
the definitions of \(\mathfrak T_{t,h,\alpha}^{\mathrm{mic}}\) and
\(\Lambda_{t,h,\alpha}^{\mathrm{mic}}\) imply that for every
\[
\Phi\in\mathcal D(\mathfrak T_{t,h,\alpha}^{\mathrm{mic}})
\]
one has
\[
\mathfrak T_{t,h,\alpha}^{\mathrm{mic}}\mathfrak p_{\varepsilon_{h,\alpha}}\Phi
=
\mathfrak T_{t,h,\alpha}^{\mathrm{mic}}\Phi,
\qquad
\Lambda_{t,h,\alpha}^{\mathrm{mic}}(\mathfrak p_{\varepsilon_{h,\alpha}}\Phi)
=
\Lambda_{t,h,\alpha}^{\mathrm{mic}}(\Phi).
\]
Since \(\mathfrak p_{\varepsilon_{h,\alpha}}\) is an orthogonal projection on \(L^2(\R_+)\),
\[
\|\mathfrak p_{\varepsilon_{h,\alpha}}\Phi\|_{\mathfrak G_{t,h,\alpha}^{\mathrm{mic}}}
\le
\|\Phi\|_{\mathfrak G_{t,h,\alpha}^{\mathrm{mic}}}.
\]
Hence the supremum defining \(\mathfrak D_{t,h,\alpha}^{\mathrm{mic}}\) may be restricted to
profiles satisfying
\[
\mathfrak p_{\varepsilon_{h,\alpha}}\Phi=\Phi.
\]

Let
\[
\Phi\in \mathcal D(\mathfrak T_{t,h,\alpha}^{\mathrm{mic}})
\]
with \(\mathfrak p_{\varepsilon_{h,\alpha}}\Phi=\Phi\), and set
\[
v:=\mathcal L_{t,h}\mathfrak p_{\varepsilon_{h,\alpha}}\Phi.
\]
Since
\[
\mathfrak l_{\varepsilon_{h,\alpha}}\Phi\in \mathfrak G_{t,h,\alpha}^{\mathrm{near}},
\]
Definition~\ref{def:local-transfer} and Proposition~\ref{prop:two-scale-factorization}(ii) give
\[
v=\mathcal L_{t,h}^{(\alpha)}\mathfrak l_{\varepsilon_{h,\alpha}}\Phi
\in
\mathfrak G_t\cap H_{t,h,\alpha}^{\mathrm{near}}.
\]
By Proposition~\ref{prop:micro-local}(iii),
\begin{equation}\label{eq:micro-near-functional-forward}
h^{-H_{XY}}\mathcal F_{t,h}(v)=\Lambda_{t,h,\alpha}^{\mathrm{mic}}(\Phi).
\end{equation}
Also, by Lemma~\ref{lem:micro-dilation},
Proposition~\ref{prop:two-scale-factorization}(ii), and
Proposition~\ref{prop:micro-local}(ii),
\[
\|v\|_{H_t}
=
\|\mathfrak p_{\varepsilon_{h,\alpha}}\Phi\|_{L^2(\R_+)}
=
\|\Phi\|_{L^2(\R_+)},
\]
while
\[
\mathfrak T_{t,h,\alpha}^{\mathrm{mic}}\Phi
=
h^\beta \mathcal D_{t,h}\mathcal S_t v.
\]
Because \(\mathcal S_t v\in H_{t,h,\alpha}^{\mathrm{near}}\), Proposition~\ref{prop:two-scale-factorization}(i)
and Lemma~\ref{lem:dilation-isometry} yield
\[
\left\|
\mathfrak T_{t,h,\alpha}^{\mathrm{mic}}\Phi
\right\|_{L^2(\R_+)}
=
\|\mathcal S_t v\|_{H_t}.
\]
Hence
\begin{equation}\label{eq:micro-near-norm-forward}
\|\Phi\|_{\mathfrak G_{t,h,\alpha}^{\mathrm{mic}}}
=
\|v\|_{\mathfrak G_t}.
\end{equation}

Conversely, let
\[
v\in\mathfrak G_t\cap H_{t,h,\alpha}^{\mathrm{near}}.
\]
Set
\[
\phi:=\mathcal D_{t,h}^{(\alpha)}v\in \mathfrak G_{t,h,\alpha}^{\mathrm{near}},
\qquad
\Phi:=\mathfrak d_{\varepsilon_{h,\alpha}}\phi\in L^2(\R_+).
\]
By Lemma~\ref{lem:micro-dilation},
\[
\mathfrak l_{\varepsilon_{h,\alpha}}\Phi=\phi,
\]
so \(\Phi\in\mathcal D(\mathfrak T_{t,h,\alpha}^{\mathrm{mic}})\). Since \(\phi\) is supported in
\([0,1]\), one has
\[
\mathfrak p_{\varepsilon_{h,\alpha}}\Phi=\Phi,
\]
and Proposition~\ref{prop:two-scale-factorization}(ii) gives
\[
\mathcal L_{t,h}\mathfrak p_{\varepsilon_{h,\alpha}}\Phi
=
\mathcal L_{t,h}^{(\alpha)}\phi
=
v.
\]
Applying Proposition~\ref{prop:micro-local}(iii) therefore yields
\begin{equation}\label{eq:micro-near-functional-backward}
h^{-H_{XY}}\mathcal F_{t,h}(v)=\Lambda_{t,h,\alpha}^{\mathrm{mic}}(\Phi).
\end{equation}
The same norm calculation gives
\begin{equation}\label{eq:micro-near-norm-backward}
\|\Phi\|_{\mathfrak G_{t,h,\alpha}^{\mathrm{mic}}}
=
\|v\|_{\mathfrak G_t}.
\end{equation}

The two constructions are inverse to one another, and
\eqref{eq:micro-near-functional-forward}--\eqref{eq:micro-near-norm-backward} identify the two
suprema. This proves \eqref{eq:micro-near-dual}.
\end{proof}

\begin{corollary}[Near-layer screening at every boundary-layer scale]
\label{cor:micro-near-screening}
Assume the hypotheses of Corollary~\ref{cor:global-collapse}. Fix \(\alpha\in(0,1)\). Then
\begin{equation}\label{eq:near-screening}
\sup\left\{
h^{-H_{XY}}\left|\mathcal F_{t,h}(v)\right|:
v\in\mathfrak G_t\cap H_{t,h,\alpha}^{\mathrm{near}},
\ \|v\|_{\mathfrak G_t}\le1
\right\}
\longrightarrow0
\qquad\text{as }h\downarrow0.
\end{equation}
\end{corollary}

\begin{proof}
Corollary~\ref{cor:global-collapse} gives
\(\mathfrak D_{t,h,\alpha}^{\mathrm{mic}}\to0\). Proposition~\ref{prop:micro-near-dual}
identifies this micro dual norm with the displayed near-layer supremum for all sufficiently small
\(h\). The supremum therefore converges to zero.
\end{proof}

\begin{definition}[\(h^\alpha\)-stable boundary cutoffs]
\label{def:alpha-cutoff-smooth}
Assume \(\eta>0\), fix \(t\in(0,T)\), and let \(\alpha\in(0,1)\). We say that \(\mathcal S_t\)
has \emph{\(h^\alpha\)-stable boundary cutoffs at time \(t\)} if there exist
\[
\delta_*\in(0,1),\qquad C_{\mathrm{scut},\alpha}<\infty,\qquad h_{0,\alpha}\in(0,t]
\]
such that for every \(0<h\le h_{0,\alpha}\) there exists a cutoff
\[
\chi_{h,\alpha}:[0,t]\to[0,1]
\]
with
\[
0\le \chi_{h,\alpha}\le 1,\qquad
\chi_{h,\alpha}\equiv1\ \text{on }[t-(1-\delta_*)h^\alpha,t],\qquad
\operatorname{supp}\chi_{h,\alpha}\subset[t-h^\alpha,t],
\]
and such that for every \(v\in\mathfrak G_t\),
\begin{equation}\label{eq:alpha-cutoff-smooth}
\chi_{h,\alpha}v\in\mathfrak G_t
\qquad\text{and}\qquad
\|\chi_{h,\alpha}v\|_{\mathfrak G_t}
\le
C_{\mathrm{scut},\alpha}\|v\|_{\mathfrak G_t}.
\end{equation}
\end{definition}

\begin{remark}[Why the cutoff lives on [0,t]]
The cutoff is used only as a multiplier on \(H_t=L^2([0,t])\) and on the graph space
\(\mathfrak G_t\). What matters in Definition~\ref{def:alpha-cutoff-smooth} is the boundary-layer
support, the plateau near \(t\), and the graph-norm stability in
\eqref{eq:alpha-cutoff-smooth}. No ambient smooth extension beyond \([0,t]\) is needed for the
screening argument.
\end{remark}

\begin{definition}[Far/strip localization on an \(h^\alpha\) layer]
\label{def:alpha-layer-projections}
Fix \(t\in(0,T)\), let \(\alpha\in(0,1)\), let \(h\in(0,t]\), and let \(\delta\in(0,1)\). Define
the orthogonal projections on \(H_t=L^2([0,t])\) by
\begin{align}
\bigl(\Pi_{t,h}^{\mathrm{far},\alpha}f\bigr)(s)
&:=f(s)\one_{[0,\ t-h^\alpha]}(s), \label{eq:Pi-far-alpha}\\
\bigl(\Pi_{t,h}^{\mathrm{strip},\alpha,\delta}f\bigr)(s)
&:=f(s)\one_{[t-h^\alpha,\ t-(1-\delta)h^\alpha]}(s).
\label{eq:Pi-strip-alpha}
\end{align}
\end{definition}

\begin{lemma}[Far-zone estimate on an \(h^\alpha\) buffer]
\label{lem:far-kernel-alpha}
Assume Assumption~\ref{ass:model}. Fix \(t\in(0,T)\), let \(i\in\{Y,XY\}\), and let
\(\alpha\in(0,1)\). Then
\begin{equation}\label{eq:far-kernel-alpha}
\|\Pi_{t,h}^{\mathrm{far},\alpha}g_{i,t,h}\|_{H_t}
=
O\!\left(h^{1-\alpha(1-H_i)}\right)
\qquad\text{as }h\downarrow0.
\end{equation}
\end{lemma}

\begin{proof}
Write \(H:=H_i\). By Lemma~\ref{lem:offdiag-pointwise}, there exist constants
\(C<\infty\) and \(h_1\in(0,t]\). After shrinking \(h_1\) if necessary, depending only on
\(t\), \(\alpha\), and the constants from Lemma~\ref{lem:offdiag-pointwise}, we may assume
\(h_1\le1\). Since \(h^\alpha\ge h\) for \(0<h\le1\) and \(\alpha\in(0,1)\), every
\(s\in[0,t-h^\alpha]\) satisfies \(s\le t-h\) for \(0<h\le h_1\). Hence
\[
|g_{i,t,h}(s)|
\le
Ch\bigl((t-s)^{H-\frac32}+(t-s)^{H-\frac12}\bigr)
\qquad\text{for }0\le s\le t-h^\alpha.
\]
Therefore
\begin{align*}
\|\Pi_{t,h}^{\mathrm{far},\alpha}g_{i,t,h}\|_{H_t}^2
&\lesssim
h^2\int_0^{t-h^\alpha}
\bigl((t-s)^{2H-3}+(t-s)^{2H-1}\bigr)\,ds.
\end{align*}
With \(u=t-s\), this becomes
\[
\|\Pi_{t,h}^{\mathrm{far},\alpha}g_{i,t,h}\|_{H_t}^2
\lesssim
h^2\int_{h^\alpha}^{t}(u^{2H-3}+u^{2H-1})\,du.
\]
Since \(H\in(0,\tfrac12)\),
\[
\int_{h^\alpha}^{t}u^{2H-3}\,du
=
O\!\left(h^{-2\alpha(1-H)}\right),
\qquad
\int_{h^\alpha}^{t}u^{2H-1}\,du
=
O(1).
\]
Hence
\[
\|\Pi_{t,h}^{\mathrm{far},\alpha}g_{i,t,h}\|_{H_t}^2
=
O\!\left(h^{2-2\alpha(1-H)}\right)+O(h^2)
=
O\!\left(h^{2-2\alpha(1-H)}\right),
\]
and taking square roots proves \eqref{eq:far-kernel-alpha}.
\end{proof}

\begin{lemma}[Old-edge strip estimate on an \(h^\alpha\) layer]
\label{lem:strip-kernel-alpha}
Assume Assumption~\ref{ass:model}. Fix \(t\in(0,T)\), let \(i\in\{Y,XY\}\), let
\(\alpha\in(0,1)\), and fix \(\delta\in(0,1)\). Then the projection
\(\Pi_{t,h}^{\mathrm{strip},\alpha,\delta}\) from
Definition~\ref{def:alpha-layer-projections} satisfies
\begin{equation}\label{eq:strip-kernel-alpha}
\|\Pi_{t,h}^{\mathrm{strip},\alpha,\delta}g_{i,t,h}\|_{H_t}
=
O\!\left(\delta^{1/2}h^{1-\alpha(1-H_i)}\right)
\qquad\text{as }h\downarrow0.
\end{equation}
\end{lemma}

\begin{proof}
Write \(H:=H_i\). Lemma~\ref{lem:offdiag-pointwise} yields
\(C<\infty\) and \(h_1\in(0,t]\). After shrinking \(h_1\) if necessary, depending only on
\(t\), \(\alpha\), \(\delta\), and the constants from Lemma~\ref{lem:offdiag-pointwise}, we may
assume \(h_1\le1\) and \(h_1^{1-\alpha}\le1-\delta\). Then for every \(0<h\le h_1\),
\[
h\le (1-\delta)h^\alpha.
\]
Hence every
\[
s\in[t-h^\alpha,\ t-(1-\delta)h^\alpha],
\]
satisfies \(s\le t-h\), so Lemma~\ref{lem:offdiag-pointwise} gives
\[
|g_{i,t,h}(s)|
\le
Ch\bigl((t-s)^{H-\frac32}+(t-s)^{H-\frac12}\bigr).
\]
Therefore
\begin{align*}
\|\Pi_{t,h}^{\mathrm{strip},\alpha,\delta}g_{i,t,h}\|_{H_t}^2
&\lesssim
h^2\int_{t-h^\alpha}^{t-(1-\delta)h^\alpha}
\bigl((t-s)^{2H-3}+(t-s)^{2H-1}\bigr)\,ds.
\end{align*}
With \(u=t-s\), this becomes
\[
\|\Pi_{t,h}^{\mathrm{strip},\alpha,\delta}g_{i,t,h}\|_{H_t}^2
\lesssim
h^2\int_{(1-\delta)h^\alpha}^{h^\alpha}(u^{2H-3}+u^{2H-1})\,du.
\]
Since \(H\in(0,\tfrac12)\), the second integral is \(O(\delta h^{2\alpha H})\), while
\[
\int_{(1-\delta)h^\alpha}^{h^\alpha}u^{2H-3}\,du
=
h^{-2\alpha(1-H)}
\int_{1-\delta}^{1}r^{2H-3}\,dr
=
O\!\left(\delta h^{-2\alpha(1-H)}\right),
\]
because \(r^{2H-3}\) is bounded on \([1-\delta,1]\) for fixed \(\delta\in(0,1)\). Hence
\[
\|\Pi_{t,h}^{\mathrm{strip},\alpha,\delta}g_{i,t,h}\|_{H_t}^2
=
O\!\left(\delta h^{2-2\alpha(1-H)}\right),
\]
and taking square roots proves \eqref{eq:strip-kernel-alpha}.
\end{proof}

\begin{theorem}[Observable screening under \(h^\alpha\)-stable boundary cutoffs]
\label{thm:observable-screening-smooth}
Assume Assumptions~\ref{ass:model} and~\ref{ass:obs-source}, assume the extra regularity from
Definition~\ref{def:model-c2}, let \(\eta>0\), and fix \(t\in(0,T)\). Suppose
\[
0<H_{XY}<H_Y<\tfrac12.
\]
Fix \(\alpha\in(0,1)\), and assume \(\mathcal S_t\) has \(h^\alpha\)-stable boundary
cutoffs at time \(t\) in the sense of Definition~\ref{def:alpha-cutoff-smooth}. Then
\begin{equation}\label{eq:observable-screening-smooth-F}
\|\mathcal F_{t,h}\|_{\mathfrak G_t^*}
=
o(h^{H_{XY}})
\qquad\text{as }h\downarrow0,
\end{equation}
and consequently
\begin{equation}\label{eq:observable-screening-smooth-G}
\widetilde G_{X\to Y}(t,h)
=
o(h^{2H_{XY}})
\qquad\text{as }h\downarrow0.
\end{equation}
\end{theorem}

\begin{remark}[Role of the cutoff hypothesis]
Theorem~\ref{thm:observable-screening-smooth} is the screening statement once
\(h^\alpha\)-stable boundary cutoffs are available.  The remaining sections verify that cutoff
hypothesis for the model classes covered by the exact-correction analysis.  The theorem is therefore
a conditional screening result, not an isolated regularity assumption.
\end{remark}

\begin{proof}
Let \(\delta_*\), \(C_{\mathrm{scut},\alpha}\), and \(h_{0,\alpha}\) be given by
Definition~\ref{def:alpha-cutoff-smooth}. For \(0<h\le h_{0,\alpha}\) and \(v\in\mathfrak G_t\)
with \(\|v\|_{\mathfrak G_t}\le1\), define
\[
w_h:=\chi_{h,\alpha}v,
\qquad
r_h:=v-w_h.
\]
Then
\[
w_h\in\mathfrak G_t\cap H_{t,h,\alpha}^{\mathrm{near}},
\qquad
\|w_h\|_{\mathfrak G_t}\le C_{\mathrm{scut},\alpha},
\]
and
\[
r_h\in\mathfrak G_t,
\qquad
\|r_h\|_{\mathfrak G_t}\le 1+C_{\mathrm{scut},\alpha}.
\]
Since \(\mathcal F_{t,h}\) is linear,
\[
|\mathcal F_{t,h}(v)|
\le
|\mathcal F_{t,h}(w_h)|
+
|\mathcal F_{t,h}(r_h)|.
\]

Define
\[
\omega_\alpha(h)
:=
\sup\left\{
h^{-H_{XY}}|\mathcal F_{t,h}(u)|:
u\in\mathfrak G_t\cap H_{t,h,\alpha}^{\mathrm{near}},
\|u\|_{\mathfrak G_t}\le1
\right\}.
\]
Corollary~\ref{cor:micro-near-screening} gives $\omega_\alpha(h)\to0$. Since
$w_h\in\mathfrak G_t\cap H_{t,h,\alpha}^{\mathrm{near}}$, homogeneity gives
\[
|\mathcal F_{t,h}(w_h)|
\le
\omega_\alpha(h)h^{H_{XY}}\|w_h\|_{\mathfrak G_t}
\le
C_{\mathrm{scut},\alpha}\omega_\alpha(h)h^{H_{XY}}.
\]

Because \(\chi_{h,\alpha}\equiv1\) on \([t-(1-\delta_*)h^\alpha,t]\) and
\(\operatorname{supp}\chi_{h,\alpha}\subset[t-h^\alpha,t]\), the remainder \(r_h\) is supported in
\[
[0,\ t-(1-\delta_*)h^\alpha]
=
[0,\ t-h^\alpha]\cup[t-h^\alpha,\ t-(1-\delta_*)h^\alpha].
\]
Hence
\[
r_h
=
\Pi_{t,h}^{\mathrm{far},\alpha}r_h
+
\Pi_{t,h}^{\mathrm{strip},\alpha,\delta_*}r_h.
\]
Using Cauchy--Schwarz, Lemmas~\ref{lem:far-kernel-alpha} and~\ref{lem:strip-kernel-alpha}, we obtain
\begin{align*}
\bigl|
\langle g_{XY,t,h},r_h\rangle_{H_t}
\bigr|
&\le
\|\Pi_{t,h}^{\mathrm{far},\alpha}g_{XY,t,h}\|_{H_t}\|r_h\|_{H_t}
+
\|\Pi_{t,h}^{\mathrm{strip},\alpha,\delta_*}g_{XY,t,h}\|_{H_t}\|r_h\|_{H_t}\\
&=
O\!\left(h^{1-\alpha(1-H_{XY})}\right),
\end{align*}
because \(\|r_h\|_{H_t}\le \|r_h\|_{\mathfrak G_t}\le 1+C_{\mathrm{scut},\alpha}\).
Likewise,
\[
\bigl|
\langle g_{Y,t,h},\mathcal S_t r_h\rangle_{H_t}
\bigr|
\le
\|g_{Y,t,h}\|_{H_t}\|\mathcal S_t r_h\|_{H_t}
=
O(h^{H_Y}),
\]
using Proposition~\ref{prop:old} and the graph norm inequality. Therefore
\[
|\mathcal F_{t,h}(r_h)|
=
O\!\left(h^{1-\alpha(1-H_{XY})}\right)
+
O(h^{H_Y}).
\]

Combining the near and remainder bounds gives
\[
\sup_{\|v\|_{\mathfrak G_t}\le1}
|\mathcal F_{t,h}(v)|
\le
C_{\mathrm{scut},\alpha}\omega_\alpha(h)h^{H_{XY}}
+
O\!\left(h^{1-\alpha(1-H_{XY})}\right)
+
O(h^{H_Y}).
\]
Because \(\alpha<1\) and \(H_Y>H_{XY}\),
\[
1-\alpha(1-H_{XY})>H_{XY},
\qquad
H_Y>H_{XY},
\]
so the two remainder terms are \(o(h^{H_{XY}})\). This proves
\eqref{eq:observable-screening-smooth-F}. Then Proposition~\ref{prop:graph-dual} gives
\eqref{eq:observable-screening-smooth-G}.
\end{proof}

\begin{corollary}[Quantitative observable screening from a quantitative near-layer profile]
\label{cor:observable-screening-rate}
Assume the hypotheses of Theorem~\ref{thm:observable-screening-smooth}.  For \(0<h\le h_{0,\alpha}\),
define the near-layer screening profile
\[
\omega_\alpha(h)
:=
\sup\left\{
h^{-H_{XY}}\left|\mathcal F_{t,h}(v)\right|:
v\in\mathfrak G_t\cap H_{t,h,\alpha}^{\mathrm{near}},
\ \|v\|_{\mathfrak G_t}\le1
\right\}.
\]
Suppose moreover that there exist \(\rho>0\), \(C_{\omega,\alpha}<\infty\), and
\(\widetilde h_{0,\alpha}\in(0,h_{0,\alpha}]\) such that
\[
\omega_\alpha(h)\le C_{\omega,\alpha}h^\rho
\qquad\text{for every }0<h\le \widetilde h_{0,\alpha}.
\]
Set
\[
\varepsilon
:=
\min\!\bigl\{\rho,\;1-\alpha(1-H_{XY})-H_{XY},\;H_Y-H_{XY}\bigr\}.
\]
If \(\varepsilon>0\), then
\begin{equation}\label{eq:observable-screening-rate-F}
\|\mathcal F_{t,h}\|_{\mathfrak G_t^*}
=
O(h^{H_{XY}+\varepsilon})
\qquad\text{as }h\downarrow0,
\end{equation}
and consequently
\begin{equation}\label{eq:observable-screening-rate-G}
\widetilde G_{X\to Y}(t,h)
=
O(h^{2H_{XY}+2\varepsilon})
\qquad\text{as }h\downarrow0.
\end{equation}
\end{corollary}

\begin{proof}
Repeat the proof of Theorem~\ref{thm:observable-screening-smooth}, but retain the explicit scale.
For \(\|v\|_{\mathfrak G_t}\le1\), the near-layer term satisfies
\[
|\mathcal F_{t,h}(w_h)|
\le
C_{\mathrm{scut},\alpha}\omega_\alpha(h)h^{H_{XY}}
\le
C_{\mathrm{scut},\alpha}C_{\omega,\alpha}h^{H_{XY}+\rho}.
\]
The same far/strip calculation used in Theorem~\ref{thm:observable-screening-smooth} gives the
unchanged remainder estimate
\[
|\mathcal F_{t,h}(r_h)|
=
O\!\left(h^{1-\alpha(1-H_{XY})}\right)
+
O(h^{H_Y}).
\]
Therefore
\[
\|\mathcal F_{t,h}\|_{\mathfrak G_t^*}
\le
C h^{H_{XY}+\rho}
+
C h^{1-\alpha(1-H_{XY})}
+
C h^{H_Y}
\]
for some finite constant \(C\), uniformly for \(0<h\le \widetilde h_{0,\alpha}\).  Factoring out
\(h^{H_{XY}}\) and using the definition of \(\varepsilon\) gives
\eqref{eq:observable-screening-rate-F}.  Proposition~\ref{prop:graph-dual} then yields
\eqref{eq:observable-screening-rate-G}.
\end{proof}

\section{Reduction of smooth graph localization to exact micro correction}
\label{sec:smooth-cutoff-reduction}

The last reduction is intrinsic to the micro analysis.  A smooth time-side cutoff near \(t\) becomes
a slowly varying cutoff on the moving micro window \([0,\varepsilon_{h,\alpha}^{-1}]\).  Observable
graph localization is reduced to producing one exact corrected micro output on that window with the
required weighted boundary-layer control.

\begin{proposition}[Exact full \(h\)-scale trace identity for a graph element]
\label{prop:full-h-trace}
Assume Assumptions~\ref{ass:model} and~\ref{ass:obs-source}, let \(\eta>0\), fix \(t\in(0,T)\),
and suppose
\[
0<H_{XY}<H_Y<\tfrac12.
\]
Set
\[
\beta:=H_Y-H_{XY}.
\]
For \(i\in\{Y,XY\}\), define the full \(h\)-scale Volterra operator
\[
\mathfrak V_{i,t,h}^{\mathrm{full}}:L^2(\R_+)\to L^2(\R_+)
\]
by
\begin{equation}\label{eq:full-h-volterra}
\bigl(
\mathfrak V_{i,t,h}^{\mathrm{full}}\Theta
\bigr)(r)
:=
\int_r^{t/h}
(q-r)^{H_i-\frac12}
L_i(t-hr,t-hq)\Theta(q)\,dq,
\qquad r\ge0.
\end{equation}
Then:
\begin{enumerate}[label=(\roman*)]
\item for every \(f\in H_t\),
\begin{equation}\label{eq:full-h-representation}
\mathcal D_{t,h}\mathcal K_{i,t}f
=
h^{H_i}
\mathfrak V_{i,t,h}^{\mathrm{full}}
\bigl(
\mathcal D_{t,h}f
\bigr)
\qquad\text{in }L^2(\R_+);
\end{equation}
\item for every \(v\in \mathfrak G_t\), if
\begin{equation}\label{eq:full-h-traces}
\Phi_{v,h}:=\mathcal D_{t,h}v,
\qquad
\Psi_{v,h}:=h^\beta \mathcal D_{t,h}\mathcal S_t v,
\end{equation}
then
\begin{equation}\label{eq:full-h-exact}
\mathfrak V_{Y,t,h}^{\mathrm{full}}\Psi_{v,h}
=
\mathfrak V_{XY,t,h}^{\mathrm{full}}\Phi_{v,h}
\qquad\text{in }L^2(\R_+);
\end{equation}
\item if \(\chi\in C^\infty([0,1])\), \(0\le \chi\le1\), and
\[
\chi\equiv1\ \text{on }[0,1-\delta_*]
\qquad\text{for some }\delta_*\in(0,1),
\]
then with \(\widetilde\chi_{h,\alpha}\) from Proposition~\ref{prop:micro-cutoff-scaling},
\begin{equation}\label{eq:full-h-time-cutoff}
\mathcal L_{t,h}
\bigl(
\widetilde\chi_{h,\alpha}\Phi_{v,h}
\bigr)
=
\chi_{h,\alpha}^{\mathrm{bdry}}v,
\qquad
\chi_{h,\alpha}^{\mathrm{bdry}}(s)
:=
\chi\!\left(\frac{t-s}{h^\alpha}\right)\one_{[t-h^\alpha,t]}(s).
\end{equation}
\end{enumerate}
\end{proposition}

\begin{proof}
Fix \(i\in\{Y,XY\}\) and \(f\in H_t\). For \(0\le r\le t/h\),
\begin{align*}
\bigl(\mathcal D_{t,h}\mathcal K_{i,t}f\bigr)(r)
&=
h^{1/2}
\int_0^{t-hr}
(t-hr-s)^{H_i-\frac12}
L_i(t-hr,s)f(s)\,ds.
\end{align*}
With the change of variables \(s=t-hq\), \(ds=-h\,dq\), this becomes
\begin{align*}
\bigl(\mathcal D_{t,h}\mathcal K_{i,t}f\bigr)(r)
&=
h^{H_i}
\int_r^{t/h}
(q-r)^{H_i-\frac12}
L_i(t-hr,t-hq)\,
h^{1/2}f(t-hq)\,dq \\
&=
h^{H_i}
\bigl(
\mathfrak V_{i,t,h}^{\mathrm{full}}(\mathcal D_{t,h}f)
\bigr)(r).
\end{align*}
For \(r>t/h\), both sides vanish. This proves \eqref{eq:full-h-representation}.

Now let \(v\in\mathfrak G_t\). Since \(\mathcal K_{Y,t}\mathcal S_t v=\mathcal K_{XY,t}v\), applying
\eqref{eq:full-h-representation} with \(f=\mathcal S_t v\) and \(f=v\) gives
\[
h^{H_Y}
\mathfrak V_{Y,t,h}^{\mathrm{full}}
\bigl(
\mathcal D_{t,h}\mathcal S_t v
\bigr)
=
h^{H_{XY}}
\mathfrak V_{XY,t,h}^{\mathrm{full}}
\bigl(
\mathcal D_{t,h}v
\bigr).
\]
Since \(\beta=H_Y-H_{XY}\), this is \eqref{eq:full-h-exact}.

Finally, for \(0\le s\le t\),
\begin{align*}
\bigl(
\mathcal L_{t,h}
(\widetilde\chi_{h,\alpha}\Phi_{v,h})
\bigr)(s)
&=
h^{-1/2}
\widetilde\chi_{h,\alpha}\!\left(\frac{t-s}{h}\right)
\Phi_{v,h}\!\left(\frac{t-s}{h}\right)\one_{[0,t]}(s) \\
&=
\widetilde\chi_{h,\alpha}\!\left(\frac{t-s}{h}\right)
v(s)\one_{[0,t]}(s).
\end{align*}
By Proposition~\ref{prop:micro-cutoff-scaling},
\[
\widetilde\chi_{h,\alpha}\!\left(\frac{t-s}{h}\right)
=
\chi\!\left(
\varepsilon_{h,\alpha}\frac{t-s}{h}
\right)
\one_{\left\{0\le \frac{t-s}{h}\le \varepsilon_{h,\alpha}^{-1}\right\}}
=
\chi\!\left(\frac{t-s}{h^\alpha}\right)\one_{[t-h^\alpha,t]}(s).
\]
This proves \eqref{eq:full-h-time-cutoff}.
\end{proof}

\begin{theorem}[Scaled exact micro-correction criterion for boundary localization]
\label{thm:smooth-cutoff-criterion}
Assume Assumptions~\ref{ass:model} and~\ref{ass:obs-source}, let \(\eta>0\), fix \(t\in(0,T)\),
and suppose
\[
0<H_{XY}<H_Y<\tfrac12.
\]
Fix \(\alpha\in(0,1)\), choose \(\chi\in C^\infty([0,1])\) with
\[
0\le \chi\le1,
\qquad
\chi\equiv1\ \text{on }[0,1-\delta_*]
\]
for some \(\delta_*\in(0,1)\), and let \(\widetilde\chi_{h,\alpha}\) be the corresponding moving
micro cutoff from Proposition~\ref{prop:micro-cutoff-scaling}. Assume in addition that there exist
\[
C_{\mathrm{mc},\alpha}<\infty,
\qquad
h_{0,\alpha}\in(0,t]
\]
such that for every \(0<h\le h_{0,\alpha}\) and every \(v\in\mathfrak G_t\), setting
\[
\Phi_{v,h}:=\mathcal D_{t,h}v,
\]
there exists
\[
\widetilde\Psi_{v,h}\in L^2(\R_+)
\]
with
\begin{equation}\label{eq:smooth-cutoff-criterion-exact}
\operatorname{supp}(\widetilde\Psi_{v,h})
\subset [0,\varepsilon_{h,\alpha}^{-1}],
\qquad
\mathfrak V_{Y,t,h,\alpha}^{\mathrm{mic}}\widetilde\Psi_{v,h}
=
\mathfrak V_{XY,t,h,\alpha}^{\mathrm{mic}}
\bigl(
\widetilde\chi_{h,\alpha}\Phi_{v,h}
\bigr)
\end{equation}
and
\begin{equation}\label{eq:smooth-cutoff-criterion-bound}
\left\|
\widetilde\chi_{h,\alpha}\Phi_{v,h}
\right\|_{L^2(\R_+)}^2
+
\frac{\eta^2}{\nu_Y^2}
\left\|
\widetilde\Psi_{v,h}
\right\|_{L^2(\R_+)}^2
\le
C_{\mathrm{mc},\alpha}^2\|v\|_{\mathfrak G_t}^2.
\end{equation}
Then, with the time-side cutoff
\[
\chi_{h,\alpha}:=\chi_{h,\alpha}^{\mathrm{bdry}}
\]
from \eqref{eq:full-h-time-cutoff}, one has for every \(0<h\le h_{0,\alpha}\) and every
\(v\in\mathfrak G_t\),
\begin{equation}\label{eq:smooth-cutoff-criterion-scaled}
\chi_{h,\alpha}v\in \mathfrak G_t\cap H_{t,h,\alpha}^{\mathrm{near}}
\end{equation}
and
\begin{equation}\label{eq:smooth-cutoff-criterion-scaled-bound}
\|\chi_{h,\alpha}v\|_{H_t}^2
+
\frac{\eta^2}{\nu_Y^2}
h^{2\beta}
\|\mathcal S_t(\chi_{h,\alpha}v)\|_{H_t}^2
\le
C_{\mathrm{mc},\alpha}^2\|v\|_{\mathfrak G_t}^2.
\end{equation}
\end{theorem}

\begin{proof}
Fix \(0<h\le h_{0,\alpha}\) and \(v\in\mathfrak G_t\), and define
\[
\widetilde\Phi_{v,h}:=\widetilde\chi_{h,\alpha}\Phi_{v,h}.
\]
Set
\[
\phi_{v,h}:=
\mathfrak l_{\varepsilon_{h,\alpha}}\widetilde\Phi_{v,h}\in E,
\qquad
\psi_{v,h}:=
\varepsilon_{h,\alpha}^{-\beta}
\mathfrak l_{\varepsilon_{h,\alpha}}\widetilde\Psi_{v,h}\in E.
\]
Applying Proposition~\ref{prop:micro-volterra} to
\eqref{eq:smooth-cutoff-criterion-exact}, and using
\(\mathfrak l_{\varepsilon_{h,\alpha}}\mathfrak d_{\varepsilon_{h,\alpha}}=I_E\) from
Lemma~\ref{lem:micro-dilation}, gives
\[
\mathfrak A_{Y,t,h,\alpha}\psi_{v,h}
=
\mathfrak A_{XY,t,h,\alpha}\phi_{v,h}.
\]
Hence
\[
(\phi_{v,h},\psi_{v,h})\in \Gamma_{t,h,\alpha}^{\mathrm{loc,ex}}.
\]
By Proposition~\ref{prop:local-exact-identification},
\[
w_{v,h}:=
\mathcal L_{t,h}^{(\alpha)}\phi_{v,h}
\in
\mathfrak G_t\cap H_{t,h,\alpha}^{\mathrm{near}}
\]
Because \(\widetilde\Phi_{v,h}\) and \(\widetilde\Psi_{v,h}\) are supported in
\([0,\varepsilon_{h,\alpha}^{-1}]\), Proposition~\ref{prop:two-scale-factorization}(ii) gives
\[
\mathcal L_{t,h}^{(\alpha)}\phi_{v,h}
=
\mathcal L_{t,h}\widetilde\Phi_{v,h},
\]
so Proposition~\ref{prop:full-h-trace} yields
\[
w_{v,h}
=
\mathcal L_{t,h}\widetilde\Phi_{v,h}
=
\chi_{h,\alpha}^{\mathrm{bdry}}v.
\]
By Proposition~\ref{prop:local-exact-identification},
\[
\mathcal S_t w_{v,h}
=
h^{-\alpha\beta}\mathcal L_{t,h}^{(\alpha)}\psi_{v,h}
=
h^{-\beta}\mathcal L_{t,h}\widetilde\Psi_{v,h},
\]
where we also used \(\varepsilon_{h,\alpha}=h^{1-\alpha}\). Lemma~\ref{lem:dilation-isometry}
therefore gives
\begin{equation}\label{eq:smooth-cutoff-criterion-norm}
\|w_{v,h}\|_{H_t}^2
+
\frac{\eta^2}{\nu_Y^2}
h^{2\beta}\|\mathcal S_t w_{v,h}\|_{H_t}^2
=
\left\|
\widetilde\Phi_{v,h}
\right\|_{L^2(\R_+)}^2
+
\frac{\eta^2}{\nu_Y^2}
\left\|
\widetilde\Psi_{v,h}
\right\|_{L^2(\R_+)}^2.
\end{equation}
Combining
\eqref{eq:smooth-cutoff-criterion-bound} with
\eqref{eq:smooth-cutoff-criterion-norm} yields
\[
\chi_{h,\alpha}^{\mathrm{bdry}}v\in \mathfrak G_t\cap H_{t,h,\alpha}^{\mathrm{near}},
\]
and
\[
\|\chi_{h,\alpha}^{\mathrm{bdry}}v\|_{H_t}^2
+
\frac{\eta^2}{\nu_Y^2}
h^{2\beta}
\|\mathcal S_t(\chi_{h,\alpha}^{\mathrm{bdry}}v)\|_{H_t}^2
\le
C_{\mathrm{mc},\alpha}^2\|v\|_{\mathfrak G_t}^2.
\]
By construction,
\[
\chi_{h,\alpha}^{\mathrm{bdry}}(s)
=
\chi\!\left(\frac{t-s}{h^\alpha}\right)\one_{[t-h^\alpha,t]}(s),
\]
so
\[
0\le \chi_{h,\alpha}^{\mathrm{bdry}}\le1,
\qquad
\chi_{h,\alpha}^{\mathrm{bdry}}\equiv1
\ \text{on }[t-(1-\delta_*)h^\alpha,t],
\qquad
\operatorname{supp}\chi_{h,\alpha}^{\mathrm{bdry}}\subset[t-h^\alpha,t].
\]
This proves \eqref{eq:smooth-cutoff-criterion-scaled} and
\eqref{eq:smooth-cutoff-criterion-scaled-bound}.
\end{proof}

\begin{remark}[Exact micro correction and smooth localization]
\label{rem:smooth-cutoff-criterion}
Theorem~\ref{thm:smooth-cutoff-criterion} supplies the scaled boundary-localization mechanism used
in the closing reduction. It already gives the correct time-side cutoff in the boundary-layer norm, and isolates
the only missing upgrade to Definition~\ref{def:alpha-cutoff-smooth}: the factor \(h^\beta\) in
front of \(\|\mathcal S_t(\chi_{h,\alpha}v)\|_{H_t}\). By Proposition~\ref{prop:full-h-trace}, the
full-\(h\) trace pair is exact before cutoff, so that upgrade is encoded entirely by the moving
exact-cutoff defect. The later sections resolve it by passing to the growing-window exact defect and
comparing that object with its frozen commutator model.
\end{remark}

\begin{definition}[Moving exact-cutoff defect for a graph element]
\label{def:moving-cutoff-defect}
Assume Assumptions~\ref{ass:model} and~\ref{ass:obs-source}, let \(\eta>0\), fix \(t\in(0,T)\),
let \(\alpha\in(0,1)\), and suppose
\[
0<H_{XY}<H_Y<\tfrac12.
\]
Set
\[
\beta:=H_Y-H_{XY}.
\]
Let \(\chi\in C^\infty([0,1])\) satisfy
\[
0\le \chi\le1,
\qquad
\chi\equiv1\ \text{on }[0,1-\delta_*]
\]
for some \(\delta_*\in(0,1)\), and let \(\widetilde\chi_{h,\alpha}\) be the associated moving micro
cutoff from Proposition~\ref{prop:micro-cutoff-scaling}. For \(v\in\mathfrak G_t\) and \(h>0\),
define the exact full-\(h\) traces
\[
\Phi_{v,h}:=\mathcal D_{t,h}v,
\qquad
\Psi_{v,h}:=h^\beta\mathcal D_{t,h}\mathcal S_t v,
\]
and the corresponding moving exact-cutoff defect by
\begin{equation}\label{eq:moving-cutoff-defect}
\mathfrak R_{v,h,\alpha}^{\mathrm{cut}}
:=
\mathfrak V_{XY,t,h,\alpha}^{\mathrm{mic}}
\bigl(
\widetilde\chi_{h,\alpha}\Phi_{v,h}
\bigr)
-
\mathfrak V_{Y,t,h,\alpha}^{\mathrm{mic}}
\bigl(
\widetilde\chi_{h,\alpha}\Psi_{v,h}
\bigr).
\end{equation}
\end{definition}

\begin{proposition}[Exact commutator-tail decomposition of the moving exact-cutoff defect]
\label{prop:moving-cutoff-decomposition}
Assume the hypotheses of Definition~\ref{def:moving-cutoff-defect}. Let \(v\in\mathfrak G_t\),
let \(0<h\le t\) satisfy \(h^\alpha\le t\), and set
\[
\varepsilon:=\varepsilon_{h,\alpha}=h^{1-\alpha}.
\]
With \(\Phi_{v,h}\), \(\Psi_{v,h}\), and \(\widetilde\chi_{h,\alpha}\) as above, define
\begin{align}
\bigl(
\mathfrak C_{XY,v,h,\alpha}^{\mathrm{cut}}
\bigr)(r)
&:=
\int_r^{\varepsilon^{-1}}
(q-r)^{H_{XY}-\frac12}
\bigl(
\widetilde\chi_{h,\alpha}(q)-\widetilde\chi_{h,\alpha}(r)
\bigr)
L_{XY}(t-hr,t-hq)\Phi_{v,h}(q)\,dq,
\label{eq:moving-cutoff-com-XY}\\
\bigl(
\mathfrak C_{Y,v,h,\alpha}^{\mathrm{cut}}
\bigr)(r)
&:=
\int_r^{\varepsilon^{-1}}
(q-r)^{H_Y-\frac12}
\bigl(
\widetilde\chi_{h,\alpha}(q)-\widetilde\chi_{h,\alpha}(r)
\bigr)
L_Y(t-hr,t-hq)\Psi_{v,h}(q)\,dq,
\label{eq:moving-cutoff-com-Y}
\end{align}
for \(r\ge0\), and define the old-tail forcing by
\begin{equation}\label{eq:moving-cutoff-old-tail}
\begin{aligned}
\bigl(
\mathfrak O_{v,h,\alpha}^{\mathrm{old}}
\bigr)(r)
&:=
\one_{[0,\varepsilon^{-1}]}(r)
\int_{\varepsilon^{-1}}^{t/h}
\Bigl[
(q-r)^{H_Y-\frac12}
L_Y(t-hr,t-hq)\Psi_{v,h}(q)
\\
&\qquad\qquad\qquad
-
(q-r)^{H_{XY}-\frac12}
L_{XY}(t-hr,t-hq)\Phi_{v,h}(q)
\Bigr]\,dq.
\end{aligned}
\end{equation}
Then:
\begin{enumerate}[label=(\roman*)]
\item one has the exact decomposition
\begin{equation}\label{eq:moving-cutoff-decomposition}
\mathfrak R_{v,h,\alpha}^{\mathrm{cut}}
=
\mathfrak C_{XY,v,h,\alpha}^{\mathrm{cut}}
-
\mathfrak C_{Y,v,h,\alpha}^{\mathrm{cut}}
+
\widetilde\chi_{h,\alpha}\mathfrak O_{v,h,\alpha}^{\mathrm{old}}
\qquad\text{in }L^2(\R_+);
\end{equation}
\item defining the unit-interval traces
\begin{equation}\label{eq:moving-cutoff-local-traces}
\phi_{v,h}^{\sharp}
:=
\mathfrak l_{\varepsilon}\Phi_{v,h}\in E,
\qquad
\psi_{v,h}^{\sharp}
:=
\varepsilon^{-\beta}\mathfrak l_{\varepsilon}\Psi_{v,h}\in E,
\end{equation}
and the rescaled old-tail source
\begin{equation}\label{eq:moving-cutoff-old-tail-local}
\omega_{v,h,\alpha}^{\mathrm{old}}
:=
\varepsilon^{H_{XY}}
\mathfrak l_{\varepsilon}
\mathfrak O_{v,h,\alpha}^{\mathrm{old}}
\in E,
\end{equation}
one has
\begin{equation}\label{eq:moving-cutoff-old-tail-local-id}
\mathfrak A_{XY,t,h,\alpha}\phi_{v,h}^{\sharp}
-
\mathfrak A_{Y,t,h,\alpha}\psi_{v,h}^{\sharp}
=
\omega_{v,h,\alpha}^{\mathrm{old}}
\qquad\text{in }E;
\end{equation}
\item with the same cutoff \(\chi\) on \([0,1]\) used to define \(\widetilde\chi_{h,\alpha}\),
\begin{equation}\label{eq:moving-cutoff-unit-reduction}
\varepsilon^{H_{XY}}
\mathfrak l_{\varepsilon}
\mathfrak R_{v,h,\alpha}^{\mathrm{cut}}
=
\mathfrak W_{XY,t,h,\alpha,\chi}^{\mathrm{loc}}
\phi_{v,h}^{\sharp}
-
\mathfrak W_{Y,t,h,\alpha,\chi}^{\mathrm{loc}}
\psi_{v,h}^{\sharp}
+
\chi\,\omega_{v,h,\alpha}^{\mathrm{old}}
\qquad\text{in }E.
\end{equation}
\end{enumerate}
\end{proposition}

\begin{proof}
For \(r>\varepsilon^{-1}\), all terms in \eqref{eq:moving-cutoff-decomposition} vanish, so it is
enough to work on \(0\le r\le \varepsilon^{-1}\). There,
\begin{align*}
\mathfrak R_{v,h,\alpha}^{\mathrm{cut}}(r)
&=
\int_r^{\varepsilon^{-1}}
(q-r)^{H_{XY}-\frac12}
L_{XY}(t-hr,t-hq)\widetilde\chi_{h,\alpha}(q)\Phi_{v,h}(q)\,dq \\
&\quad-
\int_r^{\varepsilon^{-1}}
(q-r)^{H_Y-\frac12}
L_Y(t-hr,t-hq)\widetilde\chi_{h,\alpha}(q)\Psi_{v,h}(q)\,dq.
\end{align*}
Adding and subtracting \(\widetilde\chi_{h,\alpha}(r)\) inside both integrals gives
\begin{align*}
\mathfrak R_{v,h,\alpha}^{\mathrm{cut}}(r)
&=
\mathfrak C_{XY,v,h,\alpha}^{\mathrm{cut}}(r)
-
\mathfrak C_{Y,v,h,\alpha}^{\mathrm{cut}}(r) \\
&\quad+
\widetilde\chi_{h,\alpha}(r)
\Biggl[
\int_r^{\varepsilon^{-1}}
(q-r)^{H_{XY}-\frac12}
L_{XY}(t-hr,t-hq)\Phi_{v,h}(q)\,dq \\
&\hspace{5.0em}-
\int_r^{\varepsilon^{-1}}
(q-r)^{H_Y-\frac12}
L_Y(t-hr,t-hq)\Psi_{v,h}(q)\,dq
\Biggr].
\end{align*}
By Proposition~\ref{prop:full-h-trace}(ii),
\[
\int_r^{t/h}
(q-r)^{H_{XY}-\frac12}
L_{XY}(t-hr,t-hq)\Phi_{v,h}(q)\,dq
=
\int_r^{t/h}
(q-r)^{H_Y-\frac12}
L_Y(t-hr,t-hq)\Psi_{v,h}(q)\,dq.
\]
Subtracting the common tail from \(q=\varepsilon^{-1}\) to \(q=t/h\) yields
\[
\begin{aligned}
&\int_r^{\varepsilon^{-1}}
(q-r)^{H_{XY}-\frac12}
\\
&\qquad\times
L_{XY}(t-hr,t-hq)\Phi_{v,h}(q)\,dq
\\
&\quad
-\int_r^{\varepsilon^{-1}}
(q-r)^{H_Y-\frac12}
\\
&\qquad\quad\times
L_Y(t-hr,t-hq)\Psi_{v,h}(q)\,dq
\\
&=
\mathfrak O_{v,h,\alpha}^{\mathrm{old}}(r),
\end{aligned}
\]
which proves \eqref{eq:moving-cutoff-decomposition}.

For item~(ii), Proposition~\ref{prop:micro-volterra} gives
\[
\mathfrak d_{\varepsilon}
\mathfrak A_{XY,t,h,\alpha}\phi_{v,h}^{\sharp}
=
\varepsilon^{H_{XY}}
\mathfrak V_{XY,t,h,\alpha}^{\mathrm{mic}}\Phi_{v,h},
\]
and, since \(\beta=H_Y-H_{XY}\),
\[
\mathfrak d_{\varepsilon}
\mathfrak A_{Y,t,h,\alpha}\psi_{v,h}^{\sharp}
=
\varepsilon^{H_Y-\beta}
\mathfrak V_{Y,t,h,\alpha}^{\mathrm{mic}}\Psi_{v,h}
=
\varepsilon^{H_{XY}}
\mathfrak V_{Y,t,h,\alpha}^{\mathrm{mic}}\Psi_{v,h}.
\]
Therefore
\[
\mathfrak d_{\varepsilon}
\bigl(
\mathfrak A_{XY,t,h,\alpha}\phi_{v,h}^{\sharp}
-
\mathfrak A_{Y,t,h,\alpha}\psi_{v,h}^{\sharp}
\bigr)
=
\varepsilon^{H_{XY}}
\bigl(
\mathfrak V_{XY,t,h,\alpha}^{\mathrm{mic}}\Phi_{v,h}
-
\mathfrak V_{Y,t,h,\alpha}^{\mathrm{mic}}\Psi_{v,h}
\bigr).
\]
By the full-tail identity established in the proof of item~(i), the bracket on the right is
exactly \(\mathfrak O_{v,h,\alpha}^{\mathrm{old}}\). Applying \(\mathfrak l_{\varepsilon}\) and using
\(\mathfrak l_{\varepsilon}\mathfrak d_{\varepsilon}=I_E\) from
Lemma~\ref{lem:micro-dilation} proves \eqref{eq:moving-cutoff-old-tail-local-id}.

For item~(iii), Proposition~\ref{prop:micro-cutoff-scaling} gives
\[
\mathfrak l_{\varepsilon}
\bigl(
\widetilde\chi_{h,\alpha}\Phi_{v,h}
\bigr)
=
\chi\,\phi_{v,h}^{\sharp},
\qquad
\mathfrak l_{\varepsilon}
\bigl(
\widetilde\chi_{h,\alpha}\Psi_{v,h}
\bigr)
=
\varepsilon^{\beta}\chi\,\psi_{v,h}^{\sharp}.
\]
Applying Proposition~\ref{prop:micro-volterra} to these two identities yields
\[
\varepsilon^{H_{XY}}
\mathfrak l_{\varepsilon}
\mathfrak R_{v,h,\alpha}^{\mathrm{cut}}
=
\mathfrak A_{XY,t,h,\alpha}(\chi\phi_{v,h}^{\sharp})
-
\mathfrak A_{Y,t,h,\alpha}(\chi\psi_{v,h}^{\sharp}).
\]
Now expand the right-hand side by adding and subtracting
\(\chi \mathfrak A_{XY,t,h,\alpha}\phi_{v,h}^{\sharp}\) and
\(\chi \mathfrak A_{Y,t,h,\alpha}\psi_{v,h}^{\sharp}\):
\begin{align*}
\mathfrak A_{XY,t,h,\alpha}(\chi\phi_{v,h}^{\sharp})
-
\mathfrak A_{Y,t,h,\alpha}(\chi\psi_{v,h}^{\sharp})
&=
\bigl(
\mathfrak A_{XY,t,h,\alpha}(\chi\phi_{v,h}^{\sharp})
-
\chi\,\mathfrak A_{XY,t,h,\alpha}\phi_{v,h}^{\sharp}
\bigr) \\
&\quad-
\bigl(
\mathfrak A_{Y,t,h,\alpha}(\chi\psi_{v,h}^{\sharp})
-
\chi\,\mathfrak A_{Y,t,h,\alpha}\psi_{v,h}^{\sharp}
\bigr) \\
&\quad+
\chi\bigl(
\mathfrak A_{XY,t,h,\alpha}\phi_{v,h}^{\sharp}
-
\mathfrak A_{Y,t,h,\alpha}\psi_{v,h}^{\sharp}
\bigr).
\end{align*}
By the definitions \eqref{eq:local-com-Y}--\eqref{eq:local-com-XY}, the first two brackets equal
\[
\mathfrak W_{XY,t,h,\alpha,\chi}^{\mathrm{loc}}\phi_{v,h}^{\sharp},
\qquad
\mathfrak W_{Y,t,h,\alpha,\chi}^{\mathrm{loc}}\psi_{v,h}^{\sharp},
\]
respectively. Using \eqref{eq:moving-cutoff-old-tail-local-id} for the last term gives
\eqref{eq:moving-cutoff-unit-reduction}.
\end{proof}

\begin{definition}[Unit-interval source for the moving cutoff problem]
\label{def:moving-cutoff-local-source}
Under the hypotheses and notation of Proposition~\ref{prop:moving-cutoff-decomposition}, define
the unit-interval source term by
\begin{equation}\label{eq:moving-cutoff-local-source}
\Sigma_{v,h,\alpha}^{\mathrm{cut}}
:=
\mathfrak W_{XY,t,h,\alpha,\chi}^{\mathrm{loc}}
\phi_{v,h}^{\sharp}
-
\mathfrak W_{Y,t,h,\alpha,\chi}^{\mathrm{loc}}
\psi_{v,h}^{\sharp}
+
\chi\,\omega_{v,h,\alpha}^{\mathrm{old}}
\in E.
\end{equation}
\end{definition}

\begin{theorem}[Unscaled exact-correction criterion for boundary graph localization]
\label{thm:unscaled-smooth-cutoff-criterion}
Assume Assumptions~\ref{ass:model} and~\ref{ass:obs-source}, let \(\eta>0\), fix \(t\in(0,T)\),
and suppose
\[
0<H_{XY}<H_Y<\tfrac12.
\]
Fix \(\alpha\in(0,1)\), choose \(\chi\in C^\infty([0,1])\) with
\[
0\le \chi\le1,
\qquad
\chi\equiv1\ \text{on }[0,1-\delta_*]
\]
for some \(\delta_*\in(0,1)\), and let \(\widetilde\chi_{h,\alpha}\) be the associated moving micro
cutoff from Proposition~\ref{prop:micro-cutoff-scaling}. Assume that there exist
\[
C_{\mathrm{tail},\alpha}<\infty,
\qquad
h_{0,\alpha}\in(0,t]
\]
such that for every \(0<h\le h_{0,\alpha}\) and every \(v\in\mathfrak G_t\), with
\(\Phi_{v,h}\), \(\Psi_{v,h}\), and \(\mathfrak R_{v,h,\alpha}^{\mathrm{cut}}\) as in
Definition~\ref{def:moving-cutoff-defect}, there exists
\[
\mathfrak C_{v,h}\in L^2(\R_+)
\]
satisfying
\begin{equation}\label{eq:unscaled-cutoff-correction-eq}
\operatorname{supp}(\mathfrak C_{v,h})
\subset [0,\varepsilon_{h,\alpha}^{-1}],
\qquad
\mathfrak V_{Y,t,h,\alpha}^{\mathrm{mic}}\mathfrak C_{v,h}
=
\mathfrak R_{v,h,\alpha}^{\mathrm{cut}}
\end{equation}
and
\begin{equation}\label{eq:unscaled-cutoff-correction-bound}
\|\mathfrak C_{v,h}\|_{L^2(\R_+)}
\le
C_{\mathrm{tail},\alpha}h^\beta\|v\|_{\mathfrak G_t}.
\end{equation}
Then \(\mathcal S_t\) has \(h^\alpha\)-stable boundary cutoffs at time \(t\) in the sense of
Definition~\ref{def:alpha-cutoff-smooth}. The cutoff can be chosen with constants
\[
C_{\mathrm{scut},\alpha}<\infty,
\qquad
h_{1,\alpha}\in(0,t]
\]
such that for every \(0<h\le h_{1,\alpha}\), the cutoff
\[
\chi_{h,\alpha}(s)
:=
\chi\!\left(\frac{t-s}{h^\alpha}\right)\one_{[t-h^\alpha,t]}(s)
\]
satisfies
\[
0\le \chi_{h,\alpha}\le1,
\qquad
\chi_{h,\alpha}\equiv1\ \text{on }[t-(1-\delta_*)h^\alpha,t],
\qquad
\operatorname{supp}\chi_{h,\alpha}\subset[t-h^\alpha,t],
\]
and for every \(v\in\mathfrak G_t\),
\begin{equation}\label{eq:unscaled-cutoff-graph}
\chi_{h,\alpha}v\in\mathfrak G_t
\qquad\text{and}\qquad
\|\chi_{h,\alpha}v\|_{\mathfrak G_t}
\le
C_{\mathrm{scut},\alpha}\|v\|_{\mathfrak G_t}.
\end{equation}
\end{theorem}

\begin{proof}
Set
\[
\beta:=H_Y-H_{XY}.
\]
Fix \(0<h\le h_{0,\alpha}\) and \(v\in\mathfrak G_t\), and abbreviate
\[
\widetilde\Phi_{v,h}:=\widetilde\chi_{h,\alpha}\Phi_{v,h},
\qquad
\widetilde\Psi_{v,h}:=
\widetilde\chi_{h,\alpha}\Psi_{v,h}
+
\mathfrak C_{v,h}.
\]
By \eqref{eq:unscaled-cutoff-correction-eq} and
Definition~\ref{def:moving-cutoff-defect},
\[
\mathfrak V_{Y,t,h,\alpha}^{\mathrm{mic}}\widetilde\Psi_{v,h}
=
\mathfrak V_{XY,t,h,\alpha}^{\mathrm{mic}}\widetilde\Phi_{v,h}.
\]
Moreover,
\[
\operatorname{supp}(\widetilde\Phi_{v,h})
\subset[0,\varepsilon_{h,\alpha}^{-1}],
\qquad
\operatorname{supp}(\widetilde\Psi_{v,h})
\subset[0,\varepsilon_{h,\alpha}^{-1}],
\]
because \(\widetilde\chi_{h,\alpha}\) and \(\mathfrak C_{v,h}\) are supported in
\([0,\varepsilon_{h,\alpha}^{-1}]\).

By Lemma~\ref{lem:dilation-isometry},
\[
\|\widetilde\Phi_{v,h}\|_{L^2(\R_+)}
\le
\|\Phi_{v,h}\|_{L^2(\R_+)}
=
\|v\|_{H_t}
\le
\|v\|_{\mathfrak G_t}.
\]
Likewise,
\[
\|\widetilde\chi_{h,\alpha}\Psi_{v,h}\|_{L^2(\R_+)}
\le
\|\Psi_{v,h}\|_{L^2(\R_+)}
=
h^\beta\|\mathcal S_t v\|_{H_t}
\le
\frac{\nu_Y}{\eta}h^\beta\|v\|_{\mathfrak G_t}.
\]
Combining this with \eqref{eq:unscaled-cutoff-correction-bound} yields
\begin{equation}\label{eq:unscaled-cutoff-output-bound}
\|\widetilde\Psi_{v,h}\|_{L^2(\R_+)}
\le
\left(
\frac{\nu_Y}{\eta}
+
C_{\mathrm{tail},\alpha}
\right)
h^\beta\|v\|_{\mathfrak G_t}.
\end{equation}
Hence there exists \(C_{\mathrm{mc},\alpha}<\infty\) such that
\begin{equation}\label{eq:unscaled-cutoff-micro-bound}
\|\widetilde\Phi_{v,h}\|_{L^2(\R_+)}^2
+
\frac{\eta^2}{\nu_Y^2}
h^{-2\beta}
\|\widetilde\Psi_{v,h}\|_{L^2(\R_+)}^2
\le
C_{\mathrm{mc},\alpha}^2\|v\|_{\mathfrak G_t}^2.
\end{equation}

Set
\[
\phi_{v,h}:=
\mathfrak l_{\varepsilon_{h,\alpha}}\widetilde\Phi_{v,h}\in E,
\qquad
\psi_{v,h}:=
\varepsilon_{h,\alpha}^{-\beta}
\mathfrak l_{\varepsilon_{h,\alpha}}\widetilde\Psi_{v,h}\in E.
\]
Applying Proposition~\ref{prop:micro-volterra} to the exact micro identity above, and using
\(\mathfrak l_{\varepsilon_{h,\alpha}}\mathfrak d_{\varepsilon_{h,\alpha}}=I_E\) from
Lemma~\ref{lem:micro-dilation}, gives
\[
\mathfrak A_{Y,t,h,\alpha}\psi_{v,h}
=
\mathfrak A_{XY,t,h,\alpha}\phi_{v,h}.
\]
So
\[
(\phi_{v,h},\psi_{v,h})\in\Gamma_{t,h,\alpha}^{\mathrm{loc,ex}}.
\]
By Proposition~\ref{prop:local-exact-identification},
\[
w_{v,h}:=\mathcal L_{t,h}^{(\alpha)}\phi_{v,h}\in \mathfrak G_t\cap H_{t,h,\alpha}^{\mathrm{near}}
\]
and
\[
\mathcal S_t w_{v,h}
=
h^{-\alpha\beta}\mathcal L_{t,h}^{(\alpha)}\psi_{v,h}.
\]
Because \(\widetilde\Phi_{v,h}\) and \(\widetilde\Psi_{v,h}\) are supported in
\([0,\varepsilon_{h,\alpha}^{-1}]\), Proposition~\ref{prop:two-scale-factorization}(ii) gives
\[
\mathcal L_{t,h}^{(\alpha)}\phi_{v,h}
=
\mathcal L_{t,h}\widetilde\Phi_{v,h},
\qquad
\mathcal L_{t,h}^{(\alpha)}\psi_{v,h}
=
\varepsilon_{h,\alpha}^{-\beta}\mathcal L_{t,h}\widetilde\Psi_{v,h}.
\]
Since \(\varepsilon_{h,\alpha}=h^{1-\alpha}\), it follows that
\[
\mathcal S_t w_{v,h}
=
h^{-\beta}\mathcal L_{t,h}\widetilde\Psi_{v,h}.
\]
Also, Proposition~\ref{prop:full-h-trace} gives
\[
w_{v,h}
=
\mathcal L_{t,h}\widetilde\Phi_{v,h}
=
\chi_{h,\alpha}v.
\]
Therefore \(\chi_{h,\alpha}v\in\mathfrak G_t\). Using Lemma~\ref{lem:dilation-isometry},
\begin{align*}
\|\chi_{h,\alpha}v\|_{\mathfrak G_t}^2
&=
\|w_{v,h}\|_{H_t}^2
+
\frac{\eta^2}{\nu_Y^2}
\|\mathcal S_t w_{v,h}\|_{H_t}^2 \\
&=
\|\widetilde\Phi_{v,h}\|_{L^2(\R_+)}^2
+
\frac{\eta^2}{\nu_Y^2}
h^{-2\beta}
\|\widetilde\Psi_{v,h}\|_{L^2(\R_+)}^2 \\
&\le
C_{\mathrm{mc},\alpha}^2\|v\|_{\mathfrak G_t}^2
\end{align*}
by \eqref{eq:unscaled-cutoff-micro-bound}. This proves \eqref{eq:unscaled-cutoff-graph}.

Finally, by definition,
\[
\chi_{h,\alpha}(s)
=
\chi\!\left(\frac{t-s}{h^\alpha}\right)\one_{[t-h^\alpha,t]}(s),
\]
so the stated support and plateau properties hold. Thus Definition~\ref{def:alpha-cutoff-smooth}
is satisfied with \(h_{1,\alpha}:=h_{0,\alpha}\) and \(C_{\mathrm{scut},\alpha}:=C_{\mathrm{mc},\alpha}\).
\end{proof}

\begin{corollary}[Unit-interval sufficient criterion for boundary graph localization]
\label{cor:unit-interval-unscaled-criterion}
Assume the hypotheses of Proposition~\ref{prop:moving-cutoff-decomposition}. Fix
\[
\chi\in C^\infty([0,1]),
\qquad
0\le \chi\le1,
\qquad
\chi\equiv1\ \text{on }[0,1-\delta_*]
\]
for some \(\delta_*\in(0,1)\), and let \(\widetilde\chi_{h,\alpha}\) be the associated moving micro
cutoff. Assume that there exist
\[
C_{\mathrm{loc},\alpha}<\infty,
\qquad
h_{0,\alpha}\in(0,t]
\]
such that for every \(0<h\le h_{0,\alpha}\) and every \(v\in\mathfrak G_t\), with
\(\Sigma_{v,h,\alpha}^{\mathrm{cut}}\) from Definition~\ref{def:moving-cutoff-local-source},
there exists
\[
\zeta_{v,h}^{\mathrm{cut}}\in E
\]
such that
\begin{equation}\label{eq:unit-interval-unscaled-eq}
\mathfrak A_{Y,t,h,\alpha}\zeta_{v,h}^{\mathrm{cut}}
=
\Sigma_{v,h,\alpha}^{\mathrm{cut}}
\qquad\text{in }E
\end{equation}
and
\begin{equation}\label{eq:unit-interval-unscaled-bound}
\|\zeta_{v,h}^{\mathrm{cut}}\|_E
\le
C_{\mathrm{loc},\alpha}h^{\alpha\beta}\|v\|_{\mathfrak G_t},
\qquad
\beta:=H_Y-H_{XY}.
\end{equation}
Then the conclusion of Theorem~\ref{thm:unscaled-smooth-cutoff-criterion} holds. In particular,
\(\mathcal S_t\) has \(h^\alpha\)-stable boundary cutoffs at time \(t\), and the screening
conclusions
\eqref{eq:unscaled-smooth-cutoff-screening-F}--\eqref{eq:unscaled-smooth-cutoff-screening-G} hold.
\end{corollary}

\begin{proof}
Fix \(0<h\le h_{0,\alpha}\) and \(v\in\mathfrak G_t\), and set
\[
\varepsilon:=\varepsilon_{h,\alpha}=h^{1-\alpha}.
\]
Define the micro correction
\begin{equation}\label{eq:unit-interval-unscaled-micro}
\mathfrak C_{v,h}
:=
\varepsilon^\beta
\mathfrak d_{\varepsilon}\zeta_{v,h}^{\mathrm{cut}}
\in L^2(\R_+).
\end{equation}
By Lemma~\ref{lem:micro-dilation},
\[
\operatorname{supp}(\mathfrak C_{v,h})\subset[0,\varepsilon^{-1}]
\]
and
\[
\|\mathfrak C_{v,h}\|_{L^2(\R_+)}
=
\varepsilon^\beta
\|\zeta_{v,h}^{\mathrm{cut}}\|_E
\le
\varepsilon^\beta C_{\mathrm{loc},\alpha}h^{\alpha\beta}\|v\|_{\mathfrak G_t}
=
C_{\mathrm{loc},\alpha}h^\beta\|v\|_{\mathfrak G_t}.
\]
By Proposition~\ref{prop:micro-volterra},
\begin{align*}
\mathfrak V_{Y,t,h,\alpha}^{\mathrm{mic}}\mathfrak C_{v,h}
&=
\varepsilon^{-H_Y}
\mathfrak d_{\varepsilon}
\mathfrak A_{Y,t,h,\alpha}
\mathfrak l_{\varepsilon}
\bigl(
\varepsilon^\beta \mathfrak d_{\varepsilon}\zeta_{v,h}^{\mathrm{cut}}
\bigr) \\
&=
\varepsilon^{-H_Y+\beta}
\mathfrak d_{\varepsilon}
\mathfrak A_{Y,t,h,\alpha}\zeta_{v,h}^{\mathrm{cut}} \\
&=
\varepsilon^{-H_{XY}}
\mathfrak d_{\varepsilon}
\Sigma_{v,h,\alpha}^{\mathrm{cut}},
\end{align*}
where we used \(\mathfrak l_{\varepsilon}\mathfrak d_{\varepsilon}=I_E\) and
\(\beta=H_Y-H_{XY}\). By \eqref{eq:moving-cutoff-unit-reduction},
\[
\varepsilon^{-H_{XY}}
\mathfrak d_{\varepsilon}
\Sigma_{v,h,\alpha}^{\mathrm{cut}}
=
\mathfrak R_{v,h,\alpha}^{\mathrm{cut}}.
\]
Hence
\[
\mathfrak V_{Y,t,h,\alpha}^{\mathrm{mic}}\mathfrak C_{v,h}
=
\mathfrak R_{v,h,\alpha}^{\mathrm{cut}}.
\]
So the hypotheses \eqref{eq:unscaled-cutoff-correction-eq}--\eqref{eq:unscaled-cutoff-correction-bound}
of Theorem~\ref{thm:unscaled-smooth-cutoff-criterion} are satisfied. The conclusion now follows from
that theorem together with Corollary~\ref{cor:unscaled-smooth-cutoff-screening}.
\end{proof}

\begin{corollary}[Exact fixed-window defect criterion on the unit interval]
\label{cor:unit-interval-defect-criterion}
Assume the hypotheses of Proposition~\ref{prop:moving-cutoff-decomposition} and the extra regularity
from Definition~\ref{def:model-c2}. Fix
\[
\chi\in C^\infty([0,1]),
\qquad
0\le \chi\le1,
\qquad
\chi\equiv1\ \text{on }[0,1-\delta_*]
\]
for some \(\delta_*\in(0,1)\), and let \(\widetilde\chi_{h,\alpha}\) be the associated moving micro
cutoff. Assume that there exist
\[
C_{\mathrm{def},\alpha}<\infty,
\qquad
h_{0,\alpha}\in(0,t]
\]
such that for every \(0<h\le h_{0,\alpha}\) and every \(v\in\mathfrak G_t\), with
\(\phi_{v,h}^{\sharp}\) and \(\psi_{v,h}^{\sharp}\) from
\eqref{eq:moving-cutoff-local-traces}, one has
\[
\chi\phi_{v,h}^{\sharp}\in H^\beta((0,1)),
\qquad
\beta:=H_Y-H_{XY},
\]
and
\begin{equation}\label{eq:unit-interval-defect-criterion}
\left\|
\mathfrak F_{t,h,\alpha,1}^{\mathrm{ex}}
\bigl(
\chi\phi_{v,h}^{\sharp},
\chi\psi_{v,h}^{\sharp}
\bigr)
\right\|_{E}
\le
C_{\mathrm{def},\alpha}h^{\alpha\beta}\|v\|_{\mathfrak G_t}.
\end{equation}
Then the hypotheses of Corollary~\ref{cor:unit-interval-unscaled-criterion} are satisfied. In
particular, \(\mathcal S_t\) has \(h^\alpha\)-stable boundary cutoffs at time \(t\), and the
screening conclusions \eqref{eq:unscaled-smooth-cutoff-screening-F}--
\eqref{eq:unscaled-smooth-cutoff-screening-G} hold.
\end{corollary}

\begin{proof}
Fix \(0<h\le h_{0,\alpha}\) and \(v\in\mathfrak G_t\). Since
\[
\chi\phi_{v,h}^{\sharp}\in H^\beta((0,1)),
\qquad
\chi\psi_{v,h}^{\sharp}\in E,
\]
Theorem~\ref{thm:micro-window-correction} with \(R=1\) applies to the pair
\[
\bigl(
\chi\phi_{v,h}^{\sharp},
\chi\psi_{v,h}^{\sharp}
\bigr).
\]
Define
\[
\widetilde\psi_{v,h}^{\sharp}
:=
\mathfrak T_{t,h,\alpha,1}^{\mathrm{ex}}
\bigl(
\chi\phi_{v,h}^{\sharp}
\bigr)
\in E.
\]
Then
\[
\bigl(
\chi\phi_{v,h}^{\sharp},
\widetilde\psi_{v,h}^{\sharp}
\bigr)
\in
\Gamma_{t,h,\alpha,1}^{\mathrm{ex}},
\]
and
\begin{equation}\label{eq:unit-interval-defect-correction}
\left\|
\widetilde\psi_{v,h}^{\sharp}
-\chi\psi_{v,h}^{\sharp}
\right\|_{E}
\le
C_{t,1}
\left\|
\mathfrak F_{t,h,\alpha,1}^{\mathrm{ex}}
\bigl(
\chi\phi_{v,h}^{\sharp},
\chi\psi_{v,h}^{\sharp}
\bigr)
\right\|_{E},
\end{equation}
by \eqref{eq:window-correction-bound}. Set
\[
\xi_{v,h}^{\mathrm{cut}}
:=
\widetilde\psi_{v,h}^{\sharp}
-\chi\psi_{v,h}^{\sharp}
\in E.
\]
Using \eqref{eq:unit-interval-defect-criterion} in
\eqref{eq:unit-interval-defect-correction} yields
\[
\|\xi_{v,h}^{\mathrm{cut}}\|_{E}
\le
C_{t,1}C_{\mathrm{def},\alpha}h^{\alpha\beta}\|v\|_{\mathfrak G_t}.
\]

Because \((\chi\phi_{v,h}^{\sharp},\widetilde\psi_{v,h}^{\sharp})\in\Gamma_{t,h,\alpha,1}^{\mathrm{ex}}\),
one has
\[
\mathfrak A_{Y,t,h,\alpha}\widetilde\psi_{v,h}^{\sharp}
=
\mathfrak A_{XY,t,h,\alpha}(\chi\phi_{v,h}^{\sharp})
\qquad\text{in }E.
\]
Therefore
\begin{align*}
\mathfrak A_{Y,t,h,\alpha}\xi_{v,h}^{\mathrm{cut}}
&=
\mathfrak A_{XY,t,h,\alpha}(\chi\phi_{v,h}^{\sharp})
-\mathfrak A_{Y,t,h,\alpha}(\chi\psi_{v,h}^{\sharp}).
\end{align*}
Now add and subtract \(\chi\mathfrak A_{XY,t,h,\alpha}\phi_{v,h}^{\sharp}\) and
\(\chi\mathfrak A_{Y,t,h,\alpha}\psi_{v,h}^{\sharp}\), and use
\eqref{eq:moving-cutoff-old-tail-local-id}. This gives
\begin{align*}
\mathfrak A_{Y,t,h,\alpha}\xi_{v,h}^{\mathrm{cut}}
&=
\bigl(
\mathfrak A_{XY,t,h,\alpha}(\chi\phi_{v,h}^{\sharp})
-\chi\mathfrak A_{XY,t,h,\alpha}\phi_{v,h}^{\sharp}
\bigr) \\
&\quad-
\bigl(
\mathfrak A_{Y,t,h,\alpha}(\chi\psi_{v,h}^{\sharp})
-\chi\mathfrak A_{Y,t,h,\alpha}\psi_{v,h}^{\sharp}
\bigr) \\
&\quad+
\chi\bigl(
\mathfrak A_{XY,t,h,\alpha}\phi_{v,h}^{\sharp}
-\mathfrak A_{Y,t,h,\alpha}\psi_{v,h}^{\sharp}
\bigr) \\
&=
\mathfrak W_{XY,t,h,\alpha,\chi}^{\mathrm{loc}}\phi_{v,h}^{\sharp}
-\mathfrak W_{Y,t,h,\alpha,\chi}^{\mathrm{loc}}\psi_{v,h}^{\sharp}
+\chi\,\omega_{v,h,\alpha}^{\mathrm{old}} \\
&=
\Sigma_{v,h,\alpha}^{\mathrm{cut}}
\qquad\text{in }E,
\end{align*}
where the last identity is Definition~\ref{def:moving-cutoff-local-source}. Thus
\eqref{eq:unit-interval-unscaled-eq} and \eqref{eq:unit-interval-unscaled-bound} hold with
\[
C_{\mathrm{loc},\alpha}:=C_{t,1}C_{\mathrm{def},\alpha}.
\]
Corollary~\ref{cor:unit-interval-unscaled-criterion} now gives the conclusion.
\end{proof}

\begin{corollary}[Source-only fixed-window criterion on the unit interval]
\label{cor:unit-interval-source-only}
Assume the hypotheses of Proposition~\ref{prop:moving-cutoff-decomposition} and the extra regularity
from Definition~\ref{def:model-c2}. Fix
\[
\chi\in C^\infty([0,1]),
\qquad
0\le \chi\le1,
\qquad
\chi\equiv1\ \text{on }[0,1-\delta_*]
\]
for some \(\delta_*\in(0,1)\), and let \(\widetilde\chi_{h,\alpha}\) be the associated moving micro
cutoff. Let
\[
\beta:=H_Y-H_{XY}.
\]
Assume that there exist
\[
C_{\mathrm{src},\alpha}<\infty,
\qquad
h_{0,\alpha}\in(0,t]
\]
such that for every \(0<h\le h_{0,\alpha}\) and every \(v\in\mathfrak G_t\), with
\(\phi_{v,h}^{\sharp}\) and \(\psi_{v,h}^{\sharp}\) from
\eqref{eq:moving-cutoff-local-traces}, the distribution
\begin{equation}\label{eq:unit-interval-source-only-def}
\mathfrak S_{v,h,\alpha}^{\mathrm{cut}}
:=
-\frac{L_Y(t,t)\mathbf c_t}{\Gamma(1-\beta)}
\mathfrak M_{t,h,\alpha,1}^{\mathrm{win}}
\mathfrak D_{\beta,1}^{\mathrm{dist}}
\bigl(
\chi\phi_{v,h}^{\sharp}
\bigr)
-\mathfrak K_{t,h,\alpha,1}^{\mathrm{win}}
\bigl(
\chi\phi_{v,h}^{\sharp}
\bigr)
\end{equation}
belongs to \(E=L^2([0,1])\) and satisfies
\begin{equation}\label{eq:unit-interval-source-only-bound}
\left\|
\mathfrak S_{v,h,\alpha}^{\mathrm{cut}}
\right\|_{E}
\le
C_{\mathrm{src},\alpha}h^{\alpha\beta}\|v\|_{\mathfrak G_t}.
\end{equation}
Then the hypotheses of Corollary~\ref{cor:unit-interval-defect-criterion} are satisfied. In
particular, \(\mathcal S_t\) has \(h^\alpha\)-stable boundary cutoffs at time \(t\), and the
screening conclusions \eqref{eq:unscaled-smooth-cutoff-screening-F}--
\eqref{eq:unscaled-smooth-cutoff-screening-G} hold.
\end{corollary}

\begin{proof}
After shrinking \(h_{0,\alpha}\) if needed, we may assume that
Theorem~\ref{thm:micro-window-regularity} and Proposition~\ref{prop:BY-micro} both apply with
\(R=1\) for every \(0<h\le h_{0,\alpha}\).

Fix \(0<h\le h_{0,\alpha}\) and \(v\in\mathfrak G_t\), and set
\[
\Phi_{v,h}^{\chi}:=\chi\phi_{v,h}^{\sharp}\in E.
\]
By \eqref{eq:unit-interval-source-only-def},
\[
\frac{L_Y(t,t)\mathbf c_t}{\Gamma(1-\beta)}
\mathfrak M_{t,h,\alpha,1}^{\mathrm{win}}
\mathfrak D_{\beta,1}^{\mathrm{dist}}
\Phi_{v,h}^{\chi}
=
-\mathfrak S_{v,h,\alpha}^{\mathrm{cut}}
-\mathfrak K_{t,h,\alpha,1}^{\mathrm{win}}\Phi_{v,h}^{\chi}
\qquad\text{in }\mathcal D'((0,1)).
\]
The right-hand side belongs to \(E\), because
\(\mathfrak S_{v,h,\alpha}^{\mathrm{cut}}\in E\) by assumption and
\(\mathfrak K_{t,h,\alpha,1}^{\mathrm{win}}\in\mathcal L(E)\) by
Theorem~\ref{thm:micro-window-regularity}. Since
\(\mathfrak M_{t,h,\alpha,1}^{\mathrm{win}}\) is boundedly invertible for small \(h\), it follows
that
\[
\mathfrak D_{\beta,1}^{\mathrm{dist}}\Phi_{v,h}^{\chi}\in E.
\]
Lemma~\ref{lem:window-distribution}(i) therefore gives
\[
\chi\phi_{v,h}^{\sharp}\in H^\beta((0,1)).
\]

Now define
\[
\Psi_{v,h}^{\chi}:=\chi\psi_{v,h}^{\sharp}\in E.
\]
By \eqref{eq:window-defect}, the exact fixed-window defect of the pair
\((\Phi_{v,h}^{\chi},\Psi_{v,h}^{\chi})\) is
\[
\mathfrak F_{t,h,\alpha,1}^{\mathrm{ex}}
\bigl(
\Phi_{v,h}^{\chi},
\Psi_{v,h}^{\chi}
\bigr)
=
\mathfrak B_{Y,t,h,\alpha,1}^{\mathrm{mic}}\Psi_{v,h}^{\chi}
+\mathfrak S_{v,h,\alpha}^{\mathrm{cut}}
\qquad\text{in }E.
\]
Therefore
\begin{equation}\label{eq:source-only-split}
\left\|
\mathfrak F_{t,h,\alpha,1}^{\mathrm{ex}}
\bigl(
\chi\phi_{v,h}^{\sharp},
\chi\psi_{v,h}^{\sharp}
\bigr)
\right\|_{E}
\le
\left\|
\mathfrak B_{Y,t,h,\alpha,1}^{\mathrm{mic}}
\bigl(
\chi\psi_{v,h}^{\sharp}
\bigr)
\right\|_{E}
+
\left\|
\mathfrak S_{v,h,\alpha}^{\mathrm{cut}}
\right\|_{E}.
\end{equation}

By Proposition~\ref{prop:BY-micro}(ii), there exists \(C_t<\infty\) such that
\[
\left\|
\mathfrak B_{Y,t,h,\alpha,1}^{\mathrm{mic}}
\right\|_{\mathcal L(E)}
\le
C_t
\qquad\text{for }0<h\le h_{0,\alpha}.
\]
Also,
\[
\|\chi\psi_{v,h}^{\sharp}\|_{E}
\le
\|\psi_{v,h}^{\sharp}\|_{E}
=
\varepsilon_{h,\alpha}^{-\beta}
\left\|
\Psi_{v,h}
\right\|_{L^2([0,\varepsilon_{h,\alpha}^{-1}])}.
\]
Using \(\Psi_{v,h}=h^\beta\mathcal D_{t,h}\mathcal S_t v\),
Lemma~\ref{lem:dilation-isometry}, and \(\varepsilon_{h,\alpha}=h^{1-\alpha}\), we get
\[
\|\psi_{v,h}^{\sharp}\|_{E}
\le
\varepsilon_{h,\alpha}^{-\beta}h^\beta\|\mathcal S_t v\|_{H_t}
=
h^{\alpha\beta}\|\mathcal S_t v\|_{H_t}
\le
\frac{\nu_Y}{\eta}h^{\alpha\beta}\|v\|_{\mathfrak G_t}.
\]
Hence
\[
\left\|
\mathfrak B_{Y,t,h,\alpha,1}^{\mathrm{mic}}
\bigl(
\chi\psi_{v,h}^{\sharp}
\bigr)
\right\|_{E}
\le
C_t\frac{\nu_Y}{\eta}h^{\alpha\beta}\|v\|_{\mathfrak G_t}.
\]
Combining this with \eqref{eq:source-only-split} and
\eqref{eq:unit-interval-source-only-bound} proves
\eqref{eq:unit-interval-defect-criterion}. Corollary~\ref{cor:unit-interval-defect-criterion} now
gives the conclusion.
\end{proof}

\begin{definition}[Conjugated old-tail source on the unit interval]
\label{def:unit-interval-old-tail-conj}
Under the hypotheses of Proposition~\ref{prop:moving-cutoff-decomposition}, set
\[
\lambda:=H_Y+\frac12.
\]
For \(v\in\mathfrak G_t\) and \(0<h\le t\), define the conjugated old-tail source by
\begin{equation}\label{eq:unit-interval-old-tail-conj}
\mathfrak O_{v,h,\alpha}^{\mathrm{conj}}
:=
\frac{1}{\Gamma(\lambda)}
\mathfrak D_{-,1}^{\lambda}
\bigl(
\chi\,\omega_{v,h,\alpha}^{\mathrm{old}}
\bigr)
\in \mathcal D'((0,1)),
\end{equation}
where \(\omega_{v,h,\alpha}^{\mathrm{old}}\) is given by
\eqref{eq:moving-cutoff-old-tail-local}.
\end{definition}

\begin{proposition}[Exact conjugated decomposition of the unit-interval defect]
\label{prop:unit-interval-defect-conj}
Assume the hypotheses of Proposition~\ref{prop:moving-cutoff-decomposition} and the extra regularity
from Definition~\ref{def:model-c2}. Let \(v\in\mathfrak G_t\), let \(0<h\le t\) satisfy
\(h^\alpha\le t\), and suppose
\[
\chi\phi_{v,h}^{\sharp}\in H^\beta((0,1)),
\qquad
\beta:=H_Y-H_{XY}.
\]
Then, in \(\mathcal D'((0,1))\),
\begin{equation}\label{eq:unit-interval-defect-conj}
\mathfrak F_{t,h,\alpha,1}^{\mathrm{ex}}
\bigl(
\chi\phi_{v,h}^{\sharp},
\chi\psi_{v,h}^{\sharp}
\bigr)
=
\mathfrak Q_{Y,t,h,\alpha,\chi}^{\mathrm{loc}}\psi_{v,h}^{\sharp}
-\mathfrak Q_{XY,t,h,\alpha,\chi}^{\mathrm{loc}}\phi_{v,h}^{\sharp}
-\mathfrak O_{v,h,\alpha}^{\mathrm{conj}}.
\end{equation}
In particular, if
\[
\mathfrak Q_{XY,t,h,\alpha,\chi}^{\mathrm{loc}}\phi_{v,h}^{\sharp}
+
\mathfrak O_{v,h,\alpha}^{\mathrm{conj}}
\in E,
\]
then \(\mathfrak O_{v,h,\alpha}^{\mathrm{conj}}\in E\) and
\eqref{eq:unit-interval-defect-conj} holds in \(E=L^2([0,1])\).
\end{proposition}

\begin{proof}
Because \(\chi\phi_{v,h}^{\sharp}\in H^\beta((0,1))\) and \(\chi\psi_{v,h}^{\sharp}\in E\), the exact
fixed-window defect
\[
\mathfrak F_{t,h,\alpha,1}^{\mathrm{ex}}
\bigl(
\chi\phi_{v,h}^{\sharp},
\chi\psi_{v,h}^{\sharp}
\bigr)
\]
is well-defined by \eqref{eq:window-defect}. By
\eqref{eq:moving-cutoff-unit-reduction},
\[
\mathfrak A_{XY,t,h,\alpha}(\chi\phi_{v,h}^{\sharp})
-\mathfrak A_{Y,t,h,\alpha}(\chi\psi_{v,h}^{\sharp})
=
\mathfrak W_{XY,t,h,\alpha,\chi}^{\mathrm{loc}}\phi_{v,h}^{\sharp}
-\mathfrak W_{Y,t,h,\alpha,\chi}^{\mathrm{loc}}\psi_{v,h}^{\sharp}
+\chi\,\omega_{v,h,\alpha}^{\mathrm{old}}
\]
in \(E\). Applying \((\Gamma(\lambda))^{-1}\mathfrak D_{-,1}^{\lambda}\) in the distributional sense
gives
\begin{align*}
\frac{1}{\Gamma(\lambda)}
\mathfrak D_{-,1}^{\lambda}
\bigl(
\mathfrak A_{Y,t,h,\alpha}(\chi\psi_{v,h}^{\sharp})
-\mathfrak A_{XY,t,h,\alpha}(\chi\phi_{v,h}^{\sharp})
\bigr)
&=
\mathfrak Q_{Y,t,h,\alpha,\chi}^{\mathrm{loc}}\psi_{v,h}^{\sharp}\\
&\quad-
\mathfrak Q_{XY,t,h,\alpha,\chi}^{\mathrm{loc}}\phi_{v,h}^{\sharp}
-\mathfrak O_{v,h,\alpha}^{\mathrm{conj}},
\end{align*}
where we used Proposition~\ref{prop:local-smooth-commutator} and
Definition~\ref{def:unit-interval-old-tail-conj}. On the other hand,
Proposition~\ref{prop:micro-window-conjugated} with \(R=1\) identifies the left-hand side exactly
with
\[
\mathfrak F_{t,h,\alpha,1}^{\mathrm{ex}}
\bigl(
\chi\phi_{v,h}^{\sharp},
\chi\psi_{v,h}^{\sharp}
\bigr)
\]
in \(\mathcal D'((0,1))\). This proves \eqref{eq:unit-interval-defect-conj}.

If the displayed sum belongs to \(E\), then the right-hand side of
\eqref{eq:unit-interval-defect-conj} is an \(E\)-distribution minus
\(\mathfrak O_{v,h,\alpha}^{\mathrm{conj}}\), while the left-hand side belongs to \(E\) by
\eqref{eq:window-defect}. Hence \(\mathfrak O_{v,h,\alpha}^{\mathrm{conj}}\in E\), and the identity
holds in \(E\).
\end{proof}

\begin{corollary}[Reduction to source-old-tail cancellation on the unit interval]
\label{cor:unit-interval-source-tail}
Assume the hypotheses of Proposition~\ref{prop:moving-cutoff-decomposition} and the extra regularity
from Definition~\ref{def:model-c2}. Fix
\[
\chi\in C^\infty([0,1]),
\qquad
0\le \chi\le1,
\qquad
\chi\equiv1\ \text{on }[0,1-\delta_*]
\]
for some \(\delta_*\in(0,1)\), and let \(\widetilde\chi_{h,\alpha}\) be the associated moving micro
cutoff. Assume that there exist
\[
C_{\mathrm{tail},\alpha}<\infty,
\qquad
h_{0,\alpha}\in(0,t]
\]
such that for every \(0<h\le h_{0,\alpha}\) and every \(v\in\mathfrak G_t\), with
\(\phi_{v,h}^{\sharp}\), \(\psi_{v,h}^{\sharp}\), and \(\mathfrak O_{v,h,\alpha}^{\mathrm{conj}}\)
as above, one has
\[
\chi\phi_{v,h}^{\sharp}\in H^\beta((0,1)),
\qquad
\beta:=H_Y-H_{XY},
\]
and
\begin{equation}\label{eq:unit-interval-source-tail}
\left\|
\mathfrak Q_{XY,t,h,\alpha,\chi}^{\mathrm{loc}}\phi_{v,h}^{\sharp}
+
\mathfrak O_{v,h,\alpha}^{\mathrm{conj}}
\right\|_{E}
\le
C_{\mathrm{tail},\alpha}h^{\alpha\beta}\|v\|_{\mathfrak G_t}.
\end{equation}
Then the conclusion of Corollary~\ref{cor:unit-interval-defect-criterion} holds. In particular,
\(\mathcal S_t\) has \(h^\alpha\)-stable boundary cutoffs at time \(t\), and the screening
conclusions \eqref{eq:unscaled-smooth-cutoff-screening-F}--
\eqref{eq:unscaled-smooth-cutoff-screening-G} hold.
\end{corollary}

\begin{proof}
After shrinking \(h_{0,\alpha}\) if needed, Proposition~\ref{prop:local-smooth-commutator} applies
for every \(0<h\le h_{0,\alpha}\). Fix such an \(h\) and \(v\in\mathfrak G_t\). By
Proposition~\ref{prop:unit-interval-defect-conj},
\[
\mathfrak F_{t,h,\alpha,1}^{\mathrm{ex}}
\bigl(
\chi\phi_{v,h}^{\sharp},
\chi\psi_{v,h}^{\sharp}
\bigr)
=
\mathfrak Q_{Y,t,h,\alpha,\chi}^{\mathrm{loc}}\psi_{v,h}^{\sharp}
-\Bigl(
\mathfrak Q_{XY,t,h,\alpha,\chi}^{\mathrm{loc}}\phi_{v,h}^{\sharp}
+
\mathfrak O_{v,h,\alpha}^{\mathrm{conj}}
\Bigr)
\]
in \(E\).

By \eqref{eq:local-com-bound}, there exists \(C_t<\infty\), depending only on \(t\) and the fixed
cutoff \(\chi\), such that
\[
\left\|
\mathfrak Q_{Y,t,h,\alpha,\chi}^{\mathrm{loc}}
\right\|_{\mathcal L(E)}
\le
C_t.
\]
Also,
\[
\|\psi_{v,h}^{\sharp}\|_{E}
=
\varepsilon_{h,\alpha}^{-\beta}
\|\Psi_{v,h}\|_{L^2([0,\varepsilon_{h,\alpha}^{-1}])}
\le
\varepsilon_{h,\alpha}^{-\beta}h^\beta\|\mathcal S_t v\|_{H_t}
=
h^{\alpha\beta}\|\mathcal S_t v\|_{H_t}
\le
\frac{\nu_Y}{\eta}h^{\alpha\beta}\|v\|_{\mathfrak G_t},
\]
using \(\Psi_{v,h}=h^\beta\mathcal D_{t,h}\mathcal S_t v\), Lemma~\ref{lem:dilation-isometry}, and
\(\varepsilon_{h,\alpha}=h^{1-\alpha}\). Therefore
\[
\left\|
\mathfrak Q_{Y,t,h,\alpha,\chi}^{\mathrm{loc}}\psi_{v,h}^{\sharp}
\right\|_{E}
\le
C_t\frac{\nu_Y}{\eta}h^{\alpha\beta}\|v\|_{\mathfrak G_t}.
\]
Combining this with \eqref{eq:unit-interval-source-tail} gives
\[
\left\|
\mathfrak F_{t,h,\alpha,1}^{\mathrm{ex}}
\bigl(
\chi\phi_{v,h}^{\sharp},
\chi\psi_{v,h}^{\sharp}
\bigr)
\right\|_{E}
\le
C'_{t,\alpha}h^{\alpha\beta}\|v\|_{\mathfrak G_t},
\]
for some \(C'_{t,\alpha}<\infty\). Corollary~\ref{cor:unit-interval-defect-criterion} now gives the
conclusion.
\end{proof}

\begin{definition}[Preconjugated source-old-tail discrepancy on the unit interval]
\label{def:unit-interval-source-preconj}
Under the hypotheses of Proposition~\ref{prop:moving-cutoff-decomposition}, define
\begin{equation}\label{eq:unit-interval-source-preconj}
\mathfrak Y_{v,h,\alpha}^{\mathrm{src}}
:=
\mathfrak A_{XY,t,h,\alpha}
\bigl(
\chi\phi_{v,h}^{\sharp}
\bigr)
-
\chi\,\mathfrak A_{Y,t,h,\alpha}\psi_{v,h}^{\sharp}
\in E.
\end{equation}
\end{definition}

\begin{proposition}[Exact source-old-tail recombination on the unit interval]
\label{prop:unit-interval-source-preconj}
Assume the hypotheses of Proposition~\ref{prop:moving-cutoff-decomposition} and the extra regularity
from Definition~\ref{def:model-c2}. Let \(v\in\mathfrak G_t\), let \(0<h\le t\) satisfy
\(h^\alpha\le t\), and suppose
\[
\chi\phi_{v,h}^{\sharp}\in H^\beta((0,1)),
\qquad
\beta:=H_Y-H_{XY}.
\]
Then:
\begin{enumerate}[label=(\roman*)]
\item one has the exact preconjugated identity
\begin{equation}\label{eq:unit-interval-source-preconj-id}
\mathfrak Y_{v,h,\alpha}^{\mathrm{src}}
=
\mathfrak W_{XY,t,h,\alpha,\chi}^{\mathrm{loc}}\phi_{v,h}^{\sharp}
+
\chi\,\omega_{v,h,\alpha}^{\mathrm{old}}
\qquad\text{in }E;
\end{equation}
\item in \(\mathcal D'((0,1))\),
\begin{equation}\label{eq:unit-interval-source-preconj-conj}
\frac{1}{\Gamma(\lambda)}
\mathfrak D_{-,1}^{\lambda}
\mathfrak Y_{v,h,\alpha}^{\mathrm{src}}
=
\mathfrak Q_{XY,t,h,\alpha,\chi}^{\mathrm{loc}}\phi_{v,h}^{\sharp}
+
\mathfrak O_{v,h,\alpha}^{\mathrm{conj}},
\qquad
\lambda:=H_Y+\frac12;
\end{equation}
\item consequently,
\begin{equation}\label{eq:unit-interval-defect-preconj}
\mathfrak F_{t,h,\alpha,1}^{\mathrm{ex}}
\bigl(
\chi\phi_{v,h}^{\sharp},
\chi\psi_{v,h}^{\sharp}
\bigr)
=
\mathfrak Q_{Y,t,h,\alpha,\chi}^{\mathrm{loc}}\psi_{v,h}^{\sharp}
-
\frac{1}{\Gamma(\lambda)}
\mathfrak D_{-,1}^{\lambda}
\mathfrak Y_{v,h,\alpha}^{\mathrm{src}}
\qquad\text{in }\mathcal D'((0,1)).
\end{equation}
\end{enumerate}
If the right-hand side of \eqref{eq:unit-interval-source-preconj-conj} belongs to \(E\), then
\eqref{eq:unit-interval-source-preconj-conj} and \eqref{eq:unit-interval-defect-preconj} both hold in
\(E=L^2([0,1])\).
\end{proposition}

\begin{proof}
By definition of \(\mathfrak W_{XY,t,h,\alpha,\chi}^{\mathrm{loc}}\),
\[
\mathfrak A_{XY,t,h,\alpha}
\bigl(
\chi\phi_{v,h}^{\sharp}
\bigr)
=
\mathfrak W_{XY,t,h,\alpha,\chi}^{\mathrm{loc}}\phi_{v,h}^{\sharp}
+
\chi\,\mathfrak A_{XY,t,h,\alpha}\phi_{v,h}^{\sharp}
\qquad\text{in }E.
\]
Using \eqref{eq:moving-cutoff-old-tail-local-id},
\[
\mathfrak A_{XY,t,h,\alpha}\phi_{v,h}^{\sharp}
-\mathfrak A_{Y,t,h,\alpha}\psi_{v,h}^{\sharp}
=
\omega_{v,h,\alpha}^{\mathrm{old}}
\qquad\text{in }E.
\]
Subtracting \(\chi\,\mathfrak A_{Y,t,h,\alpha}\psi_{v,h}^{\sharp}\) proves
\eqref{eq:unit-interval-source-preconj-id}.

Applying \((\Gamma(\lambda))^{-1}\mathfrak D_{-,1}^{\lambda}\) to
\eqref{eq:unit-interval-source-preconj-id}, Proposition~\ref{prop:local-smooth-commutator}, and
Definition~\ref{def:unit-interval-old-tail-conj} gives
\eqref{eq:unit-interval-source-preconj-conj}.

Finally, substituting \eqref{eq:unit-interval-source-preconj-conj} into
\eqref{eq:unit-interval-defect-conj} from Proposition~\ref{prop:unit-interval-defect-conj} yields
\eqref{eq:unit-interval-defect-preconj}. If
\eqref{eq:unit-interval-source-preconj-conj} belongs to \(E\), then
the left-hand side of \eqref{eq:unit-interval-source-preconj-conj} is an \(E\)-distribution, so
\eqref{eq:unit-interval-source-preconj-conj} holds in \(E\). Since
\(\mathfrak Q_{Y,t,h,\alpha,\chi}^{\mathrm{loc}}\psi_{v,h}^{\sharp}\in E\), the right-hand side of
\eqref{eq:unit-interval-defect-preconj} also belongs to \(E\), and therefore
\eqref{eq:unit-interval-defect-preconj} holds in \(E\) as well.
\end{proof}

\begin{proposition}[Automatic \(H^\beta\)-regularity from the preconjugated source discrepancy]
\label{prop:unit-interval-source-preconj-regularity}
Assume the hypotheses of Proposition~\ref{prop:moving-cutoff-decomposition} and the extra regularity
from Definition~\ref{def:model-c2}. Fix
\[
\chi\in C^\infty([0,1]),
\qquad
0\le \chi\le1,
\qquad
\chi\equiv1\ \text{on }[0,1-\delta_*]
\]
for some \(\delta_*\in(0,1)\), and let \(\widetilde\chi_{h,\alpha}\) be the associated moving micro
cutoff. Then there exist \(h_{0,\alpha}\in(0,t]\) and \(C_{t,\alpha}<\infty\) such that for every
\(0<h\le h_{0,\alpha}\) and every \(v\in\mathfrak G_t\), if
\begin{equation}\label{eq:unit-interval-source-preconj-regularity-hyp}
\frac{1}{\Gamma(\lambda)}
\mathfrak D_{-,1}^{\lambda}
\mathfrak Y_{v,h,\alpha}^{\mathrm{src}}
\in E,
\qquad
\lambda:=H_Y+\frac12,
\end{equation}
then
\[
\chi\phi_{v,h}^{\sharp}\in H^\beta((0,1)),
\qquad
\beta:=H_Y-H_{XY},
\]
and
\begin{equation}\label{eq:unit-interval-source-preconj-regularity-bound}
\left\|
\mathfrak D_{\beta,1}^{\mathrm{dist}}
\bigl(
\chi\phi_{v,h}^{\sharp}
\bigr)
\right\|_{E}
\le
C_{t,\alpha}
\left(
\left\|
\frac{1}{\Gamma(\lambda)}
\mathfrak D_{-,1}^{\lambda}
\mathfrak Y_{v,h,\alpha}^{\mathrm{src}}
\right\|_{E}
+
\|v\|_{\mathfrak G_t}
\right).
\end{equation}
\end{proposition}

\begin{proof}
After shrinking \(h_{0,\alpha}\) if needed, Theorem~\ref{thm:micro-window-regularity},
Proposition~\ref{prop:BY-micro}, and Proposition~\ref{prop:local-smooth-commutator} all apply with
\(R=1\) for every \(0<h\le h_{0,\alpha}\).

Fix such an \(h\) and \(v\in\mathfrak G_t\). Set
\[
\Phi_{v,h}^{\chi}:=\chi\phi_{v,h}^{\sharp}\in E,
\qquad
\Psi_{v,h}^{\chi}:=\chi\psi_{v,h}^{\sharp}\in E.
\]
By Definition~\ref{def:unit-interval-source-preconj} and the identity
\[
\mathfrak W_{Y,t,h,\alpha,\chi}^{\mathrm{loc}}\psi_{v,h}^{\sharp}
=
\mathfrak A_{Y,t,h,\alpha}\Psi_{v,h}^{\chi}
-\chi\,\mathfrak A_{Y,t,h,\alpha}\psi_{v,h}^{\sharp},
\]
one has
\[
\mathfrak Y_{v,h,\alpha}^{\mathrm{src}}
=
\mathfrak A_{XY,t,h,\alpha}\Phi_{v,h}^{\chi}
-\mathfrak A_{Y,t,h,\alpha}\Psi_{v,h}^{\chi}
+\mathfrak W_{Y,t,h,\alpha,\chi}^{\mathrm{loc}}\psi_{v,h}^{\sharp}
\qquad\text{in }E.
\]
Applying \((\Gamma(\lambda))^{-1}\mathfrak D_{-,1}^{\lambda}\), using
Proposition~\ref{prop:BY-micro}(i) for \(\Psi_{v,h}^{\chi}\) and
Proposition~\ref{prop:local-smooth-commutator} for the commutator term, yields
\begin{equation}\label{eq:unit-interval-source-preconj-regularity-id}
\frac{1}{\Gamma(\lambda)}
\mathfrak D_{-,1}^{\lambda}
\bigl(
\mathfrak A_{XY,t,h,\alpha}\Phi_{v,h}^{\chi}
\bigr)
=
\frac{1}{\Gamma(\lambda)}
\mathfrak D_{-,1}^{\lambda}
\mathfrak Y_{v,h,\alpha}^{\mathrm{src}}
+
\mathfrak B_{Y,t,h,\alpha,1}^{\mathrm{mic}}\Psi_{v,h}^{\chi}
-\mathfrak Q_{Y,t,h,\alpha,\chi}^{\mathrm{loc}}\psi_{v,h}^{\sharp}
\end{equation}
in \(\mathcal D'((0,1))\).

The first term on the right-hand side belongs to \(E\) by
\eqref{eq:unit-interval-source-preconj-regularity-hyp}. The other two terms also belong to \(E\),
because
\[
\mathfrak B_{Y,t,h,\alpha,1}^{\mathrm{mic}},
\mathfrak Q_{Y,t,h,\alpha,\chi}^{\mathrm{loc}}
\in\mathcal L(E)
\]
for small \(h\). Hence the right-hand side of
\eqref{eq:unit-interval-source-preconj-regularity-id} belongs to \(E\).

Now apply the distributional source decomposition
\eqref{eq:window-source-decompose-dist} from
Theorem~\ref{thm:micro-window-regularity} with \(R=1\) to \(\Phi_{v,h}^{\chi}\). This gives
\[
\frac{1}{\Gamma(\lambda)}
\mathfrak D_{-,1}^{\lambda}
\bigl(
\mathfrak A_{XY,t,h,\alpha}\Phi_{v,h}^{\chi}
\bigr)
=
\frac{L_Y(t,t)\mathbf c_t}{\Gamma(1-\beta)}
\mathfrak M_{t,h,\alpha,1}^{\mathrm{win}}
\mathfrak D_{\beta,1}^{\mathrm{dist}}\Phi_{v,h}^{\chi}
+
\mathfrak K_{t,h,\alpha,1}^{\mathrm{win}}\Phi_{v,h}^{\chi}
\]
in \(\mathcal D'((0,1))\). Since
\(\mathfrak K_{t,h,\alpha,1}^{\mathrm{win}}\in\mathcal L(E)\) and
\(\mathfrak M_{t,h,\alpha,1}^{\mathrm{win}}\) is boundedly invertible for small \(h\),
\eqref{eq:unit-interval-source-preconj-regularity-id} implies that
\[
\mathfrak D_{\beta,1}^{\mathrm{dist}}\Phi_{v,h}^{\chi}\in E.
\]
Lemma~\ref{lem:window-distribution}(i) therefore yields
\[
\chi\phi_{v,h}^{\sharp}\in H^\beta((0,1)).
\]

Moreover, the same identity gives
\begin{align*}
\left\|
\mathfrak D_{\beta,1}^{\mathrm{dist}}\Phi_{v,h}^{\chi}
\right\|_{E}
&\le
C_{t,\alpha}
\Biggl(
\left\|
\frac{1}{\Gamma(\lambda)}
\mathfrak D_{-,1}^{\lambda}
\mathfrak Y_{v,h,\alpha}^{\mathrm{src}}
\right\|_{E}
+
\|\Psi_{v,h}^{\chi}\|_{E}
+
\|\psi_{v,h}^{\sharp}\|_{E}
+
\|\Phi_{v,h}^{\chi}\|_{E}
\Biggr).
\end{align*}
Now \(\|\Phi_{v,h}^{\chi}\|_{E}\le \|\phi_{v,h}^{\sharp}\|_{E}=\|v\|_{H_t}\le\|v\|_{\mathfrak G_t}\),
while
\[
\|\psi_{v,h}^{\sharp}\|_{E}
\le
\varepsilon_{h,\alpha}^{-\beta}h^\beta\|\mathcal S_t v\|_{H_t}
=
h^{\alpha\beta}\|\mathcal S_t v\|_{H_t}
\le
\frac{\nu_Y}{\eta}h^{\alpha\beta}\|v\|_{\mathfrak G_t}
\le
\frac{\nu_Y}{\eta}\|v\|_{\mathfrak G_t},
\]
and the same bound holds for \(\|\Psi_{v,h}^{\chi}\|_{E}\). This proves
\eqref{eq:unit-interval-source-preconj-regularity-bound}.
\end{proof}

\begin{corollary}[Reduction to a single preconjugated source discrepancy]
\label{cor:unit-interval-source-preconj}
Assume the hypotheses of Proposition~\ref{prop:moving-cutoff-decomposition} and the extra regularity
from Definition~\ref{def:model-c2}. Fix
\[
\chi\in C^\infty([0,1]),
\qquad
0\le \chi\le1,
\qquad
\chi\equiv1\ \text{on }[0,1-\delta_*]
\]
for some \(\delta_*\in(0,1)\), and let \(\widetilde\chi_{h,\alpha}\) be the associated moving micro
cutoff. Assume that there exist
\[
C_{\mathrm{pre},\alpha}<\infty,
\qquad
h_{0,\alpha}\in(0,t]
\]
such that for every \(0<h\le h_{0,\alpha}\) and every \(v\in\mathfrak G_t\),
\begin{equation}\label{eq:unit-interval-source-preconj-bound}
\left\|
\frac{1}{\Gamma(\lambda)}
\mathfrak D_{-,1}^{\lambda}
\mathfrak Y_{v,h,\alpha}^{\mathrm{src}}
\right\|_{E}
\le
C_{\mathrm{pre},\alpha}h^{\alpha\beta}\|v\|_{\mathfrak G_t},
\qquad
\lambda:=H_Y+\frac12.
\end{equation}
Then the conclusion of Corollary~\ref{cor:unit-interval-defect-criterion} holds. In particular,
\(\mathcal S_t\) has \(h^\alpha\)-stable boundary cutoffs at time \(t\), and the screening
conclusions \eqref{eq:unscaled-smooth-cutoff-screening-F}--
\eqref{eq:unscaled-smooth-cutoff-screening-G} hold.
\end{corollary}

\begin{proof}
  After shrinking \(h_{0,\alpha}\) if needed, Proposition~\ref{prop:local-smooth-commutator} applies
  for every \(0<h\le h_{0,\alpha}\), and
  Proposition~\ref{prop:unit-interval-source-preconj-regularity} applies as well. Fix such an \(h\)
  and \(v\in\mathfrak G_t\). By
  Proposition~\ref{prop:unit-interval-source-preconj-regularity} and
  \eqref{eq:unit-interval-source-preconj-bound},
  \[
  \chi\phi_{v,h}^{\sharp}\in H^\beta((0,1)),
  \qquad
  \beta:=H_Y-H_{XY}.
  \]
  By
  Proposition~\ref{prop:unit-interval-source-preconj},
  \[
  \mathfrak F_{t,h,\alpha,1}^{\mathrm{ex}}
  \bigl(
\chi\phi_{v,h}^{\sharp},
\chi\psi_{v,h}^{\sharp}
\bigr)
=
\mathfrak Q_{Y,t,h,\alpha,\chi}^{\mathrm{loc}}\psi_{v,h}^{\sharp}
-
\frac{1}{\Gamma(\lambda)}
\mathfrak D_{-,1}^{\lambda}
\mathfrak Y_{v,h,\alpha}^{\mathrm{src}}
\]
in \(E\).

By \eqref{eq:local-com-bound}, there exists \(C_t<\infty\), depending only on \(t\) and the fixed
cutoff \(\chi\), such that
\[
\left\|
\mathfrak Q_{Y,t,h,\alpha,\chi}^{\mathrm{loc}}
\right\|_{\mathcal L(E)}
\le
C_t.
\]
Also, as in the proof of Corollary~\ref{cor:unit-interval-source-tail},
\[
\|\psi_{v,h}^{\sharp}\|_{E}
\le
\varepsilon_{h,\alpha}^{-\beta}h^\beta\|\mathcal S_t v\|_{H_t}
=
h^{\alpha\beta}\|\mathcal S_t v\|_{H_t}
\le
\frac{\nu_Y}{\eta}h^{\alpha\beta}\|v\|_{\mathfrak G_t}.
\]
Therefore
\[
\left\|
\mathfrak Q_{Y,t,h,\alpha,\chi}^{\mathrm{loc}}\psi_{v,h}^{\sharp}
\right\|_{E}
\le
C_t\frac{\nu_Y}{\eta}h^{\alpha\beta}\|v\|_{\mathfrak G_t}.
\]
Combining this with \eqref{eq:unit-interval-source-preconj-bound} gives
\[
\left\|
\mathfrak F_{t,h,\alpha,1}^{\mathrm{ex}}
\bigl(
\chi\phi_{v,h}^{\sharp},
\chi\psi_{v,h}^{\sharp}
\bigr)
\right\|_{E}
\le
C'_{t,\alpha}h^{\alpha\beta}\|v\|_{\mathfrak G_t},
\]
for some \(C'_{t,\alpha}<\infty\). Corollary~\ref{cor:unit-interval-defect-criterion} now gives the
conclusion.
\end{proof}

\begin{definition}[Separated old-tail pieces on the unit interval]
\label{def:unit-interval-old-tail-split}
Under the hypotheses of Proposition~\ref{prop:moving-cutoff-decomposition}, set
\[
\varepsilon:=\varepsilon_{h,\alpha}=h^{1-\alpha},
\qquad
\lambda:=H_Y+\frac12,
\qquad
\mu:=H_{XY}+\frac12.
\]
For \(v\in\mathfrak G_t\) and \(0<h\le t\), define
\begin{align}
\vartheta_{Y,v,h,\alpha}^{\mathrm{old}}
&:=
\varepsilon^{H_{XY}}
\mathfrak l_{\varepsilon}
\Biggl[
\one_{[0,\varepsilon^{-1}]}
\int_{\varepsilon^{-1}}^{t/h}
(q-r)^{\lambda-1}
L_Y(t-hr,t-hq)\Psi_{v,h}(q)\,dq
\Biggr],
\label{eq:unit-interval-target-old-tail}\\
\vartheta_{XY,v,h,\alpha}^{\mathrm{old}}
&:=
\varepsilon^{H_{XY}}
\mathfrak l_{\varepsilon}
\Biggl[
\one_{[0,\varepsilon^{-1}]}
\int_{\varepsilon^{-1}}^{t/h}
(q-r)^{\mu-1}
L_{XY}(t-hr,t-hq)\Phi_{v,h}(q)\,dq
\Biggr].
\label{eq:unit-interval-source-old-tail}
\end{align}
\end{definition}

\begin{proposition}[Full-\(h\) source realization of the preconjugated discrepancy]
\label{prop:unit-interval-source-full}
Assume the hypotheses of Proposition~\ref{prop:moving-cutoff-decomposition}. Let \(v\in\mathfrak G_t\)
and let \(0<h\le t\) satisfy \(h^\alpha\le t\). With \(\widetilde\chi_{h,\alpha}\) from
Definition~\ref{def:moving-cutoff-defect}, define the full-\(h\) source discrepancy by
\begin{equation}\label{eq:full-h-source-discrepancy}
\mathfrak Z_{v,h,\alpha}^{\mathrm{src}}
:=
\mathfrak V_{XY,t,h}^{\mathrm{full}}
\bigl(
\widetilde\chi_{h,\alpha}\Phi_{v,h}
\bigr)
-
\widetilde\chi_{h,\alpha}\,
\mathfrak V_{XY,t,h}^{\mathrm{full}}\Phi_{v,h}
\in L^2(\R_+).
\end{equation}
Then:
\begin{enumerate}[label=(\roman*)]
\item \(\mathfrak Z_{v,h,\alpha}^{\mathrm{src}}\) is supported in \([0,\varepsilon_{h,\alpha}^{-1}]\);
\item one has
\begin{equation}\label{eq:unit-interval-old-tail-split}
\omega_{v,h,\alpha}^{\mathrm{old}}
=
\vartheta_{Y,v,h,\alpha}^{\mathrm{old}}
-
\vartheta_{XY,v,h,\alpha}^{\mathrm{old}}
\qquad\text{in }E;
\end{equation}
\item one has the exact scaled source identity
\begin{equation}\label{eq:full-h-source-discrepancy-local}
\varepsilon_{h,\alpha}^{H_{XY}}
\mathfrak l_{\varepsilon_{h,\alpha}}
\mathfrak Z_{v,h,\alpha}^{\mathrm{src}}
=
\mathfrak W_{XY,t,h,\alpha,\chi}^{\mathrm{loc}}\phi_{v,h}^{\sharp}
-
\chi\,\vartheta_{XY,v,h,\alpha}^{\mathrm{old}}
\qquad\text{in }E;
\end{equation}
\item consequently,
\begin{equation}\label{eq:unit-interval-source-preconj-split}
\mathfrak Y_{v,h,\alpha}^{\mathrm{src}}
=
\varepsilon_{h,\alpha}^{H_{XY}}
\mathfrak l_{\varepsilon_{h,\alpha}}
\mathfrak Z_{v,h,\alpha}^{\mathrm{src}}
+
\chi\,\vartheta_{Y,v,h,\alpha}^{\mathrm{old}}
\qquad\text{in }E.
\end{equation}
\end{enumerate}
\end{proposition}

\begin{proof}
Set \(\varepsilon:=\varepsilon_{h,\alpha}\). For \(r>\varepsilon^{-1}\), one has
\(\widetilde\chi_{h,\alpha}(r)=0\), and for \(q\ge r\) also \(\widetilde\chi_{h,\alpha}(q)=0\). Hence
both terms in \eqref{eq:full-h-source-discrepancy} vanish, proving (i).

Item (ii) follows from \eqref{eq:moving-cutoff-old-tail},
\eqref{eq:moving-cutoff-old-tail-local}, and Definitions
\eqref{eq:unit-interval-target-old-tail}--\eqref{eq:unit-interval-source-old-tail}.

For \(0\le r\le \varepsilon^{-1}\), expanding \eqref{eq:full-h-source-discrepancy} gives
\begin{align*}
\mathfrak Z_{v,h,\alpha}^{\mathrm{src}}(r)
&=
\int_r^{\varepsilon^{-1}}
(q-r)^{\mu-1}
L_{XY}(t-hr,t-hq)
\bigl(
\widetilde\chi_{h,\alpha}(q)-\widetilde\chi_{h,\alpha}(r)
\bigr)
\Phi_{v,h}(q)\,dq\\
&\quad-
\widetilde\chi_{h,\alpha}(r)
\int_{\varepsilon^{-1}}^{t/h}
(q-r)^{\mu-1}
L_{XY}(t-hr,t-hq)\Phi_{v,h}(q)\,dq.
\end{align*}
Applying \(\varepsilon^{H_{XY}}\mathfrak l_{\varepsilon}\), using
Proposition~\ref{prop:micro-cutoff-scaling} and the definition
\eqref{eq:unit-interval-source-old-tail}, yields
\eqref{eq:full-h-source-discrepancy-local}.

Finally, combining \eqref{eq:full-h-source-discrepancy-local} with
\eqref{eq:unit-interval-old-tail-split} gives
\begin{align*}
\varepsilon^{H_{XY}}\mathfrak l_{\varepsilon}\mathfrak Z_{v,h,\alpha}^{\mathrm{src}}
+
\chi\,\vartheta_{Y,v,h,\alpha}^{\mathrm{old}}
&=
\mathfrak W_{XY,t,h,\alpha,\chi}^{\mathrm{loc}}\phi_{v,h}^{\sharp}
-
\chi\,\vartheta_{XY,v,h,\alpha}^{\mathrm{old}}
+
\chi\,\vartheta_{Y,v,h,\alpha}^{\mathrm{old}}\\
&=
\mathfrak W_{XY,t,h,\alpha,\chi}^{\mathrm{loc}}\phi_{v,h}^{\sharp}
+
\chi\,\omega_{v,h,\alpha}^{\mathrm{old}},
\end{align*}
which is exactly \(\mathfrak Y_{v,h,\alpha}^{\mathrm{src}}\) by
\eqref{eq:unit-interval-source-preconj-id}. This proves \eqref{eq:unit-interval-source-preconj-split}.
\end{proof}

\begin{corollary}[Reduction to full-\(h\) source discrepancy plus target old tail]
\label{cor:unit-interval-source-full}
Assume the hypotheses of Proposition~\ref{prop:moving-cutoff-decomposition} and the extra regularity
from Definition~\ref{def:model-c2}. Fix
\[
\chi\in C^\infty([0,1]),
\qquad
0\le \chi\le1,
\qquad
\chi\equiv1\ \text{on }[0,1-\delta_*]
\]
for some \(\delta_*\in(0,1)\), and let \(\widetilde\chi_{h,\alpha}\) be the associated moving micro
cutoff. Assume that there exist
\[
C_{\mathrm{src},\alpha}<\infty,
\qquad
C_{\mathrm{tgt},\alpha}<\infty,
\qquad
h_{0,\alpha}\in(0,t]
\]
such that for every \(0<h\le h_{0,\alpha}\) and every \(v\in\mathfrak G_t\),
\begin{align}
\left\|
\frac{1}{\Gamma(\lambda)}
\mathfrak D_{-,1}^{\lambda}
\Bigl(
\varepsilon_{h,\alpha}^{H_{XY}}
\mathfrak l_{\varepsilon_{h,\alpha}}
\mathfrak Z_{v,h,\alpha}^{\mathrm{src}}
\Bigr)
\right\|_{E}
&\le
C_{\mathrm{src},\alpha}h^{\alpha\beta}\|v\|_{\mathfrak G_t},
\label{eq:full-h-source-discrepancy-bound}\\
\left\|
\frac{1}{\Gamma(\lambda)}
\mathfrak D_{-,1}^{\lambda}
\bigl(
\chi\,\vartheta_{Y,v,h,\alpha}^{\mathrm{old}}
\bigr)
\right\|_{E}
&\le
C_{\mathrm{tgt},\alpha}h^{\alpha\beta}\|v\|_{\mathfrak G_t},
\qquad
\lambda:=H_Y+\frac12.
\label{eq:target-old-tail-bound}
\end{align}
Then the conclusion of Corollary~\ref{cor:unit-interval-source-preconj} holds. In particular,
\(\mathcal S_t\) has \(h^\alpha\)-stable boundary cutoffs at time \(t\), and the screening
conclusions \eqref{eq:unscaled-smooth-cutoff-screening-F}--
\eqref{eq:unscaled-smooth-cutoff-screening-G} hold.
\end{corollary}

\begin{proof}
By Proposition~\ref{prop:unit-interval-source-full},
\[
\mathfrak Y_{v,h,\alpha}^{\mathrm{src}}
=
\varepsilon_{h,\alpha}^{H_{XY}}
\mathfrak l_{\varepsilon_{h,\alpha}}
\mathfrak Z_{v,h,\alpha}^{\mathrm{src}}
+
\chi\,\vartheta_{Y,v,h,\alpha}^{\mathrm{old}}
\qquad\text{in }E.
\]
Applying \((\Gamma(\lambda))^{-1}\mathfrak D_{-,1}^{\lambda}\) and using
\eqref{eq:full-h-source-discrepancy-bound} together with \eqref{eq:target-old-tail-bound} yields
\[
\left\|
\frac{1}{\Gamma(\lambda)}
\mathfrak D_{-,1}^{\lambda}
\mathfrak Y_{v,h,\alpha}^{\mathrm{src}}
\right\|_{E}
\le
\bigl(
C_{\mathrm{src},\alpha}+C_{\mathrm{tgt},\alpha}
\bigr)
h^{\alpha\beta}\|v\|_{\mathfrak G_t}.
\]
Corollary~\ref{cor:unit-interval-source-preconj} now gives the conclusion.
\end{proof}

\begin{definition}[Flat outer-edge profile for the endpoint cutoff]
\label{def:endpoint-flat-cutoff}
Let
\[
\chi\in C^\infty([0,1]),
\qquad
0\le \chi\le1,
\qquad
\chi\equiv1\ \text{on }[0,1-\delta_*]
\]
for some \(\delta_*\in(0,1)\). We say that \(\chi\) is \emph{flat at the outer edge} if
\[
\chi^{(m)}(1)=0
\qquad\text{for every }m\in\mathbb N\cup\{0\}.
\]
\end{definition}

\begin{definition}[Full-\(h\) target old-tail commutator]
\label{def:full-h-target-old-tail}
Assume the hypotheses of Proposition~\ref{prop:moving-cutoff-decomposition}, and let \(\chi\) be as
in Definition~\ref{def:endpoint-flat-cutoff}. For \(v\in\mathfrak G_t\) and \(0<h\le t\), define
\begin{equation}\label{eq:full-h-target-old-tail}
\bigl(
\mathfrak C_{Y,v,h,\alpha}^{\mathrm{old}}
\bigr)(r)
:=
\int_{\varepsilon_{h,\alpha}^{-1}}^{t/h}
(q-r)^{\lambda-1}
\bigl(
\widetilde\chi_{h,\alpha}(q)-\widetilde\chi_{h,\alpha}(r)
\bigr)
L_Y(t-hr,t-hq)\Psi_{v,h}(q)\,dq,
\end{equation}
where
\[
\lambda:=H_Y+\frac12.
\]
\end{definition}

\begin{proposition}[Exact full-\(h\) commutator representation of the target old tail]
\label{prop:full-h-target-old-tail}
Assume the hypotheses of Definition~\ref{def:full-h-target-old-tail}. Then:
\begin{enumerate}[label=(\roman*)]
\item \(\mathfrak C_{Y,v,h,\alpha}^{\mathrm{old}}\) is supported in
\([0,\varepsilon_{h,\alpha}^{-1}]\);
\item one has the exact identity
\begin{equation}\label{eq:full-h-target-old-tail-id}
\chi\,\vartheta_{Y,v,h,\alpha}^{\mathrm{old}}
=
-\varepsilon_{h,\alpha}^{H_{XY}}
\mathfrak l_{\varepsilon_{h,\alpha}}
\mathfrak C_{Y,v,h,\alpha}^{\mathrm{old}}
\qquad\text{in }E.
\end{equation}
\end{enumerate}
\end{proposition}

\begin{proof}
Let \(\varepsilon:=\varepsilon_{h,\alpha}\). Because \(\widetilde\chi_{h,\alpha}\) is supported in
\([0,\varepsilon^{-1}]\), one has \(\widetilde\chi_{h,\alpha}(r)=0\) for \(r>\varepsilon^{-1}\), and
the integrand in \eqref{eq:full-h-target-old-tail} also vanishes for \(r>\varepsilon^{-1}\). This
proves (i).

Now fix \(0\le r\le \varepsilon^{-1}\). Since \(\widetilde\chi_{h,\alpha}(q)=0\) for
\(q\ge\varepsilon^{-1}\), the definition \eqref{eq:full-h-target-old-tail} reduces to
\[
\mathfrak C_{Y,v,h,\alpha}^{\mathrm{old}}(r)
=
-\widetilde\chi_{h,\alpha}(r)
\int_{\varepsilon^{-1}}^{t/h}
(q-r)^{\lambda-1}
L_Y(t-hr,t-hq)\Psi_{v,h}(q)\,dq.
\]
Applying \(\varepsilon^{H_{XY}}\mathfrak l_\varepsilon\) and using
\[
(\mathfrak l_\varepsilon\widetilde\chi_{h,\alpha})(u)=\chi(u)
\qquad\text{for }u\in[0,1]
\]
gives exactly \eqref{eq:full-h-target-old-tail-id}, by
\eqref{eq:unit-interval-target-old-tail}.
\end{proof}

\begin{lemma}[Crossing-region quotient bounds for the flat micro cutoff]
\label{lem:full-h-target-old-tail-quotient}
Assume the hypotheses of Definition~\ref{def:full-h-target-old-tail}, and suppose in addition that
\(\chi\) is flat at the outer edge in the sense of Definition~\ref{def:endpoint-flat-cutoff}. Set
\[
\Sigma_{h,\alpha}
:=
\left\{
0\le r\le \varepsilon_{h,\alpha}^{-1}\le q\le \frac{t}{h}
\right\}.
\]
Define on \(\Sigma_{h,\alpha}\)
\[
\delta_{Y,h,\alpha}^{\mathrm{old}}(r,q)
:=
\frac{
\widetilde\chi_{h,\alpha}(q)-\widetilde\chi_{h,\alpha}(r)
}{q-r}.
\]
Then there exists \(C_\chi<\infty\), depending only on \(\chi\), such that
\begin{equation}\label{eq:full-h-target-old-tail-quotient}
\left\|
\delta_{Y,h,\alpha}^{\mathrm{old}}
\right\|_{C^1(\Sigma_{h,\alpha})}
\le
C_\chi \varepsilon_{h,\alpha}
\qquad\text{for all }0<h\le t.
\end{equation}
\end{lemma}

\begin{proof}
Let \(\varepsilon:=\varepsilon_{h,\alpha}\). Since \(q\ge \varepsilon^{-1}\), one has
\(\widetilde\chi_{h,\alpha}(q)=0\), so
\[
\delta_{Y,h,\alpha}^{\mathrm{old}}(r,q)
=
-\frac{\widetilde\chi_{h,\alpha}(r)}{q-r}
=
-\frac{\chi(\varepsilon r)}{q-r}
\qquad\text{on }\Sigma_{h,\alpha}.
\]
Because \(\chi\) is flat at \(1\), Taylor's theorem at the outer edge gives
\[
|\chi(x)|\le C_\chi(1-x)^2,
\qquad
|\chi'(x)|\le C_\chi(1-x)
\qquad\text{for }x\in[0,1].
\]
Also, for \((r,q)\in\Sigma_{h,\alpha}\),
\[
q-r\ge \varepsilon^{-1}-r=\frac{1-\varepsilon r}{\varepsilon}.
\]
Therefore
\[
\left|
\delta_{Y,h,\alpha}^{\mathrm{old}}(r,q)
\right|
\le
C_\chi
\frac{(1-\varepsilon r)^2}{(1-\varepsilon r)/\varepsilon}
\le
C_\chi \varepsilon.
\]
Next,
\[
\partial_q\delta_{Y,h,\alpha}^{\mathrm{old}}(r,q)
=
\frac{\chi(\varepsilon r)}{(q-r)^2},
\]
so the same bounds imply
\[
\left|
\partial_q\delta_{Y,h,\alpha}^{\mathrm{old}}(r,q)
\right|
\le
C_\chi
\frac{(1-\varepsilon r)^2}{((1-\varepsilon r)/\varepsilon)^2}
\le
C_\chi \varepsilon^2
\le
C_\chi \varepsilon.
\]
Finally,
\[
\partial_r\delta_{Y,h,\alpha}^{\mathrm{old}}(r,q)
=
\frac{-\varepsilon\chi'(\varepsilon r)(q-r)-\chi(\varepsilon r)}{(q-r)^2},
\]
hence
\begin{align*}
\left|
\partial_r\delta_{Y,h,\alpha}^{\mathrm{old}}(r,q)
\right|
&\le
\frac{
\varepsilon |\chi'(\varepsilon r)|(q-r)+|\chi(\varepsilon r)|
}{(q-r)^2}\\
&\le
C_\chi
\frac{
\varepsilon(1-\varepsilon r)\cdot (1-\varepsilon r)/\varepsilon
+(1-\varepsilon r)^2
}{
((1-\varepsilon r)/\varepsilon)^2
}\\
&\le
C_\chi \varepsilon^2
\le
C_\chi \varepsilon.
\end{align*}
This proves \eqref{eq:full-h-target-old-tail-quotient}.
\end{proof}

\begin{proposition}[\(\alpha\)-scale realization of the target old tail]
\label{prop:full-h-target-old-tail-alpha}
Assume the hypotheses of Definition~\ref{def:full-h-target-old-tail}, and suppose moreover that
\(\chi\) is flat at the outer edge in the sense of Definition~\ref{def:endpoint-flat-cutoff}. Set
\[
\beta:=H_Y-H_{XY},
\qquad
\varepsilon:=\varepsilon_{h,\alpha}=h^{1-\alpha},
\qquad
R_{h,\alpha}:=th^{-\alpha}.
\]
Define the extended \(\alpha\)-scale tail trace by
\begin{equation}\label{eq:full-h-target-old-tail-alpha-trace}
\widehat\Psi_{v,h,\alpha}^{\mathrm{old}}(q)
:=
\varepsilon^{-(\beta+\frac12)}
\Psi_{v,h}\!\left(\frac{q}{\varepsilon}\right)\one_{[1,R_{h,\alpha}]}(q),
\qquad q\ge0.
\end{equation}
Then:
\begin{enumerate}[label=(\roman*)]
\item one has
\begin{equation}\label{eq:full-h-target-old-tail-alpha-trace-bound}
\left\|
\widehat\Psi_{v,h,\alpha}^{\mathrm{old}}
\right\|_{L^2(\R_+)}
\le
\varepsilon^{-\beta}
\left\|
\Psi_{v,h}
\right\|_{L^2([\varepsilon^{-1},t/h])}
\le
\frac{\nu_Y}{\eta}h^{\alpha\beta}\|v\|_{\mathfrak G_t};
\end{equation}
\item if
\begin{equation}\label{eq:full-h-target-old-tail-alpha-op}
\bigl(
\mathfrak W_{Y,t,h,\alpha,\chi}^{\mathrm{old},\alpha}\Theta
\bigr)(u)
:=
\chi(u)
\int_1^{R_{h,\alpha}}
(q-u)^{\lambda-1}
L_Y(t-h^\alpha u,t-h^\alpha q)\Theta(q)\,dq,
\qquad 0\le u\le 1,
\end{equation}
then
\begin{equation}\label{eq:full-h-target-old-tail-alpha-id}
\chi\,\vartheta_{Y,v,h,\alpha}^{\mathrm{old}}
=
\varepsilon^{-1/2}
\mathfrak W_{Y,t,h,\alpha,\chi}^{\mathrm{old},\alpha}
\widehat\Psi_{v,h,\alpha}^{\mathrm{old}}
\qquad\text{in }E;
\end{equation}
\item for every sufficiently small \(h>0\), there exists
\[
\mathfrak Q_{Y,t,h,\alpha,\chi}^{\mathrm{old},\alpha}
\in
\mathcal L\bigl(L^2([1,R_{h,\alpha}]),E\bigr)
\]
such that for every \(\Theta\in L^2([1,R_{h,\alpha}])\),
\begin{equation}\label{eq:full-h-target-old-tail-alpha-conj}
\frac{1}{\Gamma(\lambda)}
\mathfrak D_{-,1}^{\lambda}
\bigl(
\mathfrak W_{Y,t,h,\alpha,\chi}^{\mathrm{old},\alpha}\Theta
\bigr)
=
\mathfrak Q_{Y,t,h,\alpha,\chi}^{\mathrm{old},\alpha}\Theta
\qquad\text{in }E,
\end{equation}
and there exists \(C_{t,\chi}<\infty\), independent of \(h\), such that
\begin{equation}\label{eq:full-h-target-old-tail-alpha-opbound}
\left\|
\mathfrak Q_{Y,t,h,\alpha,\chi}^{\mathrm{old},\alpha}
\right\|_{\mathcal L(L^2([1,R_{h,\alpha}]),E)}
\le
C_{t,\chi}.
\end{equation}
\end{enumerate}
Consequently,
\begin{equation}\label{eq:full-h-target-old-tail-alpha-bound}
\left\|
\frac{1}{\Gamma(\lambda)}
\mathfrak D_{-,1}^{\lambda}
\bigl(
\chi\,\vartheta_{Y,v,h,\alpha}^{\mathrm{old}}
\bigr)
\right\|_{E}
\le
C_{t,\chi}\frac{\nu_Y}{\eta}
h^{\alpha\beta-\frac{1-\alpha}{2}}
\|v\|_{\mathfrak G_t}.
\end{equation}
\end{proposition}

\begin{proof}
For item~(i), the change of variables \(q=s/\varepsilon\) gives
\begin{align*}
\left\|
\widehat\Psi_{v,h,\alpha}^{\mathrm{old}}
\right\|_{L^2(\R_+)}^2
&=
\varepsilon^{-2\beta-1}
\int_1^{R_{h,\alpha}}
\left|
\Psi_{v,h}\!\left(\frac{q}{\varepsilon}\right)
\right|^2\,dq\\
&=
\varepsilon^{-2\beta}
\int_{\varepsilon^{-1}}^{t/h}
\left|
\Psi_{v,h}(s)
\right|^2\,ds
\le
\varepsilon^{-2\beta}h^{2\beta}\|\mathcal S_t v\|_{H_t}^2.
\end{align*}
Since \(\varepsilon=h^{1-\alpha}\), this gives
\eqref{eq:full-h-target-old-tail-alpha-trace-bound}.

To prove item~(ii), let \(u\in[0,1]\). Using \eqref{eq:unit-interval-target-old-tail},
\eqref{eq:full-h-traces}, and the change of variables \(q=s/\varepsilon\), we obtain
\begin{align*}
\chi(u)\vartheta_{Y,v,h,\alpha}^{\mathrm{old}}(u)
&=
\chi(u)\varepsilon^{H_{XY}-\frac12}
\int_{\varepsilon^{-1}}^{t/h}
\left(
q-\frac{u}{\varepsilon}
\right)^{\lambda-1}
L_Y\!\left(t-h\frac{u}{\varepsilon},t-hq\right)
\Psi_{v,h}(q)\,dq\\
&=
\chi(u)\varepsilon^{H_{XY}-\frac12}
\int_1^{R_{h,\alpha}}
\left(
\frac{q-u}{\varepsilon}
\right)^{\lambda-1}
L_Y(t-h^\alpha u,t-h^\alpha q)
\Psi_{v,h}\!\left(\frac{q}{\varepsilon}\right)
\frac{dq}{\varepsilon}\\
&=
\varepsilon^{H_{XY}-\lambda-\frac12}
\chi(u)
\int_1^{R_{h,\alpha}}
(q-u)^{\lambda-1}
L_Y(t-h^\alpha u,t-h^\alpha q)
\Psi_{v,h}\!\left(\frac{q}{\varepsilon}\right)\,dq.
\end{align*}
Since \(\lambda=H_Y+\tfrac12\) and \(\beta=H_Y-H_{XY}\),
\[
\varepsilon^{H_{XY}-\lambda-\frac12}
\Psi_{v,h}\!\left(\frac{q}{\varepsilon}\right)
=
\varepsilon^{-1/2}\widehat\Psi_{v,h,\alpha}^{\mathrm{old}}(q).
\]
This proves \eqref{eq:full-h-target-old-tail-alpha-id}.

For item~(iii), extend \(\chi\) by zero to a function
\[
\overline\chi\in C^\infty([0,\infty)),
\]
which is possible because \(\chi\) is flat at the outer edge. For \(q\in[1,R_{h,\alpha}]\), define
\[
a_{Y,t,h,\alpha,\chi}(u,q)
:=
\overline\chi(u)L_Y(t-h^\alpha u,t-h^\alpha q),
\qquad 0\le u\le q.
\]
Fix \(\Theta\in L^2([1,R_{h,\alpha}])\). Exactly as in the proof of
Proposition~\ref{prop:micro-tail-Y}, Fubini's theorem gives
\[
\frac{1}{\Gamma(\lambda)}
\mathfrak D_{-,1}^{\lambda}
\bigl(
\mathfrak W_{Y,t,h,\alpha,\chi}^{\mathrm{old},\alpha}\Theta
\bigr)(u)
=
\int_1^{R_{h,\alpha}}
K_{Y,t,h,\alpha,\chi}^{\mathrm{old},\alpha}(u,q)\Theta(q)\,dq
\]
for almost every \(u\in[0,1]\), where
\[
K_{Y,t,h,\alpha,\chi}^{\mathrm{old},\alpha}(u,q)
:=
-
\frac{1}{\Gamma(\lambda)\Gamma(1-\lambda)}
\int_0^1
s^{-\lambda}(1-s)^\lambda
\partial_1 a_{Y,t,h,\alpha,\chi}(u+(q-u)s,q)\,ds.
\]
This defines the operator
\(\mathfrak Q_{Y,t,h,\alpha,\chi}^{\mathrm{old},\alpha}\).

To bound its norm, split
\[
\partial_1 a_{Y,t,h,\alpha,\chi}
=
\overline\chi'(\cdot)L_Y(t-h^\alpha\cdot,t-h^\alpha q)
-h^\alpha \overline\chi(\cdot)\partial_1L_Y(t-h^\alpha\cdot,t-h^\alpha q).
\]
For the first term, the support of \(\overline\chi'\) is contained in \([0,1]\), so the integral in
\(s\) is effectively restricted to
\[
0\le s\le \rho(u,q):=\frac{1-u}{q-u}.
\]
Using the change of variables \(x=u+(q-u)s\), one obtains
\[
\left|
K_{Y,t,h,\alpha,\chi}^{(1)}(u,q)
\right|
\le
C_{t,\chi}(q-u)^{-1}
\int_0^{\rho(u,q)} s^{-\lambda}\,ds
\le
C_{t,\chi}(1-u)^{1-\lambda}(q-u)^{-(2-\lambda)}.
\]
Since \(2-\lambda>1\), this kernel is Hilbert--Schmidt on
\([0,1]\times[1,\infty)\), hence the corresponding operator is bounded from
\(L^2([1,R_{h,\alpha}])\) to \(E\), uniformly in \(h\).

For the second term,
\[
\left|
K_{Y,t,h,\alpha,\chi}^{(2)}(u,q)
\right|
\le
C_{t,\chi}h^\alpha
\int_0^{\rho(u,q)} s^{-\lambda}\,ds
\le
C_{t,\chi}h^\alpha(1-u)^{1-\lambda}(q-u)^{-(1-\lambda)}.
\]
Its Hilbert--Schmidt norm is bounded by
\begin{align*}
\left\|
K_{Y,t,h,\alpha,\chi}^{(2)}
\right\|_{L^2([0,1]\times[1,R_{h,\alpha}])}^2
&\le
C_{t,\chi}h^{2\alpha}
\int_0^1
(1-u)^{2(1-\lambda)}
\int_1^{R_{h,\alpha}}
(q-u)^{-2(1-\lambda)}\,dq\,du\\
&\le
C_{t,\chi}h^{2\alpha}R_{h,\alpha}^{\,2\lambda-1}
\le
C_{t,\chi},
\end{align*}
because \(2\lambda-1=2H_Y\in(0,1)\) and \(R_{h,\alpha}=th^{-\alpha}\). Thus the second operator is
also bounded uniformly in \(h\). This proves \eqref{eq:full-h-target-old-tail-alpha-opbound}.

Finally, combining \eqref{eq:full-h-target-old-tail-alpha-id},
\eqref{eq:full-h-target-old-tail-alpha-conj},
\eqref{eq:full-h-target-old-tail-alpha-opbound}, and
\eqref{eq:full-h-target-old-tail-alpha-trace-bound} gives
\[
\left\|
\frac{1}{\Gamma(\lambda)}
\mathfrak D_{-,1}^{\lambda}
\bigl(
\chi\,\vartheta_{Y,v,h,\alpha}^{\mathrm{old}}
\bigr)
\right\|_E
\le
\varepsilon^{-1/2}
\left\|
\mathfrak Q_{Y,t,h,\alpha,\chi}^{\mathrm{old},\alpha}
\right\|
\left\|
\widehat\Psi_{v,h,\alpha}^{\mathrm{old}}
\right\|_{L^2(\R_+)}
\le
C_{t,\chi}\frac{\nu_Y}{\eta}
h^{\alpha\beta-\frac{1-\alpha}{2}}
\|v\|_{\mathfrak G_t},
\]
which is \eqref{eq:full-h-target-old-tail-alpha-bound}.
\end{proof}

\begin{corollary}[Explicit bound for the flat-edge target commutator]
\label{cor:full-h-target-old-tail-alpha}
Under the hypotheses of Proposition~\ref{prop:full-h-target-old-tail-alpha},
\begin{equation}\label{eq:full-h-target-old-tail-alpha-criterion}
\left\|
\frac{1}{\Gamma(\lambda)}
\mathfrak D_{-,1}^{\lambda}
\Bigl(
\varepsilon_{h,\alpha}^{H_{XY}}
\mathfrak l_{\varepsilon_{h,\alpha}}
\mathfrak C_{Y,v,h,\alpha}^{\mathrm{old}}
\Bigr)
\right\|_{E}
\le
C_{t,\chi}\frac{\nu_Y}{\eta}
h^{\alpha\beta-\frac{1-\alpha}{2}}
\|v\|_{\mathfrak G_t},
\end{equation}
where \(\lambda=H_Y+\frac12\).
\end{corollary}

\begin{proof}
By Proposition~\ref{prop:full-h-target-old-tail},
\[
\chi\,\vartheta_{Y,v,h,\alpha}^{\mathrm{old}}
=
-\varepsilon_{h,\alpha}^{H_{XY}}
\mathfrak l_{\varepsilon_{h,\alpha}}
\mathfrak C_{Y,v,h,\alpha}^{\mathrm{old}}
\qquad\text{in }E.
\]
Taking \((\Gamma(\lambda))^{-1}\mathfrak D_{-,1}^{\lambda}\), then applying
\eqref{eq:full-h-target-old-tail-alpha-bound}, gives
\eqref{eq:full-h-target-old-tail-alpha-criterion}.
\end{proof}

\begin{remark}[Limitation of the isolated target estimate]
\label{rem:full-h-target-old-tail-alpha}
Proposition~\ref{prop:full-h-target-old-tail-alpha} gives an explicit \(\alpha\)-scale realization
of the target old tail, but it also exposes the obstruction. After the first
full-\(h\) dilation and the second \(\alpha\)-scale zoom, the target remainder carries an unavoidable
prefactor \(\varepsilon_{h,\alpha}^{-1/2}=h^{-(1-\alpha)/2}\). So the separate target-side route
produces, at best, the weaker scale \(h^{\alpha\beta-(1-\alpha)/2}\). This is why the argument
cannot close the endpoint by bounding the target old tail in isolation: the governing
endpoint has to preserve the source-target coupling encoded by
\(\mathfrak Y_{v,h,\alpha}^{\mathrm{src}}\) or, equivalently, by the corrected full-\(h\)
realization \(\mathfrak C_{v,h,\alpha}^{\mathrm{pair}}\).
\end{remark}

\begin{definition}[Full-\(h\) truncated-target realization of the preconjugated discrepancy]
\label{def:full-h-coupled-commutator}
Assume the hypotheses of Proposition~\ref{prop:moving-cutoff-decomposition}. For \(v\in\mathfrak G_t\),
set
\[
\lambda:=H_Y+\frac12,
\qquad
\mu:=H_{XY}+\frac12,
\qquad
\varepsilon:=\varepsilon_{h,\alpha}=h^{1-\alpha},
\]
and for \(0<h\le t\), define
\begin{equation}\label{eq:full-h-coupled-commutator}
\begin{aligned}
\mathfrak C_{v,h,\alpha}^{\mathrm{pair}}
={}&
\one_{[0,\varepsilon^{-1}]}
\Biggl[
\mathfrak V_{XY,t,h}^{\mathrm{full}}
\bigl(
\widetilde\chi_{h,\alpha}\Phi_{v,h}
\bigr)
\\
&\qquad
-
\widetilde\chi_{h,\alpha}(\cdot)
\int_{\cdot}^{\varepsilon^{-1}}
\Bigl(
(q-\cdot)^{\lambda-1}
L_Y(t-h\cdot,t-hq)\Psi_{v,h}(q)
\Bigr)\,dq
\Biggr]
\in L^2(\R_+).
\end{aligned}
\end{equation}
\end{definition}

\begin{proposition}[Exact full-\(h\) realization of the preconjugated discrepancy]
\label{prop:full-h-coupled-commutator}
Assume the hypotheses of Definition~\ref{def:full-h-coupled-commutator}. Then:
\begin{enumerate}[label=(\roman*)]
\item \(\mathfrak C_{v,h,\alpha}^{\mathrm{pair}}\) is supported in
\([0,\varepsilon_{h,\alpha}^{-1}]\);
\item one has
\begin{equation}\label{eq:full-h-coupled-commutator-id}
\mathfrak Y_{v,h,\alpha}^{\mathrm{src}}
=
\varepsilon_{h,\alpha}^{H_{XY}}
\mathfrak l_{\varepsilon_{h,\alpha}}
\mathfrak C_{v,h,\alpha}^{\mathrm{pair}}
\qquad\text{in }E.
\end{equation}
\end{enumerate}
\end{proposition}

\begin{proof}
  Let \(\varepsilon:=\varepsilon_{h,\alpha}\). Since \(\widetilde\chi_{h,\alpha}\) is supported in
  \([0,\varepsilon^{-1}]\), the second term in \eqref{eq:full-h-coupled-commutator} vanishes for
  \(r>\varepsilon^{-1}\). The first term also vanishes there because
  \(\widetilde\chi_{h,\alpha}\Phi_{v,h}\) is supported in \([0,\varepsilon^{-1}]\) and the full
  Volterra integral starts at \(q=r\). This proves (i).

  By Proposition~\ref{prop:micro-volterra},
  \[
  \mathfrak A_{XY,t,h,\alpha}
  \bigl(
  \chi\phi_{v,h}^{\sharp}
  \bigr)
  =
  \varepsilon^{H_{XY}}
  \mathfrak l_{\varepsilon}
  \mathfrak V_{XY,t,h}^{\mathrm{full}}
  \bigl(
  \widetilde\chi_{h,\alpha}\Phi_{v,h}
  \bigr),
  \]
  because \(\widetilde\chi_{h,\alpha}\Phi_{v,h}\) is supported in \([0,\varepsilon^{-1}]\).

  For the target term, using \(\psi_{v,h}^{\sharp}=\varepsilon^{-\beta}\mathfrak l_{\varepsilon}\Psi_{v,h}\),
  \(\lambda=H_Y+\tfrac12\), \(\beta=H_Y-H_{XY}\), and the change of variables \(q=\varepsilon^{-1}v\),
  one gets for \(u\in[0,1]\),
  \begin{align*}
  \chi(u)\mathfrak A_{Y,t,h,\alpha}\psi_{v,h}^{\sharp}(u)
  &=
  \chi(u)
  \int_u^1
  (v-u)^{\lambda-1}
  L_Y(t-h^\alpha u,t-h^\alpha v)\psi_{v,h}^{\sharp}(v)\,dv\\
  &=
  \varepsilon^{H_Y-\beta}
  \chi(u)
  \int_{\varepsilon^{-1}u}^{\varepsilon^{-1}}
  (q-\varepsilon^{-1}u)^{\lambda-1}
  L_Y(t-h\varepsilon^{-1}u,t-hq)\Psi_{v,h}(q)\,dq\\
  &=
  \varepsilon^{H_{XY}}
  \mathfrak l_{\varepsilon}
  \Biggl[
  \widetilde\chi_{h,\alpha}(r)
  \int_r^{\varepsilon^{-1}}
  (q-r)^{\lambda-1}
  L_Y(t-hr,t-hq)\Psi_{v,h}(q)\,dq
  \Biggr](u).
  \end{align*}
    Subtracting the two identities proves \eqref{eq:full-h-coupled-commutator-id}.
  \end{proof}

\begin{proposition}[Exact split of the corrected full-\(h\) realization]
\label{prop:full-h-coupled-split}
Assume the hypotheses of Definition~\ref{def:full-h-coupled-commutator}. Then
\begin{equation}\label{eq:full-h-coupled-split}
\mathfrak C_{v,h,\alpha}^{\mathrm{pair}}
=
\mathfrak Z_{v,h,\alpha}^{\mathrm{src}}
-\mathfrak C_{Y,v,h,\alpha}^{\mathrm{old}}
\qquad\text{in }L^2(\R_+),
\end{equation}
where \(\mathfrak Z_{v,h,\alpha}^{\mathrm{src}}\) is the full-\(h\) source discrepancy from
\eqref{eq:full-h-source-discrepancy} and \(\mathfrak C_{Y,v,h,\alpha}^{\mathrm{old}}\) is the target
old-tail commutator from \eqref{eq:full-h-target-old-tail}.
\end{proposition}

\begin{proof}
  Let \(\varepsilon:=\varepsilon_{h,\alpha}\). For \(r>\varepsilon^{-1}\), all three terms in
  \eqref{eq:full-h-coupled-split} vanish, so it is enough to work on \(0\le r\le \varepsilon^{-1}\).
  By definition,
  \begin{align*}
  \mathfrak C_{v,h,\alpha}^{\mathrm{pair}}(r)
  &=
  \mathfrak V_{XY,t,h}^{\mathrm{full}}
  \bigl(
  \widetilde\chi_{h,\alpha}\Phi_{v,h}
  \bigr)(r) \\
  &\quad-
  \widetilde\chi_{h,\alpha}(r)
  \int_r^{\varepsilon^{-1}}
  (q-r)^{\lambda-1}
  L_Y(t-hr,t-hq)\Psi_{v,h}(q)\,dq.
  \end{align*}
  Adding and subtracting
  \[
  \widetilde\chi_{h,\alpha}(r)
  \int_r^{t/h}
  (q-r)^{\mu-1}
  L_{XY}(t-hr,t-hq)\Phi_{v,h}(q)\,dq
  \]
  and using the exact full-\(h\) identity \eqref{eq:full-h-exact} gives
  \begin{align*}
  \mathfrak C_{v,h,\alpha}^{\mathrm{pair}}(r)
  &=
  \mathfrak V_{XY,t,h}^{\mathrm{full}}
  \bigl(
  \widetilde\chi_{h,\alpha}\Phi_{v,h}
  \bigr)(r) \\
  &\quad-
  \widetilde\chi_{h,\alpha}(r)
  \int_r^{t/h}
  (q-r)^{\mu-1}
  L_{XY}(t-hr,t-hq)\Phi_{v,h}(q)\,dq \\
  &\quad+
  \widetilde\chi_{h,\alpha}(r)
  \int_{\varepsilon^{-1}}^{t/h}
  (q-r)^{\lambda-1}
  L_Y(t-hr,t-hq)\Psi_{v,h}(q)\,dq.
  \end{align*}
  The first two lines are exactly \(\mathfrak Z_{v,h,\alpha}^{\mathrm{src}}(r)\), while the last line
  is \(-\mathfrak C_{Y,v,h,\alpha}^{\mathrm{old}}(r)\) by
  \eqref{eq:full-h-target-old-tail}. This proves \eqref{eq:full-h-coupled-split}.
\end{proof}

\begin{corollary}[Endpoint reduction to the corrected full-\(h\) realization]
\label{cor:full-h-coupled-commutator}
Assume the hypotheses of Proposition~\ref{prop:moving-cutoff-decomposition} and the extra
regularity from Definition~\ref{def:model-c2}. If there exist
\[
C_{\mathrm{pair},\alpha}<\infty,
\qquad
h_{0,\alpha}\in(0,t]
\]
such that for every \(0<h\le h_{0,\alpha}\) and every \(v\in\mathfrak G_t\),
\begin{equation}\label{eq:full-h-coupled-commutator-bound}
\left\|
\frac{1}{\Gamma(\lambda)}
\mathfrak D_{-,1}^{\lambda}
\Bigl(
\varepsilon_{h,\alpha}^{H_{XY}}
\mathfrak l_{\varepsilon_{h,\alpha}}
\mathfrak C_{v,h,\alpha}^{\mathrm{pair}}
\Bigr)
\right\|_{E}
\le
C_{\mathrm{pair},\alpha}h^{\alpha\beta}\|v\|_{\mathfrak G_t},
\qquad
\lambda:=H_Y+\frac12,
\end{equation}
then the conclusion of Corollary~\ref{cor:unit-interval-source-preconj} holds. In particular,
\(\mathcal S_t\) has \(h^\alpha\)-stable boundary cutoffs at time \(t\), and the screening
conclusions \eqref{eq:unscaled-smooth-cutoff-screening-F}--
\eqref{eq:unscaled-smooth-cutoff-screening-G} hold.
\end{corollary}

\begin{proof}
By Proposition~\ref{prop:full-h-coupled-commutator},
\[
\mathfrak Y_{v,h,\alpha}^{\mathrm{src}}
=
\varepsilon_{h,\alpha}^{H_{XY}}
\mathfrak l_{\varepsilon_{h,\alpha}}
\mathfrak C_{v,h,\alpha}^{\mathrm{pair}}
\qquad\text{in }E.
\]
Applying \((\Gamma(\lambda))^{-1}\mathfrak D_{-,1}^{\lambda}\) and using
\eqref{eq:full-h-coupled-commutator-bound} gives \eqref{eq:unit-interval-source-preconj-bound}.
Corollary~\ref{cor:unit-interval-source-preconj} now gives the conclusion.
\end{proof}

\begin{definition}[Coupled \(\alpha\)-scale endpoint traces]
\label{def:full-h-coupled-alpha}
Assume the hypotheses of Definition~\ref{def:full-h-coupled-commutator}. Set
\[
\beta:=H_Y-H_{XY},
\qquad
\varepsilon:=\varepsilon_{h,\alpha}=h^{1-\alpha},
\qquad
R_{h,\alpha}:=th^{-\alpha}.
\]
Define
\[
\widehat E_{h,\alpha}^{X}:=L^2([0,R_{h,\alpha}]),
\qquad
\widehat E_{h,\alpha}^{Y}:=L^2([0,R_{h,\alpha}]),
\qquad
\widehat E_{h,\alpha}^{Y,\mathrm{old}}:=L^2([1,R_{h,\alpha}]).
\]
For \(v\in\mathfrak G_t\), define the extended \(\alpha\)-scale source trace and the
\(\alpha\)-scale target traces by
\begin{align}
\widehat\Phi_{v,h,\alpha}(q)
&:=
\varepsilon^{-1/2}
\Phi_{v,h}\!\left(\frac{q}{\varepsilon}\right)
\one_{[0,R_{h,\alpha}]}(q),
\label{eq:full-h-coupled-alpha-X}\\
\widehat\Psi_{v,h,\alpha}^{\mathrm{full}}(q)
&:=
\varepsilon^{-(\beta+\frac12)}
\Psi_{v,h}\!\left(\frac{q}{\varepsilon}\right)
\one_{[0,R_{h,\alpha}]}(q),
\label{eq:full-h-coupled-alpha-Y-full}\\
\widehat\Psi_{v,h,\alpha}^{\mathrm{old}}(q)
&:=
\varepsilon^{-(\beta+\frac12)}
\Psi_{v,h}\!\left(\frac{q}{\varepsilon}\right)
\one_{[1,R_{h,\alpha}]}(q).
\label{eq:full-h-coupled-alpha-Y}
\end{align}
Finally, define the admissible coupled \(\alpha\)-scale trace class by
\begin{equation}\label{eq:full-h-coupled-alpha-class}
\widehat{\mathcal A}_{t,h,\alpha}^{\mathrm{pair}}
:=
\left\{
\bigl(
\widehat\Phi_{v,h,\alpha},
\widehat\Psi_{v,h,\alpha}^{\mathrm{old}}
\bigr)
:
v\in\mathfrak G_t
\right\}
\subset
\widehat E_{h,\alpha}^{X}\times \widehat E_{h,\alpha}^{Y,\mathrm{old}}.
\end{equation}
\end{definition}

\begin{proposition}[Exact full \(\alpha\)-scale pair on the growing window]
\label{prop:full-h-coupled-alpha-exact}
Assume the hypotheses of Definition~\ref{def:full-h-coupled-alpha}. Then for every
\[
v\in\mathfrak G_t,
\qquad
0<h\le t,
\qquad
h^\alpha\le t,
\]
the pair
\[
\bigl(
\widehat\Phi_{v,h,\alpha},
\widehat\Psi_{v,h,\alpha}^{\mathrm{full}}
\bigr)
\]
belongs to the exact fixed-window graph
\[
\Gamma_{t,h,\alpha,R_{h,\alpha}}^{\mathrm{ex}}.
\]
Equivalently,
\begin{equation}\label{eq:full-h-coupled-alpha-exact}
\mathfrak T_{t,h,\alpha,R_{h,\alpha}}^{\mathrm{ex}}
\widehat\Phi_{v,h,\alpha}
=
\widehat\Psi_{v,h,\alpha}^{\mathrm{full}}.
\end{equation}
Moreover,
\begin{equation}\label{eq:full-h-coupled-alpha-exact-bound}
\left\|
\widehat\Phi_{v,h,\alpha}
\right\|_{E_{R_{h,\alpha}}}^2
+
\frac{\eta^2}{\nu_Y^2}
\left\|
\widehat\Psi_{v,h,\alpha}^{\mathrm{full}}
\right\|_{E_{R_{h,\alpha}}}^2
\le
\|v\|_{\mathfrak G_t}^2.
\end{equation}
\end{proposition}

\begin{proof}
Let \(\varepsilon:=\varepsilon_{h,\alpha}\) and \(R:=R_{h,\alpha}=t\varepsilon/h\). By
\eqref{eq:full-h-exact},
\[
\mathfrak V_{Y,t,h}^{\mathrm{full}}\Psi_{v,h}
=
\mathfrak V_{XY,t,h}^{\mathrm{full}}\Phi_{v,h}
\qquad\text{in }L^2(\R_+).
\]
For \(0\le r\le R\), Proposition~\ref{prop:micro-volterra} with the change of variables
\(q=s/\varepsilon\) gives
\begin{align*}
\bigl(
\mathfrak V_{Y,t,h,\alpha,R}^{\mathrm{mic}}
\widehat\Psi_{v,h,\alpha}^{\mathrm{full}}
\bigr)(r)
&=
\int_r^{R}
(q-r)^{H_Y-\frac12}
L_Y(t-h^\alpha r,t-h^\alpha q)
\varepsilon^{-(\beta+\frac12)}
\Psi_{v,h}\!\left(\frac{q}{\varepsilon}\right)\,dq\\
&=
\varepsilon^{H_{XY}}
\bigl(
\mathfrak V_{Y,t,h}^{\mathrm{full}}\Psi_{v,h}
\bigr)\!\left(\frac{r}{\varepsilon}\right),
\end{align*}
because \(h^\alpha=\varepsilon h\) and \(\beta=H_Y-H_{XY}\). The same computation gives
\[
\bigl(
\mathfrak V_{XY,t,h,\alpha,R}^{\mathrm{mic}}
\widehat\Phi_{v,h,\alpha}
\bigr)(r)
=
\varepsilon^{H_{XY}}
\bigl(
\mathfrak V_{XY,t,h}^{\mathrm{full}}\Phi_{v,h}
\bigr)\!\left(\frac{r}{\varepsilon}\right).
\]
Hence
\[
\mathfrak V_{Y,t,h,\alpha,R}^{\mathrm{mic}}
\widehat\Psi_{v,h,\alpha}^{\mathrm{full}}
=
\mathfrak V_{XY,t,h,\alpha,R}^{\mathrm{mic}}
\widehat\Phi_{v,h,\alpha}
\qquad\text{in }E_R,
\]
so
\[
\bigl(
\widehat\Phi_{v,h,\alpha},
\widehat\Psi_{v,h,\alpha}^{\mathrm{full}}
\bigr)
\in
\Gamma_{t,h,\alpha,R}^{\mathrm{ex}}
\]
by Definition~\ref{def:micro-window-exact}. Proposition~\ref{prop:micro-window-solver} then gives
\eqref{eq:full-h-coupled-alpha-exact}.

Finally, by the change of variables \(q=s/\varepsilon\),
\[
\left\|
\widehat\Phi_{v,h,\alpha}
\right\|_{E_R}
=
\left\|
\Phi_{v,h}
\right\|_{L^2([0,t/h])}
=
\|v\|_{H_t},
\]
while
\[
\left\|
\widehat\Psi_{v,h,\alpha}^{\mathrm{full}}
\right\|_{E_R}
=
\varepsilon^{-\beta}
\left\|
\Psi_{v,h}
\right\|_{L^2([0,t/h])}
=
h^{\alpha\beta}\|\mathcal S_t v\|_{H_t}
\le
\|\mathcal S_t v\|_{H_t}.
\]
Combining these with the definition of \(\|v\|_{\mathfrak G_t}\) proves
\eqref{eq:full-h-coupled-alpha-exact-bound}.
\end{proof}

\begin{definition}[Coupled \(\alpha\)-scale endpoint operator]
\label{def:full-h-coupled-alpha-op}
Under the hypotheses of Definition~\ref{def:full-h-coupled-alpha}, set
\[
\lambda:=H_Y+\frac12,
\qquad
\mu:=H_{XY}+\frac12.
\]
For
\[
\Theta_X\in \widehat E_{h,\alpha}^{X},
\qquad
\Theta_Y\in \widehat E_{h,\alpha}^{Y,\mathrm{old}},
\]
define
\[
\mathfrak W_{t,h,\alpha,\chi}^{\mathrm{pair},\alpha}
\bigl(
\Theta_X,\Theta_Y
\bigr)
\in E=L^2([0,1])
\]
by
\begin{align}
\bigl(
\mathfrak W_{t,h,\alpha,\chi}^{\mathrm{pair},\alpha}
(\Theta_X,\Theta_Y)
\bigr)(u)
&:=
\mathfrak W_{XY,t,h,\alpha,\chi}^{\mathrm{loc}}
\bigl(
\Theta_X|_{[0,1]}
\bigr)(u)
\notag\\
&\quad+
\varepsilon_{h,\alpha}^{-1/2}\chi(u)
\int_1^{R_{h,\alpha}}
\Bigl[
(q-u)^{\lambda-1}
L_Y(t-h^\alpha u,t-h^\alpha q)\Theta_Y(q)
\notag\\
&\hspace{7.2em}-
(q-u)^{\mu-1}
L_{XY}(t-h^\alpha u,t-h^\alpha q)\Theta_X(q)
\Bigr]\,dq,
\label{eq:full-h-coupled-alpha-op}
\end{align}
for \(0\le u\le1\). Also define the conjugated distributional realization by
\begin{equation}\label{eq:full-h-coupled-alpha-op-conj}
\mathfrak Q_{t,h,\alpha,\chi}^{\mathrm{pair},\alpha}
\bigl(
\Theta_X,\Theta_Y
\bigr)
:=
\frac{1}{\Gamma(\lambda)}
\mathfrak D_{-,1}^{\lambda}
\mathfrak W_{t,h,\alpha,\chi}^{\mathrm{pair},\alpha}
\bigl(
\Theta_X,\Theta_Y
\bigr)
\in \mathcal D'((0,1)).
\end{equation}
\end{definition}

\begin{proposition}[Exact coupled \(\alpha\)-scale realization of the preconjugated discrepancy]
\label{prop:full-h-coupled-alpha}
Assume the hypotheses of Definition~\ref{def:full-h-coupled-alpha}. Then for every
\[
v\in\mathfrak G_t,
\qquad
0<h\le t,
\qquad
h^\alpha\le t,
\]
one has
\begin{equation}\label{eq:full-h-coupled-alpha-id}
\mathfrak Y_{v,h,\alpha}^{\mathrm{src}}
=
\mathfrak W_{t,h,\alpha,\chi}^{\mathrm{pair},\alpha}
\bigl(
\widehat\Phi_{v,h,\alpha},
\widehat\Psi_{v,h,\alpha}^{\mathrm{old}}
\bigr)
\qquad\text{in }E.
\end{equation}
Consequently, in \(\mathcal D'((0,1))\),
\begin{equation}\label{eq:full-h-coupled-alpha-conj}
\frac{1}{\Gamma(\lambda)}
\mathfrak D_{-,1}^{\lambda}
\mathfrak Y_{v,h,\alpha}^{\mathrm{src}}
=
\mathfrak Q_{t,h,\alpha,\chi}^{\mathrm{pair},\alpha}
\bigl(
\widehat\Phi_{v,h,\alpha},
\widehat\Psi_{v,h,\alpha}^{\mathrm{old}}
\bigr).
\end{equation}
\end{proposition}

\begin{proof}
Let \(\varepsilon:=\varepsilon_{h,\alpha}\) and \(R:=R_{h,\alpha}\). By
Proposition~\ref{prop:unit-interval-source-full},
\[
\mathfrak Y_{v,h,\alpha}^{\mathrm{src}}
\!=\!
\mathfrak W_{XY,t,h,\alpha,\chi}^{\mathrm{loc}}\phi_{v,h}^{\sharp}
-\chi\,\vartheta_{XY,v,h,\alpha}^{\mathrm{old}}
+\chi\,\vartheta_{Y,v,h,\alpha}^{\mathrm{old}}
\qquad\text{in }E.
\]
Since
\[
\widehat\Phi_{v,h,\alpha}|_{[0,1]}
=
\phi_{v,h}^{\sharp},
\]
it remains to identify the two old-tail pieces.
For the target contribution, \eqref{eq:unit-interval-target-old-tail},
\eqref{eq:full-h-coupled-alpha-Y}, and the change of variables \(q=s/\varepsilon\) give
\begin{equation}\label{eq:full-h-coupled-alpha-old-Y}
\chi\,\vartheta_{Y,v,h,\alpha}^{\mathrm{old}}(u)
=
\varepsilon^{-1/2}\chi(u)
\int_1^{R}
(q-u)^{\lambda-1}
L_Y(t-h^\alpha u,t-h^\alpha q)
\widehat\Psi_{v,h,\alpha}^{\mathrm{old}}(q)\,dq
\end{equation}
for \(0\le u\le1\), because \(\lambda=H_Y+\tfrac12\) and \(\beta=H_Y-H_{XY}\).

For the source contribution, \eqref{eq:unit-interval-source-old-tail},
\eqref{eq:full-h-coupled-alpha-X}, and the change of variables \(q=s/\varepsilon\) give
\begin{align*}
\chi(u)\vartheta_{XY,v,h,\alpha}^{\mathrm{old}}(u)
&=
\chi(u)\varepsilon^{H_{XY}-\frac12}
\int_{\varepsilon^{-1}}^{t/h}
\left(
q-\frac{u}{\varepsilon}
\right)^{\mu-1}
L_{XY}\!\left(t-h\frac{u}{\varepsilon},t-hq\right)
\Phi_{v,h}(q)\,dq\\
&=
\chi(u)\varepsilon^{H_{XY}-\frac12}
\int_1^{R}
\left(
\frac{q-u}{\varepsilon}
\right)^{\mu-1}
L_{XY}(t-h^\alpha u,t-h^\alpha q)
\Phi_{v,h}\!\left(\frac{q}{\varepsilon}\right)
\frac{dq}{\varepsilon}\\
&=
\varepsilon^{-1/2}\chi(u)
\int_1^{R}
(q-u)^{\mu-1}
L_{XY}(t-h^\alpha u,t-h^\alpha q)
\widehat\Phi_{v,h,\alpha}(q)\,dq,
\end{align*}
because \(\mu=H_{XY}+\tfrac12\). Combining this with
\eqref{eq:full-h-coupled-alpha-old-Y} proves
\eqref{eq:full-h-coupled-alpha-id}. Applying
\((\Gamma(\lambda))^{-1}\mathfrak D_{-,1}^{\lambda}\) gives
\eqref{eq:full-h-coupled-alpha-conj}.
\end{proof}

\begin{proposition}[Exact defect realization of the coupled \(\alpha\)-scale endpoint]
\label{prop:full-h-coupled-alpha-defect}
Assume the hypotheses of Proposition~\ref{prop:full-h-coupled-alpha-exact}, and suppose moreover
that \(\chi\) is flat at the outer edge in the sense of Definition~\ref{def:endpoint-flat-cutoff}.
Let
\[
\overline\chi(q)
:=
\begin{cases}
\chi(q), & 0\le q\le1,\\
0, & 1<q\le R_{h,\alpha}.
\end{cases}
\]
Then \(\overline\chi\in C^\infty([0,R_{h,\alpha}])\), and for every \(v\in\mathfrak G_t\) one has
\[
\overline\chi\,\widehat\Phi_{v,h,\alpha}\in H^\beta((0,R_{h,\alpha})),
\qquad
\overline\chi\,\widehat\Psi_{v,h,\alpha}^{\mathrm{full}}\in E_{R_{h,\alpha}},
\]
and
\begin{equation}\label{eq:full-h-coupled-alpha-defect}
\one_{[0,1]}
\mathfrak F_{t,h,\alpha,R_{h,\alpha}}^{\mathrm{ex}}
\bigl(
\overline\chi\,\widehat\Phi_{v,h,\alpha},
\overline\chi\,\widehat\Psi_{v,h,\alpha}^{\mathrm{full}}
\bigr)
=
\mathfrak Q_{t,h,\alpha,\chi}^{\mathrm{pair},\alpha}
\bigl(
\widehat\Phi_{v,h,\alpha},
\widehat\Psi_{v,h,\alpha}^{\mathrm{old}}
\bigr)
\qquad\text{in }E.
\end{equation}
\end{proposition}

\begin{proof}
Set
\[
R:=R_{h,\alpha},
\qquad
\Phi:=\widehat\Phi_{v,h,\alpha},
\qquad
\Psi:=\widehat\Psi_{v,h,\alpha}^{\mathrm{full}}.
\]
By Proposition~\ref{prop:full-h-coupled-alpha-exact},
\[
(\Phi,\Psi)\in \Gamma_{t,h,\alpha,R}^{\mathrm{ex}}.
\]
Since \(\overline\chi\in C^\infty([0,R])\) and \(\Phi\in H^\beta((0,R))\) by
Theorem~\ref{thm:micro-window-regularity}, one has
\[
\overline\chi\,\Phi\in H^\beta((0,R)).
\]

Because \((\Phi,\Psi)\) is an exact fixed-window pair, the proof of
Theorem~\ref{thm:quant-window-buffer} applied on the window \([0,R]\) gives
\[
\mathfrak F_{t,h,\alpha,R}^{\mathrm{ex}}
\bigl(
\overline\chi\,\Phi,
\overline\chi\,\Psi
\bigr)
=
\mathfrak Q_{Y,t,h,\alpha,\overline\chi}^{\mathrm{com}}\Psi
-
\mathfrak Q_{XY,t,h,\alpha,\overline\chi}^{\mathrm{com}}\Phi
\qquad\text{in }E_R.
\]
Restrict now to \(r\in[0,1]\). Since \(\overline\chi(r)=\chi(r)\), while
\(\overline\chi(q)=0\) for \(q\ge1\), the two commutator integrals split into the local
unit-interval commutator on \([r,1]\) and the coupled tail contribution on \([1,R]\).  Using
\(\Psi|_{[1,R]}=\widehat\Psi_{v,h,\alpha}^{\mathrm{old}}\), the resulting expression is the
distributional definition \eqref{eq:full-h-coupled-alpha-op-conj} of
\[
\mathfrak Q_{t,h,\alpha,\chi}^{\mathrm{pair},\alpha}
\bigl(
\widehat\Phi_{v,h,\alpha},
\widehat\Psi_{v,h,\alpha}^{\mathrm{old}}
\bigr)
\]
on \(E=L^2([0,1])\). This proves \eqref{eq:full-h-coupled-alpha-defect}.
\end{proof}

\begin{corollary}[Endpoint reduction to a growing-window exact defect]
\label{cor:full-h-coupled-alpha-defect}
Assume the hypotheses of Proposition~\ref{prop:moving-cutoff-decomposition} and the extra
regularity from Definition~\ref{def:model-c2}, and suppose moreover
that \(\chi\) is flat at the outer edge in the sense of Definition~\ref{def:endpoint-flat-cutoff}.
If there exist
\[
C_{\mathrm{pair},\alpha}^{\flat}<\infty,
\qquad
h_{0,\alpha}\in(0,t]
\]
such that for every \(0<h\le h_{0,\alpha}\) and every \(v\in\mathfrak G_t\),
\begin{equation}\label{eq:full-h-coupled-alpha-defect-bound}
\left\|
\one_{[0,1]}
\mathfrak F_{t,h,\alpha,R_{h,\alpha}}^{\mathrm{ex}}
\bigl(
\overline\chi\,\widehat\Phi_{v,h,\alpha},
\overline\chi\,\widehat\Psi_{v,h,\alpha}^{\mathrm{full}}
\bigr)
\right\|_{E}
\le
C_{\mathrm{pair},\alpha}^{\flat}
h^{\alpha\beta}\|v\|_{\mathfrak G_t},
\end{equation}
then the hypothesis of Corollary~\ref{cor:full-h-coupled-commutator} is satisfied. Consequently,
the conclusion of Corollary~\ref{cor:full-h-coupled-commutator} follows.
\end{corollary}

\begin{proof}
By Proposition~\ref{prop:full-h-coupled-alpha-defect},
\[
\mathfrak Q_{t,h,\alpha,\chi}^{\mathrm{pair},\alpha}
\bigl(
\widehat\Phi_{v,h,\alpha},
\widehat\Psi_{v,h,\alpha}^{\mathrm{old}}
\bigr)
=
\one_{[0,1]}
\mathfrak F_{t,h,\alpha,R_{h,\alpha}}^{\mathrm{ex}}
\bigl(
\overline\chi\,\widehat\Phi_{v,h,\alpha},
\overline\chi\,\widehat\Psi_{v,h,\alpha}^{\mathrm{full}}
\bigr)
\qquad\text{in }E.
\]
Proposition~\ref{prop:full-h-coupled-alpha} identifies the left-hand side with
\[
\frac{1}{\Gamma(\lambda)}
\mathfrak D_{-,1}^{\lambda}
\mathfrak Y_{v,h,\alpha}^{\mathrm{src}}
\qquad\text{in }\mathcal D'((0,1)),
\]
and Proposition~\ref{prop:full-h-coupled-commutator} gives
\[
\mathfrak Y_{v,h,\alpha}^{\mathrm{src}}
=
\varepsilon_{h,\alpha}^{H_{XY}}
\mathfrak l_{\varepsilon_{h,\alpha}}
\mathfrak C_{v,h,\alpha}^{\mathrm{pair}}
\qquad\text{in }E.
\]
Hence the assumed bound \eqref{eq:full-h-coupled-alpha-defect-bound} is exactly the hypothesis of
Corollary~\ref{cor:full-h-coupled-commutator}, which then yields the conclusion.
\end{proof}

\begin{definition}[Normalized growing-window endpoint traces]
\label{def:full-h-coupled-alpha-normalized}
Assume the hypotheses of Definition~\ref{def:full-h-coupled-alpha}. For \(v\in\mathfrak G_t\), set
\begin{equation}\label{eq:full-h-coupled-alpha-Z-full}
\widehat Z_{v,h,\alpha}^{\mathrm{full}}
:=
h^{-\alpha\beta}\widehat\Psi_{v,h,\alpha}^{\mathrm{full}}
\in
\widehat E_{h,\alpha}^{Y},
\end{equation}
and
\begin{equation}\label{eq:full-h-coupled-alpha-Z-old}
\widehat Z_{v,h,\alpha}^{\mathrm{old}}
:=
h^{-\alpha\beta}\widehat\Psi_{v,h,\alpha}^{\mathrm{old}}
\in
\widehat E_{h,\alpha}^{Y,\mathrm{old}}.
\end{equation}
\end{definition}

\begin{proposition}[Normalized growing-window exact pair]
\label{prop:full-h-coupled-alpha-normalized}
Assume the hypotheses of Definition~\ref{def:full-h-coupled-alpha-normalized}. Then for every
\[
v\in\mathfrak G_t,
\qquad
0<h\le t,
\qquad
h^\alpha\le t,
\]
one has
\begin{equation}\label{eq:full-h-coupled-alpha-normalized-exact}
\mathfrak T_{t,h,\alpha,R_{h,\alpha}}^{\mathrm{ex}}
\widehat\Phi_{v,h,\alpha}
=
h^{\alpha\beta}\widehat Z_{v,h,\alpha}^{\mathrm{full}},
\end{equation}
and
\begin{equation}\label{eq:full-h-coupled-alpha-normalized-bound}
\left\|
\widehat\Phi_{v,h,\alpha}
\right\|_{E_{R_{h,\alpha}}}
+
\frac{\eta}{\nu_Y}
h^{\alpha\beta}
\left\|
\widehat Z_{v,h,\alpha}^{\mathrm{full}}
\right\|_{E_{R_{h,\alpha}}}
\le
\left(1+\frac{\eta}{\nu_Y}\right)\|v\|_{\mathfrak G_t}.
\end{equation}
In particular,
\begin{equation}\label{eq:full-h-coupled-alpha-normalized-Z}
\left\|
\widehat Z_{v,h,\alpha}^{\mathrm{full}}
\right\|_{E_{R_{h,\alpha}}}
=
\|\mathcal S_t v\|_{H_t}
\le
\frac{\nu_Y}{\eta}\|v\|_{\mathfrak G_t}.
\end{equation}
\end{proposition}

\begin{proof}
Equation~\eqref{eq:full-h-coupled-alpha-normalized-exact} is
\eqref{eq:full-h-coupled-alpha-exact} together with
\eqref{eq:full-h-coupled-alpha-Z-full}. The first bound follows from
\eqref{eq:full-h-coupled-alpha-exact-bound} by taking square roots and using
\[
\sqrt{a^2+b^2}\le a+b
\qquad\text{for }a,b\ge0.
\]
Finally,
\eqref{eq:full-h-coupled-alpha-normalized-Z} is exactly the computation already used in the proof
of Proposition~\ref{prop:full-h-coupled-alpha-exact}:
\[
\left\|
\widehat Z_{v,h,\alpha}^{\mathrm{full}}
\right\|_{E_{R_{h,\alpha}}}
=
h^{-\alpha\beta}
\left\|
\widehat\Psi_{v,h,\alpha}^{\mathrm{full}}
\right\|_{E_{R_{h,\alpha}}}
=
\|\mathcal S_t v\|_{H_t}.
\]
\end{proof}

\begin{definition}[Normalized growing-window exact defect]
\label{def:full-h-coupled-alpha-defect-normalized}
Assume the hypotheses of Proposition~\ref{prop:full-h-coupled-alpha-defect}. For
\[
R\ge 1,
\qquad
\Phi\in H^\beta((0,R)),
\qquad
Z\in E_R,
\]
define
\begin{equation}\label{eq:full-h-coupled-alpha-defect-normalized}
\mathfrak E_{t,h,\alpha,R,\chi}^{\mathrm{norm}}(\Phi,Z)
:=
h^{-\alpha\beta}
\one_{[0,1]}
\mathfrak F_{t,h,\alpha,R}^{\mathrm{ex}}
\bigl(
\overline\chi\,\Phi,
h^{\alpha\beta}\overline\chi\,Z
\bigr)
\in E,
\end{equation}
where \(\overline\chi\) denotes the zero extension of \(\chi\) from \([0,1]\) to \([0,R]\).
\end{definition}

\begin{proposition}[Endpoint bound is a uniform normalized-defect estimate]
\label{prop:full-h-coupled-alpha-defect-normalized}
Assume the hypotheses of Proposition~\ref{prop:moving-cutoff-decomposition} and the extra
regularity from Definition~\ref{def:model-c2}, and suppose moreover
that \(\chi\) is flat at the outer edge in the sense of Definition~\ref{def:endpoint-flat-cutoff}.
Then \eqref{eq:full-h-coupled-alpha-defect-bound} is equivalent to the statement that there exist
\[
C_{\mathrm{pair},\alpha}^{\mathrm{norm}}<\infty,
\qquad
h_{0,\alpha}\in(0,t]
\]
such that for every \(0<h\le h_{0,\alpha}\) and every \(v\in\mathfrak G_t\),
\begin{equation}\label{eq:full-h-coupled-alpha-defect-normalized-bound}
\left\|
\mathfrak E_{t,h,\alpha,R_{h,\alpha},\chi}^{\mathrm{norm}}
\bigl(
\widehat\Phi_{v,h,\alpha},
\widehat Z_{v,h,\alpha}^{\mathrm{full}}
\bigr)
\right\|_{E}
\le
C_{\mathrm{pair},\alpha}^{\mathrm{norm}}
\|v\|_{\mathfrak G_t}.
\end{equation}
\end{proposition}

\begin{proof}
By Definition~\ref{def:full-h-coupled-alpha-defect-normalized} and
\eqref{eq:full-h-coupled-alpha-Z-full},
\[
\mathfrak E_{t,h,\alpha,R_{h,\alpha},\chi}^{\mathrm{norm}}
\bigl(
\widehat\Phi_{v,h,\alpha},
\widehat Z_{v,h,\alpha}^{\mathrm{full}}
\bigr)
=
h^{-\alpha\beta}
\one_{[0,1]}
\mathfrak F_{t,h,\alpha,R_{h,\alpha}}^{\mathrm{ex}}
\bigl(
\overline\chi\,\widehat\Phi_{v,h,\alpha},
\overline\chi\,\widehat\Psi_{v,h,\alpha}^{\mathrm{full}}
\bigr).
\]
Therefore \eqref{eq:full-h-coupled-alpha-defect-bound} and
\eqref{eq:full-h-coupled-alpha-defect-normalized-bound} are the same statement with a different
normalization of the left-hand side.
\end{proof}

\begin{definition}[Normalized frozen growing-window defect]
\label{def:full-h-coupled-alpha-frozen}
Assume
\[
0<H_{XY}<H_Y<\tfrac12,
\qquad
\beta:=H_Y-H_{XY},
\]
fix \(t\in(0,T)\), fix \(\alpha\in(0,1)\), let \(0<h\le t\), let \(R\ge1\), and let
\(\chi\in C^\infty([0,1])\). Extend \(\chi\) by zero to \([0,R]\), still denoted \(\overline\chi\).
For
\[
\Phi\in H^\beta((0,R)),
\qquad
Z\in E_R,
\]
define the normalized frozen localized defect by
\begin{equation}\label{eq:full-h-coupled-alpha-frozen}
\mathfrak C_{t,h,\alpha,R,\chi}^{\mathrm{fr}}(\Phi,Z)
:=
\one_{[0,1]}L_Y(t,t)
\Bigl(
\overline\chi\,Z
-
h^{-\alpha\beta}\mathfrak T_{t,R}^{\mathrm{hl}}(\overline\chi\,\Phi)
\Bigr)
\in E.
\end{equation}
\end{definition}

\begin{proposition}[Normalized frozen localized defect is a commutator]
\label{prop:full-h-coupled-alpha-frozen}
Assume the hypotheses of Definition~\ref{def:full-h-coupled-alpha-frozen}. Then
if
\[
\mathfrak T_{t,R}^{\mathrm{hl}}\Phi=h^{\alpha\beta}Z,
\]
then
\begin{equation}\label{eq:full-h-coupled-alpha-frozen-comm}
\mathfrak C_{t,h,\alpha,R,\chi}^{\mathrm{fr}}(\Phi,Z)
=
\one_{[0,1]}L_Y(t,t)\,
h^{-\alpha\beta}
\bigl[
\overline\chi,\mathfrak T_{t,R}^{\mathrm{hl}}
\bigr]\Phi
\qquad\text{in }E.
\end{equation}
Equivalently, if one defines the normalized frozen defect by
\[
\mathfrak F_{t,h,\alpha,R}^{\mathrm{hl}}(\Phi,\Psi)
:=
L_Y(t,t)\bigl(\Psi-h^{-\alpha\beta}\mathfrak T_{t,R}^{\mathrm{hl}}\Phi\bigr),
\qquad
\Phi\in H^\beta((0,R)),\ \Psi\in E_R,
\]
then
\begin{equation}\label{eq:full-h-coupled-alpha-frozen-defect}
\one_{[0,1]}
\mathfrak F_{t,h,\alpha,R}^{\mathrm{hl}}
\bigl(
\overline\chi\,\Phi,
\overline\chi\,Z
\bigr)
=
\mathfrak C_{t,h,\alpha,R,\chi}^{\mathrm{fr}}(\Phi,Z).
\end{equation}
\end{proposition}

\begin{proof}
Since \(\mathfrak T_{t,R}^{\mathrm{hl}}\Phi=h^{\alpha\beta}Z\), one has
\[
\overline\chi\,Z-h^{-\alpha\beta}\mathfrak T_{t,R}^{\mathrm{hl}}(\overline\chi\,\Phi)
=
h^{-\alpha\beta}
\overline\chi\,\mathfrak T_{t,R}^{\mathrm{hl}}\Phi
-
h^{-\alpha\beta}\mathfrak T_{t,R}^{\mathrm{hl}}(\overline\chi\,\Phi)
=
\,
h^{-\alpha\beta}
\bigl[
\overline\chi,\mathfrak T_{t,R}^{\mathrm{hl}}
\bigr]\Phi.
\]
Multiplying by \(L_Y(t,t)\) and restricting to \([0,1]\) gives
\eqref{eq:full-h-coupled-alpha-frozen-comm}. The identity
\eqref{eq:full-h-coupled-alpha-frozen-defect} is simply the definition of
\(\mathfrak F_{t,h,\alpha,R}^{\mathrm{hl}}\).
\end{proof}

\begin{definition}[Finite-window right fractional integral]
\label{def:window-fractional-integral}
For \(R>0\), \(\lambda\in(0,1)\), and \(f\in E_R=L^2([0,R])\), define
\begin{equation}\label{eq:window-fractional-integral}
(\mathfrak I_{-,R}^{\lambda}f)(r)
:=
\frac{1}{\Gamma(\lambda)}
\int_r^R (q-r)^{\lambda-1}f(q)\,dq,
\qquad 0\le r\le R.
\end{equation}
\end{definition}

\begin{lemma}[Finite-window right fractional integral inverts the frozen transfer]
\label{lem:window-frozen-inverse}
Assume
\[
0<H_{XY}<H_Y<\tfrac12,
\qquad
\beta:=H_Y-H_{XY},
\]
fix \(t\in(0,T)\), and let \(R>0\). Then:
\begin{enumerate}[label=(\roman*)]
\item for every \(f\in E_R\),
\begin{equation}\label{eq:window-fractional-left-inverse}
\mathfrak D_{\beta,R}^{\mathrm{dist}}
\bigl(
\mathfrak I_{-,R}^{\beta}f
\bigr)
=
f
\qquad\text{in }E_R;
\end{equation}
\item for every \(f\in E_R\), one has
\[
\mathfrak I_{-,R}^{\beta}f\in H^\beta((0,R));
\]
\item for every \(f\in E_R\),
\begin{equation}\label{eq:window-frozen-inverse}
\mathfrak T_{t,R}^{\mathrm{hl}}
\bigl(
\mathfrak I_{-,R}^{\beta}f
\bigr)
=
\mathbf c_t f
\qquad\text{in }E_R,
\end{equation}
where \(\mathbf c_t\) is defined in \eqref{eq:model-coeff};
\item consequently, if \(\Phi\in H^\beta((0,R))\) and \(Z\in E_R\) satisfy
\begin{equation}\label{eq:window-frozen-exact}
\mathfrak T_{t,R}^{\mathrm{hl}}\Phi=h^{\alpha\beta}Z,
\end{equation}
then
\begin{equation}\label{eq:window-frozen-exact-inverse}
\Phi
=
\mathbf c_t^{-1}h^{\alpha\beta}
\mathfrak I_{-,R}^{\beta}Z
\qquad\text{in }H^\beta((0,R)).
\end{equation}
\end{enumerate}
\end{lemma}

\begin{proof}
For \(\lambda,\mu\in(0,1)\) with \(\lambda+\mu\le1\), the same Fubini and Beta-function
computation as in Lemma~\ref{lem:unit-semigroup} gives the finite-window semigroup identity
\begin{equation}\label{eq:window-fractional-semigroup}
\mathfrak I_{-,R}^{\lambda}\mathfrak I_{-,R}^{\mu}
=
\mathfrak I_{-,R}^{\lambda+\mu}
\qquad\text{on }E_R.
\end{equation}
Applying \eqref{eq:window-fractional-semigroup} with \(\lambda=1-\beta\) and \(\mu=\beta\) yields
\[
\mathfrak I_{-,R}^{1-\beta}\mathfrak I_{-,R}^{\beta}f
=
\mathfrak I_{-,R}^{1}f,
\qquad
(\mathfrak I_{-,R}^{1}f)(r)=\int_r^R f(q)\,dq.
\]
Since
\[
\mathfrak D_{\beta,R}^{\mathrm{dist}}g
=
-\frac{d}{dr}\mathfrak I_{-,R}^{1-\beta}g
\]
by definition, differentiating the previous identity gives
\eqref{eq:window-fractional-left-inverse}. It also follows that
\(\mathfrak D_{\beta,R}^{\mathrm{dist}}(\mathfrak I_{-,R}^{\beta}f)\in E_R\), so
Lemma~\ref{lem:window-distribution}(i) gives
\(\mathfrak I_{-,R}^{\beta}f\in H^\beta((0,R))\), proving (ii).

Now combine \eqref{eq:window-fractional-left-inverse} with
Lemma~\ref{lem:window-distribution}(ii):
\[
\mathfrak T_{t,R}^{\mathrm{hl}}
\bigl(
\mathfrak I_{-,R}^{\beta}f
\bigr)
=
\frac{\mathbf c_t}{\Gamma(1-\beta)}
\mathfrak D_{\beta,R}^{\mathrm{dist}}
\bigl(
\mathfrak I_{-,R}^{\beta}f
\bigr)
=
\mathbf c_t f.
\]
This proves \eqref{eq:window-frozen-inverse}. Finally, if
\(\mathfrak T_{t,R}^{\mathrm{hl}}\Phi=h^{\alpha\beta}Z\), then
\eqref{eq:window-frozen-inverse} shows that
\[
\mathfrak T_{t,R}^{\mathrm{hl}}
\Bigl(
\Phi-\mathbf c_t^{-1}h^{\alpha\beta}\mathfrak I_{-,R}^{\beta}Z
\Bigr)
=
0.
\]
By Definition~\ref{def:micro-window-model} and the injectivity of \(\mathfrak J_{H_{XY},R}\)
proved in \eqref{eq:window-scaling-J}, \(\mathfrak T_{t,R}^{\mathrm{hl}}\) is injective, so
\eqref{eq:window-frozen-exact-inverse} follows.
\end{proof}

\begin{theorem}[Uniform growing-window bound for the normalized frozen defect]
\label{thm:full-h-coupled-alpha-frozen-bound}
Assume the hypotheses of Definition~\ref{def:full-h-coupled-alpha-frozen}, and suppose moreover
that \(\chi\) is flat at the outer edge in the sense of Definition~\ref{def:endpoint-flat-cutoff}.
Then there exists \(C_{t,\chi}<\infty\) such that for every
\[
0<h\le t,
\qquad
R\ge1,
\qquad
\Phi\in H^\beta((0,R)),
\qquad
Z\in E_R
\]
with
\[
\mathfrak T_{t,R}^{\mathrm{hl}}\Phi=h^{\alpha\beta}Z,
\]
one has
\begin{equation}\label{eq:full-h-coupled-alpha-frozen-bound}
\left\|
\mathfrak C_{t,h,\alpha,R,\chi}^{\mathrm{fr}}(\Phi,Z)
\right\|_{E}
\le
C_{t,\chi}\|Z\|_{E_R}.
\end{equation}
For almost every \(r\in[0,1]\),
\begin{equation}\label{eq:full-h-coupled-alpha-frozen-kernel}
\bigl(
\mathfrak C_{t,h,\alpha,R,\chi}^{\mathrm{fr}}(\Phi,Z)
\bigr)(r)
=
L_Y(t,t)\int_r^R
K_{\beta,\chi,R}^{\mathrm{fr}}(r,q)Z(q)\,dq,
\end{equation}
where
\begin{equation}\label{eq:full-h-coupled-alpha-frozen-kernel-def}
K_{\beta,\chi,R}^{\mathrm{fr}}(r,q)
:=
\frac{1}{\Gamma(1-\beta)\Gamma(\beta)}
\int_0^1
\theta^{-\beta}(1-\theta)^\beta
\overline\chi'(r+\theta(q-r))\,d\theta,
\qquad
0\le r\le1,\ r\le q\le R,
\end{equation}
and
\begin{align}
\left|
K_{\beta,\chi,R}^{\mathrm{fr}}(r,q)
\right|
&\le
C_{\beta,\chi},
\qquad 0\le r\le q\le 1,
\label{eq:full-h-coupled-alpha-frozen-kernel-local}\\
\left|
K_{\beta,\chi,R}^{\mathrm{fr}}(r,q)
\right|
&\le
C_{\beta,\chi}
\left(
\frac{1-r}{q-r}
\right)^{1-\beta},
\qquad 0\le r\le1,\ 1\le q\le R.
\label{eq:full-h-coupled-alpha-frozen-kernel-tail}
\end{align}
\end{theorem}

\begin{proof}
Because \(\chi\) is flat at the outer edge and vanishes identically on \((1,R]\), its zero
extension \(\overline\chi\) belongs to \(C^\infty([0,R])\). By
Lemma~\ref{lem:window-frozen-inverse}(iv),
\[
\Phi
=
\mathbf c_t^{-1}h^{\alpha\beta}\mathfrak I_{-,R}^{\beta}Z.
\]
Hence, using Lemma~\ref{lem:window-distribution}(ii),
\begin{align}
\mathfrak C_{t,h,\alpha,R,\chi}^{\mathrm{fr}}(\Phi,Z)
&=
\one_{[0,1]}L_Y(t,t)
\Bigl(
\overline\chi\,Z
-
\frac{\mathbf c_t}{\Gamma(1-\beta)}
\mathfrak D_{\beta,R}^{\mathrm{dist}}
\bigl(
\overline\chi\,\mathbf c_t^{-1}\mathfrak I_{-,R}^{\beta}Z
\bigr)
\Bigr)
\notag\\
&=
\one_{[0,1]}L_Y(t,t)
\Bigl(
\overline\chi\,Z
-
\frac{1}{\Gamma(1-\beta)}
\mathfrak D_{\beta,R}^{\mathrm{dist}}
\bigl(
\overline\chi\,\mathfrak I_{-,R}^{\beta}Z
\bigr)
\Bigr).
\label{eq:full-h-coupled-alpha-frozen-proof-start}
\end{align}

Fix \(r\in[0,1]\). Since \(\overline\chi\) is supported in \([0,1]\), Fubini's theorem yields
\begin{align*}
\bigl(
\mathfrak I_{-,R}^{1-\beta}
(\overline\chi\,\mathfrak I_{-,R}^{\beta}Z)
\bigr)(r)
&=
\frac{1}{\Gamma(1-\beta)\Gamma(\beta)}
\int_r^1 \int_s^R
(s-r)^{-\beta}\overline\chi(s)(q-s)^{\beta-1}Z(q)\,dq\,ds \\
&=
\int_r^R A_{\beta,\chi,R}(r,q)Z(q)\,dq,
\end{align*}
where
\[
A_{\beta,\chi,R}(r,q)
:=
\frac{1}{\Gamma(1-\beta)\Gamma(\beta)}
\int_0^1
\theta^{-\beta}(1-\theta)^{\beta-1}
\overline\chi(r+\theta(q-r))\,d\theta.
\]
Because \(\overline\chi\in C^\infty([0,R])\), the map \(A_{\beta,\chi,R}\) is \(C^1\) on
\(\{0\le r\le1,\ r\le q\le R\}\), and
\[
A_{\beta,\chi,R}(r,r)=\overline\chi(r).
\]
Differentiating under the integral sign gives
\[
\partial_r A_{\beta,\chi,R}(r,q)
=
\frac{1}{\Gamma(1-\beta)\Gamma(\beta)}
\int_0^1
\theta^{-\beta}(1-\theta)^\beta
\overline\chi'(r+\theta(q-r))\,d\theta
=
K_{\beta,\chi,R}^{\mathrm{fr}}(r,q).
\]
Therefore
\begin{align*}
\frac{1}{\Gamma(1-\beta)}
\mathfrak D_{\beta,R}^{\mathrm{dist}}
\bigl(
\overline\chi\,\mathfrak I_{-,R}^{\beta}Z
\bigr)(r)
&=
-\frac{d}{dr}
\int_r^R A_{\beta,\chi,R}(r,q)Z(q)\,dq \\
&=
\overline\chi(r)Z(r)-\int_r^R
K_{\beta,\chi,R}^{\mathrm{fr}}(r,q)Z(q)\,dq.
\end{align*}
Substituting this identity into \eqref{eq:full-h-coupled-alpha-frozen-proof-start} proves
\eqref{eq:full-h-coupled-alpha-frozen-kernel}.

Now fix \(0\le r\le q\le1\). Since \(\overline\chi'\) is bounded on \([0,1]\),
\eqref{eq:full-h-coupled-alpha-frozen-kernel-def} gives
\eqref{eq:full-h-coupled-alpha-frozen-kernel-local}. Next fix \(0\le r\le1\le q\le R\). Because
\(\overline\chi'=0\) on \([1,R]\), the integrand in
\eqref{eq:full-h-coupled-alpha-frozen-kernel-def} vanishes unless
\[
r+\theta(q-r)\le 1
\qquad\Longleftrightarrow\qquad
\theta\le \frac{1-r}{q-r}.
\]
Hence
\[
\left|
K_{\beta,\chi,R}^{\mathrm{fr}}(r,q)
\right|
\le
\frac{\|\overline\chi'\|_{L^\infty([0,1])}}{\Gamma(1-\beta)\Gamma(\beta)}
\int_0^{\frac{1-r}{q-r}}
\theta^{-\beta}(1-\theta)^\beta\,d\theta
\le
C_{\beta,\chi}
\left(
\frac{1-r}{q-r}
\right)^{1-\beta},
\]
which proves \eqref{eq:full-h-coupled-alpha-frozen-kernel-tail}.

Split the operator in \eqref{eq:full-h-coupled-alpha-frozen-kernel} into the local part
\(q\in[r,1]\) and the tail part \(q\in[1,R]\). By
\eqref{eq:full-h-coupled-alpha-frozen-kernel-local}, the local kernel is bounded on the finite
triangle \(\{0\le r\le q\le1\}\), hence Hilbert--Schmidt from \(L^2([0,1])\) to \(E\). For the
tail part, \eqref{eq:full-h-coupled-alpha-frozen-kernel-tail} gives
\[
\int_0^1\int_1^R
\left|
K_{\beta,\chi,R}^{\mathrm{fr}}(r,q)
\right|^2\,dq\,dr
\le
C_{\beta,\chi}
\int_0^1\int_1^\infty
\left(
\frac{1-r}{q-r}
\right)^{2(1-\beta)}
dq\,dr.
\]
Because \(\beta<\tfrac12\), the last integral is finite:
\[
\int_1^\infty (q-r)^{-2(1-\beta)}\,dq
=
\frac{(1-r)^{2\beta-1}}{1-2\beta},
\]
so
\[
\int_0^1\int_1^\infty
\left(
\frac{1-r}{q-r}
\right)^{2(1-\beta)}
dq\,dr
=
\frac{1}{1-2\beta}\int_0^1 (1-r)\,dr
<
\infty.
\]
Thus the tail kernel is also Hilbert--Schmidt, with norm bounded independently of \(R\ge1\). The
two Hilbert--Schmidt bounds combine to give
\[
\left\|
\mathfrak C_{t,h,\alpha,R,\chi}^{\mathrm{fr}}(\Phi,Z)
\right\|_{E}
\le
C_{t,\chi}\|Z\|_{E_R},
\]
which is \eqref{eq:full-h-coupled-alpha-frozen-bound}.
\end{proof}

\begin{corollary}[Normalized growing-window exact pairs satisfy the frozen bound]
\label{cor:full-h-coupled-alpha-frozen-bound}
Assume the hypotheses of Proposition~\ref{prop:full-h-coupled-alpha-normalized}, and suppose
moreover that \(\chi\) is flat at the outer edge in the sense of
Definition~\ref{def:endpoint-flat-cutoff}. Then there exists \(C_{t,\chi}<\infty\) such that for
every \(0<h\le t\) with \(h^\alpha\le t\) and every \(v\in\mathfrak G_t\),
\begin{equation}\label{eq:full-h-coupled-alpha-frozen-bound-v}
\left\|
\mathfrak C_{t,h,\alpha,R_{h,\alpha},\chi}^{\mathrm{fr}}
\bigl(
\widehat\Phi_{v,h,\alpha},
\widehat Z_{v,h,\alpha}^{\mathrm{full}}
\bigr)
\right\|_{E}
\le
C_{t,\chi}\|v\|_{\mathfrak G_t}.
\end{equation}
\end{corollary}

\begin{proof}
Apply Theorem~\ref{thm:full-h-coupled-alpha-frozen-bound} with
\[
R=R_{h,\alpha},
\qquad
\Phi=\widehat\Phi_{v,h,\alpha},
\qquad
Z=\widehat Z_{v,h,\alpha}^{\mathrm{full}}.
\]
The exact relation
\[
\mathfrak T_{t,R_{h,\alpha}}^{\mathrm{hl}}\widehat\Phi_{v,h,\alpha}
=
h^{\alpha\beta}\widehat Z_{v,h,\alpha}^{\mathrm{full}}
\]
is Proposition~\ref{prop:full-h-coupled-alpha-normalized}. The bound
\eqref{eq:full-h-coupled-alpha-normalized-Z} then gives
\[
\left\|
\widehat Z_{v,h,\alpha}^{\mathrm{full}}
\right\|_{E_{R_{h,\alpha}}}
\le
\frac{\nu_Y}{\eta}\|v\|_{\mathfrak G_t},
\]
which proves \eqref{eq:full-h-coupled-alpha-frozen-bound-v}.
\end{proof}

\begin{lemma}[Uniform head-window bound for the finite-window fractional inverse]
\label{lem:window-fractional-head}
Assume
\[
0<\beta<\tfrac12.
\]
Then there exists \(C_\beta<\infty\) such that for every \(R\ge1\) and every \(f\in E_R\),
\begin{equation}\label{eq:window-fractional-head}
\left\|
\one_{[0,1]}\mathfrak I_{-,R}^{\beta}f
\right\|_{E}
\le
C_\beta\|f\|_{E_R}.
\end{equation}
\end{lemma}

\begin{proof}
For \(0\le r\le1\), split
\[
\bigl(
\mathfrak I_{-,R}^{\beta}f
\bigr)(r)
=
\frac{1}{\Gamma(\beta)}
\int_r^1 (q-r)^{\beta-1}f(q)\,dq
+
\frac{1}{\Gamma(\beta)}
\int_1^R (q-r)^{\beta-1}f(q)\,dq
=:
(\mathfrak L_{\beta}f)(r)+(\mathfrak T_{\beta,R}f)(r).
\]
For the local part, extend \(f|_{[0,1]}\) by zero to \(\R\) and set
\[
k_\beta(s):=\Gamma(\beta)^{-1}s^{\beta-1}\one_{(0,1)}(s).
\]
Since \(k_\beta\in L^1(\R)\), Young's inequality gives
\[
\|\mathfrak L_\beta f\|_{E}
\le
\|k_\beta\|_{L^1(\R)}\|f\|_{L^2([0,1])}
\le
C_\beta\|f\|_{E_R}.
\]

For the tail part, the kernel
\[
K_{\beta,R}^{\mathrm{tail}}(r,q)
:=
\Gamma(\beta)^{-1}(q-r)^{\beta-1}
\one_{\{0\le r\le1,\ 1\le q\le R\}}
\]
is Hilbert--Schmidt uniformly in \(R\ge1\), because
\begin{align*}
\int_0^1\int_1^R
\left|
K_{\beta,R}^{\mathrm{tail}}(r,q)
\right|^2\,dq\,dr
&\le
C_\beta
\int_0^1\int_1^\infty
(q-r)^{2\beta-2}\,dq\,dr \\
&=
\frac{C_\beta}{1-2\beta}
\int_0^1 (1-r)^{2\beta-1}\,dr
<
\infty,
\end{align*}
since \(\beta<\tfrac12\). Hence
\[
\|\mathfrak T_{\beta,R}f\|_{E}
\le
C_\beta\|f\|_{E_R}
\]
with \(C_\beta\) independent of \(R\). Combining the local and tail estimates proves
\eqref{eq:window-fractional-head}.
\end{proof}

\begin{lemma}[Head-window compatibility for cutoff-supported inputs]
\label{lem:window-head-compat}
Assume the hypotheses of Theorem~\ref{thm:micro-window-regularity}, let \(R\ge1\), and let
\(\chi\in C^\infty([0,1])\) with zero extension \(\overline\chi\) to \([0,R]\). For
\[
\Phi\in H^\beta((0,R)),
\qquad
Z\in E_R,
\]
set
\[
\phi:=\Phi|_{[0,1]},
\qquad
z:=Z|_{[0,1]}.
\]
Then on \(E=L^2([0,1])\),
\begin{align}
\one_{[0,1]}\mathfrak B_{Y,t,h,\alpha,R}^{\mathrm{mic}}(\overline\chi Z)
&=
\mathfrak B_{Y,t,h,\alpha,1}^{\mathrm{mic}}(\chi z),
\label{eq:window-head-compat-BY}\\
\one_{[0,1]}\mathfrak T_{t,R}^{\mathrm{hl}}(\overline\chi\Phi)
&=
\mathfrak T_{t,1}^{\mathrm{hl}}(\chi\phi),
\label{eq:window-head-compat-T}\\
\one_{[0,1]}\mathfrak M_{t,h,\alpha,R}^{\mathrm{win}}
\mathfrak T_{t,R}^{\mathrm{hl}}(\overline\chi\Phi)
&=
\mathfrak M_{t,h,\alpha,1}^{\mathrm{win}}
\mathfrak T_{t,1}^{\mathrm{hl}}(\chi\phi),
\label{eq:window-head-compat-M}\\
\one_{[0,1]}\mathfrak K_{t,h,\alpha,R}^{\mathrm{win}}(\overline\chi\Phi)
&=
\mathfrak K_{t,h,\alpha,1}^{\mathrm{win}}(\chi\phi).
\label{eq:window-head-compat-K}
\end{align}
Consequently,
\begin{align}
\one_{[0,1]}
\mathfrak F_{t,h,\alpha,R}^{\mathrm{ex}}
\bigl(
\overline\chi\,\Phi,
h^{\alpha\beta}\overline\chi\,Z
\bigr)
&=
\mathfrak B_{Y,t,h,\alpha,1}^{\mathrm{mic}}(\chi z)
\notag\\
&\quad-
L_Y(t,t)\mathfrak M_{t,h,\alpha,1}^{\mathrm{win}}
\mathfrak T_{t,1}^{\mathrm{hl}}(h^{\alpha\beta}\chi\phi)
\notag\\
&\quad-
\mathfrak K_{t,h,\alpha,1}^{\mathrm{win}}(\chi\phi).
\label{eq:window-head-compat-defect}
\end{align}
\end{lemma}

\begin{proof}
Since \(\overline\chi Z\) is supported in \([0,1]\), the explicit Volterra formula from the proof of
Proposition~\ref{prop:BY-micro} shows that for \(0\le r\le1\),
\[
\bigl(
\mathfrak B_{Y,t,h,\alpha,R}^{\mathrm{mic}}(\overline\chi Z)
\bigr)(r)
\]
depends only on the values of \(\overline\chi Z\) on \([r,1]\), with the same kernel
\(L_Y(t-hr,t-hq)\) as in the window \(R=1\). This proves
\eqref{eq:window-head-compat-BY}.

Likewise, because \(\overline\chi\Phi\) is supported in \([0,1]\), one has for \(0\le r\le1\)
\[
\mathfrak J_{\beta,R}(\overline\chi\Phi)(r)
=
\int_r^R (q-r)^{-\beta}\overline\chi(q)\Phi(q)\,dq
=
\int_r^1 (q-r)^{-\beta}\chi(q)\phi(q)\,dq
=
\mathfrak J_{\beta,1}(\chi\phi)(r).
\]
Taking the distributional derivative on \((0,1)\) and using
Lemma~\ref{lem:window-distribution}(ii) gives \eqref{eq:window-head-compat-T}.

For \eqref{eq:window-head-compat-M}, note first that \eqref{eq:window-head-compat-T} identifies the
two inputs on \([0,1]\). Moreover, the definition
\[
m_{t,h,\alpha,R}(r)=\frac{g_{t,h,R}^{\mathrm{win}}(r,r)}{g_t^{\mathrm{hl}}}
\]
and the formula defining \(g_{t,h,R}^{\mathrm{win}}\) in
Proposition~\ref{prop:micro-window-decompose} imply
\[
g_{t,h,R}^{\mathrm{win}}(r,q)=g_{t,h,1}^{\mathrm{win}}(r,q)
\qquad
\text{for }0\le r\le q\le1,
\]
hence \(m_{t,h,\alpha,R}(r)=m_{t,h,\alpha,1}(r)\) for \(0\le r\le1\). This is
\eqref{eq:window-head-compat-M}.

For \eqref{eq:window-head-compat-K}, use the explicit formula from
Proposition~\ref{prop:micro-window-decompose}. On \([0,1]\), both terms defining
\(\mathfrak K_{t,h,\alpha,R}^{\mathrm{win}}(\overline\chi\Phi)\) integrate only over \(q\in[r,1]\),
because \(\overline\chi\Phi\) vanishes on \((1,R]\). Since
\[
g_{t,h,R}^{\mathrm{win}}=g_{t,h,1}^{\mathrm{win}},
\qquad
b_{t,h,\alpha,R}=b_{t,h,\alpha,1}
\qquad
\text{on }\Delta_1,
\]
the two remainder formulas agree there, which proves \eqref{eq:window-head-compat-K}.

Finally, \eqref{eq:window-head-compat-defect} is the definition \eqref{eq:window-defect}
combined with \eqref{eq:window-head-compat-BY}--\eqref{eq:window-head-compat-K}.
\end{proof}

\begin{theorem}[Exact-to-frozen perturbation on normalized growing-window exact pairs]
\label{thm:full-h-coupled-alpha-exact-to-frozen}
Assume the hypotheses of Proposition~\ref{prop:moving-cutoff-decomposition} and the extra
regularity from Definition~\ref{def:model-c2}, and suppose moreover
that \(\chi\) is flat at the outer edge in the sense of Definition~\ref{def:endpoint-flat-cutoff}.
Then there exist \(h_{0,\alpha}\in(0,t]\) and \(C_{t,\alpha,\chi}<\infty\) such that for every
\[
0<h\le h_{0,\alpha}
\qquad\text{with}\qquad
h^\alpha\le t,
\]
and every \(v\in\mathfrak G_t\),
\begin{equation}\label{eq:full-h-coupled-alpha-exact-to-frozen}
\left\|
\mathfrak E_{t,h,\alpha,R_{h,\alpha},\chi}^{\mathrm{norm}}
\bigl(
\widehat\Phi_{v,h,\alpha},
\widehat Z_{v,h,\alpha}^{\mathrm{full}}
\bigr)
-
\mathfrak C_{t,h,\alpha,R_{h,\alpha},\chi}^{\mathrm{fr}}
\bigl(
\widehat\Phi_{v,h,\alpha},
\widehat Z_{v,h,\alpha}^{\mathrm{full}}
\bigr)
\right\|_{E}
\le
C_{t,\alpha,\chi}h\|v\|_{\mathfrak G_t}.
\end{equation}
\end{theorem}

\begin{proof}
Shrink \(h_{0,\alpha}\), depending only on \(t\), \(\alpha\), \(\chi\), and the \(R=1\)
regularity thresholds entering Theorem~\ref{thm:micro-window-regularity}, if necessary so that for
every \(0<h\le h_{0,\alpha}\),
\[
h\le1,
\qquad
h^\alpha\le t,
\qquad
1<\varepsilon_{h,\alpha}^{-1},
\]
and Theorem~\ref{thm:micro-window-regularity} applies on the fixed window \(R=1\). Fix such an
\(h\), let
\[
R:=R_{h,\alpha},
\qquad
\Phi:=\widehat\Phi_{v,h,\alpha},
\qquad
Z:=\widehat Z_{v,h,\alpha}^{\mathrm{full}},
\]
and write
\[
\phi:=\Phi|_{[0,1]},
\qquad
z:=Z|_{[0,1]}.
\]
By Lemma~\ref{lem:window-head-compat},
\eqref{eq:window-defect}, \eqref{eq:full-h-coupled-alpha-defect-normalized}, and
\eqref{eq:full-h-coupled-alpha-frozen} reduce on \([0,1]\) to
\begin{align}
\mathfrak E_{t,h,\alpha,R,\chi}^{\mathrm{norm}}(\Phi,Z)
&=
\mathfrak B_{Y,t,h,\alpha,1}^{\mathrm{mic}}(\chi z)
-
L_Y(t,t)\mathfrak M_{t,h,\alpha,1}^{\mathrm{win}}
h^{-\alpha\beta}\mathfrak T_{t,1}^{\mathrm{hl}}(\chi\phi)
-
h^{-\alpha\beta}\mathfrak K_{t,h,\alpha,1}^{\mathrm{win}}(\chi\phi),
\label{eq:full-h-coupled-alpha-exact-head}\\
\mathfrak C_{t,h,\alpha,R,\chi}^{\mathrm{fr}}(\Phi,Z)
&=
L_Y(t,t)
\Bigl(
\chi z
-
h^{-\alpha\beta}\mathfrak T_{t,1}^{\mathrm{hl}}(\chi\phi)
\Bigr).
\label{eq:full-h-coupled-alpha-frozen-head}
\end{align}
Subtracting \eqref{eq:full-h-coupled-alpha-frozen-head} from
\eqref{eq:full-h-coupled-alpha-exact-head} gives
\begin{align}
\mathfrak E_{t,h,\alpha,R,\chi}^{\mathrm{norm}}(\Phi,Z)
-
\mathfrak C_{t,h,\alpha,R,\chi}^{\mathrm{fr}}(\Phi,Z)
&=
\bigl(
\mathfrak B_{Y,t,h,\alpha,1}^{\mathrm{mic}}
-
L_Y(t,t)I_E
\bigr)(\chi z)
\notag\\
&\quad-
L_Y(t,t)
\bigl(
\mathfrak M_{t,h,\alpha,1}^{\mathrm{win}}
-
I_E
\bigr)
h^{-\alpha\beta}\mathfrak T_{t,1}^{\mathrm{hl}}(\chi\phi)
\notag\\
&\quad-
h^{-\alpha\beta}\mathfrak K_{t,h,\alpha,1}^{\mathrm{win}}(\chi\phi).
\label{eq:full-h-coupled-alpha-exact-minus-frozen}
\end{align}

For the first term, \eqref{eq:BY-micro-close} with \(R=1\) gives
\begin{equation}\label{eq:full-h-coupled-alpha-exact-minus-frozen-1}
\left\|
\bigl(
\mathfrak B_{Y,t,h,\alpha,1}^{\mathrm{mic}}
-
L_Y(t,t)I_E
\bigr)(\chi z)
\right\|_E
\le
C_{t,\chi}h\|z\|_E.
\end{equation}

For the second term, \eqref{eq:full-h-coupled-alpha-frozen-head} implies
\[
h^{-\alpha\beta}\mathfrak T_{t,1}^{\mathrm{hl}}(\chi\phi)
=
\chi z
-
L_Y(t,t)^{-1}
\mathfrak C_{t,h,\alpha,R,\chi}^{\mathrm{fr}}(\Phi,Z),
\]
so by \eqref{eq:window-mult-close} with \(R=1\) and
Corollary~\ref{cor:full-h-coupled-alpha-frozen-bound},
\begin{align}
\left\|
L_Y(t,t)
\bigl(
\mathfrak M_{t,h,\alpha,1}^{\mathrm{win}}
-
I_E
\bigr)
h^{-\alpha\beta}\mathfrak T_{t,1}^{\mathrm{hl}}(\chi\phi)
\right\|_E
&\le
C_{t,\chi}h
\left(
\|z\|_E
+
\left\|
\mathfrak C_{t,h,\alpha,R,\chi}^{\mathrm{fr}}(\Phi,Z)
\right\|_E
\right)\notag\\
&\le
C_{t,\alpha,\chi}h\|v\|_{\mathfrak G_t}.
\label{eq:full-h-coupled-alpha-exact-minus-frozen-2}
\end{align}

For the third term, Proposition~\ref{prop:full-h-coupled-alpha-normalized} and
Lemma~\ref{lem:window-frozen-inverse}(iv) give
\[
\Phi
=
\mathbf c_t^{-1}h^{\alpha\beta}\mathfrak I_{-,R}^{\beta}Z
\qquad\text{in }H^\beta((0,R)).
\]
Hence
\[
h^{-\alpha\beta}\chi\phi
=
\mathbf c_t^{-1}\chi\,
\one_{[0,1]}
\mathfrak I_{-,R}^{\beta}Z.
\]
Using \eqref{eq:window-remainder-bound} with \(R=1\) and
Lemma~\ref{lem:window-fractional-head},
\begin{align}
\left\|
h^{-\alpha\beta}\mathfrak K_{t,h,\alpha,1}^{\mathrm{win}}(\chi\phi)
\right\|_E
&\le
C_t h
\left\|
h^{-\alpha\beta}\chi\phi
\right\|_E
\notag\\
&\le
C_{t,\chi}h
\left\|
\one_{[0,1]}\mathfrak I_{-,R}^{\beta}Z
\right\|_E
\notag\\
&\le
C_{t,\alpha,\chi}h\|Z\|_{E_R}
\le
C_{t,\alpha,\chi}h\|v\|_{\mathfrak G_t}.
\label{eq:full-h-coupled-alpha-exact-minus-frozen-3}
\end{align}

Finally, \eqref{eq:full-h-coupled-alpha-normalized-Z} gives
\[
\|z\|_E\le \|Z\|_{E_R}\le \frac{\nu_Y}{\eta}\|v\|_{\mathfrak G_t}.
\]
Combining this with
\eqref{eq:full-h-coupled-alpha-exact-minus-frozen},
\eqref{eq:full-h-coupled-alpha-exact-minus-frozen-1},
\eqref{eq:full-h-coupled-alpha-exact-minus-frozen-2}, and
\eqref{eq:full-h-coupled-alpha-exact-minus-frozen-3} proves
\eqref{eq:full-h-coupled-alpha-exact-to-frozen}.
\end{proof}

\begin{corollary}[Normalized growing-window endpoint bound]
\label{cor:full-h-coupled-alpha-endpoint}
Assume the hypotheses of Proposition~\ref{prop:moving-cutoff-decomposition} and the extra
regularity from Definition~\ref{def:model-c2}, and suppose moreover
that \(\chi\) is flat at the outer edge in the sense of Definition~\ref{def:endpoint-flat-cutoff}.
Then there exist \(h_{0,\alpha}\in(0,t]\) and \(C_{t,\alpha,\chi}<\infty\) such that for every
\[
0<h\le h_{0,\alpha}
\qquad\text{with}\qquad
h^\alpha\le t,
\]
and every \(v\in\mathfrak G_t\),
\begin{equation}\label{eq:full-h-coupled-alpha-endpoint}
\left\|
\mathfrak E_{t,h,\alpha,R_{h,\alpha},\chi}^{\mathrm{norm}}
\bigl(
\widehat\Phi_{v,h,\alpha},
\widehat Z_{v,h,\alpha}^{\mathrm{full}}
\bigr)
\right\|_{E}
\le
C_{t,\alpha,\chi}\|v\|_{\mathfrak G_t}.
\end{equation}
In particular, \eqref{eq:full-h-coupled-alpha-defect-normalized-bound} holds.
\end{corollary}

\begin{proof}
Theorem~\ref{thm:full-h-coupled-alpha-exact-to-frozen} splits the normalized exact defect into the
frozen commutator plus an error bounded by
\(C_{t,\alpha,\chi}h\|v\|_{\mathfrak G_t}\).  Corollary~\ref{cor:full-h-coupled-alpha-frozen-bound}
bounds the frozen commutator uniformly by \(C_{t,\alpha,\chi}\|v\|_{\mathfrak G_t}\).  After
shrinking \(h_{0,\alpha}\) if necessary, the two estimates give
\eqref{eq:full-h-coupled-alpha-endpoint}.  The final assertion is the same bound written in the
notation of \eqref{eq:full-h-coupled-alpha-defect-normalized-bound}.
\end{proof}

\begin{theorem}[Closed source-screening under endpoint-flat cutoffs]
\label{thm:observable-screening-closed}\label{cor:full-h-coupled-alpha-closed}
Assume Assumptions~\ref{ass:model} and~\ref{ass:obs-source}, assume in addition the extra
regularity from Definition~\ref{def:model-c2}, let \(\eta>0\), fix \(t\in(0,T)\) and
\(\alpha\in(0,1)\), and suppose
\[
0<H_{XY}<H_Y<\tfrac12.
\]
Choose
\[
\chi\in C^\infty([0,1]),
\qquad
0\le\chi\le1,
\qquad
\chi\equiv1\ \text{on }[0,1-\delta_*]
\]
for some \(\delta_*\in(0,1)\), and assume moreover that \(\chi\) is flat at the outer edge in the
sense of Definition~\ref{def:endpoint-flat-cutoff}. Then \(\mathcal S_t\) has
\(h^\alpha\)-stable boundary cutoffs at time \(t\), and
\begin{equation}\label{eq:observable-screening-closed-F}
\|\mathcal F_{t,h}\|_{\mathfrak G_t^*}
=
o(h^{H_{XY}})
\qquad\text{as }h\downarrow0,
\end{equation}
and consequently
\begin{equation}\label{eq:observable-screening-closed-G}
\widetilde G_{X\to Y}(t,h)
=
o(h^{2H_{XY}})
\qquad\text{as }h\downarrow0.
\end{equation}
\end{theorem}

The theorem is closed in a specific sense.  The cutoff stability required by
Theorem~\ref{thm:observable-screening-smooth} is verified by the endpoint-flat construction above.
The statement is not cutoff-free; the endpoint-flat cutoff and the \(C^2\) near-diagonal
regularity remain part of the screening mechanism.

\begin{proof}
Corollary~\ref{cor:full-h-coupled-alpha-endpoint} gives the normalized growing-window defect bound
\eqref{eq:full-h-coupled-alpha-defect-normalized-bound}.
Proposition~\ref{prop:full-h-coupled-alpha-defect-normalized} converts that bound into the
unnormalized estimate \eqref{eq:full-h-coupled-alpha-defect-bound}.
Corollary~\ref{cor:full-h-coupled-alpha-defect} therefore supplies the endpoint estimate required
by Corollary~\ref{cor:full-h-coupled-commutator}.  Hence \(\mathcal S_t\) has
\(h^\alpha\)-stable boundary cutoffs at time \(t\).  Applying
Theorem~\ref{thm:observable-screening-smooth} with these cutoffs gives
\eqref{eq:observable-screening-closed-F}; the gain estimate
\eqref{eq:observable-screening-closed-G} then follows from the graph-dual identity
\eqref{eq:graph-dual}.
\end{proof}

\begin{corollary}[Endpoint-flat screening from unscaled exact correction]
\label{cor:unscaled-smooth-cutoff-screening}
Assume the hypotheses of Theorem~\ref{thm:unscaled-smooth-cutoff-criterion}. Then
\begin{equation}\label{eq:unscaled-smooth-cutoff-screening-F}
\|\mathcal F_{t,h}\|_{\mathfrak G_t^*}
=
o(h^{H_{XY}})
\qquad\text{as }h\downarrow0,
\end{equation}
and consequently
\begin{equation}\label{eq:unscaled-smooth-cutoff-screening-G}
\widetilde G_{X\to Y}(t,h)
=
o(h^{2H_{XY}})
\qquad\text{as }h\downarrow0.
\end{equation}
\end{corollary}

\begin{proof}
Theorem~\ref{thm:unscaled-smooth-cutoff-criterion} proves that \(\mathcal S_t\) has
\(h^\alpha\)-stable boundary cutoffs at time \(t\).  Those cutoffs are the only additional
hypothesis in Theorem~\ref{thm:observable-screening-smooth}.  Substitution into that theorem gives
\eqref{eq:unscaled-smooth-cutoff-screening-F}, and the gain estimate follows from the graph-dual
identity \eqref{eq:graph-dual}.
\end{proof}

\section{Conclusion}\label{sec:conclusion}

Dynamic observability lies between threshold geometry and feasible inference.  The threshold analysis
of \cite{vidal2026latent} separates pricing visibility from path-space detectability in the
smoothing regime.  Dynamic observability is the residual predictive quotient after latent drivers
are replaced by observable histories.

The three levels have different boundaries.  At finite resolution, the canonical Gaussian experiment
generated by the relative covariance operator recovers
\[
H_{XY}=H_Y+\tfrac14
\]
as the critical boundary for evidence accumulation with resolution.  At the oracle latent-driver
level, the threshold is instead
\[
H_{XY}=H_Y,
\]
the sharp criterion for leading-order short-horizon prediction gain.  At the observable level, even
that oracle gain can be screened.  In the rougher regime \(0<H_{XY}<H_Y<\tfrac12\), the screening
theorem is first proved under an \(h^\alpha\)-stable boundary-cutoff hypothesis.  The
exact-correction closure theorem verifies that hypothesis for the model classes treated here under
the present second-order amplitude regularity assumptions.  In the endpoint-flat closure theorem,
the target asset's own past screens the dominant oracle contribution.

The conclusion is a three-level visibility structure for latent contagion.  Pricing visibility,
path-space detectability, and dynamic observability need not coincide.  Volterra covariance
operators govern the first comparison, Gaussian experiments govern the finite-resolution evidence
scale, and the hidden transfer operator governs observable screening.

The screened observable component is the estimand for projected source-screened inference.  The
oracle object from the dynamic theorem is reduced by conditioning on the target channel and
restricting information to observables; the remaining residual object is the inferential target.
The construction identifies that target and its screening boundary.

The remaining problems are local to the assumptions and quantitative to the rate.  The observable
screening theorem uses \(C^2\) amplitude regularity together with the cutoff-stability mechanism
verified in the closure argument.  The boundary of those assumptions is not settled.  A quantitative
refinement would start from Corollary~\ref{cor:observable-screening-rate}.  Once the near-layer
profile \(\omega_\alpha\) admits a rate, the global observable screening bound inherits it.  The
open problem is to derive such a rate directly from the exact-correction closure for the model
classes treated here.  In particular, if the near-layer defect satisfies
\[
\omega_\alpha(h)=O(h^\rho)
\qquad\text{for some }\rho>0,
\]
then Corollary~\ref{cor:observable-screening-rate} yields the stronger screening estimate
\[
\|\mathcal F_{t,h}\|_{\mathfrak G_t^*}
=
O(h^{H_{XY}+\varepsilon}),
\]
with
\[
\varepsilon
=
\min\!\bigl\{\rho,\;1-\alpha(1-H_{XY})-H_{XY},\;H_Y-H_{XY}\bigr\},
\]
provided this minimum is positive.  That quantitative rate problem belongs to finite-sample
analysis and remains outside the present closure theorem.  Extending the same observability problem
to nonlinear, non-Gaussian, or jump-driven rough models would require tools beyond the Gaussian
comparison and exact Volterra localization arguments used here.

\bibliographystyle{unsrt}
\bibliography{references}

\end{document}
